

% new file names from date 20180706: 
% block6_180607.tex 
% before was: NC-DC-d93.tex  

\renewcommand{\thisfile}{block6\_190328}  

\cleardoublepage
\mysection{Algebraically Graded Noncommutative Rings}
\label{sec:alg-gra-rings}

\AYcopyright

In this section we study {\em algebraically graded rings}, and the categories 
of {\em algebraically graded modules} over them. These are the graded rings 
that appear in standard texts on commutative and noncommutative algebra, and 
are distinct from the {\em cohomologically graded rings} that underlie {\em DG 
rings}. 

In the subsequent sections of the book (Sections \ref{sec:der-tors-conn} and 
\ref{sec:BDC}) we shall concentrate on {\em connected graded rings}
(traditionally known as {\em connected graded algebras}), that behave 
like ``complete local rings'' within the algebraically graded context. 

The interest in algebraically graded rings, and especially in connected graded 
rings, stems from the prominent role they played in {\em noncommutative 
algebraic geometry}, as it was developed by M. Artin and his collaborators 
since around 1990. See the papers \cite{ATV}, \cite{ArZh} and \cite{StaVdB}. 



\mysubsection{Categories of Algebraically Graded Modules} 
\label{subsec:alg-gr-mods}

Let $\K$ be a nonzero commutative base ring. An 
{\em algebraically graded $\K$-module}% 
\index{Algebraically graded module}
is a $\K$-module $M$ with a direct sum decomposition 
$M = \bigoplus_{i \in \Z} M_i$
into $\K$-submodules. The submodule $M_i$ is called the homogeneous component 
of $M$ of {\em algebraic degree $i$}.
If $m \in M_i$ is a nonzero element, then we write 
$\opn{deg}(m) = i$. 

In Subsection \ref{subsec:gr-alg} we discussed {\em cohomologically graded rings 
and  modules}. There are three features that distinguish between 
cohomologically graded $\K$-mod\-ules and algebraically graded $\K$-modules.
The first distinguishing feature is the use of upper vs.\ lower 
indices. The second is that in an algebraically graded module 
there is never a differential involved; i.e.\ it is 
not the underlying graded module of a DG module. The third and most important 
change is in {\em signs of permutations}, and thus in {\em commutativity}. In 
the cohomologically graded setting the Koszul sign rule 
dictates that for graded $\K$-modules $M$ and $N$, the braiding isomorphism
$\opn{br}_{M, N} : M \ot N \to N \ot M$,
where $\ot := \ot_{\K}$, is 
$\opn{br}_{M, N}(m \ot n) = (-1)^{i \cd j} \cd n \ot m$
for homogeneous elements $m \in M^i$ and $n \in N^j$. 
But in the algebraically graded setting there are no signs: 
$\opn{br}_{M, N}(m \ot n) := n \ot m$
for $m \in M_i$ and $n \in N_j$. 
This is reflected in Definition \ref{dfn:3738} below. 

Let $M = \bigoplus_{i \in \Z} \, M_i$ be an algebraically graded $\K$-module.
A nonzero element $m \in M$ can be expressed uniquely as a sum 
\begin{equation} \label{eqn:4260}
m = m_{1} + \cdots + m_{r} , 
\end{equation}
such that $r \geq 1$; each $m_{i}$ is a nonzero homogeneous element; and 
$\opn{deg}(m_{i}) < \opn{deg}(m_{i + 1})$. 
This is called the {\em homogeneous component decomposition} of $m$. 
Of course $m$ is homogeneous if and only if $r = 1$. 

Suppose $M$ and $N$ are algebraically graded $\K$-modules. A $\K$-linear 
homomorphism $\phi : M \to N$
is said to be of {\em algebraic degree $i$} if
$\phi(M_j) \sub N_{j + i}$ for all $j \in \Z$. 
The $\K$-module of all $\K$-linear homomorphisms of degree $i$ is denoted by 
$\opn{Hom}_{\K}(M, N)_i$. Taking the direct sum we get an algebraically graded 
$\K$-module
\begin{equation}  \label{eqn:3701}
\opn{Hom}_{\K}(M, N) := \bigoplus_{i \in \Z} \, \opn{Hom}_{\K}(M, N)_i . 
\end{equation}

\begin{dfn}  \label{dfn:3700}
Let $M$ and $N$ be algebraically graded $\K$-modules. 
A {\em strict homomorphism of algebraically graded $\K$-modules}
\index{Algebraically graded module! strict homomorphism of {\indash}s}
$\phi : M \to N$ is $\K$-linear homomorphism of algebraic degree $0$.

The {\em category of algebraically graded $\K$-modules} is the category 
$\dcat{M}(\K, \mrm{gr})$, 
\index{1-M(K,gr)@$\dcat{M}(\K, \mrm{gr})$}
whose objects are the algebraically graded 
$\K$-modules, and whose morphisms are the strict homomorphisms.
\end{dfn}

Thus 
\[ \opn{Hom}_{\dcat{M}(\K, \mrm{gr})}(M, N) = \opn{Hom}_{\K}(M, N)_0 . \]
It is easy to see that $\dcat{M}(\K, \mrm{gr})$ is a $\K$-linear abelian 
category. The kernels and cokernels are degreewise. 

Suppose $M, N \in \dcat{M}(\K, \mrm{gr})$.
Their tensor product is also algebraically graded:
\[ (M \ot N)_i = \bigoplus_{j \in \Z} \, (M_j \ot N_{i - j}) \]
and 
\[ M \ot N = \bigoplus_{i \in \Z} \, (M \ot N)_i . \]

\begin{dfn}  \label{dfn:3701}
An {\em algebraically graded central $\K$-ring}
\index{Algebraically graded ring}
is a central $\K$-ring $A$, equipped with a direct sum decomposition 
$A = \bigoplus_{i \in \Z} A_i$ 
into $\K$-submodules, such that $1_A \in A_0$, and 
$A_i \cd A_j \sub A_{i + j}$ 
for all $i, j$.

A {\em homomorphism of algebraically graded central $\K$-rings}
\index{Algebraically graded ring! homomorphism of {\indash}s}
$f : A \to B$ is a $\K$-ring homomorphism such that 
$f(A_i) \sub B_i$ for all $i$. 

The algebraically graded central $\K$-rings form a category, that we
denote by $\catt{Rng}_{\mrm{gr}} \centover \K$.
\end{dfn}

The base ring $\K$ is an algebraically graded ring, concentrated in degree $0$.
It is the initial object of the category 
$\catt{Rng}_{\mrm{gr}} \centover \K$. 

In the traditional ring theory literature (e.g.\ \cite{McRo}, \cite{Row}, 
\cite{ATV}, \cite{ArZh}, \cite{Ye1} and \cite{YeZh0}) the algebraically graded 
central $\K$-rings were called ``associative unital graded $\K$-algebras''. 

\begin{exa} \label{exa:3701}
Let $M$ be an algebraically graded $\K$-module. Then 
$\opn{End}_{\K}(M) \lb  := \opn{Hom}_{\K}(M, M)$
is an algebraically graded central $\K$-ring.
The grading is according to (\ref{eqn:3701}), and multiplication is 
composition. 
\end{exa}

\begin{dfn}  \label{dfn:3702}
Let $A$ be an algebraically graded central $\K$-ring.
\begin{enumerate}
\item An {\em algebraically graded left $A$-module} 
\index{Algebraically graded module}
is a left $A$-module $M$, equipped with a direct sum decomposition 
$M = \bigoplus_{i \in \Z} \, M_i$
into $\K$-submodules, such that 
$A_i \cd M_j \sub M_{i + j}$ for all $i, j$.

\item Suppose $M$ and $N$ are algebraically graded left $A$-modules.
A $\K$-linear homomorphism $\phi : M \to N$ of degree $i$ is said to be a
{\em homomorphism of algebraically graded $A$-modules}
\index{Algebraically graded module! homomorphism of {\indash}s}
if $\phi(a \cd m) = a \cd \phi(m)$
for all $a \in A$ and $m \in M$. 

\item The $\K$-module of all $A$-linear homomorphisms of degree $i$ is denoted 
by $\opn{Hom}_{A}(M, N)_i$. Taking the direct sum we get an algebraically 
graded 
$\K$-module
\[ \opn{Hom}_{A}(M, N) := \bigoplus_{i \in \Z} \, \opn{Hom}_{A}(M, N)_i . \]
\end{enumerate}
\end{dfn}

Note that there is an inclusion 
$\opn{Hom}_{A}(M, N) \sub \opn{Hom}_{\K}(M, N)$
in $\dcat{M}(\K, \mrm{gr})$.

\begin{dfn}  \label{dfn:3704}
Let $A$ be an algebraically graded central $\K$-ring.
\begin{enumerate}
\item Suppose $M$ and $N$ are algebraically graded left $A$-modules.
The elements of $\opn{Hom}_{A}(M, N)_0$ are called 
{\em strict homomorphisms of algebraically graded $A$-modules}. 
\index{Algebraically graded module! strict homomorphism of {\indash}s}

\item The category of algebraically graded left $A$-modules, with strict 
homomorphisms, is denoted by $\dcat{M}(A, \mrm{gr})$. 
\index{1-M(A,gr)@$\dcat{M}(A, \mrm{gr})$}
\end{enumerate}
\end{dfn}

Thus 
\begin{equation}  \label{eqn:3736}
\opn{Hom}_{\dcat{M}(A, \mrm{gr})}(M, N) = \opn{Hom}_{A}(M, N)_0 . 
\end{equation}

Clearly, to put an algebraically graded left $A$-module structure on an
algebraically graded $\K$-module $M$ is the same as to give a
homomorphism $A \to \lb \opn{End}_{\K}(M)$ in 
$\catt{Rng}_{\mrm{gr}} \centover \K$.

Algebraically graded right $A$-modules, and algebraically graded 
$A$-$B$-bi\-modules, are defined similarly to Definition \ref{dfn:3702}. There 
are no signs anywhere. 
All homomorphisms are $\K$-linear, and all bimodules are central over $\K$. 

\begin{dfn}  \label{dfn:3735}
Let $A$ be an algebraically graded ring. Given an algebraically graded 
$A$-module $M$ and an integer $i$, the {\em $i$-th degree twist of $M$} 
\index{Algebraically graded module! degree twist of}
\index{1-M(i)@$M(i)$}
is the algebraically graded $A$-module $M(i)$ whose degree $j$ homogeneous 
component is $M(i)_j := M_{i + j}$.
The action of $A$ is not changed. 
\end{dfn}

The degree twist is an automorphism of the abelian category 
$\dcat{M}(A, \mrm{gr})$. And 
\[ \opn{Hom}_{\dcat{M}(A, \mrm{gr})}(M, N(i)) = \opn{Hom}_{A}(M, N)_i . \]
There is a similar twisting for graded right modules and for graded 
bimodules. 

There is a functor 
$\opn{Ungr} : \catt{Rng}_{\mrm{gr}} \centover \K \to 
\catt{Rng}_{} \centover \K$
that forgets the grading. For a fixed ring 
$A \in  \catt{Rng}_{\mrm{gr}} \centover \K$
there is a forgetful functor 
\begin{equation} \label{eqn:3715}
\opn{Ungr} : \dcat{M}(A, \mrm{gr}) \to \dcat{M}(\opn{Ungr}(A)) . 
\end{equation}
This is an exact faithful $\K$-linear functor. 

\begin{rem} \label{rem:3795}
The functor $\opn{Ungr}$ in (\ref{eqn:3715}) is usually not full, for two 
reasons. First, there could be nonzero homomorphisms of nonzero 
degree, so that 
\[ \opn{Hom}_{\dcat{M}(A, \mrm{gr})}(M, N) = 
\opn{Hom}_A(M, N)_0 \subsetneqq \opn{Hom}_A(M, N) . \]
Second, it is easy find an example where
\[ \opn{Ungr} \bigl( \opn{Hom}_A(M, N) \bigr) \subsetneqq
\opn{Hom}_{\opn{Ungr}(A)} \bigl( \opn{Ungr}(M), \opn{Ungr}(N) \bigr) . \]
Just take 
$M := \bigoplus_{i \in \Z} \, A(i)$ and $N := A$. Then
\[ \opn{Ungr} \bigl( \opn{Hom}_A(M, N) \bigr) \cong 
\bigoplus_{i \in \Z} \, \opn{Ungr}(A)  \]
and 
\[ \opn{Hom}_{\opn{Ungr}(A)} \bigl( \opn{Ungr}(M), \opn{Ungr}(N) \bigr) \cong 
\prod_{i \in \Z} \, \opn{Ungr}(A) \,  \]
in $\dcat{M}(\K)$.
\end{rem}

\begin{rem} \label{rem:3900}
Homomorphisms between algebraically graded modules and cohomologically graded 
modules are distinct, not only in the position of the degree indices. 
Compare Definition \ref{dfn:3702}(2) to Definition \ref{dfn:1690}. 
This distinction is due to the Koszul sign rule, that is present only in the 
cohomologically graded setting. 
\end{rem}

\begin{dfn}  \label{dfn:3738}
Let $A$ and $B$ be algebraically graded central $\K$-rings. 
\begin{enumerate}
\item The {\em opposite ring} of $A$ is the central $\K$-ring 
$A^{\mrm{op}}$, that is the same algebraically graded $\K$-module as $A$, with 
an isomorphism
$\opn{op} := \opn{id} : A \iso A^{\mrm{op}}$
in $\dcat{M}(\K, \mrm{gr})$. The multiplication $\cd^{\mrm{op}}$ of 
$A^{\mrm{op}}$ is 
\[ \opn{op}(a_1) \, \cd^{\mrm{op}} \, \opn{op}(a_2) := \opn{op}(a_2 \cd a_1) . 
\]

\item The ring $A$ is called a {\em commutative graded ring} 
\index{Algebraically graded ring! commutative}
if  $a_2 \cd a_1 = a_1 \cd a_2$ for all $a_1, a_2 \in A$. 

\item The tensor product $A \ot B$ is made into an algebraically graded ring 
with multiplication 
\[ (a_1 \ot b_1) \cd (a_2 \ot b_2) :=  
(a_1 \cd a_2) \ot (b_1 \ot b_2)  . \]

\item The {\em enveloping} ring of $A$ is the  algebraically graded ring 
$A^{\mrm{en}} := A \ot A^{\mrm{op}}$.
\end{enumerate}
\end{dfn}

In terms of opposite rings, an algebraically graded ring $A$ is commutative if 
and only if $A = A^{\mrm{op}}$. 
The category of algebraically graded right $A$-modules is identified with  
$\dcat{M}(A^{\mrm{op}}, \mrm{gr})$; 
and the category of $\K$-central algebraically graded $A$-$B$-bimodules is  
identified with $\dcat{M}(A \ot B^{\mrm{op}}, \mrm{gr})$. 

There are graded ring homomorphisms 
$A \to A \ot B \lar B$,
and corresponding restriction functors 
\begin{equation} \label{eqn:3740}
\dcat{M}(A, \mrm{gr}) \xlar{\opn{Rest}_A}
\dcat{M}(A \ot B, \mrm{gr}) \xar{\opn{Rest}_B}
\dcat{M}(B, \mrm{gr})  . 
\end{equation}
We shall usually omit any explicit reference to these restriction 
functors.

Given $M \in \dcat{M}(A, \mrm{gr})$ and 
$N \in \dcat{M}(B, \mrm{gr})$, the tensor product 
$M \ot N$ belongs to 
$\dcat{M}(A \ot B, \mrm{gr})$. 

Suppose $M \in \dcat{M}(A^{\mrm{op}}, \mrm{gr})$
and $N \in \dcat{M}(A, \mrm{gr})$.
Their tensor product over $A$ is also algebraically graded:
\begin{equation} \label{eqn:3720}
M \ot_A N = \bigoplus_{i \in \Z} \, (M \ot_A N)_i 
\in \dcat{M}(\K, \mrm{gr}) , 
\end{equation}
where $(M \ot_A N)_i$ is the image of $(M \ot N)_i$ under the canonical 
surjection 
$M \ot N \surj M \ot_A N$
in $\dcat{M}(\K, \mrm{gr})$. 
See Lemma \ref{lem:1690} for the corresponding statement in the cohomologically 
graded setting. 

There is an isomorphism of algebraically graded $\K$-rings
\begin{equation} \label{3721}
(A^{\mrm{en}})^{\mrm{op}} \iso A^{\mrm{en}} , \quad 
a_1 \ot a_2 \mapsto a_2 \ot a_1 . 
\end{equation}

From here on in this section, and in Sections \ref{sec:der-tors-conn}
and \ref{sec:BDC}, we assume the following convention. 

\begin{conv}  \label{conv:3700}
We fix a nonzero commutative base ring $\K$. The symbol $\ot$ means $\ot_{\K}$. 
All graded rings are assumed to be 
algebraically graded central $\K$-rings, and all graded ring homomorphisms are 
over $\K$. I.e.\  we work inside the category 
$\catt{Rng}_{\mrm{gr}} \centover \K$.

Let $A$ be a graded ring. Its enveloping ring is the graded ring 
$A^{\mrm{en}} := A \ot A^{\mrm{op}}$. 
All graded $A$-modules are assumed by default to be 
algebraically graded left $A$-modules.
For $M, N \in \dcat{M}(A, \mrm{gr})$ the expression 
$\opn{Hom}_A(M, N)$ refers to the graded $\K$-module in Definition 
\ref{dfn:3702}(3). 
All graded bimodules are assumed to be $\K$-central;  i.e.\ a graded
$A$-$B$-bimodule is an object of
$\dcat{M}(A \ot B^{\mrm{op}}, \mrm{gr})$.
\end{conv}

\begin{prop} \label{prop:3825}
Let $A$ be a graded ring. 
\begin{enumerate}
\item The category $\dcat{M}(A, \mrm{gr})$ is abelian. 

\item The category $\dcat{M}(A, \mrm{gr})$ has inverse limits. 

\item The category $\dcat{M}(A, \mrm{gr})$ has direct limits. Moreover, direct 
limits in $\dcat{M}(A, \mrm{gr})$ commute with the ungrading functor
\tup{(\ref{eqn:3715})}.
\end{enumerate}
\end{prop}

\begin{proof}
(1) The kernels and cokernels in $\dcat{M}(A, \mrm{gr})$ are degreewise, so the 
abelian property of $\dcat{M}(A, \mrm{gr})$ is inherited from the 
abelian property of $\dcat{M}(\K)$.

\medskip \noindent
(2) Inverse limits in  $\dcat{M}(A, \mrm{gr})$ are taken degreewise. 
Namely, given an inverse system  
$\{ M_x \}_{x \in X^{\mrm{op}}}$ of objects of $\dcat{M}(A, \mrm{gr})$,
indexed by a directed set $X$, let 
$M_i := \lim_{\lar x} (M_x)_i \in \dcat{M}(\K)$.
Then 
$M := \bigoplus_{i \in \Z} M_i \in \dcat{M}(A, \mrm{gr})$
is the inverse limit of $\{ M_x \}_{x \in X^{\mrm{op}}}$.

\medskip \noindent
(3) Direct limits in $\dcat{M}(A, \mrm{gr})$ are taken degreewise. 
Namely, given a direct system  
$\{ N_y \}_{y \in Y}$ of objects of $\dcat{M}(A, \mrm{gr})$,
indexed by a directed set $Y$, let 
$N_i := \lim_{y \to} (N_y)_i \in \dcat{M}(\K)$.
Then 
$N := \bigoplus_{i \in \Z} N_i \in \dcat{M}(A, \mrm{gr})$
is the direct limit of $\{ N_y \}_{y \in Y}$.

Because direct sums and direct limits commute with each other, it follows that 
$\lim_{y \to}$ commutes with $\opn{Ungr}$. 
\end{proof}

\begin{dfn} \label{dfn:4540}
A graded ring $A$  is called {\em nonnegative} 
\index{Algebraically graded ring! nonnegative}
if $A_i = 0$ for all $i < 0$; i.e.\ if $A = \bigoplus_{i \geq 0} A_i$.

If $A$ is a nonnegative graded ring $A$, then the two-sided graded ideal
$\m :=  \bigoplus_{i \geq 1} A_i$ is called the {\em augmentation ideal}. 
\index{Algebraically graded ring! augmentation ideal of}
It is the kernel of the surjective graded ring homomorphism $A \to A_0$, that we 
call the {\em augmentation homomorphism}. 

The augmentation ideal of the graded ring $A^{\mrm{op}}$ is denoted by 
$\m^{\mrm{op}}$. 
\end{dfn}

\begin{dfn} \label{dfn:4541}
Let $A$ be a graded ring $A$ and let 
$M = \bigoplus_{i \in \Z} M_i$ be a graded $A$-module.
\begin{enumerate}
\item We call $M$ a {\em bounded above} graded module if $M_i = 0$ for 
$i \gg 0$.

\item We call $M$ a {\em bounded below} graded module if $M_i = 0$ for 
$i \ll 0$.

\item We call $M$ a {\em bounded} graded module if it is both bounded above and 
bounded below.
\end{enumerate}
\end{dfn}

\begin{exa} \label{exa:4540}
If $A$ is a nonnegative graded ring, then as a graded module $A$ is bounded 
below. For such $A$, the bounded below graded modules behave like finitely 
generated modules over a noetherian local commutative ring -- see Proposition 
\ref {prop:4540} below.
\end{exa}

The next proposition is the first version of the graded Nakayama Lemma that we 
present; the second is Proposition \ref{prop:4590}. 

\begin{prop} \label{prop:4540}
Let $A$ be a nonnegative graded ring, with augmentation ideal $\m$. Let $M$ be 
a bounded below graded $A$-module, and let $M' \sub M$ be a graded 
$A$-submodule. If $M' + \m \cd M = M$ then $M' = M$. 
\end{prop}

\begin{proof}
Say $M = \bigoplus_{i \geq i_0} M_i$ for an integer $i_0$.
We will prove that $M'_i = M_i$
as $A_0$-modules for all $i$, by induction on $i \geq i_0 - 1$.
For $i =  i_0 - 1$ this is clear because both these $A_0$-modules are zero. 

Now take $i \geq i_0$, and assume that $M'_j = M_j$ for all  $j < i$. 
Suppose $m \in M_i$ is a nonzero element. 
Because $M = M' + \m \cd M$, we can express $m$ as follows:
$m = m' + \sum_{k = 1}^r a_k \cd m_k$,
where $m' \in M'_i$; $r \geq 0$; $a_k \in \m$ and $m_k \in M$ are 
nonzero homogeneous elements; 
and $\opn{deg}(a_k) + \opn{deg}(m_k) = i$.
But then $\opn{deg}(m_k) < i$, so by the induction hypothesis
we have $m_k \in M'$. We conclude that $m \in M'_i$.
\end{proof}

\begin{dfn}  \label{dfn:3706}
Let $A$ be a graded ring. 
\begin{enumerate}
\item A graded $A$-module $M$ is called a {\em graded-finite  
$A$-module} if $M$ can be generated by finitely many homogeneous elements. 

\item The full subcategory of $\dcat{M}(A, \mrm{gr})$ on the graded-finite
$A$-modules is denoted by $\dcat{M}_{\mrm{f}}(A, \mrm{gr})$.
\index{1-Mf(A,gr)@$\dcat{M}_{\mrm{f}}(A, \mrm{gr})$}

\item $A$ is called a {\em left graded-noetherian  ring} if every graded left 
ideal $\a \sub A$ is a graded-finite  $A$-module, in the sense of item (1) 
above. 

\item $A$ is called a {\em graded-noetherian  ring} if both $A$ and 
$A^{\mrm{op}}$ are left graded-noetherian  rings. 
\end{enumerate}
\end{dfn}

An obvious modification of item (1) above gives the definition of a 
graded-finite  right $A$-module. Of course the ring $A$ is called a right 
graded-noetherian  ring if $A^{\mrm{op}}$ is a left graded-noetherian  ring.

\begin{prop} \label{prop:4260}
If $A$ is left graded-noetherian ring, then $\dcat{M}_{\mrm{f}}(A, 
\mrm{gr})$ is a thick abelian subcategory of $\dcat{M}(A, \mrm{gr})$. 
\end{prop}

\begin{exer} \label{exer:4261}
Prove this proposition. 
\end{exer}

\begin{prop} \label{prop:3745}
Let $A$ be a graded ring and $M$ a graded $A$-module. 
The following two conditions are equivalent.
\begin{itemize}
\rmitem{i} $M$ is a graded-finite  $A$-module, in the sense of Definition 
\tup{\ref{dfn:3706}(1)}.

\rmitem{ii} $\opn{Ungr}(M)$ is a finitely generated $\opn{Ungr}(A)$-module. 
\end{itemize}
\end{prop}

\begin{exer} \label{exer:3745}
Prove this proposition. 
\end{exer}

\begin{thm} \label{thm:4245}
Let $A$ be a graded ring. The following two conditions are equivalent.
\begin{itemize}
\rmitem{i} The graded ring $A$ is left graded-noetherian, in the sense of 
Definition \tup{\ref{dfn:3706}(3)}.

\rmitem{ii} The ring $\opn{Ungr}(A)$ is left noetherian.  
\end{itemize}
\end{thm}

\begin{proof}
The case when $A$ is nonnegative is proved in many textbooks, and it is a 
nice exercise for the reader. 

The general case is much harder, and the few proofs in the 
literature are quite complicated. Here is a short proof, that is due to S. 
Snigerov. 

The implication (ii) $\Rightarrow$ (i) is trivial. So let us assume that $A$ is 
left graded-noetherian, and we shall prove that $\opn{Ungr}(A)$ 
is left noetherian. The proof proceeds in a few steps. 

\medskip \noindent
Step 1. Consider an $A_0$-submodule $V \sub A_j$ for some $j$, and let  
$\b := A \cd V \sub A$, which is a graded left ideal of $A$. 
By assumption, $\b$ is generated by finitely many homogeneous elements. 
An easy calculation shows that $V = \b \cap A_j$, and that 
$V$ is finitely generated as an $A_0$-module. The consequence is that  
the ring $A_0$ is left noetherian, and that every $A_j$ is a finitely generated 
$A_0$-module. 

\medskip \noindent
Step 2. Now we fix a nonzero left ideal $L \sub A$. We are going to find a 
finite number of generators for $L$ as an $A$-module. 

For a nonzero element 
$a \in L$ we consider its homogeneous component decomposition 
(\ref{eqn:4260}), and we define its bottom degree component
$\opn{bot}(a) := a_{1}$, and its top degree component 
$\opn{top}(a) := a_{r}$.
Let $\a_{\mrm{top}}$ be the left ideal in $A$ generated by the elements
$\opn{top}(a)$, where $a$ runs over all the nonzero elements of $L$. By 
assumption the graded ideal $\a_{\mrm{top}}$ is graded-finite; so we 
can choose a finite collection $\{ a_x \}_{x \in X_{\mrm{top}}}$ of nonzero 
elements of $L$, such that their top degree components $\opn{top}(a_x)$ 
generate the left ideal $\a_{\mrm{top}}$. Now let 
\[ d_1 := \opn{max} \, \bigl\{ \opn{deg}(\opn{top}(a_x)) \mid x \in 
X_{\mrm{top}} \bigr\} \in \Z . \]

\medskip \noindent
Step 3. Define $\a_{\mrm{bot}}$ to be the left ideal in $A$ generated by the 
elements $\opn{bot}(a)$, where $a$ runs over all the nonzero elements of 
$L \cap A_{\leq d_1}$. (It is possible that there are no such elements.)
By assumption the graded left ideal $\a_{\mrm{bot}}$ is graded-finite; 
so we can choose a finite collection $\{ a_x \}_{x \in X_{\mrm{bot}}}$ of 
nonzero elements of $L \cap A_{\leq d_1}$, such that their bottom degree 
components $\opn{bot}(a_x)$ generate the left ideal $\a_{\mrm{bot}}$. 
Now we let 
\[ d_0 := \opn{min} \, \bigl( \bigl\{ \opn{deg}(\opn{bot}(a_x)) \mid x \in 
X_{\mrm{bot}} \bigr\} \cup \{ 0 \} \bigr) \in \Z . \]

\medskip \noindent
Step 4. By step 1 the $A_0$-module
$V := L \cap \bigl( \bigoplus\nolimits_{i = d_0}^{d_1} A_i \bigr)$
is finitely generated. We choose a finite collection 
 $\{ a_x \}_{x \in X_{\mrm{mid}}}$ 
that generates $V$ as an $A_0$-module.

\medskip \noindent
Step 5. Define the finite indexing set  
$X := X_{\mrm{top}} \sqcup X_{\mrm{bot}} \sqcup X_{\mrm{mid}}$.
We claim that the collection $\{ a_x \}_{x \in X}$ 
generates the left ideal $L$. Take any element $a \in L$. 
By subtracting from $a$ a finite linear combination of elements 
from $\{ a_x \}_{x \in X_{\mrm{top}}}$, with coefficients taken from 
$A_{\geq 1}$, we can assume that 
$a \in L \cap A_{\leq d_1}$. 
Next, by subtracting from $a$ a finite linear combination of elements 
from $\{ a_x \}_{x \in X_{\mrm{bot}}}$, with coefficients taken from 
$A_{\leq -1}$, we can assume that $a \in V$, the $A_0$-submodule from step 4.
So $a$ is finite linear combination of elements 
from $\{ a_x \}_{x \in X_{\mrm{mid}}}$, with coefficients taken from $A_0$. 
\end{proof}

In view of Proposition \ref{prop:3745} and Theorem \ref{thm:4245}, we can use 
this more relaxed terminology without danger of ambiguity: 

\begin{dfn} \label{dfn:4262}
Let $A$ be a graded ring. 
\begin{enumerate}
\item We call $A$ a {\em left noetherian graded ring} 
\index{Algebraically graded ring! left noetherian}
if it is left graded-noetherian. Likewise for the attributes ``right noetherian 
graded ring'' and  ``noetherian graded ring''. 

\item Let $M$ be a graded $A$-module. We call $M$ a 
{\em finite graded $A$-module}
\index{Algebraically graded module! finite}
if it is graded-finite. Likewise for right modules. 
\end{enumerate}
\end{dfn}

We do not assume the rings here are noetherian. Partly this is because of the 
facts in the remark below. Instead we resort to {\em derived 
graded-pseudo-finite} modules and complexes, see Definition \ref{dfn:4532}. 

\begin{rem} \label{rem:3901}
Unlike the commutative theory, the Hilbert Basis Theorem does not hold in the 
noncommutative setting. Thus, a finitely generated $\K$-ring $A$ need 
not be noetherian; see Example \ref{exa:3755}. Likewise, if $A$ and $B$ are 
noetherian finitely generated $\K$-rings, the tensor product $A \ot B$ need not 
be noetherian; there is a counterexample in \cite{Rog}. 
See \cite{ArSmZh} for a deep discussion of permanence properties 
of NC noetherian rings. 
\end{rem}

We now pass to complexes of algebraically graded modules, retaining Convention 
\ref{conv:3700}. 
In Subsection \ref{subsec:complexes} we introduced the category 
$\dcat{C}(\cat{M})$ of complexes with entries in an abelian category $\cat{M}$.
This is a DG category. Its strict subcategory 
$\dcat{C}_{\mrm{str}}(\cat{M})$ is abelian. Here we take the abelian category 
$\cat{M} := \dcat{M}(A, \mrm{gr})$ for a graded ring $A$. 

\begin{dfn} \label{dfn:3730}
Let $A$ be a graded ring.
\begin{enumerate}
\item The category of complexes with entries in the abelian category 
$\dcat{M}(A, \mrm{gr})$ is the DG category 
$\dcat{C}(A, \mrm{gr}) := \dcat{C}(\dcat{M}(A, \mrm{gr}))$%
\index{1-C(A,gr)@$\dcat{C}(A, \mrm{gr})$}.

\item The strict subcategory of 
$\dcat{C}(A, \mrm{gr})$ is the abelian category  
$\dcat{C}_{\mrm{str}}(A, \mrm{gr}) := 
\dcat{C}_{\mrm{str}}(\dcat{M}(A, \mrm{gr}))$% 
\index{1-Cstr(A,gr)@$\dcat{C}_{\mrm{str}}(A, \mrm{gr})$}.

\item Given a boundedness indicator $\star$, we denote by 
$\dcat{C}^{\star}(A, \mrm{gr})$ the full subcategory of 
$\dcat{C}(A, \mrm{gr})$ on the complexes $M$ having boundedness condition 
$\star$. 
\end{enumerate}
\end{dfn}

We now try to make Definition \ref{dfn:3730} more explicit. 
An object $M \in \dcat{C}(A, \mrm{gr})$ has these direct sum decompositions:
\begin{equation} \label{eqn:3730}
M = \bigoplus_{i \in \Z} \, M^i = \bigoplus_{i, j \in \Z} \, M^i_j . 
\end{equation}
Here each $M^i \in \dcat{M}(A, \mrm{gr})$, and each
$M^i_j \in \dcat{M}(\K)$. 
We call $M^i_j$ the homogeneous component of $M$ of bidegree 
$\sbmat{i \\[0em] j} \in \sbmat{\Z \\[0.1em] \Z} = \Z^2$.
The upper index $i$ is called the {\em cohomological degree}, and the lower 
index $j$ is called the {\em algebraic degree}. 
The differential $\d_{M}$ of the complex $M$ is an $A$-linear homomorphism 
$\d_M : M \to M$ of bidegree $\sbmat{1 \\[0.1em] 0}$
that satisfies $\d_M \circ \d_M = 0$. 

A morphism
$\phi : M \to N$ in the strict category 
$\dcat{C}_{\mrm{str}}(A, \mrm{gr})$
has bidegree $\sbmat{0 \\[0.1em] 0}$ and satisfies 
$\phi \circ \d_M = \d_N \circ \phi$.

Lastly, to clarify item (3) in the definition, a complex 
$M$ belongs to $\dcat{C}^{+}(A, \mrm{gr})$
if $M^i = 0$ for $i \ll 0$, etc. Indicators can be combined, so 
$\dcat{C}^{\star}_{\mrm{str}}(A, \mrm{gr})$
is a thick abelian subcategory of 
$\dcat{C}^{}_{\mrm{str}}(A, \mrm{gr})$.

Like for any abelian category, we have a homotopy category and a derived 
category. Here is the specific notation:

\begin{dfn} \label{dfn:3731} 
Let $A$ be a graded ring.
\begin{enumerate}
\item The homotopy category of complexes with entries in the abelian category 
$\dcat{M}(A, \mrm{gr})$ is the triangulated category%
\index{1-K(A,gr)@$\dcat{K}(A, \mrm{gr})$}
$\dcat{K}(A, \mrm{gr}) := \dcat{K}(\dcat{M}(A, \mrm{gr}))$.

\item The derived category of the abelian category $\dcat{M}(A, \mrm{gr})$ is 
the triangulated category% 
\index{1-D(A,gr)@$\dcat{D}(A, \mrm{gr})$}
$\dcat{D}(A, \mrm{gr}) := \dcat{D}(\dcat{M}(A, \mrm{gr}))$.

\item Given a boundedness indicator $\star$, we denote by 
$\dcat{D}^{\star}(A, \mrm{gr})$ the full triangulated subcategory of 
$\dcat{D}(A, \mrm{gr})$ on the complexes $M$ whose cohomology $\opn{H}(M)$
has boundedness condition $\star$. 
\end{enumerate}
\end{dfn}

As always for derived categories, there are functors
\begin{equation} \label{eqn:3961}
\UseTips \xymatrix @C=6ex @R=6ex {
\dcat{C}^{}_{\mrm{str}}(A, \mrm{gr})
\ar[r]^{\opn{P}}
\ar@(ur,ul)[rr]^{\til{\opn{Q}}}
&
\dcat{K}(A, \mrm{gr})
\ar[r]^{\opn{Q}}
&
\dcat{D}(A, \mrm{gr}) \, . 
} 
\end{equation}
By abuse of notation we sometimes write $\opn{Q}$ rather than 
$\til{\opn{Q}}$; or, when there is no danger of confusion, we suppress the 
categorical localization functors $\opn{Q}$ and $\til{\opn{Q}}$ altogether. 

There is yet another category associated to $A$~: it is the category 
$\dcat{G}(A, \mrm{gr}) :=  \dcat{G}(\dcat{M}(A, \mrm{gr}))$
of cohomologically graded objects in $\dcat{M}(A, \mrm{gr})$. It has 
its strict subcategory $\dcat{G}_{\mrm{str}}(A, \mrm{gr})$. 
This is the target of the cohomology functor 
\[ \opn{H} : \dcat{C}_{\mrm{str}}(A, \mrm{gr})
\to \dcat{G}_{\mrm{str}}(A, \mrm{gr}) . \]

If $A \to B$ is a graded ring homomorphism that makes $B$ into a finite left 
$A$-module, and if $A$ is left noetherian, then of course $B$ is also left 
noetherian. Here is a partial converse: 

\begin{thm} \label{thm:4240}
Let $A$ be a nonnegative graded ring, and let $a \in A$ be a nonzero 
homogeneous central element of positive degree. Define the nonnegative graded 
ring $B := A / (a)$. If $B$ is left noetherian, then $A$ is also left 
noetherian. 
\end{thm}

This result is very similar to \cite[Lemma 8.2]{ATV}. There are 
three changes: we do not assume that the element $a$ is regular (i.e.\ not a 
zero-divisor), and we do not assume that $A$ is connected graded (only 
nonnegative). On the other hand, in \cite{ATV} the element $a$ is only 
required to be normalizing, not central. Our proof is essentially copied from 
\cite{ATV}. 

\begin{proof}
Assume, for the sake of contradiction, that the graded ring $A$ is not left 
noetherian. Then (taking Theorem \ref{thm:4245} into account) there is a graded 
left ideal $L \sub A$ that is not finitely generated. By Zorn's Lemma we can 
assume that $L$ is maximal with this property. 
Thus every graded left ideal $L^{\dag}$ such that 
$L \subneq L^{\dag} \sub A$ must be finitely generated.  

Consider the graded $A$-module $M := A / L$. Every graded $A$-module 
$M^{\dag} \sub M$ is of the form $M^{\dag} = L^{\dag} / L$, and therefore 
$M^{\dag}$ is a finitely generated $A$-module. 
In other words, $M$ is a noetherian object of $\dcat{M}(A, \mrm{gr})$.

Let $i := \opn{deg}(a)$. We examine the commutative diagram (\ref{4590}) in 
$\dcat{M}(A, \mrm{gr})$. 
Viewing the columns of this diagram as complexes, whose nonzero terms are in 
cohomological degrees $0$ and $1$, this is a short exact sequence in the 
abelian category $\dcat{C}_{\mrm{str}}(A, \mrm{gr})$. 
\begin{equation} \label{4590}
\UseTips \xymatrix @C=6ex @R=6ex {
0
\ar[r]
&
L
\ar[r]
&
A
\ar[r]
&
M
\ar[r]
&
0
\\
0
\ar[r]
&
L(-i)
\ar[u]^{a \cd (-)}
\ar[r]
&
A(-i)
\ar[u]^{a \cd (-)}
\ar[r]
&
M(-i)
\ar[u]^{a \cd (-)}
\ar[r]
&
0
}
\end{equation}

Define the graded $A$-modules $\bar{M}$ and $\bar{L}$ by these exact 
sequences:
\[ 0 \to \bar{M} \to M(-i) \xar{a \cd (-)} M \]
and 
\begin{equation} \label{eqn:4241}
L(-i) \xar{a \cd (-)} L \to \bar{L} \to 0 .
\end{equation}
The long exact cohomology sequence of (\ref{4590}) gives 
rise to this exact sequence in $\dcat{M}(A, \mrm{gr})$~:
\begin{equation} \label{eqn:4240}
\bar{M} \xar{\phi} \bar{L} \xar{\psi} B . 
\end{equation}
Because $M$ is a noetherian object of $\dcat{M}(A, \mrm{gr})$,
so is its degree twist $M(-i)$, and also the subobject $\bar{M} \sub M(-i)$. 
On the other hand, we are given that the ring  $B$ 
is noetherian; so it is a noetherian object of $\dcat{M}(A, \mrm{gr})$.
The exact sequence (\ref{eqn:4240}) shows that $\bar{L}$ is also a 
noetherian object of $\dcat{M}(A, \mrm{gr})$.
We see that $\bar{L}$ is a finitely generated $A$-module. 

Finally, let us choose finitely many homogeneous generators of the graded 
$A$-module $\bar{L}$, and then lift them to homogeneous elements of $L$. 
Let $L' \sub L$ be the $A$-submodule generated by this finite collection of 
elements. The exact sequence (\ref{eqn:4241}) shows that 
$L = (a \cd L) + L'$.
The graded $A$-module $L$ is bounded below, because $L \sub A$ and $A$ is 
nonnegative. The element $a$ belongs to the augmentation ideal $\m$ of $A$. 
Therefore, by Proposition \ref{prop:4540} (the graded Nakayama Lemma),
we have $L' = L$. This contradicts the assumption that $L$ is not finitely 
generated. 
\end{proof}




\mysubsection{Properties of Algebraically Graded Modules}
\label{subsec:props-ag-mods}

In this subsection we continue with Convention \ref{conv:3700}, to which we add:

\begin{conv}  \label{conv:4560}
The base ring $\K$ is a field. 
\end{conv}

See Remark \ref{rem:3740} regarding the requirement that $\K$ be a field. 
In a sense the category $\dcat{M}(\K, \mrm{gr})$ is boring under Convention
\ref{conv:4560}~:

\begin{prop} \label{prop:3851}
For an object $M \in \dcat{M}(\K, \mrm{gr})$ there is an isomorphism
$M \cong \bigoplus_{x \in X} \K(-i_x)$,
where $X$ is an indexing set and $\{ i_x \}_{x \in X}$ is a collection of 
integers. 
\end{prop}

\begin{proof}
This is a variation on the usual proof for $\K$-modules: by Zorn's Lemma there 
is a maximal linearly independent collection 
$\{ m_x \}_{x \in X}$ of homogeneous elements of $M$. 
This is a basis of $M$. The number $i_x$ is the degree of $m_x$.  
\end{proof}

\begin{dfn} \label{dfn:3795}
Let $A$ be a graded ring. 
We define the graded bimodule 
\[ A^* := \opn{Hom}_{\K}(A, \K) \in 
\dcat{M}(A^{\mrm{en}}, \mrm{gr}) , \]
see formula (\ref{eqn:3701}). 
\end{dfn}

\begin{dfn}  \label{dfn:3740}
Let $A$ be a graded ring.  
\begin{enumerate}
\item A graded $A$-module $P$ is called a {\em graded-free $A$-module}% 
\index{Algebraically graded module! graded-free}
if
$P \cong \lb \bigoplus_{x \in X} A(-i_x)$
in $\dcat{M}(A, \mrm{gr})$
for some collection $\{ i_x \}_{x \in X}$ of integers.

\item A graded $A$-module $P$ is called a 
{\em graded-projective $A$-module}% 
\index{Algebraically graded module! graded-projective}
if it is a projective object in the abelian category 
$\dcat{M}(A, \mrm{gr})$.

\item  A graded $A$-module $I$ is called a 
{\em graded-cofree $A$-module}%
\index{Algebraically graded module! graded-cofree}
if
$I \cong \lb \prod_{x \in X} A^*(-p_x)$
in $\dcat{M}(A, \mrm{gr})$ 
for some collection $\{ p_x \}_{x \in X}$ of integers.

\item A graded $A$-module $I$ is called a 
{\em graded-injective $A$-module}% 
\index{Algebraically graded module! graded-injective}
if it is an injective object in the abelian category $\dcat{M}(A, \mrm{gr})$.

\item A graded $A$-module $P$ is called a {\em graded-flat $A$-module} 
\index{Algebraically graded module! graded-flat}
if the functor 
\[ (-) \ot_A P : \dcat{M}(A^{\mrm{op}}, \mrm{gr}) \to
\dcat{M}(\K, \mrm{gr}) \]
is exact. 
\end{enumerate}
\end{dfn}

Note that a graded-free $A$-module $P$ is of the form 
$P \cong A \ot_{} V$ for some $V \in \dcat{M}(\K, \mrm{gr})$;
and a graded-cofree $A$-module $I$ is of the form 
$I \cong \opn{Hom}_{\K}(A \ot W, \K)$ for some $W \in \dcat{M}(\K, \mrm{gr})$.

\begin{prop} \label{prop:3795}
Let $A$ and $B$ be graded rings.
\begin{enumerate}
\item There is an isomorphism 
\[ \opn{Hom}_{A}(-, A^*) \cong \opn{Hom}_{\K}(-, \K)  \]
of functors 
\[ \dcat{M}(A \ot B^{\mrm{op}}, \mrm{gr})^{\mrm{op}} \to 
\dcat{M}(B \ot A^{\mrm{op}}, \mrm{gr}) . \]

\item The graded bimodule $A^*$ is graded-injective over $A$, 
namely it is an injective object in the category
$\dcat{M}(A, \mrm{gr})$. 
\end{enumerate}
\end{prop}

\begin{proof}
This is by the adjunction isomorphism 
\[ \opn{Hom}_{A} \bigl( M, \opn{Hom}_{\K}(A, \K) \bigr) \cong
\opn{Hom}_{\K}(M, \K) \]
in $\dcat{M}(\K, \mrm{gr})$.
\end{proof}

The ungrading functor (\ref{eqn:3715}) preserves finiteness (i.e.\ 
being a finitely generated module). It also preserves freeness and 
projectivity, as the next proposition shows. 

\begin{prop} \label{prop:3715} 
Let $A$ be a graded ring, and let $P \in \dcat{M}(A, \mrm{gr})$. 
\begin{enumerate}
\item If $P$  is a graded-free $A$-module, then $\opn{Ungr}(P)$ is a 
free $\opn{Ungr}(A)$-module.

\item $P$ is a graded-projective $A$-module if and only if it is a direct 
summand, in $\dcat{M}(A, \mrm{gr})$, of a graded-free $A$-module.

\item If $P$ is a graded-projective $A$-module, then $\opn{Ungr}(P)$ is a 
projective $\opn{Ungr}(A)$-module.
\end{enumerate} 
\end{prop}

\begin{exer} \label{exer:3715} 
Prove Proposition \ref{prop:3715}.
\end{exer}

\begin{rem} \label{rem:4805}
In the situation of Proposition \ref{prop:3715}, one can show that $P$ is a 
graded-flat $A$-module if and only if  $\opn{Ungr}(P)$ is a flat 
$\opn{Ungr}(A)$-module. The proof is not easy. See 
\cite[Proposition 3.4.5]{NV}.
\end{rem}

For injectives things are much more complicated, as the next theorem and 
example show.

\begin{thm}[Van den Bergh] \label{thm:3715} 
Let $A$ be a left noetherian nonnegative graded ring, and let 
$I$ be graded-injective $A$-module. Then $\opn{Ungr}(I)$ has injective 
dimension at most $1$ as an $\opn{Ungr}(A)$-module.
\end{thm}

\begin{proof}
See \cite{Ye17}. A more general statement can be found in \cite{SoZa}.
\end{proof}

\begin{exa} \label{exa:3800}
Consider the graded ring $A := \K[t]$, the ring of polynomials in a variable 
$t$ of degree $1$. The graded $A$-module $I := \K[t, t^{-1}]$ is 
graded-injective, but $\opn{Ungr}(I)$ has injective dimension $1$ over 
$\opn{Ungr}(A)$. 
\end{exa}

\begin{prop} \label{prop:4905}  %{lem:4045} 
\mbox{}
\begin{enumerate}
\item A graded-free $A$-module $P$ is graded-projective. 

\item A graded-projective $A$-module $P$ is graded-flat.

\item A graded-cofree $A$-module $I$ is graded-injective. 
\end{enumerate}
\end{prop}

\begin{exer} \label{exer:4045}
Prove this proposition. 
\end{exer}

\begin{prop} \label{prop:3900}
Let $A$ and $B$ be graded rings.
\begin{enumerate}
\item If $P \in \dcat{M}(A \ot B, \mrm{gr})$ is projective, then it is 
projective in $\dcat{M}(A, \mrm{gr})$.

\item If $P \in \dcat{M}(A \ot B, \mrm{gr})$ is flat, then it is 
flat in $\dcat{M}(A, \mrm{gr})$.

\item If $I \in \dcat{M}(A \ot B, \mrm{gr})$ is injective, then it is 
injective in $\dcat{M}(A, \mrm{gr})$.
\end{enumerate}
\end{prop}

\begin{proof}
(1) Take $N \in \dcat{C}(A, \mrm{gr})$. 
We have isomorphisms
\[ \begin{aligned}
& \opn{Hom}_{A}(P, N) \cong
\opn{Hom}_{A \ot B}(P, \opn{Hom}_A(A \ot B, N)) 
\\ & \quad  \cong \opn{Hom}_{A \ot B}(P, \opn{Hom}_{\K}(B, N)) 
\end{aligned} \]
in $\dcat{C}_{\mrm{str}}(\K, \mrm{gr})$. 
As an object of $\dcat{M}(\K, \mrm{gr})$, $B$ is projective.
Hence, if $N$ is an acyclic complex, then so is $\opn{Hom}_{A}(P, N)$.

\medskip \noindent
(2, 3) Like item (1), but using the fact that $B$ is graded-flat over $\K$, and 
using suitable adjunction isomorphisms.
\end{proof}

\begin{dfn} \label{dfn:4260}
Let $A$ be a graded ring.
\begin{enumerate}
\item Let $M$ be a graded $A$-module, and let $N \sub M$ 
be a graded submodule. We say that $N$ is a 
{\em graded-essential submodule}% 
\index{Algebraically graded module! graded-essential submodule of}
of $M$ if for every nonzero graded submodule $M' \sub M$ the intersection 
$N \cap M'$ is nonzero. 

\item A homomorphism $\phi : N \to M$ in $\dcat{M}(A, \mrm{gr})$
is called a {\em graded-essential monomorphism}% 
\index{Algebraically graded module! graded-essential monomorphism of {\indash}s}
if $\phi$ is an injective 
homomorphism, and $\phi(N)$ is a graded-essential submodule of $M$.
\end{enumerate}
\end{dfn}

Item (2) of the definition is a special case of the usual definition of an 
essential monomorphism in an abelian category. 

\begin{exer} \label{exer}
Let $N \sub M$ be graded $A$-modules. Show that $N$ is a graded-essential 
submodule of $M$ if and only if $\opn{Ungr}(N)$ is an essential 
$\opn{Ungr}(A)$-submodule of $\opn{Ungr}(M)$.
(Hint: prove that for every nonzero element $m \in M$ there is a homogeneous 
element $a \in A$ such that $a \cd m \in N$ and $a \cd m \neq 0$. Do this by 
induction on the number $r$ appearing in the decomposition (\ref{eqn:4260}).
See the proof of \cite[Lemma 3.3.13]{NV}, which is not totally correct, but 
can be fixed using our hint.)
\end{exer}

\begin{prop} \label{prop:3735}
Let $A$ be a graded ring. 
\begin{enumerate}
\item The category $\dcat{M}(A, \mrm{gr})$ has enough projectives.

\item The category $\dcat{M}(A, \mrm{gr})$ has enough injectives. Moreover, 
every $M \in \lb \dcat{M}(A, \mrm{gr})$ has an injective hull, namely there is 
a graded-essential mono\-morphism $M \inj I$ to a graded-injective module $I$. 
\end{enumerate}
\end{prop}

\begin{proof}
(1) Given a module $M \in \dcat{M}(A, \mrm{gr})$ and a homogeneous element 
$m \in M_i$, there is a homomorphism 
$A(-i) \to M$ that sends $1_A \mapsto m$. Therefore, by taking a sufficiently 
large direct sum of degree twists of $A$, there is a surjection 
$P \surj M$ from a graded-free module $P$. And graded-free modules are 
graded-projective, by Proposition \ref{prop:4905}. 

\medskip \noindent
(2) For every nonzero homogeneous element $m \in M_q$ there is a homomorphism 
$M \to \K(-q)$ in $\dcat{M}(\K, \mrm{gr})$
that is nonzero on $m$. By Proposition \ref{prop:3795}(1)
we get a homomorphism $M \to A^*(-q)$ in $\dcat{M}(A, \mrm{gr})$
that is nonzero on $m$. Thus, by taking a sufficiently large product 
$I :=  \prod_{x \in X} A^*(-q_x)$
in $\dcat{M}(A, \mrm{gr})$ we obtain a monomorphism $M \inj I$. And according 
Proposition \ref{prop:4905},  the graded-cofree object $I$ is injective in 
$\dcat{M}(A, \mrm{gr})$.

As for injective hulls: this is the same as in the ungraded case; see 
\cite[Theorem 3.30]{Rot} or \cite[Section 3.D]{Lam}.
\end{proof}

Now a structural result on injective objects in $\dcat{M}(A, \mrm{gr})$. 

\begin{thm} \label{thm:3795}
Let $A$ be a left noetherian graded ring. 
\begin{enumerate}
\item If $\{ I_x \}_{x \in X}$ is a collection of graded-injective $A$-modules, 
then $I := \bigoplus_{x \in X} I_x$ is a  graded-injective $A$-module. 

\item Every graded-injective $A$-module $I$ is a direct sum of indecomposable
graded-injective $A$-modules. 
\end{enumerate} 
\end{thm}

\begin{proof}
This is because $\dcat{M}(A, \mrm{gr})$ is a locally noetherian Grothendieck 
abelian category. See \cite[Section V.4]{Ste}.
Item (1) can be proved directly using the graded version of the Baer criterion 
(Theorem \ref{thm:151}). 
\end{proof}

\begin{dfn}  \label{dfn:3703}
A {\em connected graded $\K$-ring} 
\index{Algebraically graded ring! connected}
is a nonnegative graded $\K$-ring 
$A$ (see Conventions \ref{conv:3700} and 
\ref{conv:4560},  and Definition \ref{dfn:4540}), such that  the ring 
homomorphism $\K \to A_0$ is bijective, 
and $A_i$ is a finitely generated $\K$-module for all $i \geq 0$. 
\end{dfn}

Here are a few examples of connected graded rings. 

\begin{exa} \label{exa:3755}
Consider the noncommutative polynomial ring 
$A := \lb \K \bra{x_1, \ldots, x_n}$ 
in $n \geq 1$ variables, all of algebraic degree $1$. This is a connected 
graded $\K$-ring. If $n \geq 2$ then $A$ is not commutative and not noetherian 
(on either side). If $n = 1$ then 
$\K \bra{x_1} = \K[x_1]$, see next example. 
\end{exa}

\begin{exa} \label{exa:3757} 
Let
$A := \K [x_1, \ldots, x_n]$,
the commutative polynomial ring 
in $n \geq 1$ variables, all of algebraic degree $1$. This is a connected 
graded $\K$-ring, commutative and noetherian.  
\end{exa}

\begin{exa} \label{exa:3756} 
Let $\K[t]$ be the commutative polynomial ring in a variable $t$ of algebraic 
degree $1$. Next let 
\[ A := \K[t] \bra{x, y} / (y \cd x - x \cd y - t^2) , \]
where $x$ and $y$ are noncommuting variables of algebraic 
degree $1$ (that commute with $t$). This is the {\em homogeneous first Weyl 
algebra}. It is a connected graded $\K$-ring, noetherian, but not commutative. 
Indeed, if $\opn{char}(\K) = 0$, then the center of $A$ is $\K[t]$. 
\end{exa}

\begin{exa} \label{exa:4075} 
Let $\g$ be a finite Lie algebra over $\K$. Choose a basis 
$v_1, \ldots, v_n$ for $\g$ as a $\K$-module, and let 
$\ga_{i, j, k} \in \K$ be the constants describing the Lie bracket, i.e.\
\[ [v_i, v_j] = \sum_k \ga_{i, j, k} \cd v_k \in \g . \]
Let $\K[t]$ be the commutative polynomial ring in a variable $t$ of algebraic 
degree $1$. Let $x_1, \ldots, x_n$ be variables of algebraic degree $1$, and 
define the graded ring 
\[ A := \K[t] \bra{x_1, \ldots, x_n} \, / \, (I) , \]
where $I$ is the two-sided ideal generated by the quadratic elements 
\[ x_i \cd x_j - x_j \cd x_i  - \sum_{k = 1}^n \ga_{i, j, k} \cd x_k \cd t . \]
This is the {\em homogeneous universal enveloping ring} of $\g$. 
The quotient $A / (t - 1)$ is the  universal enveloping ring
$\opn{U}(\g)$. The ring $A$ is connected graded, noetherian, but often 
noncommutative (since the Lie algebra $\g$ embeds into the quotient ring 
$\opn{U}(\g)$ with its commutator Lie bracket). 
\end{exa}

\begin{exa} \label{exa:3775}
Suppose $A$ and $B$ are both connected graded $\K$-rings. Then 
$A \ot B$ is a connected graded $\K$-ring.
\end{exa}

\begin{dfn} \label{dfn:3724}
A graded $\K$-module $M = \bigoplus_{i \in \Z} \, M_i$ 
is called {\em degreewise finite} 
\index{Algebraically graded module! degreewise-finite}
if the $\K$-module $M_i$ is finitely generated for every $i \in \Z$. 

The full subcategory of $\dcat{M}(\K, \mrm{gr})$ on the 
degreewise finite modules is denoted by 
$\dcat{M}_{\mrm{dwf}}(\K, \mrm{gr})$. 
\end{dfn}

It is obvious that the category $\dcat{M}_{\mrm{dwf}}(\K, \mrm{gr})$ is a thick 
abelian subcategory of $\dcat{M}(\K, \mrm{gr})$, 
closed under subobjects and quotient objects.

\begin{dfn} \label{dfn:3725}
For $M \in \dcat{M}(\K, \mrm{gr})$ we let
\[ M^* := \opn{Hom}_{\K}(M, \K) \in \dcat{M}(\K, \mrm{gr}) . \]
This is the {\em graded $\K$-linear dual of $M$}. 
\index{Algebraically graded module! graded $\K$-linear dual}
\index{1-Mstar@$M^*$}
\end{dfn}

According to formula (\ref{eqn:3701}) we have
$M^* = \bigoplus_{i \in \Z} (M^*)_i$, where 
\[ (M^*)_i = \opn{Hom}_{\K}(M, \K)_i = \opn{Hom}_{\K}(M_{-i}, \K) = 
(M_{-i})^* . \]

\begin{lem} \label{lem:3725}
If $M \in \dcat{M}_{\mrm{dwf}}(\K, \mrm{gr})$ then 
$M^* \in \dcat{M}_{\mrm{dwf}}(\K, \mrm{gr})$, 
and the Hom-evaluation homomorphism
$\opn{ev}_M : M \to M^{**} = (M^*)^*$ 
in $\dcat{M}_{\mrm{dwf}}(\K, \mrm{gr})$ is bijective. 
Thus $(-)^*$ is a duality of the category 
$\dcat{M}_{\mrm{dwf}}(\K, \mrm{gr})$.
\end{lem}

The easy proof is omitted. 

\begin{dfn} \label{dfn:3723} 
Let $A$ be a graded ring. 
A module $M \in \dcat{M}(A, \mrm{gr})$ is called {\em degreewise $\K$-finite} 
\index{Algebraically graded module! degreewise-finite}
if it is degreewise finite as a graded $\K$-module. 

The full subcategory of $\dcat{M}(A, \mrm{gr})$ on the degreewise $\K$-finite 
modules is denoted by 
$\dcat{M}_{\mrm{dwf}}(A, \mrm{gr})$.
\index{1-Mdwf(A,gr)@$\dcat{M}_{\mrm{dwf}}(A, \mrm{gr})$}
\end{dfn}

Of course $\dcat{M}_{\mrm{dwf}}(A, \mrm{gr})$ is a thick abelian 
subcategory of $\dcat{M}(A, \mrm{gr})$,
closed under subobjects and quotients.

\begin{prop} \label{prop:4045} 
Let $A$ be a graded ring.  
If $M \in \dcat{M}_{\mrm{dwf}}(A, \mrm{gr})$ then 
$M^* \in \dcat{M}_{\mrm{dwf}}(A^{\mrm{op}}, \mrm{gr})$, 
and the Hom-evaluation homomorphism 
$\opn{ev}_M : M \to M^{**}$
in \lb $\dcat{M}(A, \mrm{gr})$ is bijective. Thus 
\[ (-)^* : \dcat{M}_{\mrm{dwf}}(A, \mrm{gr})^{\mrm{op}} \to 
\dcat{M}_{\mrm{dwf}}(A^{\mrm{op}}, \mrm{gr}) \]
is an equivalence of categories.
\end{prop}

\begin{exer} \label{exer:4075}
Prove this proposition. (Hint: use Lemma \ref{lem:3725}.)
\end{exer}

Noncommutative connected graded rings are very similar to complete commutative 
local rings. Here is the second variant of the graded Nakayama Lemma (the first 
was Proposition \ref{prop:4540}). 

\begin{prop} \label{prop:4590}    
Let $A$ be a connected graded ring, and let 
$M \in \dcat{M}(A, \mrm{gr})$ be a bounded below graded module. Define  
$V := \K \ot_A M \in \dcat{M}(\K, \mrm{gr})$.
\begin{enumerate}
\item There is a surjection 
$\pi : A \ot_{} V \surj M$
in $\dcat{M}(A, \mrm{gr})$, such that the diagram 
\[ \UseTips \xymatrix @C=8ex @R=6ex {
A \ot V
\ar[r]^(0.55){\pi}
\ar[d]_{\be}
&
M
\ar[d]^{\al}
\\
V
\ar[r]^{\opn{id}}
&
V
} \]
in which the morphisms $\al$ and $\be$ are the canonical projections, is 
commutative.  

\item If $V = 0$ then $M = 0$. 

\item If $M$ is a graded-projective $A$-module, then the surjection 
$\pi$ above is bijective. Hence $M$ is a graded-free $A$-module. 
\end{enumerate}
\end{prop}

Note that the projection $\be$ is induced by the augmentation homomorphism 
$A \to \K$, so $\opn{Ker}(\be) = \m \ot V$. Likewise, 
$\opn{Ker}(\al) = \m \cd M \sub M$.

\begin{proof}  \mbox{}

\smallskip \noindent
(1) Say $M = \bigoplus_{i \geq i_0} \, M_i$. 
Choose a section $\pi_0 : V \inj M$ in $ \dcat{M}(\K, \mrm{gr})$ of the 
canonical projection $\al : M \surj V$. 
Since $A \ot V$ is a graded-free $A$-module, $\pi_0$ extends to a homomorphism 
$\pi : A \ot_{} V \to M$ in $\dcat{M}(A, \mrm{gr})$.
See the commutative diagram below, in which $\ga(v) := 1_A \ot v$. 
\[ \UseTips \xymatrix @C=8ex @R=6ex {
V
\ar[dr]^{\pi_0}
\ar[d]_{\ga}
\\
A \ot V
\ar[r]^{\pi}
&
M
} \]
By Proposition \ref{prop:4540} the homomorphism $\pi$ is surjective. 

\medskip \noindent 
(2) Clear from the existence of the surjection $\pi$.

\medskip \noindent 
(3) Let $N := \opn{Ker}(\pi) \sub A \ot V$, so there is a short exact sequence 
\[ 0 \to N \to A \ot V \xar{\pi} M \to  0 \]
in $\dcat{M}(A, \mrm{gr})$, and $N$ is a bounded below graded $A$-module. 
If $M$ is graded-projective, then this sequence is split. 
It follows that $\K \ot_{A} N = 0$. By part (2) we see that $N = 0$. 
\end{proof}

\begin{dfn} \label{dfn:4265}
Let $A$ be a connected graded $\K$-ring, and let $M \in \dcat{M}(A, \mrm{gr})$. 
The {\em socle}% 
\index{Algebraically graded module! socle of}%
\index{1-Soc(M)@$\opn{Soc}(M)$}
of $M$ is the graded $A$-submodule
$\opn{Soc}(M) := \opn{Hom}_A(\K, M) \sub M$.
\end{dfn}

In other words, 
$\opn{Soc}(M) = \{ m \in M \mid \m \cd m = 0 \}$.
The $A$-module structure of $\opn{Soc}(M)$ is through the augmentation 
homomorphism $A \to \K$; so Proposition \ref{prop:3851} applies to it. 

Here is a dual to the graded Nakayama Lemma. 

\begin{prop} \label{prop:4595}  
Let $A$ be a connected graded $\K$-ring, let 
$N \in \dcat{M}(A, \mrm{gr})$ be a bounded above graded module,
and let $W := \opn{Soc}(N)$. 
\begin{enumerate}
\item There is an injection
$\si : N \inj \opn{Hom}_{\K}(A, W)$
in $\dcat{M}(A, \mrm{gr})$,  such that the diagram 
\[ \UseTips \xymatrix @C=8ex @R=6ex {
W
\ar[r]^{\opn{id}}
\ar[d]_{\al}
&
W
\ar[d]^{\be}
\\
N
\ar[r]^(0.35){\si}
&
\opn{Hom}_{\K}(A, W)
} \]
in which $\al$ and $\be$ are the canonical embeddings, is commutative. 

\item If $W = 0$ then $N = 0$.

\item If $N$ is a graded-injective $A$-module, then the homomorphism $\si$ 
above is bijective. Hence $N$ is a graded-cofree $A$-module.
\end{enumerate}
\end{prop}

Note that $\be$ is induced from the augmentation homomorphism 
$A \to \K$, and 
$\opn{Im}(\be) = \opn{Soc} \bigl( \opn{Hom}_{\K}(A, W) \bigr)$.

\begin{proof} \mbox{} 

\smallskip \noindent 
(1) Say $N = \bigoplus_{i \leq i_1} \, N_i$ for some $i_1 \in \Z$.
Choose a splitting 
$\si_0 : N \surj W$ in $\dcat{M}(\K, \mrm{gr})$ of the inclusion 
$\al : W \inj N$. 
By adjunction there is an isomorphism 
\[ \opn{Hom}_{\K}(N, W) \cong 
\opn{Hom}_{A}(N,  \opn{Hom}_{\K}(A, W)) \]
in $\dcat{M}(\K, \mrm{gr})$, so there is a homomorphism 
$\si : N \to \opn{Hom}_{\K}(A, W)$ in $\dcat{M}(A, \mrm{gr})$ that lifts 
$\si_0$. See the commutative diagram below, where 
$\ga(\chi) := \chi(1_A)$.  
\[ \UseTips \xymatrix @C=8ex @R=6ex {
N
\ar[dr]_{\si_0}
\ar[r]^(0.35){\si}
&
\opn{Hom}_{\K}(A, W)
\ar[d]^{\ga}
\\
&
W
} \]
A descending induction on degree, starting from $i_1$, shows that 
$\si$ is an injection. 

\medskip \noindent 
(2) Clear from the existence of $\si$. 

\medskip \noindent 
(3) Let 
$M := \opn{Coker}(\si)$. 
If $N$ is graded-injective, then the exact sequence 
\[ 0 \to N \xar{\si} \opn{Hom}_{\K}(A, W) \to M \to 0 \]
in $\dcat{M}(A, \mrm{gr})$ is split. 
Applying the additive functor $\opn{Soc}$ to it, we see that 
$\opn{Soc}(M) = 0$. Since $M$ is a bounded above graded module, 
item (2) says that $M = 0$. 
\end{proof}

\begin{dfn} \label{dfn:3726} 
Let $A$ be a graded ring. 
A module $N \in \dcat{M}(A^{\mrm{op}}, \mrm{gr})$ is called a 
{\em cofinite graded $A^{\mrm{op}}$-module} 
\index{Algebraically graded module! cofinite}
if $N \cong M^*$ for some finite graded 
$A$-module $M$.

The full subcategory of $\dcat{M}(A^{\mrm{op}}, \mrm{gr})$ on the cofinite 
graded $A^{\mrm{op}}$-modules is denoted by 
\index{1-Mcof(A,gr)@$\dcat{M}_{\mrm{cof}}(A, \mrm{gr})$}
$\dcat{M}_{\mrm{cof}}(A^{\mrm{op}}, \mrm{gr})$.
\end{dfn}

Recall that an object $M$ in an abelian category $\cat{M}$ is called {\em 
noetherian} (resp.\ {\em artinian}) if it satisfies the ascending (resp.\ 
descending) chain condition on subobjects. 
Let us denote the corresponding full subcategories of $\cat{M}$ by
$\cat{M}_{\mrm{not}}$ and $\cat{M}_{\mrm{art}}$ respectively. 
Then $\cat{M}_{\mrm{not}}, \cat{M}_{\mrm{art}} \sub \cat{M}$
are thick abelian subcategories, closed under subobjects and 
quotient objects. 

Here is a variant of the Matlis Duality Theorem (see Remark \ref{rem:2245}):

\begin{thm}[NC Graded Matlis Duality] \label{thm:4045} 
\index{Matlis Duality Theorem! noncommutative graded} 
Assume $A$ is a left noetherian connected graded ring. 
\begin{enumerate}
\item The category $\dcat{M}_{\mrm{f}}(A, \mrm{gr})$ is a thick abelian 
subcategory of $\dcat{M}_{\mrm{dwf}}(A, \mrm{gr})$, 
closed under subobjects and quotient objects. 

\item The category $\dcat{M}_{\mrm{cof}}(A^{\mrm{op}}, \mrm{gr})$
is a thick abelian subcategory of 
$\dcat{M}_{\mrm{dwf}}(A^{\mrm{op}}, \mrm{gr})$, 
closed under subobjects and quotient objects. 

\item The functor 
\[ (-)^* : \dcat{M}_{\mrm{f}}(A, \mrm{gr})^{\mrm{op}} \to 
\dcat{M}_{\mrm{cof}}(A^{\mrm{op}}, \mrm{gr}) \]
is an equivalence. 

\item The modules in $\dcat{M}_{\mrm{f}}(A, \mrm{gr})$ are the noetherian 
objects in the abelian category $\dcat{M}(A, \mrm{gr})$.

\item The modules in $\dcat{M}_{\mrm{cof}}(A^{\mrm{op}}, \mrm{gr})$ are the 
artinian objects in the abelian category $\dcat{M}(A^{\mrm{op}}, \mrm{gr})$.
\end{enumerate}
\end{thm}

\begin{proof} \mbox{}

\smallskip \noindent
(1) Since $A$ is connected, as a left module it belongs to 
$\dcat{M}_{\mrm{dwf}}(A, \mrm{gr})$. Therefore every finite graded $A$-module 
belongs to $\dcat{M}_{\mrm{dwf}}(A, \mrm{gr})$, i.e.\
$\dcat{M}_{\mrm{f}}(A, \mrm{gr}) \sub \dcat{M}_{\mrm{dwf}}(A, \mrm{gr})$.

Because $A$ is left noetherian, it follows that 
$\dcat{M}_{\mrm{f}}(A, \mrm{gr})$
is a thick abelian subcategory of
$\dcat{M}(A, \mrm{gr})$, 
closed under subobjects and quotient objects.

\medskip \noindent 
(2, 3) By definition 
$\dcat{M}_{\mrm{cof}}(A^{\mrm{op}}, \mrm{gr})$
is the essential image under the functor $(-)^*$ of the category
$\dcat{M}_{\mrm{f}}(A, \mrm{gr})$.
Now use the equivalence of Proposition \ref{prop:4045} and part (1). 

\medskip \noindent 
(4) This is proved the same way as for ungraded left noetherian rings. 

\medskip \noindent 
(5) We know by (2, 3, 4) that the objects of 
$\dcat{M}_{\mrm{cof}}(A^{\mrm{op}}, \mrm{gr})$
are artinian objects in 
$\dcat{M}(A^{\mrm{op}}, \mrm{gr})$.

For the opposite direction, let 
$N \in \dcat{M}(A^{\mrm{op}}, \mrm{gr})$ be an artinian object. 
The descending chain condition forces $N$ to be bounded above, namely 
$N_i = 0$ for $i \gg 0$. 
Consider the socle
$W := \opn{Hom}_{A^{\mrm{op}}}(\K, N)$. 
This is a graded $A^{\mrm{op}}$-submodule of $N$, isomorphic to 
$\bigoplus_{x \in X}  \K(-i_x)$ for some indexing set $X$ and a collection of 
integers $\{ i_x \}_{x \in X}$. By the descending chain condition the set $X$ 
has to be finite. Hence 
\[ \bigoplus_{x \in X}  A^*(-i_x) \cong A^* \ot W \cong 
\opn{Hom}_{\K}(A, W) . \]
Proposition \ref{prop:4595} tells us that there is an injective 
homomorphism 
$N \inj \lb \bigoplus_{x \in X} A^*(-i_x)$.
We know that $A^*(-i_x) \in \dcat{M}_{\mrm{cof}}(A, \mrm{gr})$; 
and hence so is $N$. 
\end{proof}

\begin{rem} \label{rem:3903}
If $A$ is connected graded but not left noetherian, then there are cofinite 
graded $A^{\mrm{op}}$-modules that are not artinian. Take the ring 
$A := \K \bra{x_1, x_2}$ from Example \ref{exa:3755}. Then 
$A^* \in \dcat{M}(A^{\mrm{op}}, \mrm{gr})$
is cofinite but not artinian. 
\end{rem}



\mysubsection{Resolutions and Derived Functors} 
\label{gr-res-der-fun}

In this subsection we continue with Conventions \ref{conv:3700} and 
\ref{conv:4560}. Thus $\K$ is a base field and $\ot = \ot_{\K}$. 
In the next definitions we collect the algebraically graded versions of some 
definitions from Sections \ref{sec:resol} and \ref{sec:exist-resol}.

\begin{dfn} \label{dfn:4275}   
Let $A$ be a graded ring, and let $P \in \dcat{C}(A, \mrm{gr})$. 
\begin{enumerate}
\item  The complex $P$ is called a
{\em graded-free complex}% 
\index{Complex of algebraically graded modules! graded-free}
if all the graded $A$-modules $P^i$ are 
graded-free, and the differential is zero. 

\item The complex $P$ is called a
{\em semi-graded-free complex}% 
\index{Complex of algebraically graded modules! semi-graded-free}
if it admits a filtration 
$\{ F_j(P) \}_{j \geq -1}$ in $\dcat{C}_{\mrm{str}}(A, \mrm{gr})$, 
such that $F_{-1}(P) = 0$, 
$P = \bigcup_j F_j(P)$, and 
$\opn{Gr}^F_j(P)$ is a graded-free complex for every $j$. 

\item The complex $P$ is called a 
{\em K-graded-projective complex}% 
\index{Complex of algebraically graded modules! K-graded-projective}
if for every acyclic complex $N \in \dcat{C}(A, \mrm{gr})$, the complex 
$\opn{Hom}_A(P, N)$ is acyclic.
\end{enumerate}
\end{dfn}

\begin{dfn} \label{dfn:4276}   
Let $A$ be a graded ring, and let $I \in \dcat{C}(A, \mrm{gr})$. 
\begin{enumerate}
\item  The complex $I$ is called a 
{\em graded-cofree complex}% 
\index{Complex of algebraically graded modules! graded-cofree}
if all the $A$-modules $I^p$ are graded-cofree, and the differential is zero.

\item The complex $I$ is called a 
{\em semi-graded-cofree complex}% 
\index{Complex of algebraically graded modules! semi-graded-cofree}
if it admits a cofiltration $\{ F_q(I) \}_{q \geq -1}$ in 
$\dcat{C}_{\mrm{str}}(A, \mrm{gr})$, such that $F_{-1}(I) = 0$, each 
$\opn{Gr}^F_q(I)$ is graded-cofree complex, and 
$I = \opn{lim}_{\lar q} F_q(I)$.

\item The complex $I$ is called a 
{\em K-graded-injective complex}% 
\index{Complex of algebraically graded modules! K-graded-injective}
if for every acyclic complex $N \in \dcat{C}(A, \mrm{gr})$, the complex
$\opn{Hom}_A(N, I)$ is acyclic. 
\end{enumerate}
\end{dfn}

\begin{dfn} \label{dfn:4277} 
A complex $P \in \dcat{C}(A, \mrm{gr})$ is called a 
{\em K-graded-flat complex}% 
\index{Complex of algebraically graded modules! K-graded-flat}
if for every acyclic complex 
$N \in \dcat{C}(A^{\mrm{op}}, \mrm{gr})$ the complex 
$N \ot_A P$ is acyclic. 
\end{dfn}

\begin{thm} \label{thm:3740}
Let $A$ be a graded ring, and let $M \in \dcat{C}(A, \mrm{gr})$.
\begin{enumerate}
\item There exists a quasi-isomorphism 
$\rho : P \to M$ in $\dcat{C}_{\mrm{str}}(A, \mrm{gr})$,
where $P$ is a semi-graded-free complex, 
and $\opn{sup}(P) = \opn{sup}(\opn{H}(M))$. 

\item  There exists a quasi-isomorphism 
$\rho : M \to I$ in $\dcat{C}_{\mrm{str}}(A, \mrm{gr})$,
where $I$ is a semi-graded-cofree complex, and
$\opn{inf}(I) = \opn{inf}(\opn{H}(M))$. 
\end{enumerate}
\end{thm}

\begin{proof} \mbox{}

\smallskip \noindent 
(1) This is a modification of Theorem \ref{thm:3305} and Corollary 
\ref{cor:3340}. The proof is the same, after a few obvious modifications. 

\medskip \noindent 
(2) This is a modification of Theorem \ref{thm:3325} and Corollary 
\ref{cor:3345}. The proof is the same, after a few obvious modifications. 
The injective object $\K^*$ we use here is $\K^* := \K$ of course. 
\end{proof}

The various kinds of complexes mentioned in the definitions above are related 
as follows. 

\begin{thm} \label{thm:4275}  
Let $A$ be a graded ring.
\begin{enumerate}
\item A complex $P \in \dcat{C}(A, \mrm{gr})$ is K-graded-projective if and 
only if it is 
isomorphic in $\dcat{K}(A, \mrm{gr})$ to a semi-graded-free complex $Q$.

\item If $P \in \dcat{C}(A, \mrm{gr})$ is K-graded-projective, then it is 
K-graded-flat. 

\item If $P \in \dcat{C}(A, \mrm{gr})$ is K-graded-projective, then 
$P^* \in \dcat{C}(A^{\mrm{op}}, \mrm{gr})$
is K-graded-injective. 

\item A complex $I \in \dcat{C}(A, \mrm{gr})$ is K-graded-injective iff it is 
isomorphic in $\dcat{K}(A, \mrm{gr})$ to a semi-graded-cofree complex $J$.
\end{enumerate}
\end{thm}

\begin{proof} \mbox{}

\smallskip \noindent
(1) By Theorem \ref{thm:1575}, that applies to this context too, a 
semi-graded-free complex $Q$ is K-graded-projective. And by Proposition 
\ref{prop:1520}, a complex $P \in \dcat{C}(A, \mrm{gr})$ is 
K-graded-projective iff $\opn{Hom}_{\dcat{K}(A, \mrm{gr})}(P, N) = 0$ for every 
acyclic complex $N \in \dcat{C}(A, \mrm{gr})$.

First let's assume that there is an isomorphism $P \cong Q$ in 
$\dcat{K}(A, \mrm{gr})$ to a semi-graded-free complex $Q$.
Then for every acyclic complex 
$N \in \dcat{C}(A, \mrm{gr})$ we get 
\[ \opn{Hom}_{\dcat{K}(A, \mrm{gr})}(P, N) \cong 
\opn{Hom}_{\dcat{K}(A, \mrm{gr})}(Q, N) = 0 . \]
Therefore $P$ is K-graded-projective.

Conversely, let's assume that $P$ is K-graded-projective. By Theorem 
\ref{thm:3740}(1) there is a quasi-isomorphism
$\rho : Q \to P$ in $\dcat{C}_{\mrm{str}}(A, \mrm{gr})$ from a 
semi-graded-free complex $Q$. By Corollary \ref{cor:1923}, 
$\opn{P}(\rho) : Q \to P$ is an 
isomorphism in $\dcat{K}(A, \mrm{gr})$.

\medskip \noindent
(2) Take $N \in \dcat{C}(A^{\mrm{op}}, \mrm{gr})$. 
We have the adjunction isomorphism
\[ \begin{aligned}
& (N \ot_A P)^* = \opn{Hom}_{\K}(N \ot_A P, \K) 
\\
& \quad \cong
\opn{Hom}_{A}(P, \opn{Hom}_{\K}(N, \K)) = \opn{Hom}_{A}(P, N^*) 
\end{aligned} \]
in $\dcat{C}_{\mrm{str}}(\K, \mrm{gr})$. 
If $N$ is acyclic then so are $N^*$, 
$\opn{Hom}_{A}(P, N^*)$ and $(N \ot_A P)^*$.
Since the functor $(-)^*$ is faithfully exact, we see that 
$N \ot_A P$ is acyclic too. Therefore $P$ is K-graded-flat.

\medskip \noindent
(3) Take $N \in \dcat{C}(A^{\mrm{op}}, \mrm{gr})$. 
We have the adjunction isomorphism
\[ \begin{aligned}
& \opn{Hom}_{A^{\mrm{op}}}(N, P^*) =
\opn{Hom}_{A^{\mrm{op}}}(N, \opn{Hom}_{\K}(P, \K))
\\
& \quad \cong \opn{Hom}_{\K}(N \ot_A P, \K) = (N \ot_A P)^* 
\end{aligned} \]
in $\dcat{C}_{\mrm{str}}(\K, \mrm{gr})$. 
By item (2) $P$ is K-graded-flat. If $N$ is acyclic then so are
$N \ot_A P$, $(N \ot_A P)^*$ and $\opn{Hom}_{A^{\mrm{op}}}(N, P^*)$. 
Therefore $P^*$ is K-graded-injective.

\medskip \noindent
(4) By Theorem \ref{thm:1665}, that applies to this context too, a 
semi-graded-cofree complex $J$ is K-graded-injective. And by Proposition 
\ref{prop:1515}, a complex 
$I \in \dcat{C}(A, \mrm{gr})$ is K-graded-injective iff
$\opn{Hom}_{\dcat{K}(A, \mrm{gr})}(N, I) = 0$ for every acyclic complex 
$N \in \dcat{C}(A, \mrm{gr})$.

First let's assume that there is an isomorphism $I \cong J$ in 
$\dcat{K}(A, \mrm{gr})$ to a semi-graded-cofree complex $J$.
Then for every acyclic complex 
$N \in \dcat{C}(A, \mrm{gr})$ we get 
\[ \opn{Hom}_{\dcat{K}(A, \mrm{gr})}(N, I) \cong 
\opn{Hom}_{\dcat{K}(A, \mrm{gr})}(N, J) = 0 . \]
Therefore $I$ is K-graded-injective.

Conversely, let's assume that $I$ is K-graded-injective. By Theorem 
\ref{thm:3740}(2) there is a quasi-isomorphism
$\rho : I \to J$ in $\dcat{C}_{\mrm{str}}(A, \mrm{gr})$ to a 
semi-graded-cofree complex $J$. By Corollary \ref{cor:1922}, 
$\opn{P}(\rho) : I \to J$ is an 
isomorphism in $\dcat{K}(A, \mrm{gr})$.
\end{proof}

The next proposition tells us that some resolving properties of complexes are 
preserved by the restriction functor
$\dcat{C}(A \ot B, \mrm{gr}) \to \dcat{C}(A, \mrm{gr})$.

\begin{prop} \label{prop:4906}  
Let $A$ and $B$ be graded rings.
\begin{enumerate}
\item If $P \in \dcat{C}(A \ot B, \mrm{gr})$ is K-graded-projective, then it is 
K-graded-projective in $\dcat{C}(A, \mrm{gr})$.

\item If $P \in \dcat{C}(A \ot B, \mrm{gr})$ is K-graded-flat, then it is 
K-graded-flat in $\dcat{C}(A, \mrm{gr})$.

\item If $I \in \dcat{C}(A \ot B, \mrm{gr})$ is K-graded-injective, then it is 
K-graded-injective in $\dcat{C}(A, \mrm{gr})$.
\end{enumerate}
\end{prop}

\begin{proof} \mbox{}

\smallskip \noindent
(1) Take $N \in \dcat{C}(A, \mrm{gr})$. 
We have isomorphisms
\[ \begin{aligned}
& \opn{Hom}_{A}(P, N) \cong
\opn{Hom}_{A \ot B}(P, \opn{Hom}_A(A \ot B, N)) 
\\ & \quad  \cong \opn{Hom}_{A \ot B}(P, \opn{Hom}_{\K}(B, N)) 
\end{aligned} \]
in $\dcat{C}_{\mrm{str}}(\K, \mrm{gr})$. 
As an object of $\dcat{C}(\K, \mrm{gr})$, $B$ is K-graded-projective.
Hence, if $N$ is acyclic, then so is $\opn{Hom}_{A}(P, N)$.

\medskip \noindent
(2, 3) Like item (1), but using the fact that $B$ is 
K-graded-flat over $\K$, and using other adjunction isomorphisms.
\end{proof}

\begin{cor} \label{cor:3830}
Let $A$ and $B$ be graded rings, and let 
$M \in \dcat{C}(A \ot B, \mrm{gr})$.
\begin{enumerate}
\item There exists a quasi-isomorphism 
$P \to M$ in $\dcat{C}_{\mrm{str}}(A \ot B, \mrm{gr})$,
where $P$ is a K-graded-projective complex over $A$, 
every $P^i$ is a graded-projective module over $A$, and
$\opn{sup}(P) = \opn{sup}(\opn{H}(M))$. 

\item There exists a quasi-isomorphism 
$P \to M$ in $\dcat{C}_{\mrm{str}}(A \ot B, \mrm{gr})$,
where $P$ is a K-graded-flat complex over $A$, 
every $P^i$ is a graded-flat module over $A$, and
$\opn{sup}(P) = \opn{sup}(\opn{H}(M))$. 

\item There exists a quasi-isomorphism 
$M \to I$ in $\dcat{C}_{\mrm{str}}(A \ot B, \mrm{gr})$,
where $I$ is a K-graded-injective complex over $A$, 
every $I^p$ is a graded-injective module over $A$, and
$\opn{inf}(I) = \opn{inf}(\opn{H}(M))$. 
\end{enumerate}
\end{cor}

\begin{proof} \mbox{}

\smallskip \noindent
(1) Take the resolution $P \to M$ from Theorem \ref{thm:3740}(1). 
In view of Theorem \ref{thm:4275}(1), Proposition \ref{prop:4906}(1) 
and Proposition \ref{prop:3715}(2),
the complex $P$ has the required properties. 

\medskip \noindent
(2) Use item (1), plus Theorem \ref{thm:4275}(2). 

\medskip \noindent
(3) Take the resolution $M \to I$ from Theorem \ref{thm:3740}(2).
This will do, by Theorem \ref{thm:4275}(4), Proposition \ref{prop:4906}(3), 
Proposition \ref{prop:4905}(3) and Proposition \ref{prop:3900}(3).
\end{proof}

\begin{cor} \label{cor:3740}
Let $A$, $B$ and $C$ be graded rings.
\begin{enumerate}
\item The triangulated right derived bifunctor 
\[ \qquad 
\opn{RHom}_A(-, -) : \dcat{D}(A  \ot B^{\mrm{op}}, \mrm{gr})^{\mrm{op}}
\times \dcat{D}(A  \ot C^{\mrm{op}}, \mrm{gr}) \to 
\dcat{D}(B \ot C^{\mrm{op}}, \mrm{gr}) \]
exists. 

\item Let $M \in \dcat{D}(A  \ot B^{\mrm{op}}, \mrm{gr})$ and
$N \in \dcat{D}(A  \ot C^{\mrm{op}}, \mrm{gr})$.
Assume that either $M$ is K-graded-projective over $A$, or $N$
is K-graded-injective over $A$. Then the morphism 
\[ \eta^{\mrm{R}}_{M, N} : \opn{Hom}_A(M, N) \to 
\opn{RHom}_A(M, N) \]
in $\dcat{D}(B \ot C^{\mrm{op}}, \mrm{gr})$
is an isomorphism. 

\item Let $B' \to B$ and $C' \to C$ be graded $\K$-ring homomorphisms. Then 
the diagram 
\[ \UseTips \xymatrix @C=13ex @R=6ex {
\dcat{D}(A  \ot B^{\mrm{op}}, \mrm{gr})^{\mrm{op}} \times 
\dcat{D}(A  \ot C^{\mrm{op}}, \mrm{gr})
\ar[r]^(0.6){\opn{RHom}_A(-, -)}
\ar[d]_{\opn{Rest} \times}^{\opn{Rest}}
&
\dcat{D}(B \ot C^{\mrm{op}}, \mrm{gr})
\ar[d]^{\opn{Rest}}
\\
\dcat{D}(A  \ot {B'}^{\, \mrm{op}}, \mrm{gr})^{\mrm{op}} \times 
\dcat{D}(A  \ot {C'}^{\, \mrm{op}}, \mrm{gr})
\ar[r]^(0.6){\opn{RHom}_A(-, -)}
&
\dcat{D}(B' \ot C'^{\, \mrm{op}}, \mrm{gr})
} \]
is commutative up to isomorphism. 
\end{enumerate}
\end{cor}

\begin{proof}
The same arguments used in the proof of Proposition \ref{prop:1031} work here, 
because there exist enough acyclic resolutions -- see Corollary \ref{cor:3830}.
\end{proof}

\begin{cor} \label{cor:3861}
The category $\dcat{D}(A, \mrm{gr})$ has infinite direct sums. On objects, the 
direct sum is degreewise, for the $\Z^2$ grading. 
\end{cor}

\begin{proof}
The proof of Theorem \ref{thm:3140} works here. 
\end{proof}

\begin{cor} \label{cor:3741} 
Let $A$, $B$ and $C$ be graded rings.
\begin{enumerate}
\item The triangulated left derived bifunctor 
\[ \qquad (- \ot^{\mrm{L}}_{A} -) : 
\dcat{D}(B \ot A^{\mrm{op}}, \mrm{gr})
\times \dcat{D}(A  \ot C^{\mrm{op}}, \mrm{gr}) \to 
\dcat{D}(B \ot C^{\mrm{op}}, \mrm{gr}) \]
exists. 

\item Let $M \in \dcat{D}(B \ot A^{\mrm{op}}, \mrm{gr})$ and
$N \in \dcat{D}(A  \ot C^{\mrm{op}}, \mrm{gr})$.
Assume that either $M$ is K-graded-flat over $A^{\mrm{op}}$, or $N$
is K-graded-flat over $A$. Then the morphism 
\[ \eta^{\mrm{L}}_{M, N} : M \ot^{\mrm{L}}_{A} N \to M \ot_{A} N \]
in $\dcat{D}(B \ot C^{\mrm{op}}, \mrm{gr})$
is an isomorphism. 

\item Let $B' \to B$ and $C' \to C$ be graded $\K$-ring homomorphisms. Then 
the diagram 
\[ \UseTips \xymatrix @C=11ex @R=6ex {
\dcat{D}(B \ot A^{\mrm{op}}, \mrm{gr}) \times 
\dcat{D}(A \ot C^{\mrm{op}}, \mrm{gr})
\ar[r]^(0.62){ (- \ot^{\mrm{L}}_{A} -) }
\ar[d]_{\opn{Rest} \times}^{\opn{Rest}}
&
\dcat{D}(B \ot C^{\mrm{op}}, \mrm{gr})
\ar[d]^{\opn{Rest}}
\\
\dcat{D}(B' \ot A^{\mrm{op}}, \mrm{gr}) \times 
\dcat{D}(A  \ot C'^{\, \mrm{op}}, \mrm{gr})
\ar[r]^(0.62){ (- \ot^{\mrm{L}}_{A} -) }
&
\dcat{D}(B' \ot C'^{\, \mrm{op}}, \mrm{gr})
} \]
is commutative up to isomorphism. 
\end{enumerate}
\end{cor}

\begin{proof}
The same arguments used in the proof of Proposition \ref{prop:3450} work here, 
because there exist enough acyclic resolutions -- see Corollary \ref{cor:3830}.
\end{proof}

\begin{cor} \label{cor:3835}
Let $A$, $B$ and $C$ be graded rings.
\begin{enumerate}
\item Given $M \in \dcat{M}(A  \ot B^{\mrm{op}}, \mrm{gr})$ and
$N \in \dcat{M}(A  \ot C^{\mrm{op}}, \mrm{gr})$, for every $p \in \N$ there is a 
canonical isomorphism 
\[ \opn{Ext}^p_A(M, N) \cong 
\opn{H}^p \bigl( \opn{RHom}_A(M, N) \bigr) \]
in $\dcat{M}(B \ot C^{\mrm{op}}, \mrm{gr})$.
Here $\opn{Ext}^p_A(-, -)$ is the classical $p$-th right derived bifunctor of 
$\opn{Hom}_A(-, -)$. 

\item Given $M \in \dcat{M}(B  \ot A^{\mrm{op}}, \mrm{gr})$ and
$N \in \dcat{M}(A  \ot C^{\mrm{op}}, \mrm{gr})$, for every $p \in \N$ there is a 
canonical isomorphism 
\[ \opn{Tor}_p^A(M, N) \cong 
\opn{H}^{-p} \bigl( M \ot^{\mrm{L}}_{A} N \bigr) \]
in $\dcat{M}(B \ot C^{\mrm{op}}, \mrm{gr})$.
Here $\opn{Tor}_p^A(-, -)$ is the classical $p$-th left derived bifunctor of 
$(- \ot_A -)$. 
\end{enumerate}
\end{cor}

\begin{proof}
This is clear from Corollary \ref{cor:3740}(2) and Corollary \ref{cor:3741}(2) 
respectively.
\end{proof}

\begin{cor} \label{cor:3742}
Let $A$, $B$, $C$ and $D$ be graded rings.
\begin{enumerate}
\item Let $L \in \dcat{D}(A \ot B^{\mrm{op}}, \mrm{gr})$, 
$M \in \dcat{D}(B \ot C^{\mrm{op}}, \mrm{gr})$ and 
$N \in \dcat{D}(C \ot D^{\mrm{op}}, \mrm{gr})$.
There is an isomorphism 
\[ L \ot^{\mrm{L}}_{B} (M \ot^{\mrm{L}}_{C} N) \cong 
(L \ot^{\mrm{L}}_{B} M) \ot^{\mrm{L}}_{C} N \]
in 
$\dcat{D}(A \ot D^{\mrm{op}}, \mrm{gr})$. 
This isomorphism is functorial in $L, M, N$.

\item Let $L \in \dcat{D}(B \ot A^{\mrm{op}}, \mrm{gr})$, 
$M \in \dcat{D}(C \ot B^{\mrm{op}}, \mrm{gr})$ and 
$N \in \dcat{D}(C \ot D^{\mrm{op}}, \mrm{gr})$.
There is an isomorphism 
\[ \opn{RHom}_{B} \bigl( L, \opn{RHom}_{C}(M, N) \bigr) \cong
\opn{RHom}_{C}(M \ot^{\mrm{L}}_{B} L, N) \]
in 
$\dcat{D}(A \ot D^{\mrm{op}}, \mrm{gr})$. 
This isomorphism is functorial in $L, M, N$.
\end{enumerate}
\end{cor}

\begin{proof}
Both assertions are immediate consequences of the associativity of the tensor 
product for modules, and the Hom-tensor adjunction for modules. 
\end{proof}

The triangulated category $\dcat{D}(A^{\mrm{en}}, \mrm{gr})$ has a biclosed 
monoidal structure on it, with monoidal operation 
$(- \ot^{\mrm{L}}_{A} -)$ and monoidal unit $A$. See Remark \ref{rem:3450} for 
the nongraded variant. In the next definition we 
present the pertinent operations related to this structure; these will be 
needed for several constructions. We do not need to know anything about 
the concept of monoidal structure beyond this terminology. 
The interested reader can look it up in \cite{Mac2}; the previous corollaries 
(almost) guarantee that the monoidal axioms hold in this context.

\begin{dfn}[Monoidal Operations] \label{dfn:4530}
\index{Monoidal category! structure on $\dcat{D}(A^{\mrm{en}}, \mrm{gr})$}
Let $A$ and $B$ be graded rings, and let 
$M \in \dcat{D}(A \ot B^{\mrm{op}}, \mrm{gr})$ and 
$N \in \dcat{D}(B \ot A^{\mrm{op}}, \mrm{gr})$ be complexes. 
\begin{enumerate}
\item The {\em left unitor}% 
\index{Monoidal category! left unitor isomorphism}% 
\index{1-Lu@$\opn{lu}$}
isomorphism for $M$ is the obvious isomorphism 
$\opn{lu} : A \ot_A^{\mrm{L}} M \iso M$
in $\dcat{D}(A \ot B^{\mrm{op}}, \mrm{gr})$. 

\item The {\em right unitor}% 
\index{Monoidal category! right unitor isomorphism}%
\index{1-Ru@$\opn{ru}$}
isomorphism for $N$ is the obvious isomorphism 
$\opn{ru} : N \ot_A^{\mrm{L}} A \iso N$
in $\dcat{D}(B \ot A^{\mrm{op}}, \mrm{gr})$. 

\item The {\em left co-unitor}% 
\index{Monoidal category! left co-unitor isomorphism}%
\index{1-Lcu@$\opn{lcu}$}
isomorphism for $M$ is the obvious isomorphism  
$\opn{lcu} : \opn{RHom}_{A}(A, M) \iso M$
in $\dcat{D}(A \ot B^{\mrm{op}}, \mrm{gr})$.
\end{enumerate}
\end{dfn}

Here is a notion that is dual to a minimal complex of injectives
(see Definition \ref{dfn:2210}).  
Recall that for a complex $P$, the module of degree $i$ coboundaries is 
$\opn{B}^i(P) = \opn{Im} \bigl( \d_P^{i - 1} : P^{i - 1} \xar{} P^i \bigr)
\sub P^i$.

\begin{dfn} \label{dfn:4531}
Let $A$ be a connected graded ring, with augmentation ideal $\m$. 
A {\em minimal complex of graded-free $A$-modules}% 
\index{Complex of algebraically graded modules! minimal {\indash} graded-free
{\indash}}
is a bounded above complex $P$, such that for every $i$ the graded module $P^i$ 
is graded-free, and there is an inclusion 
$\opn{B}^i(P) \sub \m \cd P^i$.
\end{dfn}

\begin{prop} \label{prop:4530}
Let $A$ be a connected graded ring, and let $P$ be a bounded above complex
of graded-free $A$-modules.
The following conditions are equivalent\tup{:}
\begin{itemize}
\rmitem{i} $P$ is a minimal complex of graded-free $A$-modules.

\rmitem{ii} The complex $\K \ot_A P$ has zero differential.
\end{itemize}
\end{prop}

\begin{exer} \label{exer:4530}
Prove this proposition. 
\end{exer}

\begin{dfn} \label{dfn:4235}
Let $A$ be a connected graded ring, and let $M \in \dcat{M}(A, \mrm{gr})$. 
A {\em minimal graded-free resolution}% 
\index{Resolution! minimal graded-free}
of $M$ is a quasi-isomorphism 
$\rho : P \to M$ in $\dcat{C}_{\mrm{str}}(A, \mrm{gr})$ from a minimal
complex of graded-free $A$-modules $P$.
\end{dfn}

\begin{prop} \label{prop:4265}   
Let $A$ be a connected graded ring, and let 
$M \in \dcat{M}(A, \mrm{gr})$ be a 
bounded below graded module. Then\tup{:}
\begin{enumerate}
\item $M$ has a minimal graded-free resolution $\eta : P \to M$.

\item For every $i$ there is an isomorphism 
$P^i \cong A \ot_{} \opn{Tor}^A_{-i}(\K, M)$
in $\dcat{M}(A, \mrm{gr})$.

\item The minimal graded-free resolution $\eta : P \to M$ in item \tup{(1)} 
above is unique, up to a non-unique isomorphism in 
$\dcat{C}_{\mrm{str}}(A,\mrm{gr})$.
\end{enumerate} 
\end{prop}

\begin{proof} \mbox{}

\smallskip \noindent
(1) Let $V_0 := \K \ot_A M$ and $P^0 := A \ot V_0$. By 
Proposition \ref{prop:4590} there is a surjection 
$\eta : P^0 \to M$. Let 
$M^{-1} := \opn{Ker}(\eta) \sub P^0$. 
By construction we have 
$M^{-1} \sub \m \ot V_0 = \m \cd P^0$. 
And of course $M^{-1}$ is a bounded below graded $A$-module. 

Next let $V_{1} := \K \ot_A M^{-1}$ and $P^{-1} := A \ot V_{1}$. 
By Proposition \ref{prop:4590} there is a surjection 
$\d^{-1} : P^{-1} \to M^{-1}$. The graded module  
$M^{-2} := \opn{Ker}(\d^{-1}) \sub P^{-1}$
satisfies $M^{-2} \sub \m \ot V_1 = \m \cd P^{-1}$.
And so on. 

\medskip \noindent
(2) Due to minimality, the differential of the complex $\K \ot_A P$ is zero; 
see Proposition \ref{prop:4530}. Therefore
\[ \opn{Tor}^A_{-i}(\K, M) \cong \opn{H}^i(\K \ot^{\mrm{L}}_{A} M) \cong 
\opn{H}^i(\K \ot_{A} P) \cong \K \ot_{A} P^{i} \cong V_{-i} . \]

\medskip \noindent
(3) Suppose $\eta' : P' \to M$ is some other minimal graded-free resolution.
Because $P$ and $P'$ are K-projective in $\dcat{C}(A, \mrm{gr})$, there is a 
homotopy equivalence $\phi : P \to P'$ in $\dcat{C}_{\mrm{str}}(A, \mrm{gr})$.
Therefore 
\begin{equation} \label{eqn:4530}
\opn{id} \ot \, \phi : \K \ot_A P \to \K \ot_A P'
\end{equation}
is a quasi-isomorphism. But by Proposition \ref{prop:4530}, these complexes 
have zero differentials, so (\ref{eqn:4530}) is an isomorphism in 
$\dcat{C}_{\mrm{str}}(\K, \mrm{gr})$. The graded Nakayama Lemma (Proposition 
\ref{prop:4590}(3)) implies that $\phi : P \to P'$ is an isomorphism in 
$\dcat{C}_{\mrm{str}}(A, \mrm{gr})$.
\end{proof}

\begin{exer} \label{exer:4531}
Try to generalize Proposition \ref{prop:4265} to a complex 
$M \in \lb \dcat{C}(A, \mrm{gr})$. The trick is to find the correct
boundedness conditions on $M$.
\end{exer}

\begin{cor} \label{cor:4237}
Assume $A$ is a left noetherian connected graded ring, and $M$ is a 
finite graded $A$-module. Then in the minimal graded-free 
resolution $\eta : P \to M$, the graded-free $A$-modules $P^i$ are all finite. 
\end{cor}

\begin{proof}
In the proof of Proposition \ref{prop:4265}(1), the graded $\K$-modules $V_i$ 
are all finite. 
\end{proof}

Here is a graded version of Definitions \ref{dfn:1925} and \ref{dfn:4470}. 

\begin{dfn} \label{dfn:4532} 
Let $A$ be a graded ring. 
\begin{enumerate}
\item A complex $P \in \dcat{C}(A, \mrm{gr})$
is called {\em pseudo-finite semi-graded-free}% 
\index{Complex of algebraically graded modules! pseudo-finite semi-graded-free}
if it is a bounded above complex of finite graded-free $A$-modules.

\item A complex $L \in \dcat{D}(A, \mrm{gr})$
is called {\em derived graded-pseudo-finite}% 
\index{Complex of algebraically graded modules! derived graded-pseudo-finite}
if it belongs to the saturated 
full triangulated subcategory of $\dcat{D}(A, \mrm{gr})$ generated by the 
pseudo-finite semi-graded-free complexes.
\end{enumerate}
\end{dfn}

\begin{exa} \label{exa:4390}
If the graded ring $A$ is connected and left noetherian, then the 
derived graded-pseudo-finite complexes over $A$ are precisely the objects 
of $\dcat{D}^{-}_{\mrm{f}}(A, \mrm{gr})$.
See Theorems \ref{thm:2895} or \ref{thm:3340}. 
\end{exa}

\begin{prop} \label{prop:4266}  
Assume $A$ is a left noetherian connected graded ring.
Consider $A$ as a graded $A^{\mrm{en}}$-module. Then in the minimal graded-free 
resolution $\rho : P \to A$ of $A$ over $A^{\mrm{en}}$, the graded-free 
$A^{\mrm{en}}$-modules $P^i$ are all finite. 
\end{prop}

The subtle point is that the graded ring $A^{\mrm{en}}$ might not be left 
noetherian -- see Remark \ref{rem:3901}. 

\begin{proof}
Let us write $P^{i} \cong A^{\mrm{en}} \ot V_{-i}$, where 
$V_{-i} \in \dcat{M}(\K, \mrm{gr})$. We shall prove that all the
$V_{-i}$ are finite graded $\K$-modules. 

When we restrict the quasi-isomorphism $\rho : P \to A$
to the category $\dcat{C}_{\mrm{str}}(A^{\mrm{op}}, \mrm{gr})$, 
it becomes a homotopy equivalence. Hence 
$\tau := \rho \ot \opn{id} : P \ot_{A} \K \to A \ot_{A} \K$
is a quasi-isomorphism in $\dcat{C}_{\mrm{str}}(A, \mrm{gr})$. 
Writing 
$Q := P \ot_{A} \K \in \dcat{C}(A, \mrm{gr})$,
we get a graded-free resolution 
$\tau : Q \to \K$ over $A$. Note that 
$Q^{i} \cong A \ot V_{-i}$ in $\dcat{C}_{\mrm{str}}(A, \mrm{gr})$.

Now 
\begin{equation} \label{eqn:5433}
\K \ot_A Q \cong \K \ot_A P \ot_A \K \cong \K \ot_{A^{\mrm{en}}} P
\end{equation}
in $\dcat{C}_{\mrm{str}}(\K, \mrm{gr})$. Because $P$ is a minimal 
complex of graded-free $A^{\mrm{en}}$-modules, 
the complex $\K \ot_{A^{\mrm{en}}} P$ has zero differential; see 
Proposition \ref{prop:4530}.
But then the complex $\K \ot_A Q$ has zero differential. We conclude that $Q$ 
is a minimal complex of graded-free $A$-modules, and hence 
$\tau : Q \to \K$ is a minimal graded-free resolution over $A$. Corollary 
\ref{cor:4237} says that the graded $A$-modules $Q^i$ are finite; and hence the 
graded $\K$-modules $V_{-i}$ are finite. 
\end{proof}

\begin{cor} \label{cor:4531}
If $A$ is a left noetherian connected graded ring, then $A$ is a 
derived graded-pseudo-finite complex over $A^{\mrm{en}}$. 
\end{cor}

\begin{proof}
This is an immediate consequence of Proposition \ref{prop:4266},
which says that the complex $P$ is pseudo-finite graded-semi-free over 
$A^{\mrm{en}}$.
\end{proof}

\begin{thm} \label{thm:4535}  
Let $A$, $B$, $C$ and $D$ be graded rings. For complexes 
$L \in \dcat{D}(A \ot C^{\mrm{op}}, \mrm{gr})$,
$M \in \dcat{D}(A \ot B^{\mrm{op}}, \mrm{gr})$ and 
$N \in \dcat{D}(B \ot D^{\mrm{op}}, \mrm{gr})$
there is a morphism 
\[ \opn{ev}^{\mrm{R, L}}_{L, M, N} :
\opn{RHom}_A(L, M) \ot_B^{\mrm{L}} N \to 
\opn{RHom}_A(L, M \ot_B^{\mrm{L}} N) \]
in $\dcat{D}(C \ot D^{\mrm{op}}, \mrm{gr})$, called 
{\em derived graded tensor-evaluation}.  
\index{Tensor-evaluation morphism! derived graded}
The morphism $\opn{ev}^{\mrm{R, L}}_{L, M, N}$ 
is functorial in the objects $L, M, N$.
Moreover, if all three conditions below hold, then 
$\opn{ev}^{\mrm{R, L}}_{L, M, N}$ is an isomorphism. 
\begin{enumerate}
\rmitem{i} The complex $L$ is derived graded-pseudo-finite over $A$.

\rmitem{ii} The complex $M$ has bounded below cohomology. 

\rmitem{iii} The complex $N$ has finite graded-flat dimension over $B$. 
\end{enumerate}
\end{thm}

Just to clarify, condition (ii) says that $\opn{H}^i(M) = 0$ for $i \ll 0$. 
We do not care about the boundedness (in algebraic degree) of the graded 
modules $\opn{H}^i(M)$. 

\begin{proof}
The proof of just like that of Theorem \ref{thm:4320}, but 
working in the graded derived categories. For condition (iii), notice 
that if $N$ has finite graded-flat dimension over $B$, then there is an 
isomorphism $P \cong N$ in $\dcat{D}(B, \mrm{gr})$, where $P$ is a bounded 
complex of graded-flat $B$-modules.
\end{proof}

\begin{dfn} \label{dfn:3732} 
Let $A$ be a graded ring. 
\begin{enumerate}
\item The full subcategory of $\dcat{D}(A, \mrm{gr})$ on the complexes $M$ 
whose cohomology modules $\opn{H}^p(M)$ belong to 
$\dcat{M}_{\mrm{dwf}}(A, \mrm{gr})$
is denoted by $\dcat{D}_{\mrm{dwf}}(A, \mrm{gr})$%
\index{1-Ddwf(A,gr)@$\dcat{D}_{\mrm{dwf}}(A, \mrm{gr})$}.

\item The full subcategory of $\dcat{D}(A, \mrm{gr})$ on the complexes 
whose cohomology modules belong to 
$\dcat{M}_{\mrm{f}}(A, \mrm{gr})$
is denoted by $\dcat{D}_{\mrm{f}}(A, \mrm{gr})$%
\index{1-Df(A,gr)@$\dcat{D}_{\mrm{f}}(A, \mrm{gr})$}.

\item The full subcategory of $\dcat{D}(A, \mrm{gr})$ on the complexes 
whose cohomology modules belong to 
$\dcat{M}_{\mrm{cof}}(A, \mrm{gr})$
is denoted by $\dcat{D}_{\mrm{cof}}(A, \mrm{gr})$%
\index{1-Dcof(A,gr)@$\dcat{D}_{\mrm{cof}}(A, \mrm{gr})$}.
\end{enumerate}
\end{dfn}

Because $\dcat{M}_{\mrm{dwf}}(A, \mrm{gr})$ is a a thick abelian 
subcategory of $\dcat{M}(A, \mrm{gr})$,
it follows that $\dcat{D}_{\mrm{dwf}}(A, \mrm{gr})$
is a full triangulated subcategory of 
$\dcat{D}(A, \mrm{gr})$. Likewise for 
$\dcat{D}_{\mrm{f}}(A, \mrm{gr})$
when $A$ is left noetherian, and for 
$\dcat{D}_{\mrm{cof}}(A, \mrm{gr})$
when $A$ is right noetherian connected.

Because $(-)^*$ is exact, it induces a triangulated functor 
\begin{equation} \label{eqn:3733}
(-)^* : \dcat{D}(A, \mrm{gr})^{\mrm{op}} \to 
\dcat{D}(A^{\mrm{op}}, \mrm{gr}) . 
\end{equation}

\begin{prop} \label{prop:4591}   
Let $A$ be a graded ring.
\begin{enumerate}
\item If $M \in \dcat{D}_{\mrm{dwf}}(A, \mrm{gr})$ then 
$M^* \in \dcat{D}_{\mrm{dwf}}(A^{\mrm{op}}, \mrm{gr})$.

\item The functor 
\[ (-)^* : \dcat{D}_{\mrm{dwf}}(A, \mrm{gr})^{\mrm{op}} \to 
\dcat{D}_{\mrm{dwf}}(A^{\mrm{op}}, \mrm{gr}) \]
is an equivalence of $\K$-linear triangulated categories. 

\item If $A$ is left noetherian connected, then the functor 
\[ (-)^* : \dcat{D}_{\mrm{f}}(A, \mrm{gr})^{\mrm{op}} \to 
\dcat{D}_{\mrm{cof}}(A^{\mrm{op}}, \mrm{gr}) \]
is an equivalence of $\K$-linear triangulated categories. 
\end{enumerate}
\end{prop}

\begin{exer} \label{exer:3830}
Prove this proposition. (Hint: use Proposition \ref{prop:4045} and 
Theorem \ref{thm:4045}.) 
\end{exer} 

The next result is from \cite{VdB}. 

\begin{thm} \label{thm:3800} 
Let $A$, $B$ and $C$  be graded $\K$-rings. 
\begin{enumerate}
\item For 
$M \in \dcat{D}(A \ot B^{\mrm{op}}, \mrm{gr})$ and
$N \in \dcat{D}(A \ot C^{\mrm{op}}, \mrm{gr})$
there is a morphism 
\[ \th_{M, N} :  \opn{RHom}_A(M, N) \to \opn{RHom}_{A^{\mrm{op}}}(N^*, M^*) \]
in $\dcat{D}(B \ot C^{\mrm{op}}, \mrm{gr})$, that is functorial in $M$ and $N$.

\item If
$M \in \dcat{D}_{\mrm{dwf}}(A \ot B^{\mrm{op}}, \mrm{gr})$ 
and
$N \in \dcat{D}_{\mrm{dwf}}(A \ot C^{\mrm{op}}, \mrm{gr})$
then $\th_{M, N}$ is an isomorphism. 
\end{enumerate}
\end{thm}

\begin{proof} \mbox{}

\smallskip \noindent
(1) Choose a K-graded-projective resolution $\rho : P \to M$ in 
$\dcat{C}_{\mrm{str}}(A \ot B^{\mrm{op}}, \mrm{gr})$. 
By Proposition \ref{prop:4906} the complex $P$ is 
K-graded-projective over $A$; so there is a canonical isomorphism 
\[ \phi_{\rho} : \opn{RHom}_A(M, N) \iso \opn{Hom}_A(P, N) \]
in $\dcat{D}(B \ot C^{\mrm{op}}, \mrm{gr})$. 
$\K$-linear duality gives a homomorphism
\[ \til{\th}_{P, N} : \opn{Hom}_A(P, N) \to 
\opn{Hom}_{A^{\mrm{op}}}(N^*, P^*) \]
in $\dcat{C}_{\mrm{str}}(B \ot C^{\mrm{op}}, \mrm{gr})$.

Now $\rho^* : M^* \to P^*$ is a quasi-isomorphism
in $\dcat{C}_{\mrm{str}}(B \ot A^{\mrm{op}}, \mrm{gr})$,
and according to Proposition \ref{prop:4906}(1) and 
Theorem \ref{thm:4275}(3) the the complex $P^*$ is 
K-graded-injective over $A^{\mrm{op}}$. Therefore there is a canonical 
isomorphism 
\[ \psi_{\rho^*} : \opn{Hom}_{A^{\mrm{op}}}(N^*, P^*) \iso 
\opn{RHom}_{A^{\mrm{op}}}(N^*, M^*) \]
in $\dcat{D}(B \ot C^{\mrm{op}}, \mrm{gr})$. 
The resulting morphism 
\[ \th_{M, N} := \psi_{\rho^*} \circ \opn{Q}(\til{\th}_{P, N}) \circ \phi_{\rho}
: \opn{RHom}_A(M, N) \to \opn{RHom}_{A^{\mrm{op}}}(M^*, N^*) \]
in $\dcat{D}(B \ot C^{\mrm{op}}, \mrm{gr})$ is functorial in $M$ and $N$.

\medskip \noindent
(2) Because the restriction functors between the derived categories are 
conservative (see Proposition \ref{prop:3445}), 
we can forget the rings $B$ and $C$, and just consider 
$\th_{M, N}$ as a morphism $\dcat{D}(\K, \mrm{gr})$. 

For every integer $i$ there is a canonical isomorphism
\begin{equation} \label{eqn:3805}
\opn{H}^i \bigl( \opn{RHom}_A(M, N) \bigr)  
\cong \opn{Hom}_{\dcat{D}_{\mrm{dwf}}( A, \mrm{gr})}
\bigl( M, N[i] \bigr) .
\end{equation}
Using the isomorphism 
$N^*[-i] \cong N[i]^*$ in 
$\dcat{D}(A^{\mrm{op}}, \mrm{gr})$ we also have  
\begin{equation} \label{eqn:3806}
\begin{aligned}
& \opn{H}^i \bigl( \opn{RHom}_{A^{\mrm{op}}}(N^*, M^*) \bigr) \cong 
\opn{Hom}_{\dcat{D}_{\mrm{dwf}}(A^{\mrm{op}}, \mrm{gr})}
\bigl( N^*, M^*[i] \bigr)
\\
& \quad \cong
\opn{Hom}_{\dcat{D}_{\mrm{dwf}}(A^{\mrm{op}}, \mrm{gr})}
\bigl( N[i]^*, M^* \bigr) . 
\end{aligned}
\end{equation}
The construction of $\th_{M, N}$ in the proof of item (1) above shows that the 
$\K$-linear homomorphism 
\[ \opn{H}^i(\th_{M, N}) :  \opn{H}^i \bigl( \opn{RHom}_A(M, N) \bigr) \to
\opn{H}^i \bigl( \opn{RHom}_{A^{\mrm{op}}}(N^*, M^*) \bigr) \]
becomes, after using the isomorphisms (\ref{eqn:3805}) and 
(\ref{eqn:3806}), the duality morphism 
\[ \opn{Hom}_{\dcat{D}_{\mrm{dwf}}(A, \mrm{gr})}
\bigl( M, N[i] \bigr)  \to  
\opn{Hom}_{\dcat{D}_{\mrm{dwf}}(A^{\mrm{op}}, \mrm{gr})}
\bigl( N[i]^*, M^* \bigr) . \]
According to Proposition \ref{prop:4591}(2) this is an isomorphism.
\end{proof}

\begin{rem} \label{rem:3740}
In this section, and the two sections following it, we work over a base field 
$\K$. This simplifies a lot of the constructions. It is however possible to 
relax this condition, and to work over a noetherian commutative base ring $\K$ 
that is not a field. But then there are (at least) two complications. 
The first complication is that we will have to fix an injective cogenerator 
$\K^*$ of $\dcat{M}(\K)$; cf.\ Subsection \ref{subsec:exis-K-inj-dgmods}. 
Matlis Duality will take the form 
$M^* := \opn{Hom}_{\K}(M, \K^*)$,
and $\opn{Hom}_{\K}(-, -)$ might also involve continuity in a delicate way. 

The second complication is that for the homological algebra to be effective, 
the graded $\K$-rings would have to be {\em flat}. 
The absence of flatness poses serious difficulties. This can be tackled using 
{\em K-flat DG ring resolutions}; and these DG rings would have to be bigraded 
(with upper and lower indices, as explained in the Subsection 
\ref{subsec:alg-gr-mods}). This sort 
of resolution has not been studied yet, as far as we know; but cf.\ 
\cite{Ye16} for some preliminary ideas. 
\end{rem}





\mysubsection{Artin-Schelter Regular Graded Rings} 
\label{AS-reg}

In this subsection we talk about an important class of graded rings: the 
{\em Artin-Schelter regular graded rings}. These are connected graded 
noncommutative rings, that homologically are similar to commutative polynomial 
rings. A great deal of noncommutative ring theory and homological algebra grew 
out of the study of this class of rings. See Remark \ref{rem:4653} for 
some historical notes.

We continue with Conventions \ref{conv:3700} and \ref{conv:4560}. 
Thus all graded rings are central over the base field $\K$. 

Suppose $A$ is a connected graded ring. Recall that the augmentation ideal of 
$A$ is denoted by $\m$. The opposite ring of $A$ is $A^{\mrm{op}}$, and the 
enveloping ring is $A^{\mrm{en}}$.
We view $A / \m \cong \K$ as a graded $A^{\mrm{en}}$-module.

\begin{dfn} \label{dfn:3920} 
Let $A$ be a noetherian connected graded ring. 
\begin{enumerate}
\item The graded ring $A$ is called a {\em regular graded ring}
\index{Algebraically graded ring! regular}
if it has finite graded global cohomological dimension; namely if there is a 
natural number $d$ such that 
$\opn{Ext}^i_A(M, N) = 0$ and
$\opn{Ext}^i_{A^{\mrm{op}}}(M', N') = 0$ 
for all $i > d$,  $M, N \in \dcat{M}(A, \mrm{gr})$ and 
$M', N' \in \dcat{M}(A^{\mrm{op}}, \mrm{gr})$.
The smallest such natural number $d$ is called the 
{\em graded global dimension}
\index{Algebraically graded ring! graded global dimension of}
of the ring $A$.

\item The graded ring $A$ is called a {\em Gorenstein graded ring}
\index{Algebraically graded ring! Gorenstein} 
if the graded bimodule $A$ has finite graded injective dimension 
over $A$ and over $A^{\mrm{op}}$; namely if there is a natural number $d$ 
such that $\opn{Ext}^i_A(M, A) = 0$ and
$\opn{Ext}^i_{A^{\mrm{op}}}(M', A) = 0$ 
for all $i > d$, $M \in \dcat{M}(A, \mrm{gr})$ and 
$M' \in \dcat{M}(A^{\mrm{op}}, \mrm{gr})$.
The smallest such natural number $d$ is called the 
{\em graded injective dimension} 
\index{Algebraically graded ring! graded injective dimension of}
of the ring $A$.
\end{enumerate}
\end{dfn}

Recall that the graded projective dimension of a graded $A$-module $M$ is 
smallest generalized integer $d \in \N \cup \{ \pm \infty \}$ 
such that $\opn{Ext}^i_A(M, N) = 0$ 
for all integers $i > d$ and all $N \in \dcat{M}(A, \mrm{gr})$.

The minimal graded-free resolution of a graded module was introduced in 
Definitions  \ref{dfn:4531} and \ref{dfn:4235}. 
 
\begin{thm} \label{thm:4548} 
Let $A$ be a noetherian connected graded ring, and let $d$ be a natural 
number. The following five conditions are equivalent. 
\begin{itemize}
\rmitem{i} The graded ring $A$ is regular, of graded global dimension $d$. 

\rmitem{ii} The graded projective dimension of the graded $A$-module $\K$ is 
$d$. 

\rmitem{iii} The graded $\K$-modules 
$\opn{Tor}_i^A(\K, \K)$ vanish for all $i > d$, but not for $i = d$. 

\rmitem{iv} The graded $\K$-modules 
$\opn{Ext}^i_A(\K, \K)$ vanish for all $i > d$, but not for $i = d$. 

\rmitem{v} Let $P \to A$ be the minimal graded-free resolution of the graded
$A^{\mrm{en}}$-module $A$. Then the graded $A^{\mrm{en}}$-modules $P^{-i}$  
vanish for all $i > d$, but not for $i = d$.  
\end{itemize}
\end{thm}

Observe that in this theorem, conditions (i), (iii) and (v) are 
op-symmetric, but not conditions (ii) and (iv). 

\begin{proof}
Consider the minimal graded-free resolution
$\rho : P \to A$ of $A$ over $A^{\mrm{en}}$. 
We express the graded-free $A^{\mrm{en}}$-module $P^{-i}$, for $i \geq 0$,  as 
\begin{equation} \label{eqn:4905}
P^{-i} =  A^{\mrm{en}} \ot V_{i} \cong A \ot V_i \ot A 
\end{equation}
with $V_{i} \in \dcat{M}_{\mrm{f}}(\K, \mrm{gr})$. 
The finiteness of the graded $\K$-modules $V_{i}$ is by Proposition 
\ref{prop:4266}. In formula (\ref{eqn:4905}) the ring $A^{\mrm{en}}$ acts on 
itself from the left, and it acts on $A \ot V_i \ot A$ by the outside action,
i.e.\ 
\[ (a_1 \ot \opn{op}(a_2)) \cd (b_1 \ot v \ot b_2) = 
(a_1 \cd b_1) \ot v \ot (b_2 \cd a_2) \]
for $a_1 \ot \opn{op}(a_2) \in A^{\mrm{en}}$ and 
$b_1 \ot v \ot b_2 \in A \ot V_i \ot A$.

As explained in the proof of Proposition \ref{prop:4266}, in the category \lb
$\dcat{C}_{\mrm{str}}(A^{\mrm{op}}, \mrm{gr})$
the homomorphism $\rho : P \to A$ is a homotopy equivalence. Hence, taking 
a module $M \in \dcat{M}(A, \mrm{gr})$, and writing 
$Q_M := P \ot_{A} M$, we get a quasi-isomorphism 
\[ \tau_M := \rho \ot \opn{id}_M : Q_M = P \ot_{A} M \to A \ot_{A} M \cong M  \]
in $\dcat{C}_{\mrm{str}}(A, \mrm{gr})$. 
Because $M$ is a graded-free $\K$-module, it follows that every
\[ Q_M^{-i} = P^{-i} \ot_{A} M \cong A \ot_{} V_{i} \ot_{} M \]
is also a graded-free $A$-module.
In this way we obtain a (functorial) graded-free resolution
$\tau_M : Q_M \to  M$ 
in $\dcat{C}_{\mrm{str}}(A, \mrm{gr})$. 

Assume now that condition (v) holds. This means that 
$V_{i} = 0$ for all $i > d$, and $V_{d} \neq 0$.
The graded-free resolution $Q_M$ has 
length $d$ for every nonzero $M \in \dcat{M}(A, \mrm{gr})$.
We see that  
\[  \opn{Ext}^i_A(M, N) \cong \opn{H}^i \bigl( \opn{Hom}_A(Q_M, N) \bigr)
= 0 \]
for all $N \in  \dcat{M}(A, \mrm{gr})$ and $i > d$. By the op-symmetry 
of condition (v), it also follows that 
$\opn{Ext}^i_{A^{\mrm{op}}}(M', N') = 0$
for all $M', N' \in  \dcat{M}(A^{\mrm{op}}, \mrm{gr})$ and $i > d$.
This says that the graded global dimension of $A$ is at most $d$. 
To get equality, we test for $M = N = \K$. 
In this case we know that $Q_{\K}$ is a minimal graded-free resolution of 
$\K$ over $A$ (see the proof of Proposition \ref{prop:4266}), so the complex 
$\K \ot_A Q_{\K}$ has trivial differential. Therefore
\begin{equation} \label{eqn:4906}
\begin{aligned}
&
\opn{Ext}^i_A(\K, \K) \cong 
\opn{H}^i \bigl( \opn{Hom}_A(Q_{\K}, \K) \bigr)
\\ & \quad 
\cong \opn{H}^i \bigl( \opn{Hom}_{\K}(\K \ot_A Q_{\K}, \K) \bigr) \cong 
\opn{Hom}_{\K}(V_{i}, \K) . 
\end{aligned}
\end{equation}
This is nonzero for $i = d$. 
Likewise over $A^{\mrm{op}}$. 
We have thus verified conditions (i), (ii) and (iv). 
Condition (iii) is verified similarly: the resolution $Q_{\K}$ shows that 
$\opn{Tor}_i^A(\K, \K) \cong V_{i}$
as graded $\K$-modules.

Let us prove the implication (i) $\Rightarrow$ (v):  since  
$\opn{Ext}^i_A(\K, \K) = 0$
for all $i > d$, equations (\ref{eqn:4906}) and (\ref{eqn:4905}) say that 
$P^{-i} = 0$ in this range. On the other hand, there is some pair 
$M, N \in \dcat{M}(A, \mrm{gr})$ such 
that $\opn{Ext}^d_A(M, N) \neq 0$, and this says that $P^{-d} \neq 0$.

The other implications are proved similarly, and we leave them as an exercise. 
\end{proof}

\begin{exer} \label{exer:4550}
Finish the proof of Theorem \ref{thm:4548}.
\end{exer}

The next condition first appeared in the paper \cite{ArSchl} by M. Artin and 
W. Schelter. Throughout, ``AS'' is an abbreviation for ``Artin-Schelter''.  

\begin{dfn} \label{dfn:4545}
Let $A$ be a noetherian connected graded ring. We say that $A$ satisfies the 
{\em AS condition}% 
\index{Algebraically graded ring! AS condition on}
if there are integers $n$ and $l$ such that 
\[ \opn{Ext}^n_A(\K, A) \cong \opn{Ext}^n_{A^{\mrm{op}}}(\K, A) 
\cong \K(l) \]
as graded $\K$-modules, and  
\[ \opn{Ext}^i_A(\K, A) \cong \opn{Ext}^i_{A^{\mrm{op}}}(\K, A) \cong 0 \]
for all $i \neq n$.
The number $n$ is called the {\em AS dimension} of $A$, and the 
number $l$ is called the {\em AS index} of $A$.
\end{dfn}
 
In derived category terms, the AS condition says that 
\begin{equation} \label{eqn:4562}
\opn{RHom}_A(\K, A) \cong \opn{RHom}_{A^{\mrm{op}}}(\K, A) 
\cong \K(l)[-n]
\end{equation}
in $\dcat{D}(\K, \mrm{gr})$, 

\begin{dfn} \label{dfn:3782} 
Let $A$ be a noetherian connected graded ring. 
\begin{enumerate}
\item We say that $A$ is an {\em AS regular graded ring} 
\index{Algebraically graded ring! AS regular}
of dimension $n$ if it
is a regular graded ring (Definition \ref{dfn:3920}(1)), and it 
satisfies the {\em AS condition} (Definition \ref{dfn:4545}) with AS dimension 
$n$. 

\item We say that $A$ is an {\em AS Gorenstein graded ring} 
\index{Algebraically graded ring! AS Gorenstein}
of dimension $n$ if 
it is a Gorenstein graded ring (Definition \ref{dfn:3920}(2)), and it 
satisfies the {\em AS condition} (Definition \ref{dfn:4545}) with AS dimension 
$n$. 
\end{enumerate}
\end{dfn}

As we shall see later, in Subsection \ref{subsec:sym-der-tor}, the AS 
condition implies that the local cohomology of graded $A$-modules behaves as 
though $A$ is commutative. In the paper \cite{ArZh} by Artin and J.J. Zhang, 
subsequent to \cite{ArSchl}, this behavior was formalized as the {\em $\chi$ 
condition}; see Definition \ref{dfn:3713} below. 

Some texts do not make the assumption that the connected graded ring $A$ is 
noetherian, as we did in Definition \ref{dfn:3782}. However, without 
this assumption the theory is not as rich.  

It is natural to ask whether the AS dimension of the ring $A$ coincides 
with its graded global dimension (in the AS regular case), or with its 
graded injective dimension (in the AS Gorenstein case). The next 
propositions say that these numbers agree. 

\begin{prop} \label{prop:4560}
Let $A$ be an AS regular graded ring.
Then the AS dimension of $A$ equals its graded global dimension. 
\end{prop}

\begin{proof}
Let $n$ be the graded global dimension of $A$. 
Let $P \to A$ be the minimal graded-free resolution of the graded
$A^{\mrm{en}}$-module $A$. By Theorem \ref{thm:4548} we know that the 
graded $A^{\mrm{en}}$-modules $P^{-i}$  vanish 
for all $i > n$, but not for $i = n$.  
Let $Q_{\K} \to \K$ be the induced graded-free resolution of the $A$-module 
$\K$ from the proof of Theorem \ref{thm:4548}. Then the complex 
$Q_{\K}$ is concentrated in cohomological degrees $[-n, 0]$, and 
$Q^{-n}_{\K} \neq 0$. Also, as shown there, $Q_{\K} \to \K$ is a 
minimal graded-free resolution of $\K$ over $A$. 

Now let 
$T := \opn{Hom}_A(Q_{\K}, A) \in \dcat{C}(A^{\mrm{op}}, \mrm{gr})$,
so $\opn{RHom}_A(\K, A) \cong T$ 
in $\dcat{D}(\K, \mrm{gr})$. A quick calculation 
shows that $T$ is a minimal complex of graded-free
$A^{\mrm{op}}$-modules. This complex is concentrated in cohomological degrees 
$[0, n]$, and $T^n \neq 0$. Therefore 
$\opn{Ext}^n_A(\K, A) \cong \opn{H}^n (T) \neq 0$.
A similar calculation shows that 
$\opn{Ext}^n_{A^{\mrm{op}}}(\K, A) \neq 0$.
We see that the AS dimension of $A$ is $n$. 
\end{proof}

\begin{prop} \label{prop:4561}
Let $A$ be an AS Gorenstein graded ring.
Then the AS dimension of $A$ equals its graded injective dimension. 
\end{prop}

\begin{proof}
This is an immediate consequence of \cite[Theorem 4.5]{Jor}. 
\end{proof}

In Examples \ref{exa:3780} and \ref{exa:3990} we will present some important 
kinds of AS regular rings. 

\begin{dfn} \label{dfn:4805}
Let $A$ be a nonzero graded ring. 
A central element $a$ of $A$ is called a {\em regular central element}
\index{Regular central element! of a ring} 
if it is not a zero-divisor; i.e.\ $a \cd b = 0$ implies $b = 0$. 
\end{dfn}

If the element $a \in A$ is not central, then we have to distinguish between 
left and right regularity of $a$.

The next two theorems discuss the relation between the rings $A$ and $B$ from 
the following setup.

\begin{setup} \label{set:4565}
Let $A$ be a noetherian connected graded ring, and let $a \in A$ be a
homogeneous regular central element of positive degree. Define the connected 
graded ring $B := A / (a)$.
\end{setup}

\begin{thm} \label{thm:4546}
Assume Setup \tup{\ref{set:4565}}.
\begin{enumerate}
\item If  $B$ is a regular graded ring of graded global dimension $n - 1$, 
then $A$ is a regular graded ring of graded global dimension $n$. 

\item Moreover, if  $B$ is an AS regular graded ring of dimension $n - 1$,  
then $A$ is an AS regular graded of dimension $n$.
\end{enumerate}
\end{thm}

\begin{proof} \mbox{}

\smallskip \noindent
(1) By Theorem \ref{thm:4548}, to prove that $A$ is a regular graded ring
of graded global dimension $n$, 
it is enough to show that the cohomology of the complex
$\K \ot^{\mrm{L}}_{A} \K \in \dcat{D}(\K, \mrm{gr})$
has concentration $[-n, 0]$.
Because a graded-free resolution of $\K$ over $A$ is also a free 
resolution in the ungraded sense, and because we are only interested in the 
vanishing of cohomology, we can forget the algebraic grading, and just look at 
the complex
$\K \ot^{\mrm{L}}_{A} \K \in \dcat{D}(\K)$.
Similarly, because $B$ is a regular graded ring of graded global dimension 
$n - 1$, we know that the complex 
$\K \ot^{\mrm{L}}_{B} \K \in \dcat{D}(\K)$
satisfies 
\begin{equation} \label{eqn:4565}
\opn{con} \bigl( \opn{H}(\K \ot^{\mrm{L}}_{B} \K) \bigr) = [-n + 1, 0] . 
\end{equation}

Let $\til{B} := \opn{K}(A, a)$ be the Koszul complex on the regular central 
element $a \in A$; namely 
\[ \til{B}  = \bigl( \cdots \to 0 \to A \xar{a \cd (-)} A \to 0
\to \cdots \bigr) , \]
concentrated in cohomological degrees $-1$ and $0$. 
This is a DG ring: as a cohomologically graded ring we have 
$\til{B} := A \ot \K[x]$, where $\K[x]$ is the strongly commutative polynomial 
ring on the variable $x$ that has cohomological degree $-1$ 
(see Examples \ref{exa:4625} and \ref{exa:1099}). 
The differential is $\d(x) := a$. 
The canonical DG ring homomorphism $\til{B} \to B$ is a 
quasi-isomorphism. Also $\til{B}$ is semi-free DG module over $A$ and over 
$A^{\mrm{op}}$. According to Theorem \ref{thm:2363}(2) there is an isomorphism
\begin{equation} \label{eqn:4560}
\K \ot^{\mrm{L}}_{B} \K \cong \K \ot^{\mrm{L}}_{\til{B}} \K
\end{equation}
in $\dcat{D}(\K)$. 

By the associativity of the derived tensor product, and because $\K$ is a DG 
$A$-module via the DG ring homomorphism 
$A \to \til{B}$, there are isomorphisms
\begin{equation} \label{eqn:4550}
\K \ot^{\mrm{L}}_{A} \K \cong
\K \ot^{\mrm{L}}_{A} (\til{B} \ot^{\mrm{L}}_{\til{B}} \K) \cong
(\K \ot^{\mrm{L}}_{A} \til{B}) \ot^{\mrm{L}}_{\til{B}} \K
\end{equation}
in $\dcat{D}(\K)$. 
Now the image of the element $a$ in $\K$ is zero, so we obtain these 
isomorphisms
\[ \K \ot^{\mrm{L}}_{A} \til{B} \cong \K \ot_{A} \til{B} \cong 
\bigl( \K \xar{0 \cd (-)} \K \bigr) \cong \K[1] \oplus \K \]
in $\dcat{D}(\til{B}^{\mrm{op}})$.
Substituting this into (\ref{eqn:4550}) we get 
\[ \K \ot^{\mrm{L}}_{A} \K \cong 
\bigl( \K[1] \oplus \K \bigr) \ot^{\mrm{L}}_{\til{B}} \K
\cong (\K \ot^{\mrm{L}}_{\til{B}} \K)[1] \oplus 
(\K \ot^{\mrm{L}}_{\til{B}} \K) \]
in $\dcat{D}(\K)$. 
Hence 
\begin{equation} \label{eqn:4551}
\opn{H}^i \bigl( \K \ot^{\mrm{L}}_{A} \K \bigr) \cong 
\opn{H}^{i + 1} \bigl( \K \ot^{\mrm{L}}_{B} \K \bigr) \oplus 
\opn{H}^{i} \bigl( \K \ot^{\mrm{L}}_{B} \K \bigr) 
\end{equation}
as $\K$-modules. This, together with (\ref{eqn:4565}), imply that 
$\opn{con} \bigl( \opn{H}(\K \ot^{\mrm{L}}_{A} \K) \bigr) = [-n, 0]$,
as required. 

\medskip \noindent
(2) Now $B$ is an AS regular graded ring of dimension $n - 1$.
As explained in part (1) of the proof, we may neglect the algebraic grading, 
and we only need to prove that 
\begin{equation} \label{eqn:4570}
\opn{RHom}_{A}(\K, A) \cong \K[-n]
\end{equation}
in $\dcat{D}(\K)$.

Let $\til{B}$ be the DG ring from above. Because $\K$ is a DG $A$-module via 
the DG ring homomorphism $A \to \til{B}$, there is an adjunction isomorphism
\begin{equation} \label{eqn:4566} 
\opn{RHom}_{A}(\K, A) \cong 
\opn{RHom}_{\til{B}} \bigl( \K, \opn{RHom}_{A}(\til{B}, A) \bigr) 
\end{equation}
in $\dcat{D}(\K)$. 
Now $\til{B}$ is a semi-free DG module over $A$, so 
\[ \opn{RHom}_{A}(\til{B}, A) \cong 
\opn{Hom}_{A}(\til{B}, A) \cong 
\bigl( A \xar{(-) \cd a} A \bigr) \cong B[-1] \]
in $\dcat{D}(\til{B})$. Substituting this into (\ref{eqn:4566}), and using 
Theorem \ref{thm:2363}(2), we get 
\[ \opn{RHom}_{A}(\K, A) \cong 
\opn{RHom}_{\til{B}} \bigl( \K, B[-1] \bigr) 
\cong \opn{RHom}_{B}(\K, B)[-1] \]
in $\dcat{D}(\K)$.
The AS condition for $B$ says that 
$\opn{RHom}_{B}(\K, B) \cong \K[-n + 1]$ 
in $\dcat{D}(\K)$. Hence the isomorphism (\ref{eqn:4570}) holds. 
\end{proof}

\begin{rem} \label{rem:4560}
The proof of the theorem above can actually yield more, with some 
modifications. If we were to make the DG ring $\til{B}$ 
into a bigraded object, with both algebraic and cohomological grading, then we 
would obtain an isomorphism (\ref{eqn:4560}) in the derived category 
$\dcat{D}(\til{B}, \mrm{gr})$ of bigraded DG $\til{B}$-modules. 
The problem is that we did not develop this kind of theory in the book...

The outcome of this more refined construction would be that formula 
(\ref{eqn:4551}) is an isomorphism in $\dcat{M}(\K, \mrm{gr})$. 
This would give the ability to compare the AS indexes of $A$ and $B$. As can 
be guessed (cf.\ the easy case of Example \ref{exa:3990}, taking $a$ to be
one of the variables), if the AS index of $A$ is $l$, then the AS index of $B$ 
is $l - \opn{deg}(a)$. 
\end{rem}

\begin{thm} \label{thm:4562}
Assume Setup \tup{\ref{set:4565}}.
\begin{enumerate}
\item If $A$ is a Gorenstein graded ring of graded injective dimension 
$n$, then $B$ is a Gorenstein graded ring of graded injective dimension at most
$n - 1$. 

\item Moreover, if $A$ is an AS Gorenstein graded ring of dimension $n$, 
then $B$ is an AS Gorenstein graded ring of dimension $n - 1$.
\end{enumerate}
\end{thm}

\begin{exer} \label{exer:4560}
Prove Theorem \ref{thm:4562}. (Hint: Study the proof of Theorem 
\ref{thm:4546}.)
\end{exer}

\begin{rem} \label{rem:4653}
Here are some historical notes on the material in this subsection. 
Around 1985, M. Artin and W. Schelter \cite{ArSchl} began a systematic study of 
connected graded noncommutative rings (better known as graded algebras), with 
the hope of obtaining a classification in low dimensions, in analogy to the 
commutative projective geometry of curves and surfaces. This is how the concept 
of {\em Artin-Schelter regular graded rings} was born. 
Closely related concepts are {\em Sklyanin algebras} and {\em Koszul algebras}. 

The area of research initiated by Artin and Schelter is often called 
{\em NC algebraic geometry}. Indeed, in 1994 Artin and J.J. Zhang \cite{ArZh} 
introduced the concept of {\em NC projective scheme}. Given a noetherian 
connected graded ring $A$, let 
$\dcat{M}_{\mrm{tor}}(A, \mrm{gr}) \sub \dcat{M}(A, \mrm{gr})$
be the full subcategory on the $\m$-torsion modules. The quotient abelian 
category $\dcat{M}(A, \mrm{gr}) / \dcat{M}_{\mrm{tor}}(A, \mrm{gr})$
is viewed as the category of quasi-coherent sheaves on an imaginary scheme  
$\opn{Proj}(A)$. Note that if $A$ happens to be commutative, then, by the 
famous theorem of J.P. Serre, this quotient category is canonically equivalent 
to $\cat{QCoh} X$, where $X = \opn{Proj}(A)$ is the genuine projective scheme. 
The geometric metaphor of Artin and Zhang turned out to be very useful in 
predicting ``geometric properties'' of connected graded rings, that were later 
proved algebraically, working in $\dcat{M}(A, \mrm{gr})$ or its derived 
category $\dcat{D}(A, \mrm{gr})$. 

For AS regular graded rings of dimension $\leq 3$ generated in degree $1$, the 
classification was pretty much complete by 1991, in papers by Artin, jointly 
with J. Tate and M. Van den Bergh; see the papers \cite{ATV} and 
\cite{ATV2}, and the survey \cite{StaVdB}. Research on AS graded rings of 
higher dimensions, and with generators in various degrees, was later conducted 
by J.T. Stafford, S.P. Smith, Zhang, and their students and collaborators. 

In 2006, V. Ginzburg published the highly influential preprint \cite{Ginz} on 
CY (Calabi-Yau) algebras. Since then, the research on AS regular rings and 
their relatives has merged with the research on CY rings and categories. 
(It turns out that for connected graded rings, the twisted CY rings are 
precisely the AS regular rings.) See Remark \ref{rem:4505} for more on CY 
theory. Two of the latest papers on graded CY rings are 
\cite{ReRoZh} and \cite{ReRo}. 

In the subsequent sections of the book we shall deal with a few ideas that 
emerged as by-products of the Artin school of NC algebraic geometry: the $\chi$ 
condition, balanced dualizing complexes and rigid dualizing complexes.
\end{rem}








\cleardoublepage
\mysection{Derived Torsion over NC Graded Rings}
\label{sec:der-tors-conn}

\AYcopyright

As already mentioned, connected graded noncommutative rings are similar 
in several aspects to complete local commutative rings. 
In this section we concentrate on the derived $\m$-torsion functors 
associated to a connected graded ring $A$ with augmentation ideal $\m$. 

Some of the definitions and results here are quite old, going back to papers 
from the 1990's. Other definitions and results are adaptations to the 
connected graded setting of the work in the recent paper \cite{VyYe}. 

The {\em $\chi$ condition} on a noncommutative noetherian connected graded ring 
$A$ was introduced by M. Artin and J.J. Zhang \cite{ArZh} to guarantee that 
the {\em noncommutative projective scheme} $\opn{Proj}(A)$ will have good 
``geometric properties'', resembling the classical commutative case. In this 
section and the subsequent one, we will focus on another property that the 
$\chi$ condition implies: the {\em symmetry of derived torsion}; namely that 
from the point of view of derived torsion, the ring $A$ looks commutative. 
It is important to say that this phenomenon is {\em only seen in the derived 
category of graded bimodules} -- on the elementary level the ring $A$ is often 
terribly noncommutative (cf.\ Examples \ref{exa:4080} and \ref{exa:4083}), 
and the bimodules are far from being central!

In this section we follow Conventions \ref{conv:3700} and \ref{conv:4560}.
So $\K$ is a base field, and all rings are algebraically graded central 
$\K$-rings. The definitions and results from Section \ref{sec:alg-gra-rings} 
will be used here freely. See Remark \ref{rem:3740} regarding the possibility of 
replacing the base field $\K$ with a base commutative ring. 



\mysubsection{Quasi-Compact Finite Dimensional Functors} 
\label{subsec:QC-FD-funs}

The content of this subsection is adapted from \cite[Section 1]{VyYe}. 
We state the results in the context of algebraically graded rings and modules, 
since here we need them for the graded $\m$-torsion functor 
$F := \Ga_{\m}$; yet these results hold in greater generality (see 
\cite{VyYe}). 

Recall from Subsection \ref{subsec:alg-gr-mods} that given a graded ring $A$, 
the category of graded $A$-modules is $\dcat{M}(A, \mrm{gr})$. 
This is a $\K$-linear abelian category. The category of complexes in 
$\dcat{M}(A, \mrm{gr})$ is $\dcat{C}(A, \mrm{gr})$, 
the homotopy category is $\dcat{K}(A, \mrm{gr})$, and the derived category is 
$\dcat{D}(A, \mrm{gr})$. 

In this subsection we consider the following setup:

\begin{setup} \label{set:4075}
$A$ and $B$ are graded rings, and 
$F : \dcat{M}(A, \mrm{gr}) \to \dcat{M}(B, \mrm{gr})$
is a left exact linear functor. 
\end{setup}

The functor $F$ extends to a triangulated functor
$F : \dcat{K}(A, \mrm{gr}) \to \dcat{K}(B, \mrm{gr})$
between the homotopy categories. Because there are enough K-injective 
resolutions, $F$ has a right derived functor 
$(\mrm{R} F, \eta^{\mrm{R}})$. Let us recall what this means: 
$\mrm{R} F : \dcat{D}(A, \mrm{gr}) \to \dcat{D}(B, \mrm{gr})$
is a triangulated functor, and 
$\eta^{\mrm{R}} : \opn{Q} \circ \, F \to \mrm{R} F \circ \opn{Q}$
is a morphism of triangulated functors 
$\dcat{K}(A, \mrm{gr}) \to \dcat{D}(B, \mrm{gr})$
that has a certain universal property (see Definition \ref{dfn:131}). 
The right derived functor $(\mrm{R} F, \eta^{\mrm{R}})$
can be constructed using a K-injective presentation: for each complex $M$ 
we choose a K-injective resolution 
$\rho_M : M \to I_M$
in $\dcat{C}_{\mrm{str}}(A, \mrm{gr})$, and then we take 
$\mrm{R} F(M) := F(I_M)$ and $\eta^{\mrm{R}}_M := F(\rho_M)$. 

Recall that the functor $F$ has the classical right derived functors  
$\mrm{R}^q F : \lb \dcat{M}(A, \mrm{gr}) \to \dcat{M}(B, \mrm{gr})$.
These can be expressed as follows: 
$\mrm{R}^q F(M) = \opn{H}^q(\mrm{R} F(M))$ 
for a module $M \in \dcat{M}(A, \mrm{gr})$. 
Since $F$ is left exact, the canonical homomorphism 
$F \to \mrm{R}^0 F$ is an isomorphism. 

Here is a classical definition:

\begin{dfn} \label{dfn:3885}
The {\em right cohomological dimension of $F$} 
\index{Cohomological dimension! right {\indash} of functor}
is 
\[ d := \opn{sup} \, \{ q \in \N \mid \mrm{R}^q F \neq 0 \} \in 
\N \cup \{ \pm \infty \} . \]
\end{dfn}

If $F \neq 0$ then $d \in \N \cup \{ \infty \}$, but for $F = 0$ the 
dimension is $d = -\infty$. The generalized integer $d$ is said to be finite 
if $d < \infty$; this terminology is designed so that the zero functor will 
have finite right cohomological dimension. 

Here is another standard definition. 

\begin{dfn} \label{dfn:3887}
A module $I \in \dcat{M}(A, \mrm{gr})$ is called a {\em right $F$-acyclic 
object} if $\mrm{R}^q F(I) = 0$ for all $q > 0$. 
\end{dfn}

Of course every injective object of $\dcat{M}(A, \mrm{gr})$ is right 
$F$-acyclic. But often there are many more right $F$-acyclic objects.

\begin{dfn} \label{dfn:3890}
A complex $I \in \dcat{C}(A, \mrm{gr})$ is called a {\em right $F$-acyclic 
complex} if the morphism $\eta^{\mrm{R}}_I : F(I) \to \mrm{R} F(I)$
in $\dcat{D}(B, \mrm{gr})$ is an isomorphism. 
\end{dfn}

Of course a K-injective complex is right $F$-acyclic, but often there 
are many more.
It is easy to see that a module 
$I \in \dcat{M}(A, \mrm{gr})$ is a right $F$-acyclic object
in the sense of Definition \ref{dfn:3887} if and only if it is right 
$F$-acyclic 
as a complex, in the sense of Definition \ref{dfn:3890}. 

\begin{lem} \label{lem:3885}
Let $I \in \dcat{C}(A, \mrm{gr})$ be a complex such that each of the modules  
$I^q$ is right $F$-acyclic. If $I$ is a bounded below complex, or if $F$ has 
finite right cohomological dimension, then $I$ is a right $F$-acyclic 
complex. 
\end{lem}

\begin{proof}
This is contained in the proof of \cite[Corollary I.5.3]{RD}, but for the sake 
of completeness we give a proof here. 

We can assume that neither $F$ nor $I$ are zero. 
First let us assume that $I$ is bounded below.
We can find a quasi-isomorphism $I \to J$ in 
$\dcat{C}_{\mrm{str}}(A, \mrm{gr})$, where $J$ is a bounded below complex 
of injective objects of $\dcat{M}(A, \mrm{gr})$, and thus it is a K-injective 
complex. We must show that 
$F(I) \to F(J)$ is a quasi-isomorphism. 
This amounts to showing that the complex $F(K)$ is acyclic, where $K$ is the 
standard cone on the quasi-isomorphism $I \to J$. Thus we reduce the problem to 
proving that for an acyclic bounded below complex $I$ made up of right 
$F$-acyclic objects, the complex $F(I)$ is acyclic. 

Say $\opn{inf}(I) = d_0 \in \Z$. For any $q$ let 
$\opn{Z}^q(I) := \opn{Ker}(I^q \to I^{q + 1})$, the object of degree 
$q$ cocycles. The acyclicity of the complex $I$ says that we have 
short exact sequences
\[ 0 \to \opn{Z}^{q}(I) \to I^{q} \to \opn{Z}^{q + 1}(I) \to 0 . \]
By induction on $q \geq d_0 - 1$, using the long exact cohomology sequence, we 
prove that each $\opn{Z}^{q}(I)$ is a right $F$-acyclic object.
Therefore the sequences 
\[ 0 \to F(\opn{Z}^{q}(I)) \to F(I^{q}) \to F(\opn{Z}^{q + 1}(I)) \to 0  \]
are exact for all $q$. Splicing together three sequences like this, for  
$q - 1$, $q$ and $q + 1$, shows that the sequence 
$F(I^{q - 1}) \to F(I^q) \to F(I^{q + 1})$
is exact. Thus the complex $F(I)$ is acyclic. 

Now $I$ is no longer bounded below, but the  right cohomological dimension
of $F$ is finite, say $d \in \N$. 
Let $I \to J$ be a quasi-isomorphism in 
$\dcat{C}_{\mrm{str}}(A, \mrm{gr})$, where 
$J$ is a K-injective complex made up of injective objects
(see Corollary \ref{cor:3830}(3)). We must show that 
$F(I) \to F(J)$ is a quasi-isomorphism. As above, by replacing $I$ with the 
standard cone on $I \to J$, we reduce the problem to proving 
that for an acyclic complex $I$ made up of right $F$-acyclic objects, the 
complex $F(I)$ is acyclic. 

Fix an an integer $q$. Let $M := \opn{Z}^q(I)$, and let 
\[ J := (\cdots \to 0 \to I^q \to  I^{q + 1} \to \cdots) , \]
the complex with $I^q$ placed in degree $0$ (so $J$ is a shift of a stupid 
truncation of $I$). By the previous part of the proof 
we know that $\mrm{R} F(J) \cong F(J)$ in $\dcat{D}(B, \mrm{gr})$. 
We have a quasi-isomorphism $M \to J$. Hence, for every $p > d$, 
by our assumption on $F$ we have
\[ 0 = \opn{H}^p(\mrm{R} F(M)) \cong \opn{H}^p(\mrm{R} F(J)) \cong
\opn{H}^p(F(J)) \cong \opn{H}^{p + q}(F(I))  . \]
Since $q$ was chosen arbitrarily, we conclude that $F(I)$ is acyclic.
\end{proof}

Recall the 
{\em cohomological dimension}% 
\index{Cohomological dimension! of triangulated functor}%
\index{Additive functor! finite dimensional}
of the triangulated functor 
$\mrm{R} F : \dcat{D}(A, \mrm{gr}) \to \dcat{D}(B, \mrm{gr})$
from Definition \ref{dfn:2126}.
In what follows we use the various notions introduced in Subsection
\ref{subsec:way-out}, 
adapted to the graded setting, such as the concentration interval 
$\opn{con}(M) \sub \Z$ of an 
object $M \in \dcat{G}(A, \mrm{gr})$. 

\begin{lem} \label{lem:3891}
The cohomological dimension of the right derived functor $\mrm{R} F$
equals the right cohomological dimension of $F$. 
\end{lem}

\begin{proof}
As explained in Example \ref{exa:2135}, the cohomological dimension of 
$(\mrm{R} F)|_{\dcat{M}(A)}$ equals the right cohomological dimension of the 
functor $F$. We need to prove the reverse inequality. 

We can assume that $F \neq 0$, and it has finite right cohomological 
dimension, say $d \in \N$. We shall prove that the cohomological displacement 
of the derived functor $\mrm{R} F$ is contained in the interval
$[d_0, d_1] := [0, d]$. 
To be explicit, we shall prove that for every complex
$M \in \dcat{D}(A, \mrm{gr})$
the integer interval 
$\opn{con} \bigl( \opn{H}(\mrm{R} F(M)) \bigr)$ is contained in the interval 
$\opn{con} \bigl( \opn{H}(M) \bigr) + [0, d]$. This is done by cases, and we 
can assume that $M \neq 0$. 

\medskip \noindent 
Case 1. If $M$ is not in $\dcat{D}^-(A, \mrm{gr})$, 
i.e.\ $\opn{con}(\opn{H}(M)) = [d_0, \infty]$ for some 
$d_0 \in \Z \cup \{ -\infty \}$, then we can take a K-injective resolution 
$M \to I$ such that $\opn{inf}(I) = d_0$. Then 
$\opn{H}(\mrm{R} F(M)) = \opn{H}(F(I))$
is concentrated in $[d_0, \infty]$. 

\medskip \noindent 
Case 2. Here $M \in \dcat{D}^-(A, \mrm{gr})$, so that 
$\opn{con}(\opn{H}(M)) = [d_0, d_1]$ for some 
$d_0 \in \Z \cup \{ -\infty \}$ and $d_1 \in \Z$. 
We now take a  K-injective resolution $M \to I$ such that $\opn{inf}(I) = d_0$
and $I$ is made up of injective objects. We must prove that 
$\opn{H}^q(F(I)) = 0$ for all $q > d_1 + d$. This is done like in the 
proof of Lemma \ref{lem:3885}. 
Let $N := \opn{Z}^{d_1}(I)$, and let 
\[ J := (\cdots \to 0 \to I^{d_1} \to  I^{d_1 + 1} \to \cdots) , \]
the complex with $I^{d_1}$ placed in degree $0$. 
So there is a quasi-isomorphism $N \to J$, and it is an injective resolution of 
$N$. Hence for every $p > d$ we have
\[ 0 = \opn{H}^p(\mrm{R} F(N)) \cong \opn{H}^p(F(J)) \cong 
\opn{H}^{d_1 + p}(F(I))  . \qedhere \]
\end{proof}

\begin{dfn} \label{dfn:3888}
Let $\cat{M}$ and $\cat{N}$ be linear categories with arbitrarily direct sums. 
A linear functor 
$G : \cat{M} \to \cat{N}$ is called {\em quasi-compact} 
\index{Additive functor! quasi-compact}
if it commutes with infinite direct sums. Namely, for every collection 
$\{ M_x \}_{x \in X}$ of objects of $\cat{M}$, the canonical morphism
\[ \bigoplus\nolimits_{x \in X} G(M_x) \to 
G \left( \bigoplus\nolimits_{x \in X} M_x \right) \]
in $\cat{N}$ is an isomorphism. 
\end{dfn}

The name ``quasi-compact functor'' is inspired by the property of pushforward 
of quasi-coherent sheaves along a quasi-compact map of schemes. 

\begin{lem} \label{lem:3890}
Assume that the functors $\mrm{R}^q F$ are quasi-compact, for all $q > 0$. Let 
$\{ I_x \}_{x \in X}$ be a collection of right $F$-acyclic objects of 
$\dcat{M}(A, \mrm{gr})$.
Then the object 
$I := \bigoplus_{x \in X} I_x \in \dcat{M}(A, \mrm{gr})$
is right $F$-acyclic.
\end{lem}

\begin{proof}
Take some $q > 0$. Because $\mrm{R}^qF$ is quasi-compact, the canonical 
homomorphism 
$\bigoplus_{x \in X} \mrm{R}^q F(I_x) \to \mrm{R}^q F(I)$
in $\cat{N}$ is an isomorphism. But by assumption, $\mrm{R}^q F(I_x) = 0$ 
for all $x$. 
\end{proof}

\begin{lem} \label{lem:1058}
Assume the functors $\mrm{R}^q F$ are quasi-compact for all $q \geq 0$. Let 
$\{ I_x \}_{x \in X}$ be a collection of right $F$-acyclic 
complexes in $\dcat{C}(A, \mrm{gr})$, and define 
$I := \bigoplus_{x \in X} I_x \in \dcat{C}(A, \mrm{gr})$.
If $I$ is a bounded below complex, or if $F$ has finite right cohomological 
dimension, then $I$ is a right $F$-acyclic complex.
\end{lem}

\begin{proof}
For every index $x$ let us choose a quasi-isomorphism 
$\phi_x : I_x \to J_x$, where $J_x$ is a K-injective complex consisting of 
injective objects, and 
$\opn{inf}(J_x) \geq \opn{inf}(I_x)$. Define 
$J := \bigoplus_{x \in X} J_x$.
We get a quasi-isomorphism 
$\phi := \bigoplus_{x \in X} \phi_x : I \to J$
in $\dcat{C}_{\mrm{str}}(A, \mrm{gr})$. 
By construction, if $I$ is a bounded below complex, then so is $J$. 

For every $x \in X$ there is a commutative diagram 
\[ \UseTips \xymatrix @C=9ex @R=6ex {
F(I_x)
\ar[r]^{F(\phi_x)}
\ar[d]_{\eta^{\mrm{R}}_{I_x}}
& 
F(J_x)
\ar[d]^{\eta^{\mrm{R}}_{J_x}}
\\
\mrm{R} F(I_x)
\ar[r]^(0.5){\mrm{R} F(\phi_x)}
& 
\mrm{R} F(J_x)
} \]
in $\dcat{D}(B, \mrm{gr})$. The vertical arrows are isomorphisms because both 
$I_x$ and $J_x$ are right $F$-acyclic complexes. The morphism $\mrm{R} 
F(\phi_x)$ is 
also an isomorphism. It follows that $F(\phi_x)$ is an isomorphism in 
$\dcat{D}(B, \mrm{gr})$; and therefore it is a quasi-isomorphism in 
$\dcat{C}_{\mrm{str}}(B, \mrm{gr})$. 

Next consider this commutative diagram in 
$\dcat{C}_{\mrm{str}}(B, \mrm{gr})$~:
\[ \UseTips \xymatrix @C=14ex @R=6ex {
\bigoplus_{x \in X} \, F(I_x)
\ar[r]^{ \bigoplus_{x \in X} \, F(\phi_x) }
\ar[d]_{}
& 
\bigoplus_{x \in X} \, F(J_x)
\ar[d]
\\
F(I)
\ar[r]^{ F(\phi) }
& 
F(J)
} \]
Because each $F(\phi_x)$ is a quasi-isomorphism, it follows that 
the top horizontal arrow is a quasi-isomorphism.
Because the functor $F = \mrm{R}^0 F$ is quasi-compact, the vertical arrows 
are isomorphisms. We conclude that $F(\phi)$ is a quasi-isomorphism in 
$\dcat{C}_{\mrm{str}}(B, \mrm{gr})$.

Finally we look at this  commutative diagram in $\dcat{D}(B, \mrm{gr})$~:
\[ \UseTips \xymatrix @C=8ex @R=6ex {
F(I)
\ar[r]^{F(\phi)}
\ar[d]_{\eta^{\mrm{R}}_I}
& 
F(J)
\ar[d]^{\eta^{\mrm{R}}_J}
\\
\mrm{R} F(I)
\ar[r]^{ \mrm{R} F(\phi) }
& 
\mrm{R} F(J)
} \]
We know that the morphisms $F(\phi)$ and $\mrm{R} F(\phi)$ are isomorphisms. 
For every $p$, the object  
$J^p = \bigoplus_{x \in X} \, J_x^p$ 
in $\dcat{M}(A, \mrm{gr})$ is a direct sum of injective objects. So according 
to 
Lemma \ref{lem:3890}, $J^p$ is a right $F$-acyclic object.
By Lemma \ref{lem:3885}, in either of the two cases, the complex $J$ is a 
right $F$-acyclic complex.
This says that the morphism $\eta^{\mrm{R}}_J$ is an isomorphism.
We conclude that the morphism $\eta^{\mrm{R}}_I$ is an isomorphism too, and 
this says that $I$ is a right $F$-acyclic complex.
\end{proof}

\begin{thm} \label{thm:1001}
Under Setup \tup{\ref{set:4075}}, assume the functor $F$
has finite right cohomological dimension, and the functors $\mrm{R}^q F$, for 
all $q \geq 0$, are quasi-compact%
\index{Additive functor! quasi-compact}%
\index{Additive functor! finite dimensional}.
Then the triangulated functor 
$\mrm{R} F : \dcat{D}(A, \mrm{gr}) \to \dcat{D}(B, \mrm{gr})$
has finite cohomological dimensional and is quasi-compact.
\end{thm}

\begin{proof}
By Lemma \ref{lem:3891} the functor $\mrm{R} F$ is finite 
dimensional. We need to  prove that $\mrm{R} F$ commutes with 
infinite direct sums. Namely, consider a collection \lb $\{ M_x \}_{x \in X}$ 
of complexes in $\dcat{C}(A, \mrm{gr})$, and let 
$M := \bigoplus_{x \in X} M_x$. 
We have to prove that the canonical morphism 
$\bigoplus_{x \in X} \mrm{R} F(M_x) \to  \mrm{R} F(M)$
in $\dcat{D}(B, \mrm{gr})$ is an isomorphism. 

For every index $x$ we choose a quasi-isomorphism 
$\phi_x : M_x \to I_x$ in $\dcat{C}_{\mrm{str}}(A, \mrm{gr})$,
where $I_x$ is a K-injective complex. 
Define $I := \bigoplus_{x \in X} I_x$, so there is a quasi-isomorphism 
$\phi : M \to I$ in $\dcat{C}_{\mrm{str}}(A, \mrm{gr})$. We get a commutative 
diagram 
\[ \UseTips \xymatrix @C=14ex @R=6ex {
\bigoplus_{x \in X} \mrm{R} F(M_x)
\ar[r]^{ \bigoplus \mrm{R} F(\phi_x) } 
\ar[d]_{}
& 
\bigoplus_{x \in X} \mrm{R} F(I_x)
\ar[d]
\\
\mrm{R} F(M)
\ar[r]^{ \mrm{R} F(\phi) }
& 
\mrm{R} F(I)
} \]
in $\dcat{D}(B, \mrm{gr})$, in which the horizontal arrows are isomorphisms. 
Therefore it suffices to prove that the canonical morphism 
\begin{equation} \label{eqn:1060}
\bigoplus\nolimits_{x \in X} \mrm{R} F(I_x) \to \mrm{R} F(I) 
\end{equation}
in $\dcat{D}(B, \mrm{gr})$ is an isomorphism. 

Consider this commutative diagram in $\dcat{D}(B, \mrm{gr})$~:
\[ \UseTips \xymatrix @C=12ex @R=6ex {
\bigoplus_{x \in X} \, F(I_x)
\ar[r]^{ \bigoplus \eta^{\mrm{R}}_{I_x} }
\ar[d]_{}
& 
\bigoplus_{x \in X} \, \mrm{R} F(I_x)
\ar[d]
\\
F(I)
\ar[r]^{ \eta^{\mrm{R}}_I }
& 
\mrm{R} F(I)
} \]
For each $x$ the morphism $\eta^{\mrm{R}}_{I_x}$ is an isomorphism; and hence 
the top horizontal arrow is an isomorphism. 
Because $F = \mrm{R}^0 F$ is quasi-compact, the left vertical arrow is an 
isomorphism (in $\dcat{C}_{\mrm{str}}(B, \mrm{gr})$, and so also in 
$\dcat{D}(B, \mrm{gr})$). By Lemma \ref{lem:1058} the complex $I$ is right 
$F$-acyclic, and therefore $\eta^{\mrm{R}}_I$ is 
an isomorphism. Hence the remaining arrow is an isomorphism; but this is the 
morphism (\ref{eqn:1060}).
\end{proof}



\mysubsection{Weakly Stable and Idempotent Copointed Functors}
\label{subsec:idem-cop-fun}

The content of this subsection is adapted from \cite[Section 2]{VyYe}. 
Recall that we follow Conventions \ref{conv:3700} and \ref{conv:4560}. Like 
Subsection \ref{subsec:QC-FD-funs}, the results of this subsection apply in much 
greater generality. We state the results in the special context of 
algebraically graded rings and modules, since we need them for the graded 
$\m$-torsion functor $F := \Ga_{\m}$.

\begin{dfn} \label{dfn:3895}
Let $A$ be a graded ring and 
$F : \dcat{M}(A, \mrm{gr}) \to \dcat{M}(A, \mrm{gr})$
a left exact linear functor. 
\begin{enumerate}
\item The functor $F$ is called {\em stable} 
\index{Additive functor! stable}
if for every injective object $I \in \dcat{M}(A, \mrm{gr})$, the object $F(I)$ 
is injective.  

\item The functor $F$ is called {\em weakly stable} 
\index{Additive functor! weakly stable}
if for every injective object $I \in \dcat{M}(A, \mrm{gr})$, the object $F(I)$ 
is right $F$-acyclic, in the sense of Definition \ref{dfn:3887}.
\end{enumerate}
\end{dfn}

Clearly stable implies weakly stable. 

\begin{rem} \label{rem:3895}
The name ``stable'' comes from the 
theory of {\em torsion classes}; see \cite[Section VI.7]{Ste}.
Indeed, a torsion class is called stable if the corresponding torsion functor 
$F$ is a stable functor, as defined above. 

The name ``weakly stable'' was coined in \cite{VyYe}, where it was shown that 
for a finitely generated ideal $\a$ in a commutative ring $A$, the torsion 
functor $\Ga_{\a}$ is weakly stable iff the ideal $\a$ is {\em weakly 
proregular}. We should mention the well-known fact that if the commutative 
ring $A$ is noetherian, then for every ideal $\a \sub A$ the torsion functor 
$\Ga_{\a}$ is stable (cf.\ \cite[Lemma III.3.2]{Har}; or use the Matlis 
classification of injective $A$-modules in Subsection \ref{subsec:mr-inj-res}).
\end{rem}

\begin{dfn} \label{dfn:3896} 
Let $\cat{M}$ be a linear category (e.g.\ $\dcat{M}(A, \mrm{gr})$ or
$\dcat{D}(A, \mrm{gr})$ for a graded ring $A$), with identity functor 
$\opn{Id}_{\cat{M}}$.
\begin{enumerate}
\item A {\em copointed linear functor}%
\index{Additive functor! copointed}
on $\cat{M}$ is a pair 
$(F, \si)$, consisting of a linear functor 
$F : \cat{M} \to \cat{M}$ and a morphism of functors 
$\si : F \to \opn{Id}_{\cat{M}}$.

\item The copointed linear functor $(F, \si)$ is called 
{\em idempotent}% 
\index{Additive functor! idempotent}
if the morphisms 
\[ \si_{F(M)}, \, F(\si_M) : F(F(M)) \to F(M) \]
are isomorphisms for all objects $M \in \cat{M}$. 

\item If $\cat{M}$ is a triangulated category, $F$ is a triangulated functor, 
and $\si$ is a morphism of triangulated functors, then we call 
$(F, \si)$ a {\em copointed triangulated functor}. 
\end{enumerate}
\end{dfn}

The name ``copointed'' is explained in Remark \ref{rem:3896}. 

Weak stability and idempotence together have the following effect. 

\begin{lem} \label{lem:3895}
Let $(F, \si)$ be an idempotent copointed linear functor on 
$\dcat{M}(A, \mrm{gr})$, and assume that $F$ is left exact and weakly stable. 
If $I$ is a right $F$-acyclic object of $\dcat{M}(A, \mrm{gr})$, then $F(I)$ is 
also a right $F$-acyclic object of $\dcat{M}(A, \mrm{gr})$.
\end{lem}

\begin{proof}
Choose an injective resolution $\rho : I \to J$; i.e.\ $J$ is a complex of 
injectives concentrated in nonnegative degrees, and $\rho$ is a 
quasi-isomorphism in $\dcat{C}_{\mrm{str}}(A, \mrm{gr})$. 
Since $\opn{H}^q(F(J)) \cong \mrm{R}^q F(I)$, and since $I$ is a right 
$F$-acyclic object, we see that the homomorphism of complexes
$F(\rho) : F(I) \to F(J)$ is a quasi-isomorphism. 
Therefore both $F(\rho)$ and $\mrm{R} F(F(\rho))$ are isomorphisms in 
$\dcat{D}(A, \mrm{gr})$. 

The weak stability of $F$ implies that $F(J)$ is a bounded below complex of 
right $F$-acyclic objects. According to Lemma \ref{lem:3885} the 
complex $F(J)$ is right $F$-acyclic. This means that the morphism 
$\eta^{\mrm{R}}_{F(J)} : F(J) \to \mrm{R} F (J)$
is an isomorphism. 

The idempotence of the copointed functor $F$ says that the morphisms 
$\si_{F(I)}$ and $\si_{F(J)}$ are isomorphisms in 
$\dcat{C}_{\mrm{str}}(A, \mrm{gr})$. 
We get a commutative diagram in $\dcat{D}(A, \mrm{gr})$~: 
\[ \UseTips \xymatrix @C=10ex @R=6ex {
F(I)
\ar[d]_{F(\rho)}^{\cong}
&
F(F(I))
\ar[l]_{\si_{F(I)}}^{\cong}
\ar[r]^{\eta^{\mrm{R}}_{F(I)}}
\ar[d]_{F(F(\rho))}
&
\mrm{R} F(F(I))
\ar[d]_{\mrm{R} F(F(\rho))}^{\cong}
\\
F(J)
&
F(F(J))
\ar[l]_{\si_{F(J)}}^{\cong}
\ar[r]^{\eta^{\mrm{R}}_{F(J)}}_{\cong}
&
\mrm{R} F(F(J))
} \]
We conclude that $\eta^{\mrm{R}}_{F(I)}$ is an isomorphism; and this means that 
the object $F(I)$ is right $F$-acyclic. 
\end{proof}

The next lemma is a generalization of \cite[Proposition 3.10]{PSY}.

\begin{lem} \label{lem:1063}
Suppose we are given a copointed linear functor $(F, \si)$ on \lb
$\dcat{M}(A, \mrm{gr})$. 
Then there is a unique morphism 
$\si^{\mrm{R}} : \mrm{R} F \to \opn{Id}_{\dcat{D}(A, \mrm{gr})}$
of triangulated functors from $\dcat{D}(A, \mrm{gr})$ to itself, satisfying 
this condition\tup{:} for every complex 
$M \in \dcat{D}(A, \mrm{gr})$ there is equality 
$\si^{\mrm{R}}_M \circ \eta^{\mrm{R}}_M = \si_M$
of morphisms $F(M) \to M$ in $\dcat{D}(A, \mrm{gr})$. 
\end{lem}

In a commutative diagram:
\begin{equation} \label{eqn:1455}
\UseTips \xymatrix @C=6ex @R=6ex {
F(M)
\ar[r]^{ \eta^{\mrm{R}}_M }
\ar[dr]_{ \si_M }
&
\mrm{R} F(M)
\ar[d]^{ \si^{\mrm{R}}_M }
\\
&
M
} 
\end{equation}

\begin{proof}
The existence of the morphism $\si^{\mrm{R}}$ comes for free from the universal 
property of the right derived functor. Still, for later reference, we give the 
construction. 

For a K-injective complex $I$ the morphism 
$\eta^{\mrm{R}}_I : F(I) \to \mrm{R} F(I)$ in $\dcat{D}(A, \mrm{gr})$ is an 
isomorphism, and we define 
$\si^{\mrm{R}}_I : \mrm{R} F(I) \to I$
to be 
$\si^{\mrm{R}}_I := \si_I \circ (\eta^{\mrm{R}}_I)^{-1}$.

For an arbitrary complex $M$ we choose a quasi-isomorphism 
$\rho : M \to I$ into a K-injective complex, and then we let 
$\si^{\mrm{R}}_M := \rho^{-1} \circ \si^{\mrm{R}}_I \circ \mrm{R} F(\rho)$
in $\dcat{D}(A, \mrm{gr})$. The corresponding commutative diagram in 
$\dcat{D}(A, \mrm{gr})$ is (\ref{eqn:4830}). 
%\begin{figure}[!htb]
\begin{equation} \label{eqn:4830}
\UseTips \xymatrix @C=10ex @R=6ex {
F(M)
\ar[r]^{\eta^{\mrm{R}}_M}
\ar[d]_{F(\rho)}
\ar@(ur,ul)[rr]^{\si_M}
&
\mrm{R} F(M)
\ar[d]_{\mrm{R} F(\rho)}^{\cong}
\ar[r]^{\si^{\mrm{R}}_M}
&
M
\ar[d]_{\rho}^{\cong}
\\
F(I)
\ar[r]^{\eta^{\mrm{R}}_I}_{\cong}
\ar@(dr,dl)[rr]_{\si_I}
&
\mrm{R} F(I)
\ar[r]^{\si^{\mrm{R}}_I}
&
I
} 
\end{equation}
%\end{figure}
It is easy to see that the collection of morphisms 
$\{ \si^{\mrm{R}}_M \}_{M \in \dcat{D}(A, \mrm{gr})}$ 
has the desired properties. 
\end{proof}

In this way we obtain a copointed triangulated functor 
$(\mrm{R} F, \si^{\mrm{R}})$ on $\dcat{D}(A, \mrm{gr})$.

\begin{thm} \label{thm:1004}
Let $(F, \si)$ be an  idempotent copointed linear functor on \lb
$\dcat{M}(A, \mrm{gr})$, and 
assume that $F$ is left exact and weakly stable.
\index{Additive functor! idempotent}
\index{Additive functor! copointed}
\index{Additive functor! weakly stable}
\begin{enumerate}
\item The copointed triangulated 
functor $(\mrm{R} F, \si^{\mrm{R}})$ on 
$\dcat{D}^+(A, \mrm{gr})$
is idempotent. 

\item If $F$ has finite right cohomological dimension, then the copointed 
triangulated functor $(\mrm{R} F, \si^{\mrm{R}})$ on 
$\dcat{D}(A, \mrm{gr})$
is idempotent. 
\end{enumerate}
\end{thm}

\begin{proof}
Let $M \in \dcat{D}(A, \mrm{gr})$. Choose a K-injective resolution 
$M \to I$ in $\dcat{C}_{\mrm{str}}(A, \mrm{gr})$, such that 
$I$ is a complex consisting of injective objects of $\dcat{M}(A, \mrm{gr})$, 
and $\opn{inf}(I) = \opn{inf}(\opn{H}(M))$. 
It suffices to prove that the morphisms 
\[ \si^{\mrm{R}}_{\mrm{R} F(I)}, \, \mrm{R} F(\si^{\mrm{R}}_I) : 
\mrm{R} F(\mrm{R} F(I)) \to \mrm{R} F(I) \]
in $\dcat{D}(A, \mrm{gr})$ are isomorphisms.

Note that by our choice, if $M \in \dcat{D}^+(A, \mrm{gr})$ 
then $I$ is a bounded below complex. 
Since each $I^q$ is an injective object, it is right $F$-acyclic. 
Because $F$ is a weakly stable functor, each of the objects $F(I^q)$ 
is right $F$-acyclic too. If the functor $F$ has finite right cohomological 
dimension, or if $I$ is bounded below, Lemma \ref{lem:3885} says that the 
complexes $I$ and $F(I)$ are both right $F$-acyclic complexes.  

Consider the diagram 
\begin{equation} \label{eqn:3897}
\UseTips \xymatrix @C=10ex @R=6ex {
F(F(I))
\ar[d]^{ F(\si_I) }
\ar[r]^{ \eta^{\mrm{R}}_{F(I)} }
&
\mrm{R}F(F(I))
\ar[d]^(0.55){ \mrm{R}F(\si_I) }
\ar[r]^{ \mrm{R}F(\eta^{\mrm{R}}_I) }
&
\mrm{R}F(\mrm{R}F(I))
\ar[d]^(0.55){ \mrm{R}F(\si^{\mrm{R}}_I) }
\\
F(I)
\ar[r]^{ \eta^{\mrm{R}}_{I} }
& 
\mrm{R}F(I)
\ar[r]^{ \opn{id} }
&
\mrm{R}F(I)
}
\end{equation}
in $\dcat{D}(A, \mrm{gr})$. 
The left square is commutative: it is gotten from the vertical
morphism $\si_I : F(I) \to I$, to which we apply in the horizontal direction 
the morphism of functors $\eta^{\mrm{R}} : F \to \mrm{R} F$. The right square 
is also commutative: it comes from applying the functor $\mrm{R}F$ to the 
commutative diagram 
\[ \UseTips \xymatrix @C=8ex @R=6ex {
F(I)
\ar[r]^{\eta^{\mrm{R}}_I}
\ar[d]_{\si_I}
&
\mrm{R}F(I)
\ar[d]^{\si^{\mrm{R}}_I}
\\
I
\ar[r]^{ \opn{id} }
&
I
} \]
that characterizes $\si^{\mrm{R}}_I$. 
Because $I$ and $F(I)$ are right $F$-acyclic complexes, the morphisms 
$\eta^{\mrm{R}}_I$ and $\eta^{\mrm{R}}_{F(I)}$ are isomorphisms. Hence 
$\mrm{R}F(\eta^{\mrm{R}}_I)$ is an 
isomorphism. So the horizontal morphisms in the diagram \ref{eqn:3897} are all 
isomorphisms. We are given that $F$ is idempotent, and thus $F(\si_{I})$ is an
isomorphism.  The conclusion of this discussion is that 
$\mrm{R}F(\si^{\mrm{R}}_I)$ is an isomorphism.

Next, let $\phi : \mrm{R}F(I) \to J$ be an isomorphism in 
$\dcat{D}(A, \mrm{gr})$ 
to a K-injective complex $J$ such that 
$\opn{inf}(J) = \opn{inf}(\opn{H}(\mrm{R}F(I)))$.
We know that 
$\eta^{\mrm{R}}_I : F(I) \to \mrm{R}F(I)$ is an isomorphism, and therefore the 
composed morphism 
$\phi \circ \eta^{\mrm{R}}_I : F(I) \to J$ 
is an isomorphism in $\dcat{D}(A, \mrm{gr})$. 
Let $\til{\psi} : F(I) \to J$ be a 
quasi-isomorphism in $\dcat{C}_{\mrm{str}}(A, \mrm{gr})$ representing 
$\phi \circ \eta^{\mrm{R}}_I$. Since $F(I)$ and $J$ are both right $F$-acyclic 
complexes, it follows that 
$F(\til{\psi}) : F(F(I)) \to F(J)$ 
is a  quasi-isomorphism. Consider the following commutative diagram 
\[ \UseTips \xymatrix @C=8ex @R=6ex {
F(F(I))
\ar[r]^{F(\til{\psi})}
\ar[d]^{ \si_{F(I)} } 
&
F(J)
\ar[r]^{\eta^{\mrm{R}}_{J}}
\ar[d]^{\si_{J}}
& 
\mrm{R}F(J)
\ar[d]^(0.55){ \si^{\mrm{R}}_J }
&
\mrm{R}F(\mrm{R}F(I))
\ar[d]^(0.55){ \si^{\mrm{R}}_{\mrm{R}F(I)} }
\ar[l]_{ \mrm{R}F(\phi) }
\\
F(I)
\ar[r]^{\til{\psi}}
& 
J
\ar[r]^{\opn{id}}
& 
J
& 
\mrm{R}F(I)
\ar[l]_{\phi} 
} \]
in $\dcat{D}(A, \mrm{gr})$. All 
horizontal arrows here are isomorphisms. 
We are given that $F$ is idempotent, and thus $\si_{F(I)}$ is an
isomorphism. The conclusion is that 
$\si^{\mrm{R}}_{\mrm{R}F(I)}$ is an isomorphism.
\end{proof}

We end this subsection with a notion dual to ``copointed functor''. 

\begin{dfn} \label{dfn:3897} 
Let $\cat{M}$ be a linear category (e.g.\ $\dcat{M}(A, \mrm{gr})$ or 
$\dcat{D}(A, \mrm{gr})$), 
with identity functor $\opn{Id}_{\cat{M}}$.
\begin{enumerate}
\item A {\em pointed linear functor} 
\index{Additive functor! pointed}
on $\cat{M}$ is a pair 
$(G, \tau)$, consisting of a linear functor 
$G : \cat{M} \to \cat{M}$ and a morphism of functors 
$\tau : \opn{Id}_{\cat{M}} \to G$.

\item The pointed linear functor $(G, \tau)$ is called {\em idempotent} 
\index{Additive functor! idempotent}
if the morphisms 
\[ \tau_{G(M)}, \, G(\tau_M) : G(M) \to G(G(M)) \]
are isomorphisms for all objects $M \in \cat{M}$. 

\item If $\cat{M}$ is a triangulated category, $G$ is a triangulated functor, 
and $\tau$ is a morphism of triangulated functors, then we call 
$(G, \tau)$ a {\em pointed triangulated functor}. 
\end{enumerate}
\end{dfn}

\begin{rem} \label{rem:3896}
Idempotent copointed functors already appeared in the literature under another 
name: {\em idempotent comonads}. Another name for (nearly) the same notion is 
a (Bousfield) colocalization functor, see  e.g.\ \cite{Kr2}. 
Dually, idempotent pointed functors are the same 
thing as {\em idempotent monads}. See \cite{nLab} for a discussion of these 
concepts. 

In \cite[Section 4.1]{KaSc2}, what we call an idempotent pointed functor is 
called a {\em projector}. It is proved there that for an idempotent 
pointed functor $(G, \tau)$, and an object $M \in \cat{M}$,  
the isomorphisms $\tau_{G(M)}$ and $G(\tau_M)$ are equal. 
The same proof (with arrows reversed) shows that for an idempotent 
copointed functor $(F, \si)$, and an object $M$,  
the morphisms $\si_{F(M)}$ and $F(\si_M)$ are equal. We shall not require these 
facts. 
\end{rem}




\mysubsection{Graded Torsion: Weak Stability and Idempotence} 
\label{subsec:der-tors-graded}

In this subsection we begin the study of derived torsion over a connected 
graded ring. We adhere to Conventions \ref{conv:3700} and \ref{conv:4560}.
Thus $\K$ is a field, and all graded $\K$-rings are 
algebraically graded central $\K$-rings. The symbol $\ot$ means $\ot_{\K}$. 
The enveloping ring of a graded $\K$-ring $A$ is 
the graded $\K$-ring 
$A^{\mrm{en}} = A \ot A^{\mrm{op}}$. 
The category of left algebraically graded $A$-modules is 
$\dcat{M}(A, \mrm{gr})$, and it is a $\K$-linear abelian category.
Its derived category is $\dcat{D}(A, \mrm{gr})$.

Recall the notion of left noetherian connected graded $\K$-ring from 
Definitions \ref{dfn:4262} and \ref{dfn:3703}. In this subsection we will 
assume the following setup:

\begin{setup} \label{setup:3755}
We are given a left noetherian connected graded $\K$-ring $A$, with 
augmentation ideal 
$\m = \bigoplus_{i \geq 1} A_i \sub A$.
We are also given graded $\K$-rings $B$ and $C$.
\end{setup}

The graded rings $B$ and $C$ are auxiliary -- they are just placeholders, 
allowing us to write more general formulas. For instance, 
$\dcat{M}(A \ot B^{\mrm{op}}, \mrm{gr})$
is $\dcat{M}(A^{\mrm{en}}, \mrm{gr})$ when we take $B := A$, but it is 
$\dcat{M}(A, \mrm{gr})$ when $B := \K$. 

Since $A$ is left noetherian, the augmentation ideal $\m$ is a
finite graded left $A$-module. 

We shall often make implicit use of the canonical isomorphisms 
\begin{equation} \label{eqn:3850}
A \ot_{A} M \iso M \quad \tup{and} \quad \opn{Hom}_A(A, M) \iso M
\end{equation}
in $\dcat{M}(A \ot B^{\mrm{op}}, \mrm{gr})$. There are similar canonical 
isomorphisms
\begin{equation} \label{eqn:3790}
A \ot^{\mrm{L}}_{A} M \iso M \quad \tup{and} \quad \opn{RHom}_A(A, M) \iso M
\end{equation}
in $\dcat{D}(A \ot B^{\mrm{op}}, \mrm{gr})$. Also we view
$\K \cong A / \m$ as a graded $A^{\mrm{en}}$-module via the augmentation ring 
homomorphism $A \to \K$. 

We are interested in graded $\m$-torsion. This makes sense for bimodules, and 
hence the slightly complicated definition below. 

\begin{dfn} \label{dfn:3745}
Under Setup \ref{setup:3755}:
\begin{enumerate}
\item Let $M \in \dcat{M}(A \ot B^{\mrm{op}}, \mrm{gr})$. 
An element $m \in M$ is called an {\em $\m$-torsion element} 
\index{Torsion! element}
if $\m^j \cd m = 0$ for $j \gg 0$. 

\item The set of $\m$-torsion elements of $M$ is denoted by 
$\Ga_{\m}(M)$, 
\index{1-Gamma-m@$\Ga_{\m}$}
and it is called the {\em $\m$-torsion submodule} of $M$.

\item We call $M$ an {\em $\m$-torsion module} 
\index{Torsion! module}
if $\Ga_{\m}(M) = M$.

\item We denote by 
$\dcat{M}_{\mrm{tor}}(A \ot B^{\mrm{op}}, \mrm{gr})$
the full subcategory of 
$\dcat{M}(A \ot B^{\mrm{op}}, \mrm{gr})$
on the $\m$-torsion modules. 
\end{enumerate}
\end{dfn}

The fact that $\Ga_{\m}(M)$ is a graded $(A \ot B^{\mrm{op}})$-submodule of $M$ 
is easy to see. Indeed, we can express the torsion submodule as follows: 
\begin{equation} \label{eqn:3772}
\Ga_{\m}(M) = \lim_{j \to} \, 
\opn{Hom}_A( A / \m^{j}, M) \sub \opn{Hom}_A(A, M) , 
\end{equation}
using the identification (\ref{eqn:3850}), 
where $j \geq 1$, $\m^j$ is the $j$-fold product 
$\m \cdots \m \sub A$, 
$A \to A / \m^{j + 1} \to A / \m^{j}$ are the canonical graded ring 
surjections, and the operation $\lim_{j \to}$ is in the category 
$\dcat{M}(A \ot B^{\mrm{op}}, \mrm{gr})$. 

It is clear that 
\[ \Ga_{\m} : \dcat{M}(A \ot B^{\mrm{op}}, \mrm{gr}) \to 
\dcat{M}_{\mrm{tor}}(A \ot B^{\mrm{op}}, \mrm{gr}) \]
is a linear functor%
\index{Torsion! functor},
and the inclusions 
$\Ga_{\m}(M) \sub M$ assemble into a monomorphism 
$\si_{} : \Ga_{\m} \inj \opn{Id}$
of functors from $\dcat{M}(A \ot B^{\mrm{op}}, \mrm{gr})$ to itself. 

The ring $B$ plays an auxiliary role in Definition \ref{dfn:3745}. 
Because there is no mention of torsion for $B^{\mrm{op}}$-modules, the notation 
$\dcat{M}_{\mrm{tor}}(A \ot B^{\mrm{op}}, \mrm{gr})$
is not ambiguous. However, in Subsection \ref{subsec:sym-der-tor}, where 
torsion for $B^{\mrm{op}}$ modules does become an option, we will switch to 
the richer notation 
$\dcat{M}_{(\mrm{tor}, ..)}(A \ot B^{\mrm{op}}, \mrm{gr})$.
See Definition \ref{dfn:3815}. 

\begin{lem} \label{lem:3835}    
Under Setup \tup{\ref{setup:3755}} the following hold.
\begin{enumerate}
\item The functor $\Ga_{\m}$ commutes with graded ring homomorphisms 
$C \to B$, in the sense that the diagram 
\[ \UseTips \xymatrix @C=6ex @R=6ex {
\dcat{M}(A  \ot B^{\mrm{op}}, \mrm{gr})
\ar[r]^(0.45){\Ga_{\m}}
\ar[d]_{\opn{Rest}}
&
\dcat{M}_{\mrm{tor}}(A \ot B^{\mrm{op}}, \mrm{gr})
\ar[d]^{\opn{Rest}}
\\
\dcat{M}(A  \ot C^{\mrm{op}}, \mrm{gr})
\ar[r]^(0.45){\Ga_{\m}}
&
\dcat{M}_{\mrm{tor}}(A \ot C^{\mrm{op}}, \mrm{gr})
} \]
is commutative. 

\item The functor $\Ga_{\m}$ is idempotent, in the sense that for every 
$M \in \lb \dcat{M}(A \ot B^{\mrm{op}}, \mrm{gr})$ 
the homomorphisms 
\[ \si_{\Ga_{\m}(M)} \, , \,   \Ga_{\m}(\si_{M}) \, : \,
\Ga_{\m}(\Ga_{\m}(M)) \to \Ga_{\m}(M) \]
are bijective and are equal to each other.

\item Let $\phi : M \to N$ be a homomorphism in 
$\dcat{M}(A \ot B^{\mrm{op}}, \mrm{gr})$,
such that $M$ is $\m$-torsion. Then 
there is a unique homomorphism 
$\phi' : M \to \Ga_{\m}(N)$ satisfying $\phi = \si_{N} \circ \phi'$.

\item The functor $\Ga_{\m}$ is left exact.

\item The functor $\Ga_{\m}$ commutes with direct limits. 

\item The category $\dcat{M}_{\mrm{tor}}(A \ot B^{\mrm{op}}, \mrm{gr})$
is a thick abelian subcategory of \lb
$\dcat{M}(A \ot B^{\mrm{op}}, \mrm{gr})$,
closed under taking subobjects and quotients.
\end{enumerate}
\end{lem}

\begin{exer} \label{exer:3810}
Prove Lemma \ref{lem:3835}. (Hint: for item (6) you will need the fact that $A$ 
is left noetherian.)
\end{exer}

Recall that for a  module $M \in \dcat{M}(A, \mrm{gr})$, 
its socle is $\opn{Soc}(M) = \opn{Hom}_A(\K, M)$.
The socle is a functor, and there is a monomorphism 
$\opn{Soc} \inj \Ga_{\m}$ 
of functors from $\dcat{M}(A, \mrm{gr})$ to itself.

\begin{lem} \label{lem:5095}
Let $M \in \dcat{M}(A, \mrm{gr})$, and let 
$W := \opn{Soc}(M)$. If $M$ is an $\m$-torsion module, then $W$ is a 
graded-essential submodule of $M$.  
\end{lem}

\begin{proof}
Take any nonzero graded $A$-submodule 
$M' \sub M$. Let $m \in M'$ be some nonzero homogeneous element.
Because $M$ is a torsion module, yet $m \neq 0$, 
there is a unique $j \in \N$ such that 
$\m^{j + 1} \cd m  = 0$ but $\m^{j} \cd m \neq 0$. 
So there is a homogeneous element $a \in \m^{j}$ such that $a \cd m \neq 0$. 
The element $a \cd m$ is then a nonzero homogeneous element in 
$M' \cap W$. 
\end{proof}

The $\K$-linear dual (see Definition \ref{dfn:3725}) of the ring $A$ is the 
graded bimodule 
$A^* = \opn{Hom}_{\K}(A, \K) \in \dcat{M}(A^{\mrm{en}}, \mrm{gr})$.
The augmentation ring homomorphism $A \to \K$ induces, by dualizing, a 
canonical monomorphism 
\begin{equation} \label{eqn:4810}
\K \inj A^* 
\end{equation}
in $\dcat{M}(A^{\mrm{en}}, \mrm{gr})$.
For every $W \in \dcat{M}(\K, \mrm{gr})$ there is an induced monomorphism 
\begin{equation} \label{eqn:5095}
W \cong W \ot \K \inj W \ot A^*  
\end{equation}
in $\dcat{M}(A^{\mrm{en}}, \mrm{gr})$.

\begin{lem} \label{lem:5100}   
Under Setup \tup{\ref{setup:3755}}, let $W \in \dcat{M}(\K, \mrm{gr})$, and 
define 
$I := A^* \ot W \in \dcat{M}(A, \mrm{gr})$.
Consider the canonical monomorphism $W \inj I$ 
in $\dcat{M}(A, \mrm{gr})$ from \tup{(\ref{eqn:5095})}. 
Then\tup{:}
\begin{enumerate}
\item $I$ is a graded-injective $A$-module.

\item $I$ is an $\m$-torsion $A$-module.

\item $W = \opn{Soc}(I)$.

\item $W$ is a graded-essential submodule of $I$. 
\end{enumerate}
\end{lem}

\begin{proof} \mbox{}

\smallskip \noindent
(1) Say $W = \bigoplus_{x \in X} \K(-i_x)$. Then 
$I \cong \bigoplus_{x \in X} A^*(-i_x)$ in $\dcat{M}(A, \mrm{gr})$.
According to Proposition \ref{prop:3795}(2) and Theorem \ref{thm:3795}(1), the 
graded $A$-module $I$ is graded-injective. 

\medskip \noindent
(2) The $A$-module $A^*$ is $\m$-torsion, and a direct sum of torsion modules 
is torsion (see Lemma \ref{lem:3835}(5)). 

\medskip \noindent
(3) This is because $\K = \opn{Soc}(A^*)$, via the canonical monomorphism 
(\ref{eqn:4810}). 

\medskip \noindent
(4) Combine items (2) and (3) with Lemma \ref{lem:5095}.
\end{proof}

\begin{prop} \label{prop:4811}  
Under Setup \tup{\ref{setup:3755}}, let 
$M \in \dcat{M}(A, \mrm{gr})$, and let 
$W := \opn{Soc}(M)$. The three conditions below are equivalent. 
\begin{enumerate}
\rmitem{i} $M$ is an $\m$-torsion graded $A$-module. 

\rmitem{ii} $W$ is a graded-essential $A$-submodule of $M$. 

\rmitem{iii} There is a graded-essential monomorphism $\ep : M \inj A^* \ot W$
in $\dcat{M}(A, \mrm{gr})$, that restricts to the identity on $W$. 
\end{enumerate}
\end{prop}

\begin{proof} \mbox{}

\smallskip \noindent 
(i) $\Rightarrow$ (ii): This is Lemma \ref{lem:5095}.

\medskip \noindent 
(ii) $\Rightarrow$ (iii): By Lemma \ref{lem:5100}
the graded $A$-module $A^* \ot W$ is 
graded-injective. 
Hence there is a homomorphism 
$\ep : M \to A^* \ot W$ in $\dcat{M}(A, \mrm{gr})$
that makes the diagram 
\[ \UseTips \xymatrix @C=8ex @R=6ex {
*++{W}
\ar@{>->}[r]
\ar@{>->}[d]
&
M
\ar@{-->}[dl]^{\ep}
\\
*+{A^* \ot W}
} \]
in $\dcat{M}(A, \mrm{gr})$ commutative.
Let $M' := \opn{Ker}(\ep)$. 
If $M'$ were nonzero, then, because $W \sub M$ is a graded-essential
submodule, we would have 
$W \cap M' \neq 0$. But $W \to A^* \ot W$ is a monomorphism.
We conclude that $\ep$ is a monomorphism too. 

Finally, by Lemma \ref{lem:5100}(4) we know that $W$ is a 
graded-essential submodule of $A^* \ot W$. It follows that the monomorphism 
$\ep$ is graded-essential.

\medskip \noindent 
(iii) $\Rightarrow$ (i): The graded $A$-module $A^*$ is torsion.
We know that the subcategory 
$\dcat{M}_{\mrm{tor}}(A, \mrm{gr})$ of $\dcat{M}(A, \mrm{gr})$
is closed under taking arbitrary direct sums and subobjects. 
Hence $M$ is $\m$-torsion. 
\end{proof}

Cofinite graded modules were defined in Definition \ref{dfn:3726}. 

\begin{cor} \label{cor:4270}
Assume $A$ is a noetherian connected graded ring. Let 
$M \in \dcat{M}_{\mrm{tor}}(A, \mrm{gr})$. 
The two conditions below are equivalent.
\begin{enumerate}
\rmitem{i} $M$ is a cofinite graded $A$-module.  

\rmitem{ii} $W := \opn{Soc}(M)$ is a finite graded $\K$-module.
\end{enumerate}
\end{cor}

\begin{proof}
Graded Matlis Duality (Theorem \ref{thm:4045}) tells us that 
the cofinite graded $A$-modules are the artinian objects in the abelian 
category $\dcat{M}(A, \mrm{gr})$. 

\medskip \noindent 
(i) $\Rightarrow$ (ii): We are given that $M$ is a cofinite graded $A$-module. 
Hence so is its submodule $W$. Since $W$ is an 
$A$-module via $\K$, it must be an artinian graded $\K$-module. So 
$W$ is finite over $\K$. 

\medskip \noindent 
(ii) $\Rightarrow$ (i): Since $A^*$ is a cofinite graded $A$-module, and 
$A^* \ot W$ is a finite direct sum of degree twists of $A^*$,
we see that $A^* \ot W$ is a cofinite graded $A$-module. By Proposition 
\ref{prop:4811} we know that $M$ is isomorphic to a graded submodule of 
$A^* \ot W$; so $M$ is also a cofinite graded $A$-module.
\end{proof}

\begin{rem} \label{rem:4270}
We do not know if the corollary remains true if $A$ is just left noetherian. 
\end{rem}

Now we pass to derived categories and functors. 

\begin{dfn} \label{dfn:3755}
Under Setup \ref{setup:3755}, we denote by 
$\dcat{D}_{\mrm{tor}}(A \ot B^{\mrm{op}}, \mrm{gr})$
\index{1-Dtor(A-B,gr)@$\dcat{D}_{\mrm{tor}}(A \ot B^{\mrm{op}}, \mrm{gr})$}
the full subcategory of 
$\dcat{D}(A \ot B^{\mrm{op}}, \mrm{gr})$
on the complexes $M$ whose cohomology modules $\opn{H}^p(M)$
belong to $\dcat{M}_{\mrm{tor}}(A \ot B^{\mrm{op}}, \mrm{gr})$
for all $p$.
\end{dfn}

\begin{prop} \label{prop:3850}   
Under Setup \tup{\ref{setup:3755}}  the following hold.
\begin{enumerate}
\item The functor $\Ga_{\m}$ has a triangulated right derived functor 
\[ \mrm{R} \Ga_{\m} : \dcat{D}(A \ot B^{\mrm{op}}, \mrm{gr}) \to
\dcat{D}(A \ot B^{\mrm{op}}, \mrm{gr}) . \]
\index{1-RGamma-m@$\mrm{R} \Ga_{\m}$}
\index{Torsion! derived {\indash} functor}
If $I \in \dcat{D}(A \ot B^{\mrm{op}}, \mrm{gr})$
is K-graded-injective over $A$, then the morphism 
$\eta^{\mrm{R}}_{I} : \Ga_{\m}(I) \to \mrm{R} \Ga_{\m}(I)$
is an isomorphism.

\item There is a unique morphism  
$\si^{\mrm{R}}_{} : \mrm{R} \Ga_{\m} \to \opn{Id}$%
\index{1-SiR@$\si^{\mrm{R}}_{}$}
of triangulated functors from 
$\dcat{D}(A \ot B^{\mrm{op}}, \mrm{gr})$
to itself, such that for every 
$M \in \dcat{D}(A \ot B^{\mrm{op}}, \mrm{gr})$
there is equality 
$\si^{\mrm{R}}_{M} \circ \eta^{\mrm{R}}_{M} = \opn{Q}(\si_M)$
of morphisms $\Ga_{\m}(M) \to M$ in $\dcat{D}(A \ot B^{\mrm{op}}, \mrm{gr})$. 

\item The category $\dcat{D}_{\mrm{tor}}(A \ot B^{\mrm{op}}, \mrm{gr})$
is a full triangulated subcategory of \nl
$\dcat{D}(A \ot B^{\mrm{op}}, \mrm{gr})$, and it contains the image of the 
functor 
$\mrm{R} \Ga_{\m}$. 

\item The functor $\mrm{R} \Ga_{\m}$ commutes with graded $\K$-ring 
homomorphisms $C \to B$, in the sense that the diagram 
\[ \UseTips \xymatrix @C=6ex @R=6ex {
\dcat{D}(A  \ot B^{\mrm{op}}, \mrm{gr})
\ar[r]^(0.45){ \mrm{R}\Ga_{\m}}
\ar[d]_{\opn{Rest}}
&
\dcat{D}_{\mrm{tor}}(A \ot B^{\mrm{op}}, \mrm{gr})
\ar[d]^{\opn{Rest}}
\\
\dcat{D}(A  \ot C^{\mrm{op}}, \mrm{gr})
\ar[r]^(0.45){\mrm{R} \Ga_{\m}}
&
\dcat{D}_{\mrm{tor}}(A \ot C^{\mrm{op}}, \mrm{gr})
} \]
is commutative up to isomorphism. 
Likewise for the morphism of triangulated functors $\si^{\mrm{R}}_{}$.

\item For every $M \in \dcat{D}(A \ot B^{\mrm{op}}, \mrm{gr})$ and $p \in \Z$ 
there is an isomorphism 
\[ \opn{H}^p(\mrm{R} \Ga_{\m}(M)) \cong 
\lim_{j \to} \, \opn{Ext}^p_A( A / \m^{j}, M) \]
in $\dcat{M}(A \ot B^{\mrm{op}}, \mrm{gr})$. This isomorphism is functorial in 
$M$. 
\end{enumerate}
\end{prop}

\begin{proof}
(1) By Corollary \ref{cor:3740}(3) every 
$M \in \dcat{C}_{}(A \ot B^{\mrm{op}}, \mrm{gr})$
has a resolution $\rho : M \to I$ 
in $\dcat{C}_{\mrm{str}}(A \ot B^{\mrm{op}}, \mrm{gr})$
by a complex $I$ that is K-graded-injective over $A$. If 
$\psi : I \to I'$ is a quasi-isomorphism in 
$\dcat{C}_{\mrm{str}}(A \ot B^{\mrm{op}}, \mrm{gr})$ 
between such complexes, 
then it is a homotopy equivalence in 
$\dcat{C}_{\mrm{str}}(A, \mrm{gr})$, 
and hence $\Ga_{\m}(\psi)$ is a quasi-isomorphism. 
We see that the full subcategory 
$\cat{J} \sub \dcat{K}(A \ot B^{\mrm{op}}, \mrm{gr})$
on these complexes $I$ satisfies the conditions of Theorem 
\ref{thm:1470}. Therefore the right derived functor 
$\mrm{R} \Ga_{\m}$ exists, and $\eta^{\mrm{R}}_{I}$ is an isomorphism for 
all $I \in \cat{J}$. 

\medskip \noindent
(2) This is a special case of Lemma \ref{lem:1063}. 

\medskip \noindent
(3) Clear from Lemma \ref{lem:3835}(6). 

\medskip \noindent
(4) This is immediate from items (1) and (2). 

\medskip \noindent
(5) Use formula (\ref{eqn:3772}), and the fact that cohomology commutes with 
direct limits.  
\end{proof}

\begin{dfn} \label{dfn:4045}
Under Setup \tup{\ref{setup:3755}}, the {\em $i$-th local cohomology functor} 
is the functor 
\[ \opn{H}^i_{\m} := \mrm{R}^i \Ga_{\m} : 
\dcat{M}(A \ot B^{\mrm{op}}, \mrm{gr}) \to 
\dcat{M}(A \ot B^{\mrm{op}}, \mrm{gr}) , \]
\index{1-Him@$\opn{H}^i_{\m}$}
the $i$-th right derived functor of $\Ga_{\m}$. 
\end{dfn}

\begin{cor} \label{cor:3836}
For every $i \in \N$ there is an isomorphism 
$\opn{H}^i_{\m} \cong \opn{H}^i \circ \, \mrm{R} \Ga_{\m}$
of functors from $\dcat{M}(A  \ot B^{\mrm{op}}, \mrm{gr})$ to itself. 
\end{cor}

\begin{proof}
This is immediate from Proposition \ref{prop:3850}(1). 
\end{proof}

\begin{prop} \label{prop:4050}   
For every $i \in \N$ the functor 
\[ \opn{H}^i_{\m} : \dcat{M}(A \ot B^{\mrm{op}}, \mrm{gr}) \to 
\dcat{M}(A \ot B^{\mrm{op}}, \mrm{gr}) \]
commutes with direct limits. 
\end{prop}

\begin{proof}
In view of Proposition \ref{prop:3850}(4), with $C := \K$, we can assume that 
$B = \K$, so $A \ot B^{\mrm{op}} = A$. 
Since direct limits commute with each other, by Proposition \ref{prop:3850}(5) 
it suffices to prove that for every $i, j$ the functor 
$\opn{Ext}^i_A( A / \m^{j}, -)$
commutes with direct limits. Now $A$ is left noetherian, so we can find a 
resolution $P \to A / \m^{j}$ by a nonpositive complex, where each $P^i$ is a 
finite graded-free $A$-module. And the functor 
\[ \opn{Ext}^i_A( A / \m^{j}, -) \cong 
\opn{H}^i \circ  \opn{Hom}_A( P, -) : \dcat{M}(A, \mrm{gr}) \to
\dcat{M}(A, \mrm{gr}) \]
commutes with infinite direct limits. 
\end{proof}

Definition \ref{dfn:3885}, when used in our context, says that 
the right cohomological dimension of the functor $\Ga_{\m}$ is 
\begin{equation} \label{eqn:3962}
d := \opn{sup} \, \{ p \in \N \mid \mrm{H}^p_{\m} \neq 0 \} \in 
\N \cup \{ \pm \infty \} .
\end{equation}

The cohomological dimension of a triangulated functor was introduced in 
Definition \ref{dfn:2126}. A triangulated functor is called 
{\em quasi-compact} if it commutes with infinite direct sums (Definition 
\ref{dfn:3888}). 

\begin{thm} \label{thm:4050}  
Under Setup \tup{\ref{setup:3755}}, assume the torsion functor
$\Ga_{\m} : \lb \dcat{M}(A, \mrm{gr}) \to \dcat{M}(A, \mrm{gr})$ 
has finite right cohomological dimension. Then the derived torsion functor 
\[ \mrm{R} \Ga_{\m} : \dcat{D}(A \ot B^{\mrm{op}}, \mrm{gr}) \to 
\dcat{D}(A \ot B^{\mrm{op}}, \mrm{gr}) \]
is quasi-compact and has finite cohomological dimension. 
\index{Additive functor! quasi-compact}
\index{Additive functor! finite dimensional}
\end{thm}

\begin{proof}
According to Proposition \ref{prop:3850}(4), the functor 
\[ \Ga_{\m} : \dcat{M}(A \ot B^{\mrm{op}}, \mrm{gr}) \to 
\dcat{M}(A \ot B^{\mrm{op}}, \mrm{gr}) \]
also has finite right cohomological dimension.
The right derived functors
\[ \mrm{R}^q \Ga_{\m} = \opn{H}^q_{\m} : 
\dcat{M}(A \ot B^{\mrm{op}}, \mrm{gr}) \to 
\dcat{M}(A \ot B^{\mrm{op}}, \mrm{gr}) \]
are quasi-compact by Proposition \ref{prop:4050}. So Theorem \ref{thm:1001} 
applies. 
\end{proof}

\begin{lem} \label{lem:3838}
Let $I \in \dcat{M}(A, \mrm{gr})$, and let 
$W := \opn{Soc}(I)$. The following three conditions are equivalent.
\begin{enumerate}
\rmitem{i} $I$ is graded-injective and $\m$-torsion. 

\rmitem{ii} $I$ is graded-injective and $W$ is an essential submodule of $I$. 

\rmitem{iii} There is an isomorphism 
$I \cong A^* \ot W$ in $\dcat{C}(A, \mrm{gr})$ which is the identity on $W$. 
\end{enumerate}
\end{lem}

In condition (iii) we view $W$ as a submodule of $A^* \ot W$ using the 
canonical embedding (\ref{eqn:5095}).

\begin{proof} \mbox{}

\smallskip \noindent
(i) $\Rightarrow$ (ii): This is Lemma \ref{lem:5095}.

\medskip \noindent
(ii) $\Rightarrow$ (iii): 
By Proposition \ref{prop:4811} there is an essential monomorphism 
$\ep : I \to A^* \ot W$ in $\dcat{M}(A, \mrm{gr})$
that extends the identity of $W$. Because $I$ is graded-injective, the 
monomorphism $\ep$ is split, by some epimorphism $\si : A^* \ot W \to I$. 
But $W$ is a graded-essential submodule of $A^* \ot W$, and therefore $\si$ is 
also a monomorphism. We conclude that both $\si$ and $\ep$ are 
isomorphisms. 

\medskip \noindent
(iii) $\Rightarrow$ (i): 
According to Lemma \ref{lem:5100}, the graded $A$-module $A^* \ot W$ is 
graded-injective and $\m$-torsion.
\end{proof}

\begin{lem} \label{lem:3900}
Let $M \in \dcat{M}(A, \mrm{gr})$ be an $\m$-torsion module, and let $I$ be an 
injective hull of $M$ in $\dcat{M}(A, \mrm{gr})$. Then 
$I \cong A^* \ot W$, where $W$ is the socle of $M$. 
In particular, $I$ is an  $\m$-torsion module.
\end{lem}

\begin{proof}
By Proposition \ref{prop:4811} there is a graded-essential monomorphism 
$\ep : M \inj A^* \ot W$ in $\dcat{M}(A, \mrm{gr})$.
From Lemma \ref{lem:5100}(1) we know that $A^* \ot W$ is a 
graded-injective module over $A$. So by the uniqueness (up to isomorphism) 
of graded-injective hulls, we get an isomorphism 
$I \cong A^* \ot W$. Finally, by Lemma \ref{lem:5100}(2),
$A^* \ot W$ is $\m$-torsion. 
\end{proof}

\begin{lem} \label{lem:3839}
Let $I$ be a graded-injective $A$-module. Then there is an isomorphism
$\Ga_{\m}(I) \cong A^* \ot W$ in $\dcat{M}(A, \mrm{gr})$,
where $W$ is the socle of $I$. 
\end{lem}

\begin{proof}
Consider the essential monomorphism $\ep : W \inj A^* \ot W$ 
from Lemma \ref{lem:5100}(4), and 
the monomorphism $\tau : W \inj I$.
Because $I$ is graded-injective, there is a homomorphism 
$\psi : A^* \ot W \to I$ in $\dcat{M}(A, \mrm{gr})$
such that $\psi \circ \ep = \tau$. Since $\ep$ is a graded-essential 
monomorphism, the usual calculation (see the proof of the implication 
(ii) $\Rightarrow$ (iii) in Proposition \ref{prop:4811}) shows that 
$\psi$ is a monomorphism. Let $J \sub I$ be the image of $\psi$,
so $\psi : A^* \ot W \iso J$. We know (from Lemma \ref{lem:5100}) 
that $J$ is graded-injective and $\m$-torsion. Hence there is a direct sum 
decomposition $I = J \oplus J'$
in $\dcat{M}(A, \mrm{gr})$, and $\Ga_{\m}(J) = J$.

By Proposition \ref{prop:4811} we know that 
$W \cap J' = \opn{Soc}(J')$ is a graded-essential submodule of $\Ga_{\m}(J')$; 
but $W \sub J$, so $W \cap J' = 0$, and therefore $\Ga_{\m}(J') = 0$.
The conclusion is that $\Ga_{\m}(I) = J$. 
\end{proof}

\begin{thm} \label{thm:3990}   
Under Setup \tup{\ref{setup:3755}}, the torsion functor
\[ \Ga_{\m} : \dcat{M}(A, \mrm{gr}) \to \dcat{M}(A, \mrm{gr}) \]
is stable, namely for every graded-injective $A$-module $I$
the $\m$-torsion submodule $\Ga_{\m}(I)$ is graded-injective. 
\end{thm}

\begin{proof}
This is immediate from Lemma \ref{lem:3839} and
Lemma \ref{lem:5100}(1). 
\end{proof}

Here is the graded NC version of Definition \ref{dfn:2210}.

\begin{dfn} \label{dfn:3837} \mbox{}
\begin{enumerate}
\item A {\em minimal complex of graded-injective $A$-modules}% 
\index{Complex of algebraically graded modules! minimal {\indash}
graded-injective {\indash}}
is a bounded below complex $I$ of graded-injective $A$-modules, such 
that for every $p$ the submodule $\opn{Z}^p(I) \sub I^p$ is graded-essential.

\item Let $M \in \dcat{D}^+(A, \mrm{gr})$. A 
{\em minimal graded-injective resolution}% 
\index{Resolution! minimal graded-injective}
of $M$ is a  quasi-isomorphism $\rho : M \to I$ 
in $\dcat{C}_{\mrm{str}}(A, \mrm{gr})$, where $I$ is a minimal complex of 
graded-injective $A$-modules.
\end{enumerate}
\end{dfn}

\begin{lem} \label{lem:3837}
Let $M \in \dcat{D}^+(A, \mrm{gr})$. Then $M$ has a minimal 
graded-injective resolution.
\end{lem}

\begin{proof}
The proof of Proposition \ref{prop:2210} works here too. 
\end{proof}

There is uniqueness for minimal graded-injective resolutions, but we won't need 
it. 

For the next proposition, and for later results, it will be convenient to use 
this notation:

\begin{notation} \label{not:3965}
If the torsion functor $\Ga_{\m}$ has finite right cohomological dimension, 
then we let $\star := \bra{\tup{empty}}$, the empty boundedness indicator; 
otherwise we let $\star := +$. Thus $\dcat{D}^{\star}$ is either 
$\dcat{D}$ or $\dcat{D}^+$.
\end{notation}

\begin{prop} \label{prop:3965}
With Notation \tup{\ref{not:3965}}, let 
$M \in \dcat{D}^{\star}(A \ot B^{\mrm{op}}, \mrm{gr})$. The following 
conditions are equivalent\tup{:}
\begin{enumerate}
\rmitem{i} $M \in \dcat{D}^{\star}_{\mrm{tor}}(A \ot B^{\mrm{op}}, \mrm{gr})$.

\rmitem{ii} The morphism 
$\si^{\mrm{R}}_{M} : \mrm{R} \Ga_{\m}(M) \to M$
in $\dcat{D}(A \ot B^{\mrm{op}}, \mrm{gr})$
is an isomorphism.
\end{enumerate}
\end{prop}

\begin{proof}
The implication (ii) $\Rightarrow$ (i) is due to Proposition 
\ref{prop:3850}(3).  

For the other implication, we can forget the ring $B$; it is not relevant. 
First let us assume that 
$M \in \dcat{D}^{+}_{\mrm{tor}}(A, \mrm{gr})$.
By Lemma \ref{lem:3900} the injective hull in $\dcat{M}(A, \mrm{gr})$ of 
a torsion module is torsion. 
Then, according to Theorem \ref{thm:2880}, there is a quasi-isomorphism 
$M \to I$ in $\dcat{C}^+_{\mrm{str}}(A, \mrm{gr})$,
where $I$ is a complex of torsion graded-injective $A$-modules.
It follows that homomorphism 
$\si_I : \Ga_{\m}(I) \to I$
in $\dcat{C}^+_{\mrm{str}}(A, \mrm{gr})$, that represents $\si^{\mrm{R}}_{M}$,
is an isomorphism. 

When $\opn{H}(M)$ is not bounded below, but $\Ga_{\m}$ has finite right 
cohomological dimension, we can use smart truncation, as in the proof of 
Theorem \ref{thm:2135}(2).
\end{proof}

\begin{prop} \label{prop:3990}   
With Notation \tup{\ref{not:3965}},
Let $M \in \dcat{D}^{\star}_{\mrm{tor}}(A \ot B^{\mrm{op}}, \mrm{gr})$ and 
$N \in \dcat{D}^+(A \ot C^{\mrm{op}}, \mrm{gr})$.
Then there is an isomorphism 
\[ \opn{RHom}_{A}(M, N) \iso 
\opn{RHom}_{A} \bigl( M, \mrm{R} \Ga_{\m}(N) \bigr) \]
in $\dcat{D}(B \ot C^{\mrm{op}}, \mrm{gr})$, 
and it is functorial in $M$ and $N$. 
\end{prop}

\begin{proof}
Choose a resolution $M \to I_M$ in 
$\dcat{C}^{\star}_{\mrm{str}}(A \ot B^{\mrm{op}}, \mrm{gr})$
where $I_M$ is K-graded-injective over $A$, each $I_M^p$ is 
graded-injective over $A$, and 
$\opn{inf}(I_M) = \opn{inf}(\opn{H}((M))$; this is possible by Corollary 
\ref{cor:3830}(3). Choose the same sort of resolution 
$N \to I_N$ in $\dcat{C}^+_{\mrm{str}}(A \ot C^{\mrm{op}}, \mrm{gr})$.
Then 
\begin{equation} \label{eqn:3990}
\opn{RHom}_{A}(M, N) \cong \opn{Hom}_{A}(I_M, I_N) 
\end{equation}
in $\dcat{D}(B \ot C^{\mrm{op}}, \mrm{gr})$,
and $\mrm{R} \Ga_{\m}(N) \cong \Ga_{\m}(I_N)$ in 
$\dcat{D}(A \ot C^{\mrm{op}}, \mrm{gr})$. By Theorem \ref{thm:3990} each 
$\Ga_{\m}(I^p_N)$ is graded-injective as an $A$-module, so 
$\Ga_{\m}(I_N)$ is a bounded below complex of graded-injective $A$-modules, and 
therefore it is a K-graded-injective complex over $A$. 
We see that  
\begin{equation} \label{eqn:3991}
\opn{RHom}_{A} \bigl( M, \mrm{R} \Ga_{\m}(N) \bigr) \cong
\opn{Hom}_{A} \bigl( I_M, \Ga_{\m}(I_N) \bigr)
\end{equation}
in $\dcat{D}(B \ot C^{\mrm{op}}, \mrm{gr})$. 

By Proposition \ref{prop:3965} (in both cases of Notation \ref{not:3965}) 
the homomorphism $\si_{I_M} : \Ga_{\m}(I_M) \to I_M$ is a quasi-isomorphism. 
Hence in the commutative diagram 
\[ \UseTips \xymatrix @C=12ex @R=6ex {
\opn{Hom}_{A}(I_M, I_N) 
\ar[r]^(0.45){ \opn{Hom}(\si_{I_M}, \opn{id}) }
&
\opn{Hom}_{A} \bigl( \Ga_{\m}(I_M), I_N \bigr)
\\
\opn{Hom}_{A} \bigl( I_M, \Ga_{\m}(I_N) \bigr)
\ar[r]^(0.45){ \opn{Hom}(\si_{I_M}, \opn{id}) }
\ar[u]^{ \opn{Hom}(\opn{id}, \si_{I_N}) }
&
\opn{Hom}_{A} \bigl( \Ga_{\m}(I_M), \Ga_{\m}(I_N) \bigr) 
\ar[u]_{ \opn{Hom}(\opn{id}, \si_{I_N}) }
} \]
in $\dcat{C}_{\mrm{str}}(B \ot C^{\mrm{op}}, \mrm{gr})$
the two horizontal arrows are quasi-isomorphisms.
The right vertical arrow is an isomorphism 
in $\dcat{C}_{\mrm{str}}(B \ot C^{\mrm{op}}, \mrm{gr})$ by Lemma 
\ref{lem:3835}(3). 
It follows that the left vertical arrow is a quasi-isomorphism.
Combining this with the isomorphisms (\ref{eqn:3990}) and 
(\ref{eqn:3991}) we obtain the desired isomorphism above. 
\end{proof}

The next definition is a specialization of Definitions \ref{dfn:3887}, 
\ref{dfn:3890} and \ref{dfn:3895}.

\begin{dfn} \label{dfn:3748} 
Under Setup \tup{\ref{setup:3755}}:
\begin{enumerate}
\item A module $N \in \dcat{M}(A \ot B^{\mrm{op}}, \mrm{gr})$
is called {\em graded-$\m$-flasque}% 
\index{Algebraically graded module! graded-$\m$-flasque}
if it is a right $\Ga_{\m}$-acyclic object, i.e.\ if 
$\mrm{H}^q_{\m}(N) = 0$ for all $q > 0$.

\item The functor 
\[ \Ga_{\m} : \dcat{M}(A \ot B^{\mrm{op}}, \mrm{gr}) \to
\dcat{M}(A \ot B^{\mrm{op}}, \mrm{gr}) \]
is called {\em weakly stable}% 
\index{Torsion! weakly stable {\indash} functor}
if for every graded-injective module
$I \in \lb \dcat{M}(A \ot B^{\mrm{op}}, \mrm{gr})$, 
the module  
$\Ga_{\m}(I) \in \dcat{M}(A \ot B^{\mrm{op}}, \mrm{gr})$
is graded-$\m$-flasque. 

\item A complex  $N \in \dcat{D}(A \ot B^{\mrm{op}}, \mrm{gr})$
is called a {\em K-graded-$\m$-flasque complex}% 
\index{Complex of algebraically graded modules! K-graded-$\m$-flasque}
if it is a right $\Ga_{\m}$-acyclic complex, i.e.\ if the morphism 
$\eta^{\mrm{R}}_{N} : \Ga_{\m}(N) \to \mrm{R} \Ga_{\m}(N)$
in $\dcat{D}(A \ot B^{\mrm{op}}, \mrm{gr})$ is an isomorphism. 
\end{enumerate}
\end{dfn}

The term ``flasque'' in this meaning seems to have been first used in 
\cite{YeZh5}.

\begin{exa} \label{exa:3745}
Suppose $M \in \dcat{M}(A, \mrm{gr})$ is an $\m$-torsion module. Consider its 
minimal graded-injective resolution 
$0 \to M \to I^0 \to I^1 \to  \cdots$. 
By Lemma \ref{lem:3900} and induction on $q$, we see that all the graded 
modules $I^q$ are $\m$-torsion; in fact 
$I^q \cong A^* \ot W^q$ for some $W^q \in \dcat{M}(\K, \mrm{gr})$. 
Thus
$\mrm{R} \Ga_{\m}(M) \cong \Ga_{\m}(I) = I \cong M$,
so $\mrm{R}^q \Ga_{\m}(M) = 0$ for all $q > 0$, and $M$ is a 
graded-$\m$-flasque $A$-module. 
\end{exa}

\begin{thm} \label{thm:3900}
Under Setup \tup{\ref{setup:3755}}, the torsion functor
\[ \Ga_{\m} : \dcat{M}(A \ot B^{\mrm{op}}, \mrm{gr}) \to 
\dcat{M}(A \ot B^{\mrm{op}}, \mrm{gr}) \]
is weakly stable.
\end{thm}

\begin{proof}
Take an injective object $I \in \dcat{M}(A \ot B^{\mrm{op}}, \mrm{gr})$. 
According to Proposition \ref{prop:3900}(3), the module $I$ is 
graded-injective over $A$. By Lemma \ref{lem:3839} we know that 
there is an isomorphism 
$\Ga_{\m}(I) \cong A^* \ot W$
in $\dcat{M}(A, \mrm{gr})$ for some graded $\K$-module $W$. This is a 
graded-injective $A$-module, and 
therefore $\Ga_{\m}(I)$ is graded-$\m$-flasque as an object of 
$\dcat{M}(A, \mrm{gr})$. 
Finally, by Proposition \ref{prop:3850}(4), the property of being 
graded-$\m$-flasque is insensitive to the restriction functor 
$\dcat{M}(A \ot B^{\mrm{op}}, \mrm{gr}) \to \dcat{M}(A, \mrm{gr})$,
so $\Ga_{\m}(I)$ is graded-$\m$-flasque as an object of 
$\dcat{M}(A \ot B^{\mrm{op}}, \mrm{gr})$. 
\end{proof}

In Proposition \ref{prop:3850}(2) we presented a morphism 
$\si^{\mrm{R}}_{} : \mrm{R} \Ga_{\m} \to \opn{Id}$
of triangulated functors from 
$\dcat{D}^{\star}(A \ot B^{\mrm{op}}, \mrm{gr})$
to itself. Here we are using Notation \ref{not:3965}. 
According to Definition \ref{dfn:3896} the pair 
$(\mrm{R} \Ga_{\m}, \si^{\mrm{R}})$ 
is a copointed triangulated functor. Recall that 
$(\mrm{R} \Ga_{\m}, \si^{\mrm{R}})$ is an 
idempotent copointed triangulated functor 
if for every complex
$M \in \dcat{D}^{\star}(A \ot B^{\mrm{op}}, \mrm{gr})$
the morphisms 
\[ \si^{\mrm{R}}_{\mrm{R} \Ga_{\m}(M)} \, , \,   
\mrm{R} \Ga_{\m}(\si^{\mrm{R}}_{M}) \, : \,
\mrm{R} \Ga_{\m}(\mrm{R} \Ga_{\m}(M)) \to \mrm{R} \Ga_{\m}(M) \]
in $\dcat{D}(A \ot B^{\mrm{op}}, \mrm{gr})$
are isomorphisms. 

\begin{cor}[\cite{VyYe}] \label{cor:3966}
With Notation \tup{\ref{not:3965}}, the copointed triangulated functor 
$(\mrm{R} \Ga_{\m}, \si^{\mrm{R}})$
on $\dcat{D}^{\star}(A \ot B^{\mrm{op}}, \mrm{gr})$ is idempotent.
\end{cor}

\begin{proof}
By Theorem \ref{thm:3900} the functor $\Ga_{\m}$ is weakly stable, and 
by Lemma \ref{lem:3835}(2) the copointed functor $(\Ga_{\m}, \si)$ is 
idempotent. Thus Theorem \ref{thm:1004} applies. 
\end{proof}

We end the subsection with an example showing a torsion functor as in 
Theorem \ref{thm:3900} which is weakly stable but not stable.
The example is a slight modification of an example in \cite{VyYe}. 

\begin{exa} \label{exa:4077}
Let $A := \K[t]$, the polynomial ring in a variable $t$ of degree $1$. It is 
connected noetherian, and its augmentation ideal $\m$ is generated by $t$. 
Let $B := \K[s_0, s_1, \ldots]$, the commutative polynomial ring in countably 
many variables $s_0, s_1, \ldots$ of degree $1$. So $B$ is a graded ring, 
and it is not noetherian. Theorem \ref{thm:3900} tells us that the functor 
\[ \Ga_{\m} : \dcat{M}(A \ot B, \mrm{gr}) \to \dcat{M}(A \ot B, \mrm{gr}) \]
is weakly stable. (Since $B$ is commutative, it doesn't matter whether we write 
$B$ or $B^{\mrm{op}}$). 
We will prove that $\Ga_{\m}$ is not stable. 

Consider a countable collection $\{ I_p \}_{p \in \N}$ of graded-injective 
$B$-modules. Then 
$I := \prod_{p \in \N} I_p(-p)$
is a graded-injective $B$-module, and 
therefore $J := \opn{Hom}_B(A \ot B, I)$ is a graded-injective 
$(A \ot B)$-module. If the functor $\Ga_{\m}$ on 
$\dcat{M}(A \ot B, \mrm{gr})$ were stable, this would imply 
that $\Ga_{\m}(J)$ is also a graded-injective $(A \ot B)$-module. Since 
$B \to A \ot B$ is flat, $\Ga_{\m}(J)$ is then a graded-injective $B$-module
(see Proposition \ref{prop:3900}). 

As graded $B$-modules there are isomorphisms
\[ A \ot B \cong \bigoplus\nolimits_{p \in \N} B \cd t^p 
\cong \bigoplus\nolimits_{p \in \N} B(-p) . \]
Therefore 
\[ J = \opn{Hom}_B(A \ot B, I) \cong \prod\nolimits_{p \in \N} I(p)  \]
as graded $B$-modules, and 
\[ \Ga_{\m}(J) \cong \lim_{p \to} \, 
\Bigl( \prod\nolimits_{q = 0}^p I(q) \Bigr) \cong 
\bigoplus\nolimits_{p \in \N} I(p) . \]
Now for each $p$ the graded-module $I_p$ is a direct summand of $I(p)$, and 
hence $\bigoplus_{p \in \N} I_p$ is a direct summand of $\Ga_{\m}(J)$, as 
graded $B$-modules. Therefore $\bigoplus_{p \in \N} I_p$ is a graded-injective 
$B$-module. 

The conclusion is that a countable direct sum of graded-injective $B$-modules 
is graded-injective. According to the Bass-Papp Theorem 
\cite[Theorem 3.46]{Lam}, that holds also in the graded sense (cf.\ Proposition 
\ref{prop:3745}), the graded ring $B$ is 
noetherian. This is a contradiction. 
\end{exa}



\mysubsection{Representability of Derived Torsion} 
\label{subsec:repr-der-tors}

The results of this subsection are an adaptation of results from 
\cite[Section 7]{VyYe} to the connected graded setting. 

In Section \ref{sec:BDC} we are going to use the representability of derived 
torsion to give an alternative proof of M. Van den Bergh's theorem on 
existence of balanced dualizing complexes (see Theorem \ref{thm:3713}). 

Recall that Conventions \ref{conv:3700} and \ref{conv:4560}
are in force. In this subsection we assume the following strengthening of Setup 
\ref{setup:3755}.

\begin{setup} \label{setup:3900}
We are given a left noetherian connected graded $\K$-ring $A$, with 
augmentation ideal 
$\m = \bigoplus_{i \geq 1} A_i \sub A$,
such that the torsion functor $\Ga_{\m}$ has finite right cohomological 
dimension. We are also given a graded $\K$-ring $B$.
\end{setup}

As before, the graded ring $B$ plays an auxiliary role, a placeholder for $A$ 
or $\K$. 

\begin{dfn} \label{dfn:3750}
Under Setup \ref{setup:3900}, the 
{\em dedualizing complex of $A$}% 
\index{Dedualizing complex}
is the complex 
$P_A := \mrm{R} \Ga_{\m}(A) \in \dcat{D}(A^{\mrm{en}}, \mrm{gr})$.
\end{dfn}

The name ``dedualizing'' is due to L. Positselsky \cite{Pos} . Regarding the  
lack of left-right symmetry in this definition, see Corollary \ref{cor:3875} 
and Proposition \ref{prop:3712}.

As a special case of Proposition \ref{prop:3850}, there is a morphism 
$\si^{\mrm{R}}_{A} : P_A = \mrm{R} \Ga_{\m}(A) \lb \to A$
in $\dcat{D}(A^{\mrm{en}}, \mrm{gr})$. 

\begin{thm}[Representability of Derived Torsion, \cite{VyYe}] \label{thm:3750} 
\index{Torsion! representability of derived {\indash}}
Under Setup \tup{\ref{setup:3900}} there is a unique isomorphism 
\[ \ttev^{\mrm{R, L}}_{\m, (-)} \, : 
\, P_A \ot^{\mrm{L}}_{A} (-) \, \iso \, \mrm{R} \Ga_{\m}(-) \]
of triangulated functors from 
$\dcat{D}(A \ot B^{\mrm{op}}, \mrm{gr})$
to itself, called {\em derived $\m$-torsion tensor-evaluation},
\index{Tensor-evaluation morphism! derived $\m$-torsion}
\index{1-EvRLm@$\ttev^{\mrm{R, L}}_{\m, (-)}$}
such that for every complex 
$M \in \dcat{D}(A \ot B^{\mrm{op}}, \mrm{gr})$
the diagram 
\[ \tag{$\heartsuit$}
\UseTips \xymatrix @C=10ex @R=6ex {
P_A \ot^{\mrm{L}}_{A} M
\ar[r]^{ \ttev^{\mrm{R, L}}_{\m, M} }_{\cong}
\ar[d]_{\si^{\mrm{R}}_{A} \, \ot^{\mrm{L}}_{A} \, \opn{id}_M}
& 
\mrm{R} \Ga_{\m}(M)
\ar[d]^{\si^{\mrm{R}}_{M}}
\\
A \ot^{\mrm{L}}_{A} M
\ar[r]^{\opn{lu}}_{\cong}
& 
M
} \]
is commutative.
\end{thm}

The isomorphism marked ``$\opn{lu}$'' in diagram ($\heartsuit$) is the left 
unitor isomorphism of the closed monoidal structure, see equation 
(\ref{eqn:3790}). 

The proof of the theorem requires the next lemma.

\begin{lem} \label{lem:3905}
Let 
$F, G : \dcat{D}(A, \mrm{gr}) \to \dcat{D}(B, \mrm{gr})$
be quasi-compact triangulated functors, and let $\eta : F \to G$
be a morphism of triangulated functors. Assume that
$\eta_A : F (A) \to G (A)$
is an isomorphism. Then $\eta$ is an isomorphism. 
\end{lem}

\begin{proof}
We are given that $\eta_A$ is an isomorphism. Because the algebraic degree 
twist $M \mapsto M(i)$ is an automorphism of $\dcat{D}(A, \mrm{gr})$,
it follows that $\eta_{A(i)}$ is an isomorphism for every integer $i$.
Likewise for the cohomological degree shift (i.e.\ the translation 
$M \mapsto M[i]$ on complexes), so $\eta_{A(i)[j]}$ is an isomorphism for every 
integer $j$. Both functors are quasi-compact, and therefore 
$\eta_P$ is an isomorphism for every graded-free complex of $A$-modules 
$P \cong \bigoplus_{x \in X} A(i_x)[j_x]$. 

Suppose we are given a distinguished triangle
$M' \to M \to M'' \xar{\triangle}$
in $\dcat{D}(A, \mrm{gr})$, such that two of the three morphisms 
$\eta_{M'}$, $\eta_{M}$ and $\eta_{M''}$ are isomorphisms. Because $\eta$ is a 
morphism of triangulated functors, it follows that third morphism is
also an isomorphism. 

Next consider a semi-graded-free DG module $P$. Take a filtration \lb 
$\{ F_j(P) \}_{j \geq -1}$ of $P$ as in Definition \ref{dfn:4275}(2).
For every $j$ we have a distinguished triangle
\[ F_{j - 1}(P) \xar{\th_j} F_j(P) \to \opn{Gr}^F_j(P) \xar{\triangle} \]
in $\dcat{D}(A, \mrm{gr})$, where 
$\th_j : F_{j - 1}(P) \to F_j(P)$ is the inclusion. 
Since $\opn{Gr}^F_j(P)$ is a graded-free complex, by
induction on $j$ we conclude that $\eta_{F_j(P)}$ is an isomorphism for every 
$j \geq 0$. The homotopy colimit construction
(see Definition \ref{dfn:3410})
gives a distinguished triangle
\[ \bigoplus_{j \in \N} F_j(P) \xar{\ \Theta \ } \bigoplus_{j \in \N} F_j(P) 
\to P \xar{\triangle} \]
in $\dcat{D}(A, \mrm{gr})$, where 
\[ \Theta|_{F_{j - 1}(P)} := 
(\opn{id}_{}, - \th_j) : F_{j - 1}(P) \to F_{j - 1}(P) \oplus F_{j}(P) . \]
By quasi-compactness we know that 
$\eta_{\bigoplus F_j(P)}$ is an isomorphism. Therefore $\eta_P$ 
is an isomorphism. 

Finally, every $M \in \dcat{D}(A, \mrm{gr})$ admits an  
isomorphism $M \cong P$ with $P$ semi-graded-free. Therefore $\eta_M$ is an 
isomorphism.
\end{proof}

\begin{proof}[Proof of Theorem \tup{\ref{thm:3750}}] \mbox{}

\smallskip \noindent 
Step 1. We begin by constructing the morphism of triangulated functors 
$\ttev^{\mrm{R, L}}_{\m, (-)}$.
For each complex $M \in \dcat{C}(A \ot B^{\mrm{op}}, \mrm{gr})$
we choose a K-graded-injective resolution 
$\rho_M : M \to I_M$ and a K-graded-flat resolution 
$\th_M : Q_M \to M$,
both in $\dcat{C}_{\mrm{str}}(A \ot B^{\mrm{op}}, \mrm{gr})$.
Note that $I_M$ is K-graded-$\m$-flasque over $A$, and $Q_M$ is K-graded-flat 
over $A$. We use these choices for presentations of the right derived functor 
\[ (\mrm{R} \Ga_{\m}, \eta^{\mrm{R}}) : 
\dcat{D}(A \ot B^{\mrm{op}}, \mrm{gr}) \to 
\dcat{D}(A \ot B^{\mrm{op}}, \mrm{gr}) \] 
and the left derived bifunctor 
\[ (- \ot_{A}^{\mrm{L}} -, \eta^{\mrm{L}}) : 
\dcat{D}(A^{\mrm{en}}, \mrm{gr}) \times 
\dcat{D}(A \ot B^{\mrm{op}}, \mrm{gr}) \to 
\dcat{D}(A \ot B^{\mrm{op}}, \mrm{gr}) . \]
Let us also choose a K-graded-injective resolution 
$\psi : A \to J$ in $\dcat{C}_{\mrm{str}}(A^{\mrm{en}}, \mrm{gr})$. 
With these choices we have the following presentations:
$P_A = \Ga_{\m}(J)$,
\begin{equation} \label{eqn:3910}
\mrm{R} \Ga_{\m}(M) = \Ga_{\m}(I_M)
\end{equation}
and
\begin{equation} \label{eqn:3911}
P_A \ot_A^{\mrm{L}} M = \Ga_{\m}(J) \ot_{A} Q_M .
\end{equation}

Suppose 
$N \in \dcat{C}(A \ot B^{\mrm{op}}, \mrm{gr})$
is some complex. Given homogeneous elements 
$x \in \Ga_{\m}(J)$ and $n \in N$, the tensor 
$x \ot n$ belongs to $\Ga_{\m}(J \ot_{A} N)$.
In this way we obtain a homomorphism 
\begin{equation} \label{eqn:3905} 
\ttev_{\m, J, N} : \Ga_{\m}(J) \ot_{A} N \to \Ga_{\m}(J 
\ot_{A} N) 
\end{equation}
in $\dcat{C}_{\mrm{str}}(A \ot B^{\mrm{op}}, \mrm{gr})$, 
which is functorial in $N$. 

Consider a complex $M \in \dcat{C}(A \ot B^{\mrm{op}}, \mrm{gr})$. The choices 
we made give rise to this solid diagram 
\begin{equation} \label{eqn:3906}  
\UseTips \xymatrix @C=8ex @R=6ex {
A \ot_{A} Q_M
\ar[r]^(0.6){\opn{lu}}
\ar[d]_{ \psi \, \ot_A \, \opn{id} }
&
Q_M
\ar[r]^{\th_{M}}
&
M
\ar[d]^{ \rho_M }
\\
J \ot_{A} Q_M 
\ar@{-->}[rr]^(0.6){\chi_M}
&
&
I_M
} 
\end{equation}
in $\dcat{C}_{\mrm{str}}(A \ot B^{\mrm{op}}, \mrm{gr})$.
The homomorphisms $\psi \, \ot_A \, \opn{id}$, $\th_{M} \circ \opn{lu}$  and
$\rho_M$ are quasi-isomorphisms.
Because $I_M$ is K-injective in $\dcat{C}(A \ot B^{\mrm{op}}, \mrm{gr})$, 
there is a quasi-iso\-morphism
$\chi_M : J \ot_{A} Q_M \to I_M$ 
that makes this diagram commutative up to homotopy. 

We now form diagram (\ref{eqn:3908})
in $\dcat{C}_{\mrm{str}}(A \ot B^{\mrm{op}}, \mrm{gr})$. 
%
\begin{figure}[!htb]
\begin{equation} \label{eqn:3908}  
 \UseTips \xymatrix @C=12ex @R=6ex {
\Ga_{\m}(J) \ot_{A} Q_M
\ar[r]^{ \ttev_{\m, J, Q_M} }
\ar[d]_{ \si_{J} \, \ot_A \, \opn{id} }
&
\Ga_{\m}(J \ot_{A} Q_M)
\ar[r]^(0.55){ \Ga_{\m}(\chi_M) } 
\ar[d]_{ \si_{(J \ot_{A} Q_M)} }
&
\Ga_{\m}(I_M)
\ar[d]_{ \si_{I_M} }
\\
J \ot_{A} Q_M
\ar[r]^{ \opn{id} }
&
J \ot_{A} Q_M
\ar[r]^(0.55){ \chi_M }
&
I_M
\\
A \ot_{A} Q_M 
\ar[u]^{ \psi \, \ot_A \, \opn{id} }
\ar[r]^{ \opn{id} }
&
A \ot_{A} Q_M 
\ar[u]^{ \psi \, \ot_A \, \opn{id} }
\ar[r]^(0.55){\th_{M} \circ \, \opn{lu}}
&
M
\ar[u]^{ \rho_M } 
} 
\end{equation}
\end{figure}
%
Here $\ttev_{\m, J, Q_M}$ is the homomorphism from (\ref{eqn:3905}). The 
diagram (\ref{eqn:3908}) is commutative up to homotopy. (Actually all small 
squares, except the bottom right one, are commutative in the strict sense.)
The vertical arrows $\psi \, \ot_A \, \opn{id}$ and $\rho_M$ are 
quasi-isomorphisms. Passing to $\dcat{D}(A \ot B^{\mrm{op}}, \mrm{gr})$
we get a commutative diagram, with vertical isomorphisms between the second and 
third rows. The diagram with the four extreme objects only is the one we are 
looking for. By construction it is a commutative diagram
in $\dcat{D}(A \ot B^{\mrm{op}}, \mrm{gr})$, 
and it is functorial in $M$. The morphism 
\[ \ttev^{\mrm{R, L}}_{\m, M} : P_A \ot_A^{\mrm{L}} M \to 
\mrm{R} \Ga_{\m}(M) \] 
that we want is represented by 
\[ \Ga_{\m}(\chi_M) \circ \ttev_{\m, J, Q_M}  : 
\Ga_{\m}(J) \ot_{A} Q_M \to \Ga_{\m}(I_M) . \]
See (\ref{eqn:3910}) and (\ref{eqn:3911}). 

\medskip \noindent
Step 2. In this step we prove that $\ttev^{\mrm{R, L}}_{\m, M}$
is an isomorphism for every complex 
$M \in \dcat{D}(A \ot B^{\mrm{op}}, \mrm{gr})$.
Because the functor $\opn{Rest}$ is conservative, it suffices to prove that the 
morphism 
\[ \opn{Rest}(\ttev^{\mrm{R, L}}_{\m, M}) : 
\opn{Rest}(P_A \ot_A^{\mrm{L}} M) \to \opn{Rest}(\mrm{R} \Ga_{\m}(M)) \]
in $\dcat{D}(A, \mrm{gr})$ is an isomorphism. Going over all the details of the 
construction above, and noting that 
$\rho_M : M \to I_M$ and $\th_M : Q_M \to M$ 
are K-flat and K-injective resolutions, respectively, also in the category
$\dcat{C}_{\mrm{str}}(A, \mrm{gr})$, we might as well forget about the ring 
$B$. 

So now we are in the case $B = \K$, $A \ot B^{\mrm{op}} = A$, and we want to 
prove that $\ttev^{\mrm{R, L}}_{\m, M}$ is an isomorphism for every 
$M \in \dcat{D}(A, \mrm{gr})$.
By Theorem \ref{thm:4050} the functor $\mrm{R} \Ga_{\m}$
on $\dcat{D}(A, \mrm{gr})$ is quasi-compact. The functor 
$P \ot_{A}^{\mrm{L}} (-)$ is also quasi-compact. This means that we can use 
Lemma \ref{lem:3905}, and it tells us 
that it suffices to prove that $\ttev^{\mrm{R, L}}_{\m, M}$ is an 
isomorphism for $M = A$. 

Let us examine the morphism $\ttev^{\mrm{R, L}}_{\m, A}$, 
i.e.\ $\ttev^{\mrm{R, L}}_{\m, M}$ for $M = A$. 
We can choose the K-flat resolution
$\th_A : Q_A \to A$ in $\dcat{C}_{\mrm{str}}(A, \mrm{gr})$ to be the identity 
of 
$A$. Also, we can choose the K-injective resolution 
$\rho_A : A \to I_A$ in $\dcat{C}_{\mrm{str}}(A, \mrm{gr})$ to be the 
restriction of $\psi : A \to J$. Then the homomorphism 
$\chi_A : J \ot_A Q_A \to I_A$
in diagram (\ref{eqn:3906}) can be chosen to be 
$\chi_A = \opn{id} \ot_A \opn{id}$. 
We get a commutative diagram 
\[ \UseTips \xymatrix @C=12ex @R=6ex {
\Ga_{\m}(J) \ot_{A} Q_A
\ar[r]^{ \ttev_{\m, J, Q_A} }
\ar[d]_{ \opn{id} \ot_A \th_A }
&
\Ga_{\m}(J \ot_{A} Q_A)
\ar[r]^(0.6){ \Ga_{\m}(\chi_A) } 
\ar[d]_{  \Ga_{\m}(\opn{id} \ot_A \th_A) }
&
\Ga_{\m}(I_A)
\ar[d]_{\opn{id}}
\\
\Ga_{\m}(J) \ot_{A} A
\ar[r]^{ \ttev_{\m, J, A} }
&
\Ga_{\m}(J \ot_{A} A)
\ar[r]^(0.6){\Ga_{\m}(\opn{ru})} 
&
\Ga_{\m}(J)
} \]
in $\dcat{C}_{\mrm{str}}(A, \mrm{gr})$. 
Here $\opn{ru} : J \ot_A A \iso J$ is the canonical isomorphism (the right 
unitor). The horizontal arrows in the second row, and the vertical arrows, are 
all bijective. We conclude that 
$\Ga_{\m}(\chi_A) \circ \ttev_{\m, J, Q_A}$
is bijective. Hence $\ttev^{\mrm{R, L}}_{\m, A}$ 
is an isomorphism in $\dcat{D}(A, \mrm{gr})$.


\medskip \noindent
Step 3. It remains to prove the uniqueness of 
$\ttev^{\mrm{R, L}}_{\m, (-)}$. By applying the functor 
$\mrm{R} \Ga_{\m}$ to the diagram ($\heartsuit$), and then making use of the 
morphism of functors
$\si^{\mrm{R}} : \mrm{R} \Ga_{\m} \to \opn{Id}$, 
we obtain the commutative diagram (\ref{eqn:4820})
in $\dcat{D}(A \ot B^{\mrm{op}}, \mrm{gr})$. 
%
\begin{figure}[!htb]
\begin{equation} \label{eqn:4820}
\UseTips \xymatrix @C=15ex @R=6ex {
P_A \ot^{\mrm{L}}_{A} M
\ar[r]^{ \ttev^{\mrm{R, L}}_{\m, M} }_{\cong}
& 
\mrm{R} \Ga_{\m}(M)
\\
\mrm{R} \Ga_{\m}(P_A \ot^{\mrm{L}}_{A} M)
\ar[u]^{\si^{\mrm{R}}_{(P_A \ot^{\mrm{L}}_{A} M)}}
\ar[r]^{ \mrm{R} \Ga_{\m}(\ttev^{\mrm{R, L}}_{\m, M}) }_{\cong}
\ar[d]_{\mrm{R} \Ga_{\m}(\si^{\mrm{R}}_{A} \, \ot^{\mrm{L}}_{A} \, \opn{id}_M)}
& 
\mrm{R} \Ga_{\m}(\mrm{R} \Ga_{\m}(M))
\ar[u]_{\si^{\mrm{R}}_{\mrm{R} \Ga_{\m}(M)}}
\ar[d]^{\mrm{R} \Ga_{\m}(\si^{\mrm{R}}_{M})}
\\
\mrm{R} \Ga_{\m}(A \ot^{\mrm{L}}_{A} M)
\ar[r]^{\mrm{R} \Ga_{\m}(\opn{lu})}_{\cong}
& 
\mrm{R} \Ga_{\m}(M)
}
\end{equation}
\end{figure}
%
The vertical morphisms 
$\si^{\mrm{R}}_{\mrm{R} \Ga_{\m}(M)}$
and 
$\mrm{R} \Ga_{\m}(\si^{\mrm{R}}_{M})$
are isomorphisms because of the idempotence
(Corollary \ref{cor:3966}).
Therefore the other two vertical arrows are isomorphisms. We see that the 
isomorphism $\ttev^{\mrm{R, L}}_{\m, M}$ can be expressed as the 
composition of other isomorphisms; so it is unique. 
\end{proof}

The next result is due to Van den Bergh \cite{VdB}. We give another 
proof, based on Theorem \ref{thm:3750}. 
Recall the $\K$-linear duality functor 
\[ (-)^* : \dcat{D}(A \ot B^{\mrm{op}}, \mrm{gr})^{\mrm{op}} \to 
\dcat{D}(B \ot A^{\mrm{op}}, \mrm{gr}) . \]
Its formula is  
\[  (-)^* := \opn{Hom}_{\K}(-, \K) \cong \opn{Hom}_{A}(-, A^*) . \]

\begin{cor}[Van den Bergh's Local Duality, \cite{VdB}] \label{cor:3901}
\index{Local Duality Theorem}
Under Setup \tup{\ref{setup:3900}}, for every complex 
$M \in \dcat{D}(A \ot B^{\mrm{op}},  \mrm{gr})$ there is an 
isomorphism 
\[ \mrm{R} \Ga_{\m}(M)^* \cong 
\opn{RHom}_A \bigl( M, (P_A)^* \bigr) \]
in $\dcat{D}(B \ot A^{\mrm{op}}, \mrm{gr})$.
This isomorphism is functorial in $M$. 
\end{cor}

\begin{proof}
We have the following functorial isomorphisms
\[ \begin{aligned}
& \mrm{R} \Ga_{\m}(M)^* \cong^{\dag} (P_A \ot^{\mrm{L}}_{A} M)^* 
= \opn{Hom}_{\K}(P_A \ot^{\mrm{L}}_{A} M, \K )
\\ & 
\quad \cong^{} \opn{RHom}_A(P_A \ot^{\mrm{L}}_{A} M, A^* )
\\ & 
\quad \cong^{\heartsuit} \opn{RHom}_A \bigl( M, 
\opn{RHom}_A(P_A , A^* ) \bigr) 
\\ & 
\quad \cong^{} \opn{RHom}_A \bigl( M, (P_A)^* \bigr) 
\end{aligned} \]
in $\dcat{D}(B \ot A^{\mrm{op}}, \mrm{gr})$.
The isomorphism $\cong^{\dag}$ is according to Theorem \ref{thm:3750}.
The isomorphism $\cong^{\heartsuit}$ comes from derived Hom-tensor adjunction.
\end{proof}




\mysubsection{Symmetry of Derived Torsion} 
\label{subsec:sym-der-tor}

The results of this subsection are an adaptation of results from 
\cite[Section 8]{VyYe} to the connected graded setting. 
Among other things, we prove that the $\chi$ condition of Artin-Zhang implies 
symmetry of derived torsion. 

In Section \ref{sec:BDC} below we are going to use the symmetry of derived 
torsion to give an alternative proof of Van den Bergh's theorem on existence of 
balanced dualizing complexes (see Theorem \ref{thm:3713}). 

We continue with Conventions \ref{conv:3700}  and \ref{conv:4560}.
Thus $\K$ is a field, and all graded rings are algebraically graded central 
$\K$-rings. The following setup is assumed in this subsection: 

\begin{setup} \label{setup:3855} 
We are given a noetherian connected graded $\K$-ring $A$, with augmentation 
ideal $\m$. The augmentation ideal of the opposite ring $A^{\mrm{op}}$ is 
$\m^{\mrm{op}}$.
\end{setup}

The assumption that $A$ is noetherian means that   
both $A$ and $A^{\mrm{op}}$ are left noetherian. 
The enveloping ring of $A$ is 
$A^{\mrm{en}} = A \ot A^{\mrm{op}}$, and it is often not noetherian.

As explained in the previous subsection, there is a torsion functor 
%\begin{equation} \label{eqn:3815}
$\Ga_{\m} : \dcat{M}(A, \mrm{gr}) \lb \to \dcat{M}(A, \mrm{gr})$.
%\end{equation}
But by the same token, there is also a torsion functor 
%\begin{equation} \label{eqn:3816}
$\Ga_{\m^{\mrm{op}}} : \lb \dcat{M}(A^{\mrm{op}}, \mrm{gr}) \to 
\dcat{M}(A^{\mrm{op}}, \mrm{gr})$.
%\end{equation}
On the category $\dcat{M}(A^{\mrm{en}}, \mrm{gr})$ we have two 
torsion functors:
%\begin{equation} \label{eqn:3811}
\[ \Ga_{\m}, \Ga_{\m^{\mrm{op}}} : \dcat{M}(A^{\mrm{en}}, \mrm{gr}) \to 
\dcat{M}(A^{\mrm{en}}, \mrm{gr}) . \]
%\end{equation}
These functors commute with each other. Indeed, for any 
$M \in \dcat{M}(A^{\mrm{en}}, \mrm{gr})$ there is equality 
%\begin{equation} \label{eqn:3812}
$\Ga_{\m}(\Ga_{\m^{\mrm{op}}}(M)) = \Ga_{\m^{\mrm{op}}}(\Ga_{\m}(M))$
%\end{equation}
of graded $A^{\mrm{en}}$-submodules of $M$. 

\begin{dfn} \label{dfn:3766} 
Under Setup \tup{\ref{setup:3855}}, we say that $A$ has 
{\em finite local cohomological dimension}% 
\index{Algebraically graded ring! connected {\indash} of finite local 
cohomological dimension}
if the torsion functors $\Ga_{\m}$ and $\Ga_{\m^{\mrm{op}}}$,
on $\dcat{M}(A, \mrm{gr})$ and 
$\dcat{M}(A^{\mrm{op}}, \mrm{gr})$ respectively, have finite right 
cohomological dimensions, 
in the sense of Definition \ref{dfn:3885}. Namely if there is a natural number 
$d$ such that $\mrm{H}^p_{\m} = 0$ and $\mrm{H}^p_{\m^{\mrm{op}}} = 0$ 
for all $p > d$. See also formula (\ref{eqn:3962}). 
\end{dfn}

Note that the vanishing of the derived functors  
$\mrm{H}^p_{\m} = \mrm{R}^p \Ga_{\m}$, and thus the cohomological dimension of 
the functor $\Ga_{\m}$, is the same on bimodules and on modules, by Proposition 
\ref{prop:3850}(4). Likewise for $\mrm{H}^p_{\m^{\mrm{op}}}$.

As explained in Proposition \ref{prop:3850}(1), there are derived torsion 
functors 
\begin{equation} \label{eqn:3817}
\mrm{R} \Ga_{\m}, \mrm{R} \Ga_{\m^{\mrm{op}}} : 
\dcat{D}(A^{\mrm{en}}, \mrm{gr}) \to \dcat{D}(A^{\mrm{en}}, \mrm{gr}) .
\end{equation}
There is no reason for these derived functors to commute with each other in 
general. 

\begin{dfn} \label{dfn:3717} 
Under Setup \tup{\ref{setup:3855}}, the morphisms from Proposition \lb 
\ref{prop:3850}(2) associated to the derived functors 
in formula (\ref{eqn:3817}) are denoted by%
\index{1-RGamma-m@$\mrm{R} \Ga_{\m}$}% 
\index{1-SiRm@$\si^{\mrm{R}}_{\m}$}
$\si^{\mrm{R}}_{\m} : \mrm{R} \Ga_{\m} \to \opn{Id}$
and 
$\si^{\mrm{R}}_{\m^{\mrm{op}}} : \mrm{R} \Ga_{\m^{\mrm{op}}} \to \opn{Id}$.
\end{dfn}

\begin{dfn} \label{dfn:3815} 
Under Setup \tup{\ref{setup:3855}}: 
\begin{enumerate}
\item We denote by 
$\dcat{M}_{(\mrm{tor}, ..)}(A^{\mrm{en}})$ the 
full subcategory of $\dcat{M}(A^{\mrm{en}})$ on the bimodules $M$ that are 
$\m$-torsion as $A$-modules; i.e.\ 
such that $\opn{Rest}_A(M) \in \dcat{M}_{\mrm{tor}}(A)$.

\item We denote by 
$\dcat{M}_{(.., \mrm{tor})}(A^{\mrm{en}})$ the 
full subcategory of $\dcat{M}(A^{\mrm{en}})$ on the bimodules $M$ that are 
$\m^{\mrm{op}}$-torsion as $A^{\mrm{op}}$-modules.

\item We let 
\[ \dcat{M}_{(\mrm{tor}, \mrm{tor})}(A^{\mrm{en}}) := 
\dcat{M}_{(\mrm{tor}, ..)}(A^{\mrm{en}}) \cap 
\dcat{M}_{(.., \mrm{tor})}(A^{\mrm{en}}) . \]

\item Let $\star, \diamond$ be torsion indicators, i.e.\ ``$\mrm{tor}$'' or 
``$..$''. We denote by 
$\dcat{D}_{(\star, \diamond)}(A^{\mrm{en}})$
the full subcategory of 
$\dcat{D}(A^{\mrm{en}})$ on the complexes $M$ such that 
$\opn{H}^p(M) \in \dcat{M}_{(\star, \diamond)}(A^{\mrm{en}})$ for 
every integer $p$.
\end{enumerate}
\end{dfn}

Trivially, for $M \in \dcat{D}(A^{\mrm{en}}, \mrm{gr})$ 
the $A$-modules $\opn{H}^p(\mrm{R} \Ga_{\m}(M))$ are $\m$-torsion, 
and the $A^{\mrm{op}}$-modules 
$\opn{H}^p(\mrm{R} \Ga_{\m^{\mrm{op}}}(M))$ are $\m^{\mrm{op}}$-torsion. 
Sometimes more happens:

\begin{dfn} \label{dfn:3720}
Under Setup \tup{\ref{setup:3855}}, a
complex $M \in \dcat{D}(A^{\mrm{en}}, 
\mrm{gr})$ is said to have 
{\em weakly symmetric derived $\m$-torsion}% 
\index{Complex of algebraically graded modules! with weakly symmetric derived 
$\m$-torsion}
if these two conditions hold:
\begin{itemize}
\item For every $p \in \Z$ the bimodule  $\opn{H}^p(\mrm{R} \Ga_{\m}(M))$ 
is $\m^{\mrm{op}}$-torsion.

\item For every $p \in \Z$ the bimodule 
$\opn{H}^p(\mrm{R} \Ga_{\m^{\mrm{op}}}(M))$ is $\m$-torsion.
\end{itemize} 
\end{dfn}

In terms of Definition \ref{dfn:3815}, a complex 
$M \in \dcat{D}(A^{\mrm{en}}, \mrm{gr})$
has weakly symmetric derived $\m$-torsion if 
\index{1-Dtortor(A)@$\dcat{D}_{(\mrm{tor}, \mrm{tor})}(A^{\mrm{en}}, \mrm{gr})$}
\[ \mrm{R} \Ga_{\m}(M) \, , \, \mrm{R} \Ga_{\m^{\mrm{op}}}(M) \, \in \,
\dcat{D}_{(\mrm{tor}, \mrm{tor})}(A^{\mrm{en}}, \mrm{gr}) . \]

\begin{dfn} \label{dfn:3779}
Under Setup \tup{\ref{setup:3855}}, a complex 
$M \in \dcat{D}(A^{\mrm{en}}, \mrm{gr})$ is said to have 
{\em symmetric derived $\m$-torsion}% 
\index{Complex of algebraically graded modules! with symmetric derived 
$\m$-torsion}
if there is an isomorphism 
\[ \ep_{M} : \mrm{R} \Ga_{\m}(M) \iso 
\mrm{R} \Ga_{\m^{\mrm{op}}}(M) \]
in $\dcat{D}(A^{\mrm{en}}, \mrm{gr})$, 
such that the diagram
\[ \UseTips \xymatrix @C=8ex @R=6ex {
\mrm{R} \Ga_{\m}(M)
\ar[r]^{\ep_{M}}_{\cong}
\ar[dr]_{\si^{\mrm{R}}_{\m, M}}
&
\mrm{R} \Ga_{\m^{\mrm{op}}}(M)
\ar[d]^{\si^{\mrm{R}}_{\m^{\mrm{op}}, M}}
\\
&
M
} \]
in $\dcat{D}(A^{\mrm{en}}, \mrm{gr})$ is commutative.
Such an isomorphism $\ep_M$ is called a 
{\em symmetry isomorphism}% 
\index{Complex of algebraically graded modules! symmetry isomorphism of}.
\end{dfn}

Of course if $M$ has symmetric derived $\m$-torsion, then it has 
weakly symmetric derived $\m$-torsion.

\begin{thm}[Symmetric Derived Torsion, \cite{VyYe}] 
\label{thm:3720}
\index{Algebraically graded ring! connected {\indash} of finite local 
cohomological dimension}%
\index{Complex of algebraically graded modules! with weakly symmetric derived 
$\m$-torsion}%
\index{Complex of algebraically graded modules! with symmetric derived 
$\m$-torsion}
Let $A$ be a noetherian connected graded $\K$-ring of finite local
cohomological dimension. 
If $M \in \dcat{D}(A^{\mrm{en}}, \mrm{gr})$
has weakly symmetric derived $\m$-torsion, then $M$ has 
symmetric derived $\m$-torsion.

Moreover, the symmetry isomorphism $\ep_{M}$ is unique, and it is functorial in 
such complexes $M$. 
\end{thm}

\begin{proof}
In Theorem \ref{thm:3750}, with $B := A$,
we have an isomorphism 
\[ \ttev^{\mrm{R, L}}_{\m, (-)} \, : 
\, P_A \ot^{\mrm{L}}_{A} (-) \, \iso \, \mrm{R} \Ga_{\m}(-) \]
of triangulated functors from 
$\dcat{D}(A^{\mrm{en}}, \mrm{gr})$ to itself. The same theorem, but with the 
roles of $A$ and $A^{\mrm{op}}$ exchanged, gives an isomorphism 
\[ \ttev^{\mrm{R, L}}_{\m^{\mrm{op}}, (-)}  
\, : \, (-) \ot^{\mrm{L}}_{A} P_{A^{\mrm{op}}} \, 
\iso \, \mrm{R} \Ga_{\m^{\mrm{op}}}(-) \]
of triangulated functors from 
$\dcat{D}(A^{\mrm{en}}, \mrm{gr})$ to itself. Here 
$P_{A^{\mrm{op}}} := \mrm{R} \Ga_{\m^{\mrm{op}}}(A) \in 
\dcat{D}(A^{\mrm{en}}, \mrm{gr})$.

Consider diagram (\ref{eqn:3871})  in $\dcat{D}(A^{\mrm{en}}, \mrm{gr})$.
In this diagram, $\al$ is the associativity isomorphism of the derived tensor 
product, and $\opn{ru}, \opn{lu}$ are the monoidal unitor isomorphisms -- all 
belonging to the monoidal structure of $\dcat{D}(A^{\mrm{en}}, \mrm{gr})$.
A quick check, using K-flat resolutions and elements in them, shows that this 
is a commutative diagram. 

\begin{figure}[!htb]
\begin{equation} \label{eqn:3871} 
\UseTips \xymatrix @C=8ex @R=6ex {
(P_A \ot^{\mrm{L}}_{A} M) \ot^{\mrm{L}}_{A} P_{A^{\mrm{op}}}
\ar[d]_{ \opn{id}  \ot^{\mrm{L}}_{A} \si^{\mrm{R}}_{\m^{\mrm{op}}, A} }
\ar[r]^{\al}_{\cong}
&
P_A \ot^{\mrm{L}}_{A} (M \ot^{\mrm{L}}_{A} P_{A^{\mrm{op}}})
\ar[d]^{ \si^{\mrm{R}}_{\m, A}  \ot^{\mrm{L}}_{A} \opn{id} }
\\
(P_A \ot^{\mrm{L}}_{A} M) \ot^{\mrm{L}}_{A} A
\ar[d]_{\opn{ru}}
&
A \ot^{\mrm{L}}_{A} (M \ot^{\mrm{L}}_{A} P_{A^{\mrm{op}}})
\ar[d]^{\opn{lu}}
\\
P_A \ot^{\mrm{L}}_{A} M 
\ar[dr]_{ \si^{\mrm{R}}_{\m, A}  \ot^{\mrm{L}}_{A} \opn{id} \ }
&
M \ot^{\mrm{L}}_{A} P_{A^{\mrm{op}}}
\ar[d]^{ \opn{id} \ot^{\mrm{L}}_{A} \si^{\mrm{R}}_{\m^{\mrm{op}}, A} }
\\
&
M
} 
\end{equation}
\end{figure}

By Theorem \ref{thm:3750}, diagram (\ref{eqn:3871}) gives rise to 
diagram (\ref{eqn:3872}), with its solid arrows only. 
This new diagram in $\dcat{D}(A^{\mrm{en}}, \mrm{gr})$ (solid arrows 
only) is isomorphic to the previous one (without its second row), via the 
various isomorphisms 
$\ttev^{\mrm{R, L}}_{\m, (-)}$ and 
$\ttev^{\mrm{R, L}}_{\m^{\mrm{op}}, (-)}$.

\begin{figure}[!htb]
\begin{equation} \label{eqn:3872}
\UseTips \xymatrix @C=8ex @R=6ex {
\mrm{R} \Ga_{\m^{\mrm{op}}}(\mrm{R} \Ga_{\m}(M))
\ar[r]^{\al'}_{\cong}
\ar[d]_{ \si^{\mrm{R}}_{\m^{\mrm{op}}, \mrm{R} \Ga_{\m}(M)} }
&
\mrm{R} \Ga_{\m}(\mrm{R} \Ga_{\m^{\mrm{op}}}(M))
\ar[d]^{ \si^{\mrm{R}}_{\m, \mrm{R} \Ga_{\m^{\mrm{op}}}(M)} }
\\
\mrm{R} \Ga_{\m}(M)
\ar[dr]_{ \si^{\mrm{R}}_{\m, M} }
\ar@{-->}[r]^{\ep_M}
&
\mrm{R} \Ga_{\m^{\mrm{op}}}(M)
\ar[d]^{ \si^{\mrm{R}}_{\m^{\mrm{op}}, M} }
\\
&
M
} 
\end{equation}
\end{figure}

Because 
$\mrm{R} \Ga_{\m}(M) \in 
\dcat{D}_{(\mrm{tor}, \mrm{tor})}(A^{\mrm{en}}, \mrm{gr})$,
Proposition \ref{prop:3965}, applied with $A^{\mrm{op}}$ instead 
of with $A$, says that 
$\si^{\mrm{R}}_{\m^{\mrm{op}}, \mrm{R} \Ga_{\m}(M)}$
is an isomorphism. Likewise, because 
$\mrm{R} \Ga_{\m^{\mrm{op}}}(M) \in 
\dcat{D}_{(\mrm{tor}, \mrm{tor})}(A^{\mrm{en}}, \mrm{gr})$,
Proposition \ref{prop:3965} says that 
$\si^{\mrm{R}}_{\m, \mrm{R} \Ga_{\m^{\mrm{op}}}(M)}$
is an isomorphism. We define $\ep_M$ to be the unique isomorphism (the dashed 
arrow) that makes diagram (\ref{eqn:3872}) commutative.

Diagram (\ref{eqn:3872}) shows that $\ep_M$ is unique: it has to be 
\[ \ep_M = \si^{\mrm{R}}_{\m, \mrm{R} \Ga_{\m^{\mrm{op}}}(M)} \circ 
\al' \circ (\si^{\mrm{R}}_{\m^{\mrm{op}}, \mrm{R} \Ga_{\m}(M)})^{-1} . \]
The functoriality of $\ep_M$ is a consequence of the functoriality of diagram 
(\ref{eqn:3872}). 
\end{proof}

\begin{cor} \label{cor:3875}
Under the assumptions of Theorem \tup{\ref{thm:3720}}, if the bimodule $A$ has 
has weakly symmetric derived $\m$-torsion, then there is a unique isomorphism 
$\ep_A : P_A \iso P_{A^{\mrm{op}}}$
in $\dcat{D}(A^{\mrm{en}}, \mrm{gr})$
such that 
$\si^{\mrm{R}}_{\m, A} = \si^{\mrm{R}}_{\m^{\mrm{op}}, A} \circ \ep_A$
as morphisms $P_A \to A$. 
\end{cor}

\begin{proof}
Take $M = A$ in the theorem. 
\end{proof}

The socle functor from Definition \ref{dfn:4265} extends to complexes. 

\begin{lem} \label{lem:3840}
\index{Complex of algebraically graded modules! minimal {\indash}
graded-injective {\indash}}
Let $I$ be a minimal complex of graded-injective $A$-modules, and let 
$W := \opn{Soc}(I)$. 
Then the differential of the complex $W$ is zero. 
\end{lem}

\begin{proof}
Take some nonzero homogeneous element $w \in W^p$, say $w \in W^p_q$. 
Let $U := \K \cd w \sub W^p_q$, so $U$ is a nonzero graded 
$A$-submodule of $I^p$. Because $\opn{Z}^p(I)$ is an essential graded 
$A$-submodule of $I^p$ we must have 
$\opn{Z}^p(I) \cap U \neq 0$. Thus $w \in \opn{Z}^p(I)$, so 
$\d(w) = 0$.
\end{proof}

Cofinite graded $A$-modules were defined in Definition \ref{dfn:3726}. 
By graded Matlis duality (Theorem \ref{thm:4045}) the cofinite graded 
$A$-modules are the artinian objects of $\dcat{M}(A, \mrm{gr})$. 
The full subcategory $\dcat{M}_{\mrm{cof}}(A, \mrm{gr})$ of cofinite modules 
is a thick abelian subcategory of $\dcat{M}(A, \mrm{gr})$, closed under 
subobjects and quotients.

\begin{lem} \label{lem:3856}
Consider a complex $M \in \dcat{D}^+(A, \mrm{gr})$ and an integer $p$. 
\begin{enumerate}
\item If $\opn{Ext}^p_A(\K, M)$ is a finite graded $\K$-module, then 
$\mrm{H}^p_{\m}(M)$ is a cofinite graded $A$-module. 

\item If $\mrm{H}^{q}_{\m}(M)$ is a cofinite graded $A$-module for every 
$q \leq p$, then $\opn{Ext}^{q}_A(\K, M)$ is a finite graded $\K$-module
for every $q \leq p$.
\end{enumerate}
\end{lem}

\begin{proof} \mbox{}

\smallskip \noindent
(1) Let $M \to I$ be a minimal graded-injective resolution, and let 
$W := \opn{Soc}(I)$. Then $\opn{Ext}^p_A(\K, M) \cong \opn{H}^p(W)$. 
But by Lemma \ref{lem:3840} we know that the differential of $W$ is zero, so in 
fact $\opn{Ext}^p_A(\K, M) \cong W^p$. Therefore $W$ is a finite graded 
$\K$-module. 

Next, let $J := \Ga_{\m}(I)$. This complex has the property that 
$\mrm{H}^p_{\m}(M) \cong \mrm{H}^p(J)$. According to Lemma \ref{lem:3839} we 
have an isomorphism $J^p \cong A^* \ot W^p$. 
Since $A^*$ is a cofinite graded $A$-module, and $J^p$ is a finite direct sum 
of degree twists of $A^*$, it is also a cofinite graded $A$-module. But 
$\mrm{H}^p_{\m}(M)$ is a subquotient of $J^p$, so it too is a cofinite graded 
$A$-module. 

\medskip \noindent 
(2) We continue with the minimal graded-injective resolution $M \to I$. 
We may assume that $\opn{H}(M) \neq 0$. 
Let $p_0 := \opn{inf}(\opn{H}(M))$; so 
$p_0 = \opn{inf}(I)$. Take the subcomplexes $W \sub J \sub I$ as above. We know 
that $J^q \cong A^* \ot W^q$ for every $q$. We will prove that 
$W^q$ is a finite graded $\K$-module for every $q \leq p$, by induction on $q$, 
starting with $q = p_0$. 

So take an integer $q$ in the range $p_0 \leq q \leq p$, and assume that 
$W^{q - 1}$ is a finite graded $\K$-module. 
There is an exact sequence 
\[ J^{q - 1} \xar{\d} \opn{Z}^q(J) \to \mrm{H}^q_{\m}(M) \to  0  \]
in $\dcat{M}(A, \mrm{gr})$. 
Since $J^{q - 1} \cong A^* \ot W^{q - 1}$, this is a cofinite graded 
$A$-module, as we already noticed. We also know that $\mrm{H}^q_{\m}(M)$ is 
cofinite. It follows that $\opn{Z}^q(J)$ is a cofinite graded $A$-module.
But by Lemma \ref{lem:3840} we have 
$W^q \sub \opn{Z}^q(J)$, so $W^q$ is a cofinite graded $A$-module annihilated 
by $\m$. Therefore $W^q$ is a finite graded $\K$-module.
\end{proof}

\begin{dfn}[Artin-Zhang, \cite{ArZh}] \label{dfn:3713}
Assume Setup \tup{\ref{setup:3855}}.
\begin{enumerate}
\item We say that the ring $A$ satisfies the {\em left $\chi$ condition} 
if for every $M \in \dcat{M}_{\mrm{f}}(A, \mrm{gr})$
and every integer $p$, the graded $\K$-module 
$\opn{Ext}^p_{A}(\K, M)$ is finite. 

\item The ring $A$ is said to satisfy the 
{\em $\chi$ condition}% 
\index{Algebraically graded ring! chi@$\chi$ condition on connected}
if both graded rings $A$ and $A^{\mrm{op}}$ satisfy the left $\chi$ condition.
\end{enumerate}
\end{dfn}

Another way to state the left $\chi$ condition is this: 
$M \in \dcat{M}_{\mrm{f}}(A, \mrm{gr})$ implies 
$\opn{RHom}_{A}(\K, M) \in \dcat{D}_{\mrm{f}}(\K, \mrm{gr})$.

\begin{rem} \label{rem:3855}
What we call the ``left $\chi$ condition'' was actually called the ``left 
$\chi^{\circ}$ condition'' in \cite[Definition 3.2]{ArZh}. 
Their ``left $\chi$ condition'' 
(in \cite[Definition 3.7]{ArZh}) is much more complicated to state. 
However, for a left noetherian connected graded ring $A$ (as we have here) 
these two conditions are equivalent, by \cite[Proposition 3.11(2)]{ArZh}. 
\end{rem}

\begin{dfn} \label{dfn:3711} 
Assume Setup \tup{\ref{setup:3855} }.
\begin{enumerate}
\item We say that the ring $A$ satisfies the {\em left special $\chi$ 
condition} if for every integer $p$, the graded $\K$-module 
$\opn{Ext}^p_{A}(\K, A)$ is finite. 

\item The ring $A$ is said to satisfy the 
{\em special $\chi$ condition}% 
\index{Algebraically graded ring! special $\chi$ condition on connected}
if both graded rings $A$ and $A^{\mrm{op}}$ satisfy the left special $\chi$ 
condition.
\end{enumerate}
\end{dfn}

\begin{prop} \label{prop:3855}
The following two conditions are equivalent.
\begin{enumerate}
\rmitem{i} $A$ satisfies the left $\chi$ condition. 

\rmitem{ii} For every $M \in \dcat{M}_{\mrm{f}}(A, \mrm{gr})$ and every $p$ 
the graded $A$-module $\mrm{H}^p_{\m}(M)$ is cofinite. 
\end{enumerate}
\end{prop}

\begin{proof}
This is immediate from Lemma \ref{lem:3856}. 
\end{proof}

\begin{dfn} \label{dfn:4050} 
We denote by 
$\dcat{D}^{\mrm{b}}_{(\mrm{f}, \mrm{f})}(A^{\mrm{en}}, \mrm{gr})$ the 
full subcategory of \lb 
$\dcat{D}^{\mrm{b}}(A^{\mrm{en}}, \mrm{gr})$ on the 
complexes of bimodules $M$ whose cohomology bimodules $\opn{H}^p(M)$ are 
finite graded modules over $A$ and over $A^{\mrm{op}}$. 
\end{dfn}


\begin{prop} \label{prop:3712}
\mbox{}
\begin{enumerate}
\item If $A$ satisfies the $\chi$ condition, then it satisfies 
the special $\chi$ condition. 

\item If $A$ satisfies the $\chi$ condition, then every
$M \in \dcat{D}^{\mrm{b}}_{(\mrm{f}, \mrm{f})}(A^{\mrm{en}}, \mrm{gr})$
has weakly symmetric derived $\m$-torsion. 

\item If $A$ satisfies the special $\chi$ condition, 
then the bimodule $A$ has weakly symmetric derived $\m$-torsion. 
\end{enumerate}
\end{prop}

\begin{proof} \mbox{}

\smallskip \noindent 
(1) This is trivial.

\medskip \noindent
(2) Take some 
$M \in \dcat{D}^{\mrm{b}}_{(\mrm{f}, \mrm{f})}(A^{\mrm{en}}, \mrm{gr})$.
Using smart truncation and induction on the amplitude of cohomology, 
we can assume that 
$M \in \dcat{M}_{(\mrm{f}, \mrm{f})}(A^{\mrm{en}}, \mrm{gr})$.

The $\chi$ condition tells us that $\opn{Ext}^p_{A}(\K, M)$ and 
$\opn{Ext}^p_{A^{\mrm{op}}}(\K, M)$ 
are finite as graded $\K$-modules, for all $p$. For every $q \geq 1$ the graded 
bimodule $A / \m^q$ is gotten from $\K$ by finitely many degree twists and 
extensions in $\dcat{M}(A^{\mrm{en}}, \mrm{gr})$. Therefore 
$\opn{Ext}^p_{A}(A / \m^q, M)$ and 
$\opn{Ext}^p_{A^{\mrm{op}}}(A / \m^q, M)$
are finite as graded $\K$-modules, for all $p$ and all $q \geq 1$. 
Thus as left and as right graded $A$-modules, 
$\opn{Ext}^p_{A}(A / \m^q, M)$ and 
$\opn{Ext}^p_{A^{\mrm{op}}}(A / \m^q, M)$
are of finite length, and thus torsion on both sides. Passing to the direct 
limit using Proposition \ref{prop:3850}(5), we see that $\mrm{H}^p_{\m}(M)$ and 
$\mrm{H}^p_{\m^{\mrm{op}}}(M)$ are torsion $A$-modules on both sides. 

\medskip \noindent
(3) This is just a special case of the proof of item (2). 
\end{proof}

\begin{lem} \label{lem:3870}
Assume the functor $\Ga_{\m}$ has finite cohomological dimension, and  $A$ has 
the left special $\chi$ condition. Then for every
$M \in \dcat{D}_{\mrm{f}}(A, \mrm{gr})$
and every integer $i$, the  graded $A$-module $\opn{H}^i_{\m}(M)$ is cofinite. 
\end{lem}

\begin{proof}
Fix an integer $i$. We will prove that the graded $A$-module 
$\opn{H}^i_{\m}(M)$ is cofinite. Say $\Ga_{\m}$ has cohomological dimension 
$\leq d$ for some natural number $d$. 

\medskip \noindent
Step 1. Consider 
$M' := \opn{smt}^{\leq i + 1}(M)$,
the smart truncation of the complex $M$ below $i + 1$. 
There is a distinguished triangle 
$M' \to M \to M'' \xar{{\,} \triangle {\,}}$
in $\dcat{D}(A, \mrm{gr})$, 
and hence there is a distinguished triangle 
\[ \mrm{R} \Ga_{\m}(M') \to \mrm{R} \Ga_{\m}(M) \to 
\mrm{R} \Ga_{\m}(M'') \xar{ \, \triangle \, } . \]
Because $\mrm{H}^j_{\m}(M'') = 0$ for $j \leq i + 1$, we see that 
$\mrm{H}^i_{\m}(M) \cong \mrm{H}^i_{\m}(M')$. 
Therefore, by replacing $M$ with $M'$, we can assume that 
$M \in \dcat{D}^-_{\mrm{f}}(A, \mrm{gr})$.

\medskip \noindent
Step 2. Now $M \in \dcat{D}^-_{\mrm{f}}(A, \mrm{gr})$.
Choose a resolution $P \to M$, where $P$ is a bounded above complex of finite 
graded-free $A$-modules. Let 
$P' := \opn{stt}^{\geq i - d -2}(P)$, the stupid truncation of $P$ above 
$i - d -2$. Then 
$\mrm{H}^i_{\m}(M) \cong \opn{H}^i(P')$.
But $P'$ is gotten from $A$ by finitely many cones, finite direct sums, 
translations and degree twists. Since $A$ has the left special $\chi$ 
condition, and using Lemma \ref{lem:3856}(1), the graded $A$-module 
$\mrm{H}^j_{\m}(A)$ is cofinite for every $j$. Therefore  
$\mrm{H}^j_{\m}(P')$ is cofinite for every $j$; including $j = i$. 
\end{proof}

Here is a converse to the trivial implication in Proposition 
\ref{prop:3712}(1), under an extra finiteness condition. 

\begin{prop} \label{prop:3870}
If $A$ satisfies the special $\chi$ condition and has finite local 
cohomological 
dimension, then $A$ satisfies the $\chi$ condition. 
\end{prop}

\begin{proof}
This is an immediate consequence of Lemma \ref{lem:3870} and Proposition 
\ref{prop:3855}, applied to both $A$ and $A^{\mrm{op}}$. 
\end{proof}

AS regular graded rings were introduced in Definition \ref{dfn:3782}. 

\begin{cor} \label{cor:3870}
Let $A$ be an AS regular graded ring. 
Then $A$ satisfies the $\chi$ condition and has finite local 
cohomological dimension.
\end{cor}

\begin{proof}
The special $\chi$ condition is an immediate consequence of the AS condition
(Definition \ref{dfn:4545}). 
The graded global cohomological dimension of $A$ bounds its local cohomological 
dimension, so the latter is finite.
By Proposition \ref{prop:3870}, $A$ satisfies the $\chi$ condition.
\end{proof}

Actually this is true also for AS Gorenstein rings, but the proof is harder. 
See Corollary \ref{cor:3780} and Corollary \ref{cor:4655}.

The $\chi$ condition, and the finiteness of local cohomological dimension, 
occur in many important examples of noncommutative graded rings. 
The main results that ensure it involve a lifting argument -- 
Theorems \ref{thm:4546} and \ref{thm:4250}.
Here are several examples. 

\begin{exa} \label{exa:4081}
Let $A := \K [x_1, \ldots, x_n]$, the commutative polynomial ring in 
$n \geq 1$ variables from Example \ref{exa:3757}. 
This is an AS regular graded ring of dimension $n$.
By Corollary \ref{cor:3870} the graded ring $A$ has the $\chi$ condition and 
finite local cohomological dimension.
\end{exa}

\begin{exa} \label{exa:4082}
If $B$ is a commutative noetherian connected graded $\K$-ring, then there is a 
finite homomorphism $A \to B$ from a polynomial ring $A$ as in the previous 
example. By Corollaries \ref{cor:4030} and \ref{cor:4031} the ring $B$ 
satisfies the $\chi$ condition and has finite local cohomological dimension.
\end{exa}

\begin{exa} \label{exa:4080}
Let $A$ be the homogeneous Weyl algebra from Example \ref{exa:3756}.
The element $t \in A$ is central regular of degree $1$, and 
$A / (t) \cong \K[x, y]$,
the commutative polynomial ring in two variables of degree $1$. 
The ring $\K[x, y]$ is AS regular of dimension $2$, so by Theorem 
\ref{thm:4546} the ring $A$ is AS regular of dimension $3$.
By Corollary \ref{cor:3870} the graded ring $A$ has the $\chi$ condition and 
finite local cohomological dimension.

As noted before, in characteristic $0$ the ring $A$ is very noncommutative. 
\end{exa}

\begin{exa} \label{exa:4083}
Let $A$ be the homogeneous universal enveloping algebra from Example 
\ref{exa:4075}. The element $t \in A$ is central regular of degree $1$, and 
$A / (t) \cong \K[x_1, \ldots, x_n]$,
the commutative polynomial ring in $n$ variables. 
The polynomial ring is AS regular of dimension $n$, so 
by Theorem \ref{thm:4546} the ring $A$ is AS regular of dimension $n + 1$.
By Corollary \ref{cor:3870} the graded ring $A$ has the $\chi$ condition and 
finite local cohomological dimension.

As noted in Example \ref{exa:4075}, the ring $A$ could be very 
noncommutative, if $\g$ is very nonabelian (e.g.\ semisimple). 
\end{exa}

We end this subsection with a result that ties the $\chi$ condition to symmetry 
of derived torsion. Recall that 
$\dcat{D}_{(\mrm{f}, \mrm{f})}(A^{\mrm{en}}, \mrm{gr})$
is the full subcategory of 
$\dcat{D}(A^{\mrm{en}}, \mrm{gr})$
on the complexes $M$ such that for every $p$ the graded bimodule $\opn{H}^p(M)$ 
is finite over $A$ and $A^{\mrm{op}}$. 

\begin{thm} \label{thm:4051}
\index{Algebraically graded ring! connected {\indash} of finite local 
cohomological dimension}
\index{Algebraically graded ring! chi@$\chi$ condition on connected}
\index{Complex of algebraically graded modules! with symmetric derived 
$\m$-torsion}
Under Setup \tup{\ref{setup:3855}}, assume $A$ satisfies the $\chi$ condition 
and has finite local cohomological dimension. Then there is a unique 
isomorphism 
\[ \ep : \mrm{R} \Ga_{\m} \iso \mrm{R} \Ga_{\m^{\mrm{op}}} \]
of triangulated functors 
$\dcat{D}^{\mrm{b}}_{(\mrm{f}, \mrm{f})}(A^{\mrm{en}}, \mrm{gr}) \to 
\dcat{D}(A^{\mrm{en}}, \mrm{gr})$, 
such that the diagram 
\[ \UseTips \xymatrix @C=8ex @R=6ex {
\mrm{R} \Ga_{\m}
\ar[r]^{\ep}_{\cong}
\ar[dr]_{\si^{\mrm{R}}_{\m}}
&
\mrm{R} \Ga_{\m^{\mrm{op}}}
\ar[d]^{\si^{\mrm{R}}_{\m^{\mrm{op}}}}
\\
&
\opn{Id}
} \]
is commutative.
\end{thm}

\begin{proof}
According to Proposition \ref{prop:3712}(2) every complex 
$M \in \dcat{D}^{\mrm{b}}_{(\mrm{f}, \mrm{f})}(A^{\mrm{en}}, \mrm{gr})$
has weakly symmetric derived $\m$-torsion. By Theorem \ref{thm:3720} such a 
complex $M$ has symmetric derive torsion, and the symmetry isomorphism $\ep_M$ 
is unique and functorial. 
\end{proof}




\mysubsection{NC MGM Equivalence}
\label{subsec:NC-MGM}

The {\em commutative MGM Equivalence} was \lb explained in Example 
\ref{exa:4595}.
A noncommutative variant of this theory was recently developed by 
R. Vyas and A. Yekutieli \cite{VyYe}, and in this subsection we give an 
adaptation of it to the case of a NC connected graded ring $A$. 
In Section \ref{sec:BDC} below we shall use the NC MGM Equivalence to prove the
existence of balanced trace morphisms (Theorem \ref{thm:4030}), and in Section 
\ref{sec:rigid-DC-NC} the NC MGM Equivalence will help us prove that a balanced 
dualizing complex is graded rigid (Theorem \ref{thm:4400}). 

We continue with Conventions \ref{conv:3700} and \ref{conv:4560}, so $\K$ is a 
base field, and $A$ and $B$ are graded $\K$-rings.

Recall (from Remark \ref{rem:3450}) that the derived category 
$\dcat{D}(A^{\mrm{en}}, \mrm{gr})$ has a biclosed monoidal structure on it, 
with monoidal operation $(- \ot^{\mrm{L}}_{A} -)$ and monoidal unit $A$.
\index{Monoidal category} 
For complexes 
$M \in \dcat{D}(A \ot B^{\mrm{op}}, \mrm{gr})$
and 
$N \in \dcat{D}(B \ot A^{\mrm{op}}, \mrm{gr})$
the left and right unitor isomorphisms are%
\index{Monoidal category! left unitor isomorphism}% 
\index{1-Lu@$\opn{lu}$}%
\index{Monoidal category! right unitor isomorphism}% 
\index{1-Ru@$\opn{ru}$} 
$\opn{lu} : A \ot^{\mrm{L}}_{A} M \iso M$ 
and
$\opn{ru} : M \ot^{\mrm{L}}_{A} A \iso M$
respectively. The left co-unitor isomorphism%
\index{Monoidal category! left co-unitor isomorphism}% 
\index{1-Lcu@$\opn{lcu}$} 
is $\opn{lcu} : \opn{RHom}_{A}(A, M) \iso M$.

\begin{dfn}[\cite{VyYe}] \label{dfn:3995} \mbox{}
\begin{enumerate}
\item A {\em copointed object}% 
\index{Monoidal category! copointed object in}
in the monoidal category 
$\dcat{D}(A^{\mrm{en}}, \mrm{gr})$ 
is a pair $(P, \si)$, consisting of an object 
$P \in \dcat{D}(A^{\mrm{en}}, \mrm{gr})$ and a morphism 
$\si : P \to A$ in $\dcat{D}(A^{\mrm{en}}, \mrm{gr})$.

\item The copointed object $(P, \si)$ is called 
{\em idempotent}% 
\index{Monoidal category! idempotent copointed object in}
if the morphisms 
\[  \opn{lu} \circ \, (\si \ot^{\mrm{L}}_A \opn{id}) , \ 
\opn{ru} \circ \, (\opn{id}  \ot^{\mrm{L}}_A \, \si) : 
\ P \ot_A^{\mrm{L}} P \to P \]
in $\dcat{D}(A^{\mrm{en}}, \mrm{gr})$ are both isomorphisms.
\end{enumerate}
\end{dfn}

\begin{rem} \label{rem:3995}
Here is an explanation of the name ``copointed''. In a monoidal category, with 
unit object $A$, a {\em point} of an object $P$ is a morphism $A \to P$. 
It thus makes sense to refer to a morphism in the dual direction,  
say $\si : P \to A$, as a {\em copoint} of $P$. 
\end{rem}

\begin{dfn} \label{dfn:3996}
Let  $(P, \si)$ be a copointed object in $\dcat{D}(A^{\mrm{en}}, \mrm{gr})$. 
\begin{enumerate}
\item Define the triangulated functors 
\[ F , G : \dcat{D}(A \ot B^{\mrm{op}}, \mrm{gr}) \to 
\dcat{D}(A \ot B^{\mrm{op}}, \mrm{gr}) \]
to be
$F := P \ot_{A}^{\mrm{L}} (-)$
and
$G := \opn{RHom}_{A}(P, - )$.
 
\item Let $\si : F \to \opn{Id}$ 
and $\tau : \opn{Id} \to G$
be the morphisms of triangulated functors from 
$\dcat{D}(A \ot B^{\mrm{op}}, \mrm{gr})$ to itself 
that are induced by the morphism $\si : P \to A$. Namely
\[ \si_M  := \opn{lu} \circ \, (\si \ot_{A}^{\mrm{L}} \opn{id}_M) : 
F(M) = P \ot_{A}^{\mrm{L}} M \to M  \]
and 
\[ \tau_M := \opn{RHom}_A(\si, \opn{id}_M)  \circ \opn{lcu}^{-1} : 
M \to \opn{RHom}_{A}(P, M) = G(M)  . \]
\end{enumerate}
We refer to $(F, \si)$ and $(G, \tau)$ as the 
{\em copointed and pointed triangulated functors induced by the copointed 
object $(P, \si)$}%
\index{Additive functor! copointed}%
\index{Additive functor! pointed}. 
\end{dfn}

Item (2) of the definition is shown in the 
commutative diagrams below 
in the category $\dcat{D}(A \ot B^{\mrm{op}}, \mrm{gr})$. 

\[ \UseTips \xymatrix @C=8ex @R=6ex {
P \ot_{A}^{\mrm{L}} M
\ar[d]_{\si \ot_{A}^{\mrm{L}} \opn{id}_M}
\ar[dr]^{\si_M}
\\
A \ot_{A}^{\mrm{L}} M
\ar[r]_(0.6){\opn{lu}}
&
M
}
\qquad
\UseTips \xymatrix @C=8ex @R=8ex {
M
\ar[dr]_{\tau_M}
&
\opn{RHom}_{A}(A, M)
\ar[d]^{\opn{RHom}_A(\si, \opn{id}_M)}
\ar[l]_(0.6){\opn{lcu}}
\\
&
\opn{RHom}_{A}(P, M)
} \]

\begin{dfn} \label{dfn:3997}
Let $(F, \si)$ and $(G, \tau)$ be the copointed and pointed triangulated 
functors on $\dcat{D}(A \ot B^{\mrm{op}}, \mrm{gr})$ from Definition 
\tup{\ref{dfn:3996}}. 
\begin{enumerate}
\item Let $\dcat{D}(A \ot B^{\mrm{op}}, \mrm{gr})_F$ be the full 
triangulated subcategory of $\dcat{D}(A \ot B^{\mrm{op}}, \mrm{gr})$ 
on the set of objects  
\[ \bigr\{ M \mid \si_M : F(M) \to M \tup{ is an isomorphism} \, \bigl\} . \]

\item  Let $\dcat{D}(A \ot B^{\mrm{op}}, \mrm{gr})_G$ be the full 
triangulated subcategory of $\dcat{D}(A \ot B^{\mrm{op}}, \mrm{gr})$ 
on the set of objects 
\[ \bigr\{ M \mid \tau_M : M \to G(M) \tup{ is an isomorphism} \, \bigl\} . \]
\end{enumerate}
\end{dfn}

Idempotent (co)pointed triangulated functors were introduced in Definitions 
\ref{dfn:3896} and \ref{dfn:3897}.

\begin{lem} \label{lem:3995}
If the copointed object  $(P, \si)$ is idempotent, then the copointed 
triangulated functor $(F, \si)$ and the 
pointed triangulated functor $(G, \tau)$ on 
$\dcat{D}(A \ot B^{\mrm{op}}, \mrm{gr})$ are idempotent.
\end{lem}

\begin{proof}
For $M \in \dcat{D}(A \ot B^{\mrm{op}}, \mrm{gr})$ there are equalities
(up to the associativity isomorphism of $(- \ot_A^{\mrm{L}} -)$, 
that should be inserted in the locations marked by ``$\dag$''): 
\begin{equation} \label{eqn:4006}
F(\si_M) = \opn{id}_P \ot_A^{\mrm{L}} \, 
\bigl( \opn{lu} \circ \, (\si \ot_{A}^{\mrm{L}} \opn{id}_M) \bigr) 
=^{\dag} \, 
\bigl( \opn{ru} \circ \, (\opn{id}_P \ot_A^{\mrm{L}} \, \si) \bigr)
\ot_{A}^{\mrm{L}} \opn{id}_M 
\end{equation}
and 
\begin{equation} \label{eqn:3996}
\si_{F(M)} = \opn{lu} \circ \, \bigl( \si \ot_A^{\mrm{L}} 
(\opn{id}_P \ot_A^{\mrm{L}} \opn{id}_M) \bigr) 
=^{\dag} \, 
\bigl( \opn{lu} \circ (\si \ot_A^{\mrm{L}} \opn{id}_P) \bigr) \ot_A^{\mrm{L}} 
\opn{id}_M 
\end{equation}
of morphisms
\[ F(F(M)) = P \ot_A^{\mrm{L}} P \ot_A^{\mrm{L}} M \to 
F(M) =  P \ot_A^{\mrm{L}} M  \]
in $\dcat{D}(A \ot B^{\mrm{op}}, \mrm{gr})$.
Because both $\opn{lu} \circ (\si \ot_A^{\mrm{L}} \opn{id}_P)$ and 
$\opn{ru} \circ (\opn{id}_P \ot_A^{\mrm{L}} \, \si)$ 
are isomorphisms in $\cat{D}(A^{\mrm{en}})$, it follows that the morphisms
$F(\si_M)$ and $\si_{F(M)}$ are isomorphisms in 
$\cat{D}(A \ot B^{\mrm{op}})$. 

There are also  equalities (up to the associativity and adjunction isomorphisms 
of $(- \ot_A^{\mrm{L}} -)$ and $\opn{RHom}_A(-,-)$, 
that should be inserted in the locations marked by by ``$\ddag$''): 
\begin{equation} \label{eqn:3997}
\begin{aligned}
& G(\tau_M) = 
\opn{RHom}_A \bigl( \opn{id}_P,  \opn{RHom}_A(\si, \opn{id}_M) \bigr) 
\circ 
\opn{RHom}_A(\opn{id}_P, \opn{lcu}) \\
& \quad =^{\ddag} \, 
\opn{RHom}_A(\si \ot_A^{\mrm{L}} \opn{id}_P, \opn{id}_M) \circ  
\opn{RHom}_A(\opn{lu}, \opn{id}_M) 
\end{aligned}
\end{equation}
and 
\begin{equation} \label{eqn:3998}
\begin{aligned}
& \tau_{G(M)} = \opn{RHom}_A(\si, \opn{id}_{\opn{RHom}_A(P, M)}) \circ 
\opn{lcu} \\
& \quad =^{\ddag} \, 
\opn{RHom}_A(\opn{id}_P \ot_A^{\mrm{L}} \, \si, \opn{id}_M) 
\circ  \opn{RHom}_A(\opn{ru}, \opn{id}_M) 
\end{aligned}
\end{equation}
of morphisms
\[ \begin{aligned}
& G(M) = \opn{RHom}_A(P, M) \to G(G(M))
\\
&
\quad = \opn{RHom}_A \bigl( P, \opn{RHom}_A(P, M) \bigr) \cong
\opn{RHom}_A(P \ot_A^{\mrm{L}} P, M) 
\end{aligned}   \]
in $\dcat{D}(A \ot B^{\mrm{op}}, \mrm{gr})$. 
Because both 
$\opn{lu} \circ \, (\si \ot_A^{\mrm{L}} \opn{id}_P)$ and 
$\opn{ru} \circ \, (\opn{id}_P \ot_A^{\mrm{L}} \, \si)$ 
are isomorphisms  in $\dcat{D}(A^{\mrm{en}}, \mrm{gr})$, it follows that the 
morphisms $G(\tau_M)$ and $\tau_{G(M)}$ are isomorphisms in 
$\dcat{D}(A \ot B^{\mrm{op}}, \mrm{gr})$. 
\end{proof}

\begin{lem} \label{lem:3996}
Consider the functors $F$ and $G$ on $\dcat{D}(A \ot B^{\mrm{op}}, \mrm{gr})$ 
from Definition \tup{\ref{dfn:3996}}. For every 
$M, N \in \dcat{D}(A \ot B^{\mrm{op}}, \mrm{gr})$ there is a bijection 
\[ \opn{Hom}_{\dcat{D}(A \ot B^{\mrm{op}}, \mrm{gr})} \bigl( F(M), N \bigr) 
\cong 
\opn{Hom}_{\dcat{D}(A \ot B^{\mrm{op}}, \mrm{gr})} \bigl( M, G(N) \bigr) , \]
and it is functorial in $M$ and $N$.  
\end{lem}

\begin{proof}
Choose a K-injective resolution $N \to J$ in 
$\dcat{C}(A \ot B^{\mrm{op}}, \mrm{gr})$,
and a K-flat resolution $\til{P} \to P$ in $\dcat{C}(A^{\mrm{en}}, \mrm{gr})$.
The usual Hom-tensor adjunction gives rise to an isomorphism 
\begin{equation} \label{eqn:4000}
\opn{Hom}_{A \ot B^{\mrm{op}}}(\til{P} \ot_A M, J) \cong 
\opn{Hom}_{A \ot B^{\mrm{op}}} \bigl( M, \opn{Hom}_{A}(\til{P}, J) \bigr)
\end{equation}
in $\dcat{C}_{\mrm{str}}(\K, \mrm{gr})$. From this we deduce that 
$\opn{Hom}_{A}(\til{P}, J)$
is K-injective in \lb $\dcat{C}(A \ot B^{\mrm{op}}, \mrm{gr})$.
We see that the isomorphism (\ref{eqn:4000}) represents an isomorphism 
\begin{equation} \label{eqn:4001}
\opn{RHom}_{A \ot B^{\mrm{op}}}(P \ot_A^{\mrm{L}} M, N) \cong 
\opn{RHom}_{A \ot B^{\mrm{op}}} \bigl( M, \opn{RHom}_{A}(P, N) \bigr)
\end{equation}
in $\dcat{D}(\K, \mrm{gr})$. Taking $\opn{H}^0$ in (\ref{eqn:4001}) 
gives us the isomorphism 
\[ \opn{Hom}_{\dcat{D}(A \ot B^{\mrm{op}}, \mrm{gr})}(P \ot_A^{\mrm{L}} M, N) 
\cong 
\opn{Hom}_{\dcat{D}(A \ot B^{\mrm{op}}, \mrm{gr})}
\bigl( M, \opn{RHom}_{A}(P, N) \bigr) \]
in $\dcat{M}(\K, \mrm{gr})$. This is what we want. 
\end{proof}

\begin{lem} \label{lem:4000}
Consider the functors $F$ and $G$ on $\dcat{D}(A \ot B^{\mrm{op}}, \mrm{gr})$ 
from Definition \tup{\ref{dfn:3996}}.  
Assume that the copointed object $(P, \rho)$ is idempotent. 
Then the kernel of $F$ equals the kernel 
of $G$. Namely for every $M \in \dcat{D}(A \ot B^{\mrm{op}}, \mrm{gr})$ we have 
$F(M) = 0$ if and only if $G(M) = 0$. 
\end{lem}

\begin{proof}
We shall use the adjunction formula from Lemma \ref{lem:3996}, with $N = M$. 

First assume $F(M) = 0$. Then 
$\opn{Hom}_{\dcat{D}(A \ot B^{\mrm{op}}, \mrm{gr})} \bigl( F(M), M \bigr)$
is zero, and by Lemma \ref{lem:3996} we see that 
$\opn{Hom}_{\dcat{D}(A \ot B^{\mrm{op}}, \mrm{gr})} \bigl( M, G(M) \bigr)$ 
is zero too. This implies that the morphism $\tau_M : M \to G(M)$ is zero.
Applying $G$ to it we deduce that the morphism 
$G(\tau_M) : G(M) \to G(G(M))$
is zero. But by Lemma \ref{lem:3995} the pointed functor 
$(G, \tau)$ is idempotent, and this means that $G(\tau_M)$ is an isomorphism. 
Therefore $G(M) = 0$. 

Now assume that $G(M) = 0$. Again using Lemma \ref{lem:3996}, but now in 
the reverse direction, we see that the morphism
$\si_M : F(M) \to M$ is zero. Therefore the morphism
$F(\si_M) : F(F(M)) \to F(M)$
is zero. But by Lemma \ref{lem:3995} the copointed functor 
$(F, \si)$ is idempotent, and this means that $F(\si_M)$ is an isomorphism. 
Therefore $F(M) = 0$. 
\end{proof}

Recall that Convention \ref{conv:3700} is in effect. 

\begin{thm}[Abstract Equivalence, \cite{VyYe}] \label{thm:4000}
\index{Monoidal category! idempotent copointed object in}
\index{Additive functor! copointed}
\index{Additive functor! pointed}
\index{Additive functor! idempotent}
Let $A$ and $B$ be graded $\K$-rings, and let $(P, \si)$ be an 
idempotent copointed object in $\dcat{D}(A^{\mrm{en}}, \mrm{gr})$.
Consider the triangulated functors 
\[ F, G : \dcat{D}(A \ot B^{\mrm{op}}, \mrm{gr}) \to 
\dcat{D}(A \ot B^{\mrm{op}}, \mrm{gr}) \]
and the categories 
$\dcat{D}(A \ot B^{\mrm{op}}, \mrm{gr})_F$ and 
$\dcat{D}(A \ot B^{\mrm{op}}, \mrm{gr})_G$ 
from Definitions \tup{\ref{dfn:3996}}
and \tup{\ref{dfn:3997}}. The following hold\tup{:}
\begin{enumerate}
\item The functor $G$ is a right adjoint to $F$. 

\item The copointed triangulated functor $(F, \si)$ and the 
pointed triangulated functor $(G, \tau)$ are idempotent.

\item The categories  $\dcat{D}(A \ot B^{\mrm{op}}, \mrm{gr})_F$ and 
$\dcat{D}(A \ot B^{\mrm{op}}, \mrm{gr})_G$ are the essential 
images of the functors $F$ and $G$ respectively. 

\item The functor 
\[ F : \dcat{D}(A \ot B^{\mrm{op}}, \mrm{gr})_G \to 
\dcat{D}(A \ot B^{\mrm{op}}, \mrm{gr})_F \]
is an equivalence of triangulated categories, with quasi-inverse $G$. 
\end{enumerate}
\end{thm}

\begin{proof}
(1) This is Lemma \ref{lem:3996}. 

\medskip \noindent
(2) This is Lemma \ref{lem:3995}. 

\medskip \noindent
(3) Take an object $M \in \dcat{D}(A \ot B^{\mrm{op}}, \mrm{gr})_G$. Then 
$M \cong G(M)$, so that $M$ is in the essential image of $G$. Conversely, 
suppose there is an isomorphism $\phi : M \iso G(N)$ for some $
N \in \dcat{D}(A \ot B^{\mrm{op}}, \mrm{gr})$. 
We have to prove that $\tau_M$ is an isomorphism. There is a commutative 
diagram 
\[ \UseTips \xymatrix @C=8ex @R=6ex {
M
\ar[r]^{ \phi }
\ar[d]_{ \tau_M }
& 
G(N)
\ar[d]^{ \tau_{G(N)} }
\\
G(M)
\ar[r]^(0.44){ G(\phi) }
&
G(G(N))
} \]
in $\dcat{D}(A \ot B^{\mrm{op}}, \mrm{gr})$ with horizontal isomorphisms.
The idempotence of $(G, \tau)$ says that the morphism $\tau_{G(N)}$ is an 
isomorphism. Therefore $\tau_M$ is an isomorphism. 

A similar argument (with reversed arrows, and using the idempotence of 
$(F, \si)$) tells us that the essential image of 
$F$ is $\dcat{D}(A \ot B^{\mrm{op}}, \mrm{gr})_F$.

\medskip \noindent 
(4) The morphism $\si : P \to A$  sits inside a 
distinguished triangle 
\begin{equation} \label{eqn:4005}
P \xar{\si} A \to N \xar{ \, \triangle \, }  
\end{equation}
in $\dcat{D}(A^{\mrm{en}}, \mrm{gr})$. Let us apply the functor 
$P \ot_A^{\mrm{L}} (-)$ to (\ref{eqn:4005}). We get a distinguished triangle 
\[ P \ot_A^{\mrm{L}} P 
\xar{\opn{ru} \circ \, (\opn{id} \ot_A^{\mrm{L}} \, \si)}
P \to P \ot_A^{\mrm{L}} N \xar{ \, \triangle \, } \]
in $\dcat{D}(A^{\mrm{en}}, \mrm{gr})$.
By the idempotence of $(P, \si)$, the first morphism above is an isomorphism; 
and hence $P \ot_A^{\mrm{L}} N = 0$. Therefore for every 
$M \in \dcat{D}(A \ot B^{\mrm{op}}, \mrm{gr})$ the complex
$F(N \ot_A^{\mrm{L}} M) = P \ot_A^{\mrm{L}} N \ot_A^{\mrm{L}} M$
is zero. Lemma \ref{lem:4000} tells us that 
\begin{equation} \label{eqn:4821}
G(N \ot_A^{\mrm{L}} M) = \opn{RHom}_{A}(P, N \ot_A^{\mrm{L}} M) 
\end{equation}
is zero. 

Now we go back to the distinguished triangle (\ref{eqn:4005}) and we apply to 
it the functor $(-) \ot_A^{\mrm{L}} M$, and then the functor
$\opn{RHom}_{A}(P, - )$. The result is the distinguished triangle 
\[ \opn{RHom}_{A}(P, P \ot_A^{\mrm{L}} M) \xar{\al_M} 
\opn{RHom}_{A}(P, M) \to 
\opn{RHom}_{A}(P, N \ot_A^{\mrm{L}} M)  \xar{ \, \triangle \, } \]
in $\dcat{D}(A \ot B^{\mrm{op}}, \mrm{gr})$. 
By (\ref{eqn:4821}) the third object in this triangle is zero, and it follows 
that $\al_M : G(F(M)) \to G(M)$ is an isomorphism. 
If moreover $M \in \dcat{D}(A \ot B^{\mrm{op}}, \mrm{gr})_G$, then $\tau_M$ is 
an isomorphism too, and thus we have an isomorphism 
\begin{equation} \label{eqn:4822}
\tau_M^{-1} \circ \al_M : G(F(M)) \to M
\end{equation}
in $\dcat{D}(A \ot B^{\mrm{op}}, \mrm{gr})_G$ that's functorial in $M$. 

Similarly, if we apply the functor 
$(-) \ot_A^{\mrm{L}} P$ to (\ref{eqn:4005}), we get a distinguished triangle 
\[ P \ot_A^{\mrm{L}} P 
\xar{\opn{lu} \circ \, (\si \, \ot_A^{\mrm{L}} \opn{id})}
P \to N \ot_A^{\mrm{L}} P \xar{\vartriangle} \]
in $\dcat{D}(A^{\mrm{en}}, \mrm{gr})$.
By the idempotence condition, the first morphism above is an isomorphism; and 
hence $N \ot_A^{\mrm{L}} P = 0$. Therefore for every 
$M \in \dcat{D}(A \ot B^{\mrm{op}}, \mrm{gr})$ the complex
\[ G \bigl( \opn{RHom}_{A}(N, M) \bigr) = 
\opn{RHom}_{A} \bigl( P, \opn{RHom}_{A}(N, M) \bigr) \cong 
\opn{RHom}_{A}(N \ot_A^{\mrm{L}} P, M) \]
is zero. Lemma \ref{lem:4000} tells us that 
\begin{equation} \label{eqn:4823}
F \bigl( \opn{RHom}_{A}(N, M) \bigr) = 
P \ot_A^{\mrm{L}} \opn{RHom}_{A}(N, M) 
\end{equation}
is zero. 

Next we apply the functor $\opn{RHom}_{A}(-, M)$, 
and then the functor $P \ot_A^{\mrm{L}} (-)$, to the distinguished triangle 
(\ref{eqn:4005}). We obtain distinguished triangle 
\[ P \ot_A^{\mrm{L}} \opn{RHom}_{A}(N, M)  \to 
P \ot_A^{\mrm{L}} M \xar{\be_M}
P \ot_A^{\mrm{L}} \opn{RHom}_{A}(P, M)  \xar{ \, \triangle \, } \]
in $\dcat{D}(A \ot B^{\mrm{op}}, \mrm{gr})$. By (\ref{eqn:4823}) the first 
object in this triangle is zero, and so 
$\be_M : F(M) \to F(G(M))$ is an isomorphism. 
If moreover $M \in \dcat{D}(A \ot B^{\mrm{op}}, \mrm{gr})_F$, then $\si_M$ 
is an isomorphism too, and thus we have an isomorphism 
\begin{equation} \label{eqn:4824}
\be_M \circ \si_M^{-1} : M \to F(G(M))
\end{equation}
in $\dcat{D}(A \ot B^{\mrm{op}}, \mrm{gr})_F$ that's functorial in $M$.

The isomorphisms (\ref{eqn:4822}) and (\ref{eqn:4824}) tell us that 
$F$ and $G$ are equivalences between 
$\dcat{D}(A \ot B^{\mrm{op}}, \mrm{gr})_G$ and
$\dcat{D}(A \ot B^{\mrm{op}}, \mrm{gr})_F$, quasi-inverse to each other. 
\end{proof}

We now leave the abstract setting, and return to $\m$-torsion. 
So we are in the situation of Setup \ref{setup:3900}. 
The dedualizing complex 
$P_A := \Ga_{\m}(A) \in \dcat{D}(A^{\mrm{en}}, \mrm{gr})$
from Definition \ref{dfn:3750} is equipped with the morphism 
$\si^{\mrm{R}}_{A} : P_A = \mrm{R} \Ga_{\m}(A) \to A$
from Proposition \ref{prop:3850}. 

\begin{dfn} \label{dfn:4825}
The pair $(P_A, \si^{\mrm{R}}_{A})$ is called the 
{\em dedualizing copointed object}% 
\index{Monoidal category! dedualizing copointed object of the {\indash} 
$\dcat{D}(A^{\mrm{en}}, \mrm{gr})$}
of the monoidal category $\dcat{D}(A^{\mrm{en}}, \mrm{gr})$. 
\end{dfn}

\begin{thm}[\cite{VyYe}] \label{thm:4010}
Under Setup \tup{\ref{setup:3900}}, the dedualizing copointed object 
$(P_A, \si^{\mrm{R}}_{A})$ in the monoidal category 
$\dcat{D}(A^{\mrm{en}}, \mrm{gr})$ is idempotent. 
\index{Monoidal category! idempotent copointed object in}
\end{thm}

\begin{proof}
Let's write $\si := \si^{\mrm{R}}_{A}$ and $P := P_A$. 
We shall start by proving that 
\begin{equation}  \label{eqn:4010}
\opn{lu} \circ \, (\si \, \ot_A^{\mrm{L}} \, \opn{id}) : 
P \ot_{A}^{\mrm{L}} P \to P 
\end{equation}
is an isomorphism in $\dcat{D}(A^{\mrm{en}}, \mrm{gr})$. 
Because the forgetful functor 
$\opn{Rest} : \dcat{D}(A^{\mrm{en}}, \mrm{gr}) \lb \to \dcat{D}(A, \mrm{gr})$
is conservative, it is enough if we prove that 
$\opn{Rest}(\opn{lu} \circ (\si \, \ot_A^{\mrm{L}} \, \opn{id}))$
is an isomorphism. Let us introduce the temporary notation 
$P' := \opn{Rest}(P) \in \dcat{D}(A, \mrm{gr})$. 
With this notation, what we have to show is that 
\begin{equation}  \label{eqn:1372}
\opn{lu} \circ \, (\si \, \ot_A^{\mrm{L}} \, \opn{id}) : 
P \ot_{A}^{\mrm{L}} P' \to P' 
\end{equation}
is an isomorphism in $\dcat{D}(A, \mrm{gr})$.

Consider Theorem \ref{thm:3750} with $B = \K$ and 
$M = P' \in \dcat{D}(A, \mrm{gr})$. There is a commutative diagram 
\[ \UseTips \xymatrix @C=12ex @R=6ex {
P \ot_{A}^{\mrm{L}} P'
\ar[r]^{ \ttev^{\mrm{R, L}}_{\m, P'} }
\ar[d]_{ \si \, \ot_A^{\mrm{L}} \, \opn{id} }
&
\mrm{R} \Ga_{\m}(P')
\ar[d]^{ \si^{\mrm{R}}_{P'} }
\\
A \ot^{\mrm{L}}_{A} P'
\ar[r]^{ \opn{lu} }
&
P'
} \]
in $\dcat{D}(A, \mrm{gr})$, and the horizontal arrows are isomorphisms. 
It suffices to prove that 
$\si^{\mrm{R}}_{P'} : \mrm{R} \Ga_{\m}(P') \to P'$
is an isomorphism in $\dcat{D}(A, \mrm{gr})$. 
But there is an isomorphism 
$P' \cong \mrm{R} \Ga_{\m}(A')$,
where $A' := \opn{Rest}(A) \in \dcat{D}(A, \mrm{gr})$. 
So what we need to prove is that 
\[ \si^{\mrm{R}}_{\mrm{R} \Ga_{\m}(A')} : 
\mrm{R} \Ga_{\m}(\mrm{R} \Ga_{\m}(A')) \to \mrm{R} \Ga_{\m}(A') 
\]
is an isomorphism in $\dcat{D}(A, \mrm{gr})$. This is true because the 
copointed triangulated functor 
$(\mrm{R} \Ga_{\m}, \si^{\mrm{R}})$ on $\dcat{D}(A, \mrm{gr})$
is idempotent; see Corollary \ref{cor:3966} with $B := \K$. 

Now we are going to prove that 
\begin{equation} \label{eqn:4011}
\opn{ru} \circ \, (\opn{id}  \ot_A^{\mrm{L}} \, \si)  : 
P \ot_{A}^{\mrm{L}} P \to P 
\end{equation}
is an isomorphism in $\dcat{D}(A^{\mrm{en}}, \mrm{gr})$. 

We have this commutative (up to a canonical isomorphism) diagram in 
$\dcat{D}(A^{\mrm{en}}, \mrm{gr})$
\begin{equation} \label{eqn:4012}
\UseTips \xymatrix @C=12ex @R=6ex {
P \ot_{A}^{\mrm{L}} P
\ar[d]_{ \opn{id}  \ot_A^{\mrm{L}} \, \si }
&
P \ot_{A}^{\mrm{L}} P \ot_{A}^{\mrm{L}} A
\ar[l]_{ \opn{id}  \ot_A^{\mrm{L}} \opn{ru} }
\ar[d]^{ \opn{id}  \ot_A^{\mrm{L}} \, \si \, \ot_A^{\mrm{L}} \opn{id} }
\\
P \ot_{A}^{\mrm{L}} A
&
P \ot_{A}^{\mrm{L}} A \ot^{\mrm{L}}_{A} A
\ar[l]_{ \opn{id}  \ot_A^{\mrm{L}} \opn{ru} }
} 
\end{equation}

Using Theorem \ref{thm:3750} with $B = A$ and 
$M = A \in \dcat{D}(A^{\mrm{en}}, \mrm{gr})$, 
we have this commutative diagram in $\dcat{D}(A^{\mrm{en}}, \mrm{gr})$~:
\begin{equation} \label{eqn:4013}
\UseTips \xymatrix @C=12ex @R=6ex {
P_A \ot_{A}^{\mrm{L}} A
\ar[r]^{ \ttev^{\mrm{R, L}}_{\m, A} }
\ar[d]_{ \si \, \ot_A^{\mrm{L}} \, \opn{id} }
&
\mrm{R} \Ga_{\m}(A)
\ar[d]^{ \si^{\mrm{R}}_{A} }
\\
A \ot^{\mrm{L}}_{A} A
\ar[r]^{ \opn{lu} }
&
A
} 
\end{equation}
Applying the functor $P \ot_A^{\mrm{L}} (-)$ to this diagram, we obtain this 
commutative diagram 
\begin{equation} \label{eqn:4014}
\UseTips \xymatrix @C=15ex @R=6ex {
P \ot_{A}^{\mrm{L}} P \ot_{A}^{\mrm{L}} A
\ar[d]_{ \opn{id}  \ot_A^{\mrm{L}} \, \si \, \ot_A^{\mrm{L}} \, \opn{id} }
\ar[r]^{ \opn{id} \ot_{A}^{\mrm{L}} \, \ttev^{\mrm{R, L}}_{\m, A} }
&
P \ot_{A}^{\mrm{L}} \mrm{R} \Ga_{\m}(A)
\ar[d]_{ \opn{id}  \ot_A^{\mrm{L}} \, \si^{\mrm{R}}_A }
\\
P \ot_{A}^{\mrm{L}} A \ot^{\mrm{L}}_{A} A
\ar[r]^{ \opn{id} \ot_{A}^{\mrm{L}} \opn{lu} }
&
P \ot_{A}^{\mrm{L}} A 
}
\end{equation}
in $\dcat{D}(A^{\mrm{en}}, \mrm{gr})$.
The last move is using the fact that 
$\ttev^{\mrm{R, L}}_{\m, (-)}$ is an isomorphism of functors; this 
yields the next commutative diagram: 
\begin{equation} \label{eqn:4015}
\UseTips \xymatrix @C=15ex @R=6ex {
P \ot_{A}^{\mrm{L}} \mrm{R} \Ga_{\m}(A)
\ar[d]_{ \opn{id}  \ot_A^{\mrm{L}} \, \si^{\mrm{R}}_A }
\ar[r]^{ \ttev^{\mrm{R, L}}_{\m, \mrm{R} \Ga_{\m}(A)}} 
&
\mrm{R} \Ga_{\m}
(\mrm{R} \Ga_{\m}(A))
\ar[d]_{\mrm{R} \Ga_{\m}(\si^{\mrm{R}}_A)}
\\
P \ot_{A}^{\mrm{L}} A 
\ar[r]^{ \ttev^{\mrm{R, L}}_{\m, A} }
&
\mrm{R} \Ga_{\m}(A)
}
\end{equation}
All horizontal arrows in diagrams (\ref{eqn:4012}), (\ref{eqn:4013}), 
(\ref{eqn:4014}) and (\ref{eqn:4015}) are isomorphisms. 
Since $\opn{ru}$ is an isomorphism, to prove that (\ref{eqn:4011}) is an 
isomorphism, it is enough to prove that 
the morphism $\opn{id}  \ot_A^{\mrm{L}} \, \si$,
which is the left vertical arrow in diagram (\ref{eqn:4012}), is an 
isomorphism. Now this last morphism coincides with the 
left vertical arrow in diagram (\ref{eqn:4015}). 
Therefore it is enough to prove that the right vertical arrow in 
diagram (\ref{eqn:4015}) is an isomorphism. This is 
\begin{equation} \label{eqn:4016}
\mrm{R} \Ga_{\m}(\si^{\mrm{R}}_A) : 
\mrm{R} \Ga_{\m} (\mrm{R} \Ga_{\m}(A)) \to
\mrm{R} \Ga_{\m}(A) . 
\end{equation}
Because the functor $\opn{Rest}$ is conservative, it suffices to prove that 
\[ \mrm{R} \Ga_{\m}(\si^{\mrm{R}}_{A'}) : 
\mrm{R} \Ga_{\m}(\mrm{R} \Ga_{\m}(A')) \to
\mrm{R} \Ga_{\m}(A') \]
is an isomorphism in $\dcat{D}(A, \mrm{gr})$, where, as before, we write  
$A' := \opn{Rest}(A) \in \dcat{D}(A, \mrm{gr})$.
This is true because the copointed triangulated functor 
$(\mrm{R} \Ga_{\m}, \si^{\mrm{R}})$ on $\dcat{D}(A, \mrm{gr})$
is idempotent. 
\end{proof}

\begin{dfn} \label{dfn:3790}
Under Setup \tup{\ref{setup:3900}}, consider the dedualizing copointed object 
$(P_A, \si^{\mrm{R}}_{A})$ in the monoidal category 
$\dcat{D}(A^{\mrm{en}}, \mrm{gr})$; see Definition \ref{dfn:4825}.
\begin{enumerate}
\item The {\em abstract derived $\m$-adic completion functor}% 
\index{Completion! abstract derived {\indash} functor}%
\index{1-ADCm@$\opn{ADC}_{\m}$}
is the triangulated \lb pointed functor $(\opn{ADC}_{\m}, \tau)$ 
on $\dcat{D}(A \ot B^{\mrm{op}}, \mrm{gr})$ 
that's induced by the copointed object 
$(P_A, \si^{\mrm{R}}_{A})$, in the sense of Definition \ref{dfn:3996}

\item The full subcategory of $\dcat{D}(A \ot B^{\mrm{op}}, \mrm{gr})$ on the 
complexes $M$ such that 
$\tau_{M} : M \to \opn{ADC}_{\m}(M)$
is an isomorphism is denoted by 
$\dcat{D}(A \ot B^{\mrm{op}}, \mrm{gr})_{\mrm{com}}$.
\end{enumerate}
\end{dfn}

To put item (1) in explicit terms,  
$\opn{ADC}_{\m}(M) = \opn{RHom}_A(P_A, M)$,
and the morphism 
$\tau_M : M \to \opn{ADC}_{\m}(M)$ is 
$\tau_M = \opn{RHom}_A(\si^{\mrm{R}}_{A}, \opn{id}_M)$. 

We already have the copointed triangulated functor 
$(\mrm{R} \Ga_{\m}, \si^{\mrm{R}})$ on the category 
$\dcat{D}(A \ot B^{\mrm{op}}, \mrm{gr})$. 
In analogy with Definition \ref{dfn:3790}(2) we make the next definition. 

\begin{dfn} \label{dfn:4020}
Under Setup \tup{\ref{setup:3900}}, we denote by 
$\dcat{D}(A \ot B^{\mrm{op}}, \mrm{gr})_{\mrm{tor}}$
the full subcategory of $\dcat{D}(A \ot B^{\mrm{op}}, \mrm{gr})$ on the 
complexes $M$ such that
$\si^{\mrm{R}}_{M} : \mrm{R} \Ga_{\m}(M) \to M$ 
is an isomorphism.
\end{dfn}

Proposition \ref{prop:3965} says that 
\[ \dcat{D}(A \ot B^{\mrm{op}}, \mrm{gr})_{\mrm{tor}} = 
\dcat{D}_{\mrm{tor}}(A \ot B^{\mrm{op}}, \mrm{gr}) , \]
where the latter is the full subcategory of 
$\dcat{D}(A \ot B^{\mrm{op}}, \mrm{gr})$ on the complexes $M$ whose cohomology 
modules are $\m$-torsion. 

\begin{thm}[NC MGM Equivalence, \cite{VyYe}] \label{thm:3792} 
\index{Noncommutative MGM Equivalence}
Under Setup \tup{\ref{setup:3900}} the following hold.
\begin{enumerate}
\item The categories  
$\dcat{D}(A \ot B^{\mrm{op}}, \mrm{gr})_{\mrm{com}}$ and 
$\dcat{D}(A \ot B^{\mrm{op}}, \mrm{gr})_{\mrm{tor}}$
are the essential images of the functors $\opn{ADC}_{\m}$ and 
$\mrm{R} \Ga_{\m}$ respectively.  

\item The pointed triangulated functor
$(\opn{ADC}_{\m}, \tau)$ and the copointed triangulated functor
$(\mrm{R} \Ga_{\m}, \si^{\mrm{R}})$ are idempotent. 

\item The functor 
\[ \mrm{R} \Ga_{\m} : \dcat{D}(A \ot B^{\mrm{op}}, \mrm{gr})_{\mrm{com}} \to 
\dcat{D}(A \ot B^{\mrm{op}}, \mrm{gr})_{\mrm{tor}} \]
is an equivalence of triangulated categories, with quasi-inverse 
$\opn{ADC}_{\m}$.
\end{enumerate}
\end{thm}

\begin{proof}
By Definition \ref{dfn:3790}, the pointed triangulated functor 
$(\opn{ADC}_{\m}, \tau)$
is the one induced by the copointed object $(P_A, \si^{\mrm{R}}_{A})$.
And by Theorem \ref{thm:3750}, the copointed triangulated functor 
$(\mrm{R} \Ga_{\m}, \si^{\mrm{R}})$
is the one induced by the copointed object $(P_A, \si^{\mrm{R}}_{A})$.
According to Theorem \ref{thm:4010}, the copointed object 
$(P_A, \si^{\mrm{R}}_{A})$ is idempotent.
This means that we can use Theorem \ref{thm:4000}  on abstract equivalence. 
Item (1) here is then a special case of item (3) of Theorem \ref{thm:4000};
item (2) here is a special case of item (2) of Theorem \ref{thm:4000};
and item (3) here is a special case of item (4) of Theorem \ref{thm:4000}.
\end{proof}

\begin{exa} \label{exa:4575}
In Example \ref{exa:4595} we presented the commutative MGM Equivalence. 
Recall the setting: a commutative ring $A$, and an ideal $\a \sub A$ 
generated by a weakly proregular sequence $\bsym{a}$. The module category there 
was $\dcat{M}(A)$. 
If we want to transfer this to the connected graded setting 
of the current section, where the module category is $\dcat{M}(A, \mrm{gr})$,
then we should take $A$ to be a commutative connected 
graded $\K$-ring, and $\bsym{a}$ is a sequence of homogeneous elements that 
generates the augmentation ideal $\m$. Then $A = \K[\bsym{a}]$ is a finitely 
generated commutative $\K$-ring, and thus noetherian. 
Because $A$ is commutative, left $A$-modules are the same as central 
$A$-bimodules. This is the reason we work with the category
$\dcat{M}(A, \mrm{gr})$ and not with $\dcat{M}(A^{\mrm{en}}, \mrm{gr})$.
% we may prefer to work with central $A$-bimodules. 
% Since central $A$-bimodules are the same as left $A$-modules, this means that 
% we will have to replace the  .

Regardless of which setting we choose (graded or ungraded), the 
dedualizing complex $P_A$ over the ring $A$ (as an object of 
$\dcat{D}(A, \mrm{gr})$ or $\dcat{D}(A)$ respectively) can be made explicit:
$P_A \cong \opn{K}^{\vee}_{\infty}(A; \bsym{a}) \cong 
\opn{Tel}(A; \bsym{a})$, 
where $\opn{K}^{\vee}_{\infty}(A; \bsym{a})$ is the {\em infinite dual Koszul 
complex}, and $\opn{Tel}(A; \bsym{a})$ is the {\em telescope complex}. 
Since the telescope complex is a bounded complex of free $A$-modules (of 
countable rank), the derived torsion and completion functors have  
particularly nice presentations:
$\mrm{R} \Ga_{\m} = \opn{Tel}(A; \bsym{a}) \ot_A (-)$
and 
$\opn{ADC}_{\m} = \opn{Hom}_A \bigl( \opn{Tel}(A; \bsym{a}), - \bigr)$.
Also the abstract derived completion functor $\opn{ADC}_{\m}$ is canonically 
isomorphism to $\mrm{L} \Lambda_{\m}$, the left derived functor of the 
$\m$-adic completion functor $\Lambda_{\m}$. 
See \cite{PSY} for details. 
\end{exa}






\cleardoublepage
\mysection{Balanced Dualizing Complexes over NC Graded Rings}
\label{sec:BDC}

\AYcopyright

Let $A = \bigoplus_{i \geq 0} A_i$ be a noncommutative noetherian connected 
graded  $\K$-ring, with augmentation ideal $\m = \bigoplus_{i \geq 1} A_i$. 
A {\em balanced dualizing complex} over $A$ is a dualizing complex $R$ over $A$,
in the noncommutative graded sense, that satisfies the 
{\em Noncommutative Graded Local Duality Theorem} (Theorem \ref{thm:4616} 
below). Balanced dualizing complexes were introduced in 1992 by A. Yekutieli 
\cite{Ye1}.

In this section we define balanced dualizing complexes, and then we prove their 
uniqueness and existence, the local duality theorem, and the trace 
functoriality. One of the main results of this section is Corollary 
\ref{cor:4585}, which says that the following two properties are equivalent:
\begin{itemize}
\item[(i)] The ring $A$ satisfies the $\chi$ condition, and it has 
finite local cohomological dimension.

\item[(ii)] The ring $A$ has a balanced dualizing complex.
\end{itemize}
This result is the product of the combined efforts of 
Yekutieli, J.J. Zhang and M. Van den Bergh.
The $\chi$ condition was already discussed in Subsection 
\ref{subsec:sym-der-tor}. M. Artin and Zhang had introduced the $\chi$ 
condition and the finiteness of local cohomological dimension in their 1994 
paper \cite{ArZh}, to ensure that the {\em noncommutative projective scheme} 
$\opn{Proj}(A)$ has ``good geometric properties''. On the other hand, as 
mentioned above, balanced dualizing complexes were designed to satisfy the 
Noncommutative Local Duality Theorem. It is 
remarkable that in the end, these two properties turned out to be equivalent. 

Throughout this section we adhere to Conventions \ref{conv:3700} and 
\ref{conv:4560}. Thus $\K$ is a base field, and all graded rings are 
algebraically graded central $\K$-rings, 
as defined in Section \ref{sec:alg-gra-rings}.
See Remark \ref{rem:4082} regarding the complete case ($A$ is adically complete 
instead of graded) and the arithmetic case (the base ring $\K$ is not a field). 


\mysubsection{Graded NC Dualizing Complexes}
\label{subsec:graded-NCDC}

In this subsection we introduce {\em graded noncommutative dualizing 
complexes}. This is a generalization of Grothendieck's original commutative 
definition from \cite{RD} (see Section \ref{sec:dual-cplx-comm-rng}). 

Consider a graded $\K$-ring $A$. 
Complexes of graded $A$-modules were studied in Subsections 
\ref{subsec:alg-gr-mods} and \ref{gr-res-der-fun}, and here we use these 
constructions and results. 
Let us recall that $\dcat{M}(A, \mrm{gr})$ is the abelian category of graded 
$A$-modules, and $\dcat{D}(A, \mrm{gr})$ is its derived category. 

Suppose $B$ is a second graded ring. Given complexes
$M \in \dcat{D}(A \ot B^{\mrm{op}}, \mrm{gr})$ and 
$N \in \dcat{D}(A^{\mrm{en}}, \mrm{gr})$, 
there is the 
{\em noncommutative derived Hom-evaluation morphism}%
\index{Hom-evaluation morphism! NC derived}%
\index{1-EvRRMN@$\opn{ev}^{\mrm{R, R}}_{M, N}$}
\begin{equation} \label{eqn:3915}
\opn{ev}^{\mrm{R, R}}_{M, N} : M \to 
\opn{RHom}_{A^{\mrm{op}}} \bigl( \opn{RHom}_{A}(M, N), N \bigr) 
\end{equation}
in $\dcat{D}(A \ot B^{\mrm{op}}, \mrm{gr})$.
For a choice of a K-injective resolution 
$N \to I$ in $\dcat{C}_{\mrm{str}}(A^{\mrm{en}}, \mrm{gr})$,
the morphism $\opn{ev}^{\mrm{R, R}}_{M, N}$ 
is represented by the homomorphism 
\[ \opn{ev}_{M, I} : 
M \to \opn{Hom}_{A^{\mrm{op}}} \bigl( \opn{Hom}_{A}(M, I), I \bigr) \]
in $\dcat{C}_{\mrm{str}}(A \ot B^{\mrm{op}}, \mrm{gr})$, 
whose formula is 
$\opn{ev}_{M, I}(m)(\phi) := (-1)^{p \cd q} \cd \phi(m)$
for $m \in M^p$ and $\phi \in \opn{Hom}_{A}(M, I)^q$. 

In the special case $B = A$ and 
$M = A \in \dcat{D}(A^{\mrm{en}}, \mrm{gr})$
we have the canonical isomorphisms 
$N \cong \opn{RHom}_{A}(A, N) \cong \opn{RHom}_{A^{\mrm{op}}}(A, N)$,
i.e.\ the left and right co-unitor isomorphisms of the biclosed monoidal 
structure. Then the derived NC Hom-evaluation morphism 
$\opn{ev}^{\mrm{R, R}}_{M, N}$ from (\ref{eqn:3915}) specializes to the 
{\em NC derived homothety morphism through $A$}%
\index{Homothety morphism! noncommutative derived}%
\index{1-HmRNA@$\opn{hm}^{\mrm{R}}_{N, A}$}~:
\begin{equation} \label{eqn:3916}
\opn{hm}^{\mrm{R}}_{N, A}
= \opn{ev}^{\mrm{R, R}}_{A, N} : A \to \opn{RHom}_{A^{\mrm{op}}}(N, N) 
\end{equation}
in $\dcat{D}(A^{\mrm{en}}, \mrm{gr})$.
Similarly, after exchanging $A$ and $A^{\mrm{op}}$, we recover 
the {\em NC derived homothety morphism through $A^{\mrm{op}}$}~:
\begin{equation} \label{eqn:3917}
\opn{hm}^{\mrm{R}}_{N, A^{\mrm{op}}} =
\opn{ev}^{\mrm{R, R}}_{A, N} : A \to \opn{RHom}_{A}(N, N) ,
\end{equation}
also in $\dcat{D}(A^{\mrm{en}}, \mrm{gr})$.
This should be compared to Definition \ref{dfn:3565}, which refers to the 
ungraded NC setting. 

\begin{dfn}[\cite{Ye1}] \label{dfn:3716} 
Let $A$ be a noetherian graded ring. A 
{\em graded NC dualizing complex}% 
\index{Dualizing complex! graded NC}
over $A$ is a complex 
$R \in \dcat{D}^{\mrm{b}}(A^{\mrm{en}}, \mrm{gr})$
with the following three properties\tup{:}
\begin{enumerate}
\rmitem{i} {\em Finiteness of cohomology}: for every integer $p$ the 
$A$-bimodule 
$\opn{H}^p(R)$ is a finite graded module over $A$ and over 
$A^{\mrm{op}}$. 

\rmitem{ii} {\em Finite injective dimension}: the complex $R$ has finite 
graded-injective dimension over $A$ and over $A^{\mrm{op}}$. 

\rmitem{iii} {\em NC Derived Morita property}%
\index{Derived Morita property! noncommutative}:
the noncommutative derived homothety morphisms 
\[ \opn{hm}^{\mrm{R}}_{R, A} : A \to \opn{RHom}_{A^{\mrm{op}}}(R, R) \]
and
\[ \opn{hm}^{\mrm{R}}_{R, A^{\mrm{op}}} : A \to \opn{RHom}_A(R, R) \]
in $\dcat{D}(A^{\mrm{en}}, \mrm{gr})$
are both isomorphisms.
\end{enumerate}
\end{dfn}

Condition (i) can be restated as 
$R \in \dcat{D}_{(\mrm{f}, \mrm{f})}(A^{\mrm{en}}, \mrm{gr})$.
It can be shown (using Corollary \ref{cor:3830}(3) and smart truncation, as in 
the proof of Proposition \ref{prop:2115}) 
that condition (ii) is equivalent to the existence of an 
isomorphism $R \cong I$ in $\dcat{D}(A^{\mrm{en}}, \mrm{gr})$,
where $I$ is a bounded complex, and every bimodule $I^p$ is 
graded-injective over $A$ and over $A^{\mrm{op}}$.

\begin{dfn} \label{dfn:3820}
In the situation of Definition \ref{dfn:3716}, 
given another graded ring $B$, the 
{\em duality functors}%
\index{Duality functor}
associated to the dualizing complex $R$ are the triangulated functors 
\[ D_A : \dcat{D}(A \ot B^{\mrm{op}}, \mrm{gr})^{\mrm{op}} \to 
\dcat{D}(B \ot A^{\mrm{op}}, \mrm{gr}), \quad 
D_A := \opn{RHom}_A(-, R)   \]
and 
\[ D_{A^{\mrm{op}}} : 
\dcat{D}(B \ot A^{\mrm{op}}, \mrm{gr})^{\mrm{op}} 
\to \dcat{D}(A \ot B^{\mrm{op}}, \mrm{gr}), \quad 
D_{A^{\mrm{op}}} := \opn{RHom}_{A^{\mrm{op}}}(-, R) .  \]
\end{dfn}

Since the dualizing complex $R$ is usually clear from the context, we can omit 
it from the notation of the duality functors 
$D_A$ and $D_{A^{\mrm{op}}}$. We can also omit $R$ from the notation of 
the related derived Hom-evaluation morphisms, and just write 
$\opn{ev}^{\mrm{R, R}}_{M} : M \to D_{A^{\mrm{op}}}(D_{A}(M))$
and
$\opn{ev}^{\mrm{R, R}}_{M'} : M' \to D_{A}(D_{A^{\mrm{op}}}(M'))$
for complexes
$M \in  \dcat{D}(A \ot B^{\mrm{op}}, \mrm{gr})$
and 
$M' \in  \dcat{D}(B \ot A^{\mrm{op}}, \mrm{gr})$. 

\begin{thm} \label{thm:3820}
Let $A$ and $B$ be graded rings, with $A$ noetherian.
Let $R \in \dcat{D}(A^{\mrm{en}}, \mrm{gr})$ be a graded NC 
dualizing complex over $A$, with associated duality functors $D_A$ and 
$D_{A^{\mrm{op}}}$. Let $\star$ be a boundedness indicator, and let 
$M \in \lb \dcat{D}^{\star}_{(\mrm{f}, ..)}(A \ot B^{\mrm{op}}, \mrm{gr})$.
Then the following hold\tup{:}
\begin{enumerate}
\item The complex $D_A(M)$ belongs to 
$\dcat{D}^{-\star}_{(.., \mrm{f})}(B \ot A^{\mrm{op}}, \mrm{gr})$, 
where $-\star$ is the reversed boundedness indicator.

\item The derived Hom-evaluation morphism 
$\opn{ev}^{\mrm{R, R}}_{M} : M \to D_{A^{\mrm{op}}}(D_{A}(M))$
in $\dcat{D}(A \ot B^{\mrm{op}}, \mrm{gr})$ is an isomorphism. 

\item The functor
\[ D_{A} : 
\dcat{D}^{\star}_{(\mrm{f}, ..)}(A \ot B^{\mrm{op}}, \mrm{gr})^{\mrm{op}} \to
\dcat{D}^{-\star}_{(.., \mrm{f})}(B \ot A^{\mrm{op}}, 
\mrm{gr}) \]
is an equivalence of triangulated categories, with quasi-inverse 
$D_{A^{\mrm{op}}}$. 
\end{enumerate}
\end{thm}

\begin{proof} \mbox{}

\smallskip \noindent
(1) We can forget the ring $B$. Since the functor 
$D_A : \dcat{D}(A, \mrm{gr})^{\mrm{op}} \to \dcat{D}(A^{\mrm{op}}, \mrm{gr})$
has finite cohomological dimension, and since
$D_A(A) = R \in \dcat{M}_{\mrm{f}}(A^{\mrm{op}}, \mrm{gr})$,
this is a consequence of Theorem \ref{thm:2160}(2), slightly 
modified to handle the algebraically graded situation. 

\medskip \noindent 
(2) Because the restriction functor 
$\dcat{D}(A \ot B^{\mrm{op}}, \mrm{gr}) \to \dcat{D}(A, \mrm{gr})$
is conservative, we can forget about the ring $B$. The derived Morita 
property says that 
$\opn{ev}^{\mrm{R, R}}_{M} : M \to D_{A^{\mrm{op}}}(D_{A}(M))$
is an isomorphism for $M = A$. The functors $\opn{Id}$ and 
$D_{A^{\mrm{op}}} \circ D_{A}$
have finite cohomological dimensions. The assertion is then a consequence of 
Theorem \ref{thm:2135}(2), slightly modified to handle 
the algebraically graded situation. 

\medskip \noindent 
(3) This is clear from items (2) and (3).
\end{proof}

We shall require a notion of twisting in the NC graded setting. 

\begin{dfn} \label{dfn:3780}
Let $\phi$ be an automorphism of $A$ in the category 
$\catt{Rng}_{\mrm{gr}} \centover \K$, and let $i$ be an integer. 
The {\em $(\phi, i)$-twist of the bimodule $A$}% 
\index{Twisted bimodule}%
\index{1-Afi@$A(\phi, i)$}
is the object
$A(\phi, i) \in \dcat{M}(A^{\mrm{en}}, \mrm{gr})$
defined as follows: as a left graded $A$-module it is the graded-free 
$A$-module $A(i)$ with basis element $e$ in degree $-i$. The right $A$-module 
action is given by the formula 
$e \cd^{\phi} a := \phi(a) \cd e$ for $a \in A$. 
\end{dfn}

To be more explicit, an element $m \in A(\phi, i)$ can be expressed uniquely as 
$m = b \cd e$ with $b \in A$. Then 
$m \cd^{\phi} a = (b \cd e) \cd^{\phi} a = (b \cd \phi(a)) \cd e 
\in A(\phi, i)$.

\begin{dfn} \label{dfn:3925}
An {\em invertible graded $A$-bimodule} is a bimodule 
$L \in \lb \dcat{M}(A^{\mrm{en}}, \mrm{gr})$
such that there exists some 
$L^{\vee} \in \dcat{M}(A^{\mrm{en}}, \mrm{gr})$
satisfying 
$L \ot_A L^{\vee} \cong L^{\vee} \ot_A L \cong A$
in $\dcat{M}(A^{\mrm{en}}, \mrm{gr})$.
The bimodule $L^{\vee}$ is called a quasi-inverse of $L$. 
\end{dfn}

\begin{prop} \label{prop:9325}  
Assume $A$ is connected graded. 
Let $L$ be an invertible graded $A$-bimodule.
Then $L \cong A(\phi, i)$, for a unique automorphism $\phi$ and a unique 
integer $i$. 
\end{prop}

\begin{proof}
Take a quasi-inverse $L^{\vee}$ of $L$. 
Since the tensor product commutes with the ungrading functor, we see that 
$L \ot_A L^{\vee} \cong L^{\vee} \ot_A L \cong A$ 
in $\dcat{M}(\opn{Ungr}(A)^{\mrm{en}})$. Classical Morita theory (see 
e.g.\ \cite[Chapter 4]{Row}) says that $L$ 
and $L^{\vee}$ are finite projective $\opn{Ungr}(A)$-modules on both sides. 
Hence they are flat over $A$ and $A^{\mrm{op}}$. 

Consider the obvious surjection 
\[ \K \cong \K \ot_A L^{\vee} \ot_A L \surj 
(\K \ot_A L^{\vee} \ot_{A} \K) \ot_{} (\K \ot_A L)  \]
in $\dcat{M}(A^{\mrm{en}}, \mrm{gr})$.
Since $L$ is a nonzero finite graded $A$-module and $A^{\mrm{op}}$-module, by 
the graded Nakayama Lemma (Proposition \ref{prop:4590}) the $\K$-modules  
$\K \ot_A L^{\vee} \ot_{A} \K$ and $\K \ot_A L$ are 
nonzero. Therefore $\K \ot_A L$ is a graded-free $\K$-module of rank $1$, so  
$\K \ot_A L \cong \K(i)$ in $\dcat{M}(\K, \mrm{gr})$ for some 
integer $i$. But $L$ is a flat $A$-module, and therefore, as in the classical 
commutative proof, we get an isomorphism
$L \cong A(i)$ in $\dcat{M}(A, \mrm{gr})$. 

By symmetry we also have an isomorphism 
$L \cong A(i)$ in $\dcat{M}(A^{\mrm{op}}, \mrm{gr})$.
Thus there is an element $e \in L_{-i}$ which is a basis of $L$ on both sides. 
For $a \in A$ let $\phi(a) \in A$ be the unique element such that 
$e \cd a = \phi(a) \cd e$. Then $\phi$ is a $\K$-ring automorphism of $A$, and 
$L \cong  A(\phi, i)$ in $\dcat{M}(A^{\mrm{en}}, \mrm{gr})$.

Because the only invertible elements of $A$ are the nonzero elements of 
$A_0 \cong \K$, and these are central elements, $A$ does not have nontrivial 
inner automorphisms. Thus the automorphism $\phi$ is unique. 
\end{proof}

The next theorem is a variant of Theorem \ref{thm:2175}. 

\begin{thm}[Uniqueness of NC Graded DC, \cite{Ye1}] \label{thm:3780} 
Let $A$ be a noetherian connected graded ring, and let  $R$ and $R'$ be graded 
NC dualizing complexes over $A$. Then there is an isomorphism
\[ R' \cong R \ot_A A(\phi, i)[j] \]
in $\dcat{D}(A^{\mrm{en}}, \mrm{gr})$, for a unique automorphism $\phi$ and 
unique integers $i, j$. 
\end{thm}

For the proof we shall need a few lemmas. 
Recall that a module $N \in \dcat{M}(A, \mrm{gr})$ is called bounded below 
if $N_i = 0$ for $i \ll 0$. 

\begin{lem} \label{lem:3916}
Let 
$T, T' \in \dcat{D}^-(A^{\mrm{en}}, \mrm{gr})$ 
satisfy these conditions\tup{:}
\begin{itemize}
\item For every $k$ the graded bimodules $\opn{H}^k(T)$ and $\opn{H}^k(T')$
are bounded below. 

\item There are isomorphisms
$T \ot^{\mrm{L}}_{A} T' \cong T' \ot^{\mrm{L}}_{A} T \cong A$
in $\dcat{D}(A^{\mrm{en}}, \mrm{gr})$.
\end{itemize}
Then $T \cong A(\phi, i)[j]$, for a unique automorphism $\phi$ and unique 
integers $i, j$. 
\end{lem}

\begin{proof}
This is similar to the proof of Lemma \ref{lem:2175}. Define 
$j_1 := \opn{sup}(\opn{H}(T))$ and 
$j'_1 := \opn{sup}(\opn{H}(T'))$. 
By smart truncation we can assume that 
$j_1 = \opn{sup}(T)$ and $j'_1 := \opn{sup}(T')$.
The K\"unneth trick (see Lemma \ref{lem:2177}; it works also in the 
graded setting) we have 
\[ \opn{H}^{j_1}(T) \ot_A \opn{H}^{j_1'}(T') \cong 
\opn{H}^{j_1 + j'_1}(T \ot^{\mrm{L}}_{A} T') \cong \opn{H}^{j_1 + j'_1}(A) \]
in $\dcat{M}(A^{\mrm{en}}, \mrm{gr})$. 
The graded Nakayama Lemma (Proposition \ref{prop:4590}) implies that
$\opn{H}^{j_1}(T) \ot_A \opn{H}^{j'_1}(T') \neq 0$.
Hence we must have $j_1 + j_1' = 0$ and 
$\opn{H}^{j_1}(T) \ot_A \opn{H}^{j'_1}(T') \cong A$.
By symmetry, we also have 
$\opn{H}^{j'_1}(T') \ot_A \opn{H}^{j_1}(T) \cong A$.
According to Proposition \ref{prop:9325} we know that 
$\opn{H}^{j_1}(T) \cong L$, where $L := A(\phi, i)$ for some $\phi$ and $i$. 
Therefore $\opn{H}^{j'_1}(T') \cong L'$,
where $L' := A(\phi^{-1}, -i)$. 

We now forget the left $A$-module structure of $T$, and the right $A$-module 
structure of $T'$. By the projective truncation trick (Lemma \ref{lem:2189})
we get an isomorphism $T \cong L[-j_1] \oplus N$ in 
$\dcat{D}(A^{\mrm{op}}, \mrm{gr})$, where 
$\opn{sup}(N) \leq j_1 - 1$. Similarly there is a an isomorphism 
$T' \cong L'[-j'_1] \oplus N'$ in 
$\dcat{D}(A, \mrm{gr})$, where $\opn{sup}(N') \leq j'_1 - 1$.
Then there is an isomorphism 
\begin{equation} \label{eqn:5125}
\begin{aligned}
& A \cong \opn{H}(T \ot^{\mrm{L}}_{A} T') \cong (L \ot_{A} L') 
\\ 
& \quad 
\oplus \bigl( L[-j_1] \ot_{A} \opn{H}(N') \bigr)
\oplus \bigl( \opn{H}(N) \ot_{A} L'[-j'_1] \bigr)  \oplus 
\opn{H}(N \ot^{\mrm{L}}_{A} N')
\end{aligned} 
\end{equation}
in $\dcat{G}_{\mrm{str}}(\K, \mrm{gr})$. But $A$ is concentrated in 
cohomological degree $0$, and this forces the three summands 
in the second line of (\ref{eqn:5125}) to be 
zero. Since $L'$ is graded-free of rank $1$ over $A$, it follows that 
$\opn{H}(N) = 0$. Therefore $\opn{H}^{k}(T) = 0$ for all $k \neq j_1$.
Letting $j := -j_1$ we get 
$T \cong \opn{H}^{-j}(T)[j] \cong L[j] = A(\phi, i)[j]$
in $\dcat{D}(A^{\mrm{en}}, \mrm{gr})$. 
\end{proof}

\begin{lem} \label{lem:3917}
Let $R, R' \in \dcat{D}(A^{\mrm{en}}, \mrm{gr})$ 
be graded NC dualizing complexes, and let 
$M \in \dcat{D}_{}(A^{\mrm{en}},  \mrm{gr})$.
Then there a morphism 
\[ \psi_M : M \ot^{\mrm{L}}_{A} \opn{RHom}_A(R, R') \iso 
\opn{RHom}_A \bigl( \opn{RHom}_{A^{\mrm{op}}}(M, R), R' \bigr) \]
in $\dcat{D}(A^{\mrm{en}}, \mrm{gr})$, that is functorial in $M$. 
If $M \in \dcat{D}^-_{(.., \mrm{f})}(A^{\mrm{en}},  \mrm{gr})$
then $\psi_M$ is an isomorphism. 
\end{lem}

\begin{proof}
This is the NC graded version of the isomorphism that was used in the proof of 
Theorem \ref{thm:2175}. To define $\psi_M$ we choose a K-projective 
resolution $P \to M$, and K-injective resolutions 
$R \to I$ and $R' \to I'$, all in 
$\dcat{C}_{\mrm{str}}(A^{\mrm{en}}, \mrm{gr})$.
Then $\psi_M$ is represented by the obvious homomorphism 
\[ \til{\psi}_{P, I, I'} : P \ot_{A} \opn{Hom}_A(I, I') \to 
\opn{Hom}_A \bigl( \opn{Hom}_{A^{\mrm{op}}}(P, I), I' \bigr) \]
in $\dcat{C}_{\mrm{str}}(A^{\mrm{en}}, \mrm{gr})$.

To show that $\psi_M$ is an isomorphism when 
$M \in \dcat{D}^-_{(.., \mrm{f})}(A^{\mrm{en}},  \mrm{gr})$,
we can forget the $A$-module structure on $M$, and view $\psi_M$ as a morphism 
in $\dcat{D}(\K, \mrm{gr})$. So we have a morphism $\psi_M$ between 
triangulated functors 
$\dcat{D}^-_{\mrm{f}}(A^{\mrm{op}}, \mrm{gr}) \to \dcat{D}(\K, \mrm{gr})$.
It is clear that $\psi_A$ is an isomorphism. 
By Theorem \ref{thm:2135}(1), 
which is valid also in the algebraically graded 
context, $\psi_M$ is an isomorphism for every 
$M \in \dcat{D}^-_{\mrm{f}}(A^{\mrm{op}}, \mrm{gr})$. 
\end{proof}

\begin{proof}[Proof of Theorem \tup{\ref{thm:3780}}]
Define the duality functors 
\[ D_A := \opn{RHom}_A(-, R) , \quad 
D_{A^{\mrm{op}}} := \opn{RHom}_{A^{\mrm{op}}}(-, R) , \]
\[ D'_A := \opn{RHom}_A(-, R') , \quad 
D'_{A^{\mrm{op}}} := \opn{RHom}_{A^{\mrm{op}}}(-, R') . \]
These are contravariant triangulated functors from 
$\dcat{D}(A^{\mrm{en}}, \mrm{gr})$ to itself. 
Then define the complexes 
$T := (D'_A \circ D_{A^{\mrm{op}}})(A) = \opn{RHom}_{A}(R, R')$
and 
$T' := \lb (D_A \circ D'_{A^{\mrm{op}}})(A) = \opn{RHom}_{A}(R', R)$
in $\dcat{D}(A^{\mrm{en}}, \mrm{gr})$. 
Note that by Theorem \ref{thm:3820} we have 
$T, T' \in \dcat{D}^{\mrm{b}}_{(\mrm{f}, ..)}(A^{\mrm{en}}, \mrm{gr})$.
Hence the graded bimodules $\opn{H}^k(T)$ and $\opn{H}^k(T')$ are all bounded 
below. 

According to Lemma \ref{lem:3917} and Theorem \ref{thm:3820} (applied twice) 
there are isomorphisms 
\[ \begin{aligned}
& T' \ot^{\mrm{L}}_{A} T \cong (D'_A \circ D_{A^{\mrm{op}}})(T') = 
(D'_A \circ D_{A^{\mrm{op}}} \circ D_A \circ D'_{A^{\mrm{op}}})(A) 
\\ & \quad 
\cong (D'_A \circ D'_{A^{\mrm{op}}})(A) \cong A 
\end{aligned} \]
in $\dcat{D}(A^{\mrm{en}}, \mrm{gr})$. 
By symmetry, there is also an isomorphism 
$T \ot^{\mrm{L}}_{A} T' \cong A$ 
in $\dcat{D}(A^{\mrm{en}}, \mrm{gr})$. 
By Lemma \ref{lem:3916} there is an isomorphism
$T \cong A(\phi, i)[j]$
in $\dcat{D}(A^{\mrm{en}}, \mrm{gr})$, for a unique automorphism $\phi$ and 
unique integers $i, j$. 

Finally we obtain these isomorphisms 
\[ \begin{aligned}
& R \ot_A A(\phi, i)[j] \cong 
R \ot^{\mrm{L}}_{A} T \cong^{\dag}
(D'_A \circ D_{A^{\mrm{op}}})(R) 
\\ & \quad 
\cong (D'_A \circ D_{A^{\mrm{op}}} \circ D_A)(A) \cong D'_A(A) = R'
\end{aligned} \]
in $\dcat{D}(A^{\mrm{en}}, \mrm{gr})$.
The isomorphism $\cong^{\dag}$ is from Lemma \ref{lem:3917}. 
\end{proof}






\mysubsection{Balanced DC: Definition, Uniqueness and Local Duality}
\label{subsec:bal-NCDC-LocDu}

Our goal in this subsection is to relate graded NC dualizing complexes with 
derived torsion. We continue with Conventions \ref{conv:3700} and 
\ref{conv:4560}, and we assume the following setup:

\begin{setup} \label{setup:3925}
$\K$ is a base field, and $A$ is a noetherian connected graded 
$\K$-ring, with augmentation ideal $\m$. 
The augmentation ideal of the opposite ring $A^{\mrm{op}}$ is 
$\m^{\mrm{op}}$. 
\end{setup}

Recall the graded bimodule $A^* = \opn{Hom}_{\K}(A, \K)$. 

\begin{dfn}[\cite{Ye1}] \label{dfn:3714} 
Under Setup \tup{\ref{setup:3925}}, a {\em balanced dualizing complex} 
\index{Dualizing complex! balanced}
over $A$ is a pair $(R, \be)$, where:
\begin{itemize}
\item[(B1)] $R \in \dcat{D}(A^{\mrm{en}}, \mrm{gr})$ is a graded NC 
dualizing 
complex over $A$ (Definition \ref{dfn:3716}), with symmetric derived 
$\m$-torsion (Definition \ref{dfn:3779}). 

\item[(B2)] $\be : \mrm{R} \Ga_{\m}(R) \iso A^*$
is an isomorphism in $\dcat{D}(A^{\mrm{en}}, \mrm{gr})$, called a 
{\em balancing isomorphism}.
\end{itemize}
\end{dfn}

The lack of left-right symmetry in item (B2) of this 
definition will be removed in Corollary \ref{cor:3762} below. 

\begin{rem} \label{rem:4030}
Definition \ref{dfn:3714} is a bit more sophisticated than the 
original definition in \cite{Ye1}. There the condition was that 
$\mrm{R} \Ga_{\m}(R) \cong \mrm{R} \Ga_{\m^{\mrm{op}}}(R) \cong A^*$
in $\dcat{D}(A^{\mrm{en}}, \mrm{gr})$, but a balancing isomorphism 
$\be$ was not specified. The current improved definition is influenced by later 
research, especially \cite{YeZh1} and \cite{VyYe}.
\end{rem}

The next theorem is implicit in \cite{Ye1} and \cite{YeZh1}. 

\begin{thm}[Uniqueness of BDC] \label{thm:3710}
Under Setup \tup{\ref{setup:3925}}, suppose that $(R, \be)$ and $(R', \be')$ 
are balanced dualizing complexes over $A$. 
Then there is a unique isomorphism 
$\psi : R' \iso R$
in $\dcat{D}(A^{\mrm{en}}, \mrm{gr})$, 
such that 
$\be \circ \mrm{R} \Ga_{\m}(\psi) = \be'$
as isomorphisms 
$\mrm{R} \Ga_{\m}(R') \iso A^*$
in $\dcat{D}(A^{\mrm{en}}, \mrm{gr})$.
\end{thm}

\begin{proof}
By Theorem \ref{thm:3780} there is an isomorphism 
$\psi^{\dag} : R'\iso R \ot_A L$
in $\dcat{D}(A^{\mrm{en}}, \mrm{gr})$, where 
$L = A(\phi, i)[j]$ is a translate of a twisted bimodule. 
Applying the functor $\mrm{R} \Ga_{\m}$ to $\psi^{\dag}$ we get a diagram of 
isomorphisms
\begin{equation} \label{eqn:3920}
\UseTips \xymatrix @C=10ex @R=6ex {
\mrm{R} \Ga_{\m}(R')
\ar[r]^(0.40){ \mrm{R} \Ga_{\m}(\psi^{\dag}) }_(0.40){\cong}
\ar[d]_{\be'}^{\cong}
&
\mrm{R} \Ga_{\m} \bigl( R \ot_A L \bigr)
\ar[r]^{\ga}_{\cong}
&
\mrm{R} \Ga_{\m}(R) \ot_A L
\ar[d]_{\be \, \ot \, \opn{id}}^{\cong}
\\
A^*
&
&
A^* \ot_A L
}
\end{equation}
in $\dcat{D}(A^{\mrm{en}}, \mrm{gr})$.
The isomorphism $\ga$ is the obvious one. 
Therefore $A^* \cong A^* \ot_A L$ in $\dcat{M}(A^{\mrm{en}}, \mrm{gr})$.
Due to the NC Graded Matlis Duality
(Theorem \ref{thm:4045}) we have $L \cong A$ in 
$\dcat{M}(A^{\mrm{en}}, \mrm{gr})$. 

Let us now rewrite diagram (\ref{eqn:3920}) with $L = A$. This is the solid 
diagram
\begin{equation} \label{eqn:3921}
\UseTips \xymatrix @C=10ex @R=6ex {
\mrm{R} \Ga_{\m}(R')
\ar[r]^(0.50){ \mrm{R} \Ga_{\m}(\psi^{\dag}) }_(0.50){\cong}
\ar[d]_{\be'}^{\cong}
&
\mrm{R} \Ga_{\m}(R)
\ar[d]_{\be}^{\cong}
\\
A^*
\ar@{-->}[r]^{c \cd (-)}_{\cong}
&
A^*
}
\end{equation}
in $\dcat{D}(A^{\mrm{en}}, \mrm{gr})$.
Because the automorphisms of $A^*$ in $\dcat{D}(A^{\mrm{en}}, \mrm{gr})$ are 
multiplication by nonzero elements of $\K$, there is a unique 
$c \in \K^{\times}$ for which the diagram (\ref{eqn:3921}) is commutative.
Hence $\psi := c^{-1} \cd \psi^{\dag}$
is the unique isomorphism $R' \iso R$ satisfying 
$\be \circ \mrm{R} \Ga_{\m}(\psi) = \be'$. 
\end{proof}

The next theorem is a noncommutative version of Grothendieck's Local Duality 
Theorem \cite[Theorem V.6.2]{RD}. Balanced dualizing complexes were 
invented in order to make this theorem hold. 
Theorem \ref{thm:4616} first appeared in \cite{Ye1}, in a slightly different 
formulation, and with a complicated proof. The current formulation of the 
theorem, as well as the proof, are taken from 
\cite[Proposition 3.4]{CWZ}. We thank R. Vyas for pointing out this proof to 
us. 

\begin{thm}[Local Duality] \label{thm:4616}
\index{Local Duality Theorem}  
Under Setup \tup{\ref{setup:3925}}, let 
$(R, \be)$ be a balanced dualizing complex over $A$.
There is a morphism
\[ \xi : \opn{RHom}_A(-, R) \to \bigl( \mrm{R} \Ga_{\m} (-) \bigr)^* \] 
of triangulated functors 
$\dcat{D}(A, \mrm{gr})^{\mrm{op}} \to \dcat{D}(A^{\mrm{op}}, \mrm{gr})$,
such that for every 
$M \in \lb \dcat{D}^{+}_{\mrm{f}}(A, \mrm{gr})$ 
the morphism 
\[ \xi_M : \opn{RHom}_A(M, R) \to \bigl( \mrm{R} \Ga_{\m} (M) \bigr)^* \]
in $\dcat{D}(A^{\mrm{op}}, \mrm{gr})$
is an isomorphism. 
\end{thm}

Note that this result resembles Van den Bergh's Local Duality, Corollary 
\ref{cor:3901} -- but this is misleading: the assumptions are not the same. 

\begin{proof} 
The proof is divided into three steps. 

\smallskip \noindent 
Step 1. Choose K-injective resolutions $R \to J$ and
$\rho : \Ga_{\m}(J) \to K$ in $\dcat{C}_{\mrm{str}}(A^{\mrm{en}}, \mrm{gr})$.
For each $M \in \dcat{D}(A, \mrm{gr})$ choose a K-injective resolution 
$M \to I_M$ in $\dcat{C}_{\mrm{str}}(A, \mrm{gr})$. 
There is an obvious homomorphism 
\[ \phi_M : \opn{Hom}_A(I_M, J) \to 
\opn{Hom}_A \bigl( \Ga_{\m}(I_M), \Ga_{\m}(J) \bigr)  \] 
in $\dcat{C}_{\mrm{str}}(A^{\mrm{op}}, \mrm{gr})$. 
Next we have the homomorphism 
\[ \opn{Hom}_A(\opn{id}, \rho) :  
\opn{Hom}_A \bigl( \Ga_{\m}(I_M), \Ga_{\m}(J) \bigr) \ \to \ 
\opn{Hom}_A \bigl( \Ga_{\m}(I_M), K \bigr) \] 
in $\dcat{C}_{\mrm{str}}(A^{\mrm{op}}, \mrm{gr})$. 
Composing these homomorphisms, and then going to the derived category, we 
obtain a morphism 
\begin{equation} \label{eqn:4835}
\begin{aligned}
& 
\th_M := \opn{Q} \bigl( \opn{Hom}_A(\opn{id}, \rho) \circ \phi_M \bigr) :
\\ & \quad 
\opn{RHom}_A(M, R) \to \opn{RHom}_A \bigl( \mrm{R} \Ga_{\m}(M), 
\mrm{R} \Ga_{\m}(R) \bigr)
\end{aligned}
\end{equation}
in $\dcat{D}(A^{\mrm{op}}, \mrm{gr})$. 
We have the balancing isomorphism 
\begin{equation} \label{eqn:4616}
\be : \mrm{R} \Ga_{\m}(R) \iso A^*
\end{equation}
in $\dcat{D}(A^{\mrm{en}}, \mrm{gr})$.
According to Proposition \ref{prop:3795} there is an isomorphism 
\[ \al_M : 
\opn{RHom}_A \bigl( \mrm{R} \Ga_{\m}(R), A^* \bigr) \iso 
\bigl( \mrm{R} \Ga_{\m}(R) \bigr)^* \]
in $\dcat{D}(A^{\mrm{op}}, \mrm{gr})$, and it is functorial in $M$.  
Let us define the morphism 
\begin{equation} \label{eqn:4615}
\xi_M :=  \al_M \circ \opn{RHom}(\opn{id}, \be) \circ \th_M :
\opn{RHom}_A(M, R) \to \bigl( \mrm{R} \Ga_{\m} (M) \bigr)^*
\end{equation}
in $\dcat{D}(A^{\mrm{op}}, \mrm{gr})$.
By construction the morphism $\xi_M$ is functorial in $M$. 

\medskip \noindent
Step 2. Take the complex $M := R$. 
By definition $R$ has the NC derived Morita property on the $A$ side.
This implies that the element 
$\opn{id}_R \in \opn{H}^0 \bigl( \opn{RHom}_A(R, R) \bigr)$
is a basis of this rank $1$ graded-free $A^{\mrm{op}}$-module. 
Because of the isomorphism (\ref{eqn:4616}) and by graded Matlis Duality,
we know that $\mrm{R} \Ga_{\m}(R)$ has the NC derived Morita property on the 
$A$ side. Therefore the element 
\[ \opn{id}_{\mrm{R} \Ga_{\m}(R)} \in 
\opn{H}^0 \bigl( \opn{RHom}_A \bigl( \mrm{R} \Ga_{\m}(R),
\mrm{R} \Ga_{\m}(R) \bigr) \bigr) \]
is a basis of this rank $1$ graded-free $A^{\mrm{op}}$-module. 
The construction of the morphism $\th_R$ in (\ref{eqn:4835}) shows that 
% \begin{equation} \label{eqn:4837}
$\opn{H}^0(\th_R)(\opn{id}_R) = \opn{id}_{\mrm{R} \Ga_{\m}(R)}$.
% \end{equation}
We conclude that $\opn{H}^0(\th_R)$ is an isomorphism. Since all other 
cohomologies of the complexes below vanish, it follows that 
\begin{equation} \label{eqn:4838}
\begin{aligned}
& 
\th_R  : \opn{RHom}_A(R, R) \to 
\opn{RHom}_A \bigl( \mrm{R} \Ga_{\m}(R), \mrm{R} \Ga_{\m}(R) \bigr)
\end{aligned}
\end{equation}
is an isomorphism in 
$\dcat{D}(A^{\mrm{op}}, \mrm{gr})$.
Therefore, in view of (\ref{eqn:4615}), $\xi_R$ is also an isomorphism. 

\medskip \noindent
Step 3.
Define the functor 
\[ F := \bigl( \mrm{R} \Ga_{\m} (-) \bigr)^* : 
\dcat{D}^{}(A, \mrm{gr})^{\mrm{op}} \to 
\dcat{D}^{}(A^{\mrm{op}}, \mrm{gr}) . \]
Thus $\xi : D_{A} \to F$
is a morphism of triangulated functors, and we must prove that 
$\xi_M : D_{A}(M) \to F(M)$
is an isomorphism for every 
$M \in \dcat{D}^{+}_{\mrm{f}}(A, \mrm{gr})$.
We know that 
$D_{A^{\mrm{op}}} : \dcat{D}^{-}_{\mrm{f}}(A^{\mrm{op}}, \mrm{gr})^{\mrm{op}} 
\to \dcat{D}^{+}_{\mrm{f}}(A, \mrm{gr})$
is an equivalence. Thus it suffices to prove that the morphism  
\begin{equation} \label{eqn:4839}
\ze := \xi \circ \opn{id}_{D_{A^{\mrm{op}}}} : 
D_{A} \circ D_{A^{\mrm{op}}} \to F \circ D_{A^{\mrm{op}}} 
\end{equation}
of functors
$\dcat{D}^{-}_{\mrm{f}}(A^{\mrm{op}}, \mrm{gr}) \to 
\dcat{D}(A^{\mrm{op}}, \mrm{gr})$
is an isomorphism. 

Consider the object 
$A \in \dcat{D}^{-}_{\mrm{f}}(A^{\mrm{op}}, \mrm{gr})$.
By formula (\ref{eqn:4839}) the diagram 
\[ \UseTips \xymatrix @C=8ex @R=6ex {
(D_A \circ D_{A^{\mrm{op}}})(A)
\ar[d]_{\ze_A}
\ar[r]^{\opn{id}}
&
\opn{RHom}_A(R, R) 
\ar[d]^{\xi_R}
\\
(F \circ D_{A^{\mrm{op}}})(A)
\ar[r]^{\opn{id}}
&
\bigl( \mrm{R} \Ga_{\m} (R) \bigr)^* 
} \]
in $\dcat{D}(A^{\mrm{op}}, \mrm{gr})$ is commutative. 
In step 2 we proved that $\xi_R$ is an isomorphism. Hence $\ze_A$ is an 
isomorphism. 

The functors $D_{A}$ and $D_{A^{\mrm{op}}}$ have bounded cohomological 
displacements, and hence so does their composition 
$D_{A} \circ D_{A^{\mrm{op}}}$. (In fact the composition is isomorphic to the 
identity functor, so it has cohomological displacement $[0, 0]$.) 
The functor $\mrm{R} \Ga_{\m}$ has bounded below 
cohomological displacement, and the contravariant functor $(-)^*$ has 
cohomological displacement $[0, 0]$, so their composition $F$ has 
bounded above cohomological displacement.
We conclude that the functors 
$D_{A} \circ D_{A^{\mrm{op}}}$ and $F \circ D_{A^{\mrm{op}}}$ have 
bounded above cohomological displacements.
We have already established that $\ze_A$ is an isomorphism. 
Theorem \ref{thm:2135}(1), which is valid also in the algebraically graded 
setting, says that $\ze_N$ is an isomorphism for every 
$N \in \dcat{D}^{-}_{\mrm{f}}(A^{\mrm{op}}, \mrm{gr})$.
\end{proof}

\begin{cor}[\cite{YeZh0}] \label{cor:4655}  
Under Setup \tup{\ref{setup:3925}}, assume $A$ has a balanced dualizing 
complex. Then $A$ satisfies the $\chi$ condition, and it has finite local
cohomological dimension. 
\end{cor}

\begin{proof}
Given $M \in \dcat{M}_{\mrm{f}}(A, \mrm{gr})$, Theorem \ref{thm:4616}
provides an isomorphism 
\[ \xi_M : \opn{RHom}_A(M, R) \iso \bigl( \mrm{R} \Ga_{\m}(M) \bigr)^* \] 
in $\dcat{D}(A^{\mrm{op}}, \mrm{gr})$. 
This gives an isomorphism 
\begin{equation} \label{eqn:4655}
\opn{H}^{-p}(\xi_M) : \opn{Ext}^{-p}_A(M, R) \iso \bigl( \mrm{H}^p_{\m}(M) 
\bigr)^* 
\end{equation}
in $\dcat{M}(A^{\mrm{op}}, \mrm{gr})$
for every $p$. By Theorem \ref{thm:3820}(1) we know that 
$\opn{Ext}^{-p}_A(M, R) \in \dcat{M}_{\mrm{f}}(A^{\mrm{op}}, \mrm{gr})$.
Using NC graded Matlis Duality (Theorem \ref{thm:4045})  we get
\begin{equation} \label{eqn:3840}
\mrm{H}^p_{\m}(M) \cong \bigl( \opn{Ext}^{-p}_A(M, R) \bigr)^* \in 
\dcat{M}_{\mrm{cof}}(A, \mrm{gr}) . 
\end{equation}
By Proposition \ref{prop:3855} we conclude that $A$ satisfies the left $\chi$ 
condition. The op-symmetry of the situation implies that $A^{\mrm{op}}$ also 
satisfies the left $\chi$ condition. 

Because $R \in \dcat{D}^{\mrm{b}}(A^{\mrm{en}}, \mrm{gr})$
and is nonzero, we have 
$q_0 := \opn{inf}(\opn{H}(R)) \in \Z$. 
So $\opn{Ext}^{q}_A(M, R) = 0$
for all $M \in \dcat{M}(A, \mrm{gr})$ and $q < q_0$.
Taking $p := -q$ and using (\ref{eqn:3840}) we get 
$\mrm{H}^p(\mrm{R} \Ga_{\m}(M)) = \mrm{H}^p_{\m}(M) = 0$
for all $p > -q_0$. So the functor $\mrm{R} \Ga_{\m}$ has finite 
cohomological dimension. The op-symmetry of the situation implies that the 
functor $\mrm{R} \Ga_{\m^{\mrm{op}}}$ also has finite cohomological 
dimension.
\end{proof}

\begin{cor} \label{cor:3762}
Let $(R, \be)$ be a balanced dualizing complex over $A$. Then 
there is a unique symmetry isomorphism
$\ep_{R} : \mrm{R} \Ga_{\m}(R) \iso \mrm{R} \Ga_{\m^{\mrm{op}}}(R)$
in $\dcat{D}(A^{\mrm{en}}, \mrm{gr})$, in the sense of Definition 
\tup{\ref{dfn:3779}}. Hence there is a unique isomorphism
$\be^{\mrm{op}} : \mrm{R} \Ga_{\m^{\mrm{op}}}(R) \iso A^*$ 
$\dcat{D}(A^{\mrm{en}}, \mrm{gr})$ such that 
$\be^{\mrm{op}} \circ \ep_{R} = \be$.
\end{cor}

\begin{proof}
By Corollary \ref{cor:4655} the ring $A$ has finite local
cohomological dimension. So Theorem \ref{thm:3720} applies. 
\end{proof}



\mysubsection{Balanced DC: Existence}
\label{subsec:bal-NCDC-exis}

In this subsection we are going to prove the existence of balanced dualizing 
complexes in two important situations. We shall adhere to Conventions 
\ref{conv:3700} and \ref{conv:4560}, and to Setup \ref{setup:3925}.

The next lemma is \cite[Proposition 4.4]{Ye1}. 

\begin{lem} \label{lem:3855}
Let $R$ be a graded NC dualizing complex over $A$, with duality functors 
$D_A, D_{A^{\mrm{op}}} : \dcat{D}(A^{\mrm{en}}, \mrm{gr}) \to 
\dcat{D}(A^{\mrm{en}}, \mrm{gr})$.
The following four conditions are equivalent.  
\begin{enumerate}
\rmitem{i} There is an isomorphism 
$D_A(\K) \cong \K$ in $\dcat{D}(A^{\mrm{en}}, \mrm{gr})$. 

\rmitem{i$'$} There is an isomorphism 
$D_{A^{\mrm{op}}}(\K) \cong \K$ in $\dcat{D}(A^{\mrm{en}}, \mrm{gr})$. 

\rmitem{ii} There is an isomorphism 
$\mrm{R} \Ga_{\m}(R) \cong A^* \ot_A L$
in $\dcat{D}(A^{\mrm{en}}, \mrm{gr})$, for some 
graded invertible $A$-bimodule $L$ generated in algebraic degree $0$. 

\rmitem{ii$'$} There is an isomorphism 
$\mrm{R} \Ga_{\m^{\mrm{op}}}(R) \cong L' \ot_A A^*$
in $\dcat{D}(A^{\mrm{en}}, \mrm{gr})$, for some 
graded invertible $A$-bimodule $L'$ generated in algebraic degree $0$. 
\end{enumerate}
\end{lem}

Note that by Proposition \ref{prop:9325}, a graded invertible $A$-bimodule $L$ 
generated in algebraic degree $0$ is isomorphic to $A(\phi, 0)$ for some 
automorphism $\phi$. 

\begin{proof} \mbox{}

\smallskip \noindent
(i) $\Rightarrow$ (i$'$): 
It is given that $D_A(\K) \cong \K$. 
By Theorem  \ref{thm:3820} we know that
$\K \cong D_{A^{\mrm{op}}}(D_A(\K))$. 
Together we get $D_{A^{\mrm{op}}}(\K) \cong \K$. 

\medskip \noindent 
(i$'$) $\Rightarrow$ (i): The same, by op-symmetry (replacing $A$ with 
$A^{\mrm{op}}$).

\medskip \noindent 
$\bigl($(i) and (i$'$)$\bigr)$ $\Rightarrow$ (ii):
Recall that a graded invertible bimodule $L$ is free of rank $1$ on both sides. 
So it suffices to prove separately that 
\begin{equation} \label{eqn:3922}
\mrm{R} \Ga_{\m}(R) \cong A^* \ \  \tup{in} \ \dcat{D}(A, \mrm{gr}) 
\end{equation}
and
\begin{equation} \label{eqn:3923}
\mrm{R} \Ga_{\m}(R) \cong A^* \ \ \tup{in} \ 
\dcat{D}(A^{\mrm{op}}, \mrm{gr}) .
\end{equation}

Let $R \to I$ be a minimal graded-injective resolution of $R$ over $A$, so that 
$\mrm{R} \Ga_{\m}(R) \cong \Ga_{\m}(I)$
in $\dcat{D}(A, \mrm{gr})$. 
By Lemma \ref{lem:3840} the subcomplex 
$W := \opn{Soc}(I) = \opn{Hom}_{A}(\K, I) \sub I$
has zero differential. But  
$D_A(\K) = \opn{RHom}_A(\K, R) \lb \cong \opn{Hom}_{A}(\K, I) = W$ 
in $\dcat{D}(A, \mrm{gr})$.
From condition (i) we see that $W \cong \K$ in 
$\dcat{D}(A, \mrm{gr})$; and therefore also 
$W \cong \K$ in $\dcat{C}_{\mrm{str}}(A, \mrm{gr})$.
According to Lemma \ref{lem:3839} and Proposition \ref{prop:4811}, 
$W^p \sub \Ga_{\m}(I^p)$ is an essential submodule.
Again using Lemma \ref{lem:3839} we conclude that 
$\Ga_{\m}(I) \cong A^*$ in $\dcat{C}_{\mrm{str}}(A^{\mrm{en}}, \mrm{gr})$. 
Hence the isomorphism (\ref{eqn:3922}) holds. 

For the other isomorphism, let 
$\dcat{M}_{\mrm{f}/ \K}(A, \mrm{gr}) \sub \dcat{M}(A, \mrm{gr})$
be the full subcategory on the modules that are finite over $\K$, or in other 
words, the finite length $A$-modules. Conditions (i) and (i$'$) imply, using 
induction on length of modules, that the functor 
$\opn{H}^0 \circ \, D_A : \dcat{M}_{\mrm{f}/ \K}(A, \mrm{gr})^{\mrm{op}} \to
\dcat{M}_{\mrm{f}/ \K}(A^{\mrm{op}}, \mrm{gr})$
is an equivalence, with quasi-inverse $\opn{H}^0 \circ \, D_{A^{\mrm{op}}}$. 
For every $j \geq 1$ we have 
$A / \m^j \in \dcat{M}_{\mrm{f}/ \K}(A, \mrm{gr})$. 
The duality provides an isomorphism 
\[ \opn{Hom}_{A^{\mrm{op}}} 
\bigl( \K, \opn{H}^0(D_{A}(A / \m^j)) \bigr)
\cong \opn{Hom}_{A} (A / \m^j, \K) \cong \K . \]
This we rewrite as 
\[ \opn{Hom}_{A^{\mrm{op}}} 
\bigl( \K, \opn{Ext}^0_A(A / \m^j, R) \bigr) \cong \K . \]
By functoriality we can pass to the limit. Using Proposition \ref{prop:3850}(5)
we obtain an isomorphism 
$\opn{Hom}_{A^{\mrm{op}}} \bigl( \K, \opn{H}^0_{\m}(R) \bigr) \cong \K$.
I.e.\ the socle of the graded $A^{\mrm{op}}$-module 
$\opn{H}^0_{\m}(R)$ is $\K$. 
By Proposition \ref{prop:4811} there is an essential 
monomorphism 
$\psi : \opn{H}^0_{\m}(R) \inj A^*$
in $\dcat{M}(A^{\mrm{op}}, \mrm{gr})$. 
But from the previous calculation we already know that 
$\opn{H}^0_{\m}(R) \cong A^*$ in $\dcat{M}(A, \mrm{gr})$, and hence also in 
$\dcat{M}(\K, \mrm{gr})$. As these are degreewise finite graded $\K$-modules,
$\psi$ must be bijective. 
We conclude that the isomorphism (\ref{eqn:3923}) holds.

\medskip \noindent 
$\bigl($(i) and (i$'$)$\bigr)$ $\Rightarrow$ (ii$'$): 
The same, by op-symmetry.

\medskip \noindent 
(ii) $\Rightarrow$ (i): 
By Proposition \ref{prop:3990} there is an isomorphism 
\[ \opn{RHom}_A(\K, R) \cong
\opn{RHom}_A \bigl( \K, \mrm{R} \Ga_{\m}(R) \bigr)  \]
in $\dcat{D}(A^{\mrm{en}}, \mrm{gr})$. 
Now condition (ii) implies that 
$\mrm{R} \Ga_{\m}(R) \cong A^*$
in $\dcat{D}(A, \mrm{gr})$. So we get 
\[ \opn{RHom}_A(\K, R) \cong \opn{RHom}_A(\K, A^*) 
\cong \opn{Hom}_A(\K, A^*) \cong \K \]
in $\dcat{D}(\K, \mrm{gr})$.
This implies that  
$\opn{RHom}_A(\K, R) \cong \K$
in $\dcat{D}(A^{\mrm{en}}, \mrm{gr})$.

\medskip \noindent 
(ii$'$) $\Rightarrow$ (i$'$): The same, by op-symmetry. 
\end{proof}

\begin{dfn}[\cite{Ye1}] \label{dfn:3781} 
Under Setup \tup{\ref{setup:3925}}, a {\em prebalanced dualizing complex} over 
$A$ is a graded NC dualizing complex $R$ over $A$
that satisfies the equivalent conditions of Lemma \ref{lem:3855}.
\end{dfn}

Here is the first of two existence theorems for balanced dualizing complexes. 
The second is Theorem \ref{thm:3713}. 

\begin{thm}[Existence of BDC, \cite{Ye1}] \label{thm:3781}
Under Setup \tup{\ref{setup:3925}}, assume that $A$ has a prebalanced 
dualizing complex. Then $A$ has a balanced dualizing complex.
\index{Dualizing complex! balanced}
\end{thm}

The proof comes after the next lemma. Let 
$\dcat{M}_{\mrm{f}/ \K}(A^{\mrm{en}}, \mrm{gr}) \sub 
\dcat{M}(A^{\mrm{en}}, \mrm{gr})$
be the full subcategory on the bimodules that are finite over $\K$.

\begin{lem} \label{lem:3925}
Suppose $R$ is a prebalanced dualizing complex over $A$ such that 
$\mrm{R} \Ga_{\m}(R) \cong A^*$ in 
$\dcat{D}(A^{\mrm{en}}, \mrm{gr})$. 
Then there is an isomorphism 
$D_A \cong \opn{Hom}_{\K}(-, \K)$
of functors 
$\dcat{M}_{\mrm{f}/ \K}(A^{\mrm{en}}, \mrm{gr})^{\mrm{op}} \to
\dcat{M}_{\mrm{f}/ \K}(A^{\mrm{en}}, \mrm{gr})$.
\end{lem}

\begin{proof}
Let us choose an isomorphism 
$\be : \mrm{R} \Ga_{\m}(R) \iso A^*$
in $\dcat{D}(A^{\mrm{en}}, \mrm{gr})$. 
Take $M \in \dcat{M}_{\mrm{f}/ \K}(A^{\mrm{en}}, \mrm{gr})$. 
There is a sequence of functorial isomorphisms 
\[ \begin{aligned}
& D_A(M) = \opn{RHom}_A(M, R)
\cong^{\dag} \opn{RHom}_A \bigl( M, \mrm{R} \Ga_{\m}(R) \bigr)
\\ & \quad 
\cong^{\ddag} \opn{RHom}_A(M, A^*) 
\cong \opn{Hom}_{\K}(M, \K) 
\end{aligned} \]
in $\dcat{D}(A^{\mrm{en}}, \mrm{gr})$. 
The isomorphism $\cong^{\dag}$ is by Proposition \ref{prop:3990}, and the 
isomorphism $\cong^{\ddag}$ comes from $\be$. 
\end{proof}

\begin{proof}[Proof of Theorem \tup{\ref{thm:3781}}]
The proof is in two steps. 

\smallskip \noindent 
Step 1. Let $R'$ be a prebalanced dualizing complex over $A$. 
This means that there is an isomorphism 
\begin{equation} \label{eqn:4575}
\mrm{R} \Ga_{\m}(R') \cong A^* \ot_A L 
\end{equation}
in $\dcat{D}(A^{\mrm{en}}, \mrm{gr})$, for some 
graded invertible $A$-bimodule $L$ generated in algebraic degree $0$. Define 
$L^{\vee} := \opn{Hom}_A(L, A)$. Thus, if $L \cong A(\phi, i)$, then 
$L^{\vee} \cong A(\phi^{-1}, -i)$ in $\dcat{M}(A^{\mrm{en}}, \mrm{gr})$.

Next let 
\begin{equation} \label{eqn:3929}
R := R'  \ot_A L^{\vee} \in \dcat{D}(A^{\mrm{en}}, \mrm{gr}) ,
\end{equation}
which is a graded dualizing complex. 
It is clear from (\ref{eqn:4575}) that there is an isomorphism 
\begin{equation} \label{eqn:3927}
\be : \mrm{R} \Ga_{\m}(R) \iso A^* \ \ \tup{in} \ \  
\dcat{D}(A^{\mrm{en}}, \mrm{gr}) ,
\end{equation}
and this is the balancing isomorphism that we take. 
To complete the proof that $(R, \be)$ is a balanced dualizing complex, it 
remains to prove that 
\begin{equation} \label{eqn:3925}
\mrm{R} \Ga_{\m^{\mrm{op}}}(R) \cong A^*
\ \ \tup{in} \ \ \dcat{D}(A^{\mrm{en}}, \mrm{gr}) .
\end{equation}
This will be done in the second step. 
(We can't use Theorem \ref{thm:3720} on symmetric derived torsion, since  
that would involve circular reasoning.)

\medskip \noindent 
Step 2. The graded dualizing complex $R$ is prebalanced, by (\ref{eqn:3927}). 
So according to Lemma \ref{lem:3855} there is an isomorphism 
\begin{equation} \label{eqn:3928}
\mrm{R} \Ga_{\m^{\mrm{op}}}(R) \cong L^{\diamond} \ot_A A^* 
\end{equation}
in $\dcat{D}(A^{\mrm{en}}, \mrm{gr})$, for some 
invertible graded $A$-bimodule $L^{\diamond}$ generated in algebraic degree 
$0$. 

For every $j \geq 1$ the bimodule 
$A / \m^j$ belongs to $\dcat{M}_{\mrm{f}/ \K}(A^{\mrm{en}}, \mrm{gr})$,
and it is $\m$-torsion and $\m^{\mrm{op}}$-torsion. 
Therefore we get isomorphisms 
\begin{equation} \label{eqn:3992}
\begin{aligned}
& D_{A^{\mrm{op}}}(A / \m^j) = \opn{RHom}_{A^{\mrm{op}}}(A / \m^j, R)
\\ & \quad 
\cong^{\mrm{(a)}} 
\opn{RHom}_{A^{\mrm{op}}} 
\bigl( A / \m^j, \mrm{R} \Ga_{\m^{\mrm{op}}}(R) \bigr)
\\ & \quad 
\cong^{\mrm{(b)}} 
\opn{RHom}_{A^{\mrm{op}}}(A / \m^j, L^{\diamond} \ot_A  A^*) 
\\ & \quad 
\cong L^{\diamond} \ot_A \opn{RHom}_{A^{\mrm{op}}}(A / \m^j, A^*) 
\\ & \quad 
\cong L^{\diamond} \ot_A  \opn{Hom}_{\K}(A / \m^j, \K) 
= L^{\diamond} \ot_A (A / \m^j)^* 
\end{aligned} 
\end{equation}
in $\dcat{D}(A^{\mrm{en}}, \mrm{gr})$.
Explanation: the isomorphism $\cong^{\mrm{(a)}}$ is by Proposition 
\ref{prop:3990} (transcribed to the ring $A^{\mrm{op}}$); and the isomorphism 
$\cong^{\mrm{(b)}}$ is from equation (\ref{eqn:3928}).

There are also isomorphisms 
\begin{equation} \label{eqn:3993}
\begin{aligned}
&
A / \m^j \cong^1 D_A \bigl( D_{A^{\mrm{op}}}(A / \m^j) \bigr)
\cong^2 D_A \bigl( L^{\diamond} \ot_A (A / \m^j)^* \bigr)
\\ & \quad 
\cong^3 \bigl( L^{\diamond} \ot_A (A / \m^j)^* \bigr)^*
\cong^4 (A / \m^j) \ot_A (L^{\diamond})^{\vee}
\end{aligned} 
\end{equation}
in $\dcat{M}(A^{\mrm{en}}, \mrm{gr})$, where
$(L^{\diamond})^{\vee} := \opn{Hom}_A(L^{\diamond}, A)$.
Here is how they arise: the isomorphism $\cong^1$ is by Theorem \ref{thm:3820};
the isomorphism $\cong^2$ is from equation (\ref{eqn:3992}); 
and the isomorphism $\cong^3$ is due to Lemma \ref{lem:3925}.
For the isomorphism $\cong^4$, say 
$L^{\diamond} \cong A(\nu, k)$; then 
$(L^{\diamond})^{\vee} \cong A(\nu^{-1}, -k)$ and 
$(L^{\diamond})^{*} \cong A^*(\nu^{-1}, -k) = A^* \ot_{A} A(\nu^{-1}, -k)$.

Because the isomorphisms (\ref{eqn:3993}) hold for every $j \geq 1$, it follows 
that $(L^{\diamond})^{\vee} \cong A$ in 
$\dcat{M}(A^{\mrm{en}}, \mrm{gr})$,
and hence 
$L^{\diamond} \cong A$ in 
$\dcat{M}(A^{\mrm{en}}, \mrm{gr})$.
Then the isomorphism (\ref{eqn:3928}) becomes (\ref{eqn:3925}). 
\end{proof}

Recall the AS Gorenstein graded rings from Definition 
\ref{dfn:3782}. For an AS Gorenstein graded ring $A$ there are isomorphisms 
\[ \opn{RHom}_{A}(\K, A) \cong  \opn{RHom}_{A^{\mrm{op}}}(\K, A) \cong
\K(l)[-n] \]
in $\dcat{D}(\K, \mrm{gr})$. The numbers $n$ and $l$ are called the dimension 
and AS index of $A$, respectively. 

\begin{cor}[\cite{Ye1}] \label{cor:3780}
Let $A$ be an AS Gorenstein graded ring, of dimension $n$ and 
with AS index $l$. Then there is a unique automorphism $\phi$ of $A$ 
in $\catt{Rng}_{\mrm{gr}} \centover \K$ such that the complex 
$R_A := A(\phi, -l)[n]  \in \dcat{D}(A^{\mrm{en}}, \mrm{gr})$
is a balanced dualizing complex over $A$. 
\end{cor}

\begin{proof}
The complex 
$R' := A(-l)[n] \in \dcat{D}(A^{\mrm{en}}, \mrm{gr})$
is a prebalanced dualizing complex. Theorem \ref{thm:3781} asserts that $A$ has 
a balanced dualizing complex $R_A$. The construction of $R_A$ in formula 
(\ref{eqn:3929}) in the proof of  Theorem \ref{thm:3781} shows that 
$R_A = R' \ot_A A(\phi, 0) \cong A(\phi, -l)[n]$ 
for some automorphism $\phi$ of $A$. 

The uniqueness of $n, l$ and $\phi$ is a consequence of Theorems \ref{thm:3710} 
and \ref{thm:3780}. 
\end{proof}

Here are two important examples of AS regular graded rings and their balanced 
dualizing complexes. 

\begin{exa} \label{exa:3780}
Suppose $A$ is an {\em elliptic $3$-dimensional Artin-Schelter regular graded 
ring}; see \cite{ATV}. There is a normalizing regular element $g \in A_3$, and 
it defines an automorphism $\phi_g$ of $A$ by the formula 
$\phi_g(a) \cd g = g \cd a$ for $a \in A$. There is a constant 
$\la \in \K^{\times}$
arising from the elliptic curve associated to $A$, and it  defines 
an automorphism $\phi_{\la}$ of $A$ by the formula 
$\phi_{\la}(a) := \la^i \cd a$ for $a \in A_i$. According to 
\cite[Theorem 7.18]{Ye1} the balanced dualizing complex of $A$ is 
$R_A = A(\phi_g \circ \phi_{\la}, -3)[3]$.
\end{exa}

\begin{exa} \label{exa:3990}
Let $A := \K[t_1, \ldots, t_n]$, the commutative polynomial ring in $n$ 
variables, of algebraic degrees $\opn{deg}(t_i) = l_i \geq 1$. This is an 
Artin-Schelter regular graded ring, and a calculation with a Koszul complex 
shows that $\opn{RHom}_A(\K, A) \cong \K(l)[-n]$
in $\dcat{D}(A^{\mrm{en}}, \mrm{gr})$, where 
$l := \sum_{i = 1}^n l_i$. 
Because 
$\mrm{R} \Ga_{\m}(A) \cong A^*(l)[-n]$
is a central $A$-bimodule, it follows that the balanced dualizing dualizing 
complex of $A$ is $R_A := A(-l)[n]$. 

If we want to tie this with the commutative theory from Section
\ref{sec:dual-cplx-comm-rng}, and to algebraic geometry,
then we should take 
$R_A := \Om^n_{A / \K}[n]$,
where $\Om^n_{A / \K}$ is the graded-free $A$-module generated by the 
differential form 
$\d(t_1) \wedge \cdots \wedge \d(t_n)$ 
that has algebraic degree $l$. 
\end{exa}

Now an example that relies on Example \ref{exa:3990}. 

\begin{exa} \label{exa:3991}
Suppose $B$ is a commutative noetherian connected graded $\K$-ring. 
We can find a finite homomorphism $A \to B$ in 
$\catt{Rng}_{\mrm{gr}} \centover \K$ from a commutative polynomial ring $A$. 
Let $R_A := A(-l)[n]$ be the balanced dualizing complex of $A$, as in the 
previous example. Define 
$R_B := \opn{RHom}_{A}(B, R_A) \in \dcat{D}(B^{\mrm{en}}, \mrm{gr})$,
where we use a commutative graded-injective resolution $R_A \to I$, so that 
$R_B$ is a complex of central graded $B$-modules. 
By the arguments of Proposition \ref{prop:2200}, $R_B$ is graded NC 
dualizing complex over $B$. A calculation that we won't perform (but that is an 
easy variant of the proof of Theorem \ref{thm:4030} on the balanced 
trace morphism) shows that $R_B$ is balanced. 

The moral is that the commutative algebraically graded duality theory embeds 
within the noncommutative algebraically graded duality theory.  
\end{exa}

The next theorem is an important converse to Corollary \ref{cor:4655}.

Recall that for a graded NC dualizing complex $R$, the derived homothety 
morphism through $A^{\mrm{op}}$, namely the morphism 
$\opn{hm}^{\mrm{R}}_{R, A^{\mrm{op}}} : A \to \opn{RHom}_{A}(R, R)$ 
in \lb $\dcat{D}(A^{\mrm{en}}, \mrm{gr})$, is an isomorphism. See condition 
(iii) of Definition  \ref{dfn:3716}. 
By Van den Bergh's Local Duality (Corollary \ref{cor:3901}), for every complex 
$M \in \dcat{D}(A^{\mrm{en}},  \mrm{gr})$ there is an isomorphism 
\begin{equation} \label{eqn:4580}
\de_M : \opn{RHom}_A \bigl( M, (P_A)^* \bigr) \iso \mrm{R} \Ga_{\m}(M)^* 
\end{equation}
in $\dcat{D}(A^{\mrm{en}}, \mrm{gr})$.
Here 
$P_A = \mrm{R} \Ga_{\m}(A) \in \dcat{M}(A^{\mrm{en}}, \mrm{gr})$
is the dedualizing complex of $A$. 

\begin{thm}[Existence of BDC, Van den Bergh \cite{VdB}] \label{thm:3713} 
Under Setup \tup{\ref{setup:3925}}, assume the ring $A$ satisfies the special 
$\chi$ condition, and it has finite local cohomological dimension. 
Then $A$ has a balanced dualizing complex $(R_A, \be_A)$. 
\index{Dualizing complex! balanced}

More explicitly, the complex 
$R_A := (P_A)^* \in \dcat{D}(A^{\mrm{en}}, \mrm{gr})$
is a graded dualizing complex over $A$ with symmetric derived $\m$-torsion.
It has a balancing isomorphism
$\be_A : \mrm{R} \Ga_{\m}(R_A) \iso A^*$, 
which is the unique isomorphism 
in $\dcat{D}(A^{\mrm{en}}, \mrm{gr})$
such that the diagram 
\[ \UseTips \xymatrix @C=10ex @R=8ex {
A
\ar[r]_(0.3){\opn{hm}^{\mrm{R}}_{R, A^{\mrm{op}}}}^(0.3){\cong}
\ar@(u,u)[rr]^{(\be_A)^*}
&
\opn{RHom}_{A}(R_A, R_A)
\ar[r]_(0.57){\de_{R_A}}^(0.6){\cong}
&
\mrm{R} \Ga_{\m}(R_A)^*
} \]
in $\dcat{D}(A^{\mrm{en}}, \mrm{gr})$ is commutative.
\end{thm}

\begin{proof} 
The proof proceeds in 5 steps. 

\smallskip \noindent 
Step 1. We are given that $A$ has finite local cohomological dimension and it 
satisfies the special $\chi$ condition. By Proposition \ref{prop:3712}, the 
bimodule $A$ has weakly symmetric derived $\m$-torsion. Theorem \ref{thm:3720} 
says that the bimodule $A$ has symmetric derived $\m$-torsion; namely that there 
is an isomorphism 
\begin{equation} \label{eqn:3980}
\ep_A : P_A = \mrm{R} \Ga_{\m}(A) \iso \mrm{R} \Ga_{\m^{\mrm{op}}}(A) 
\end{equation}
in $\dcat{D}^{\mrm{b}}(A^{\mrm{en}}, \mrm{gr})$. 

The left $\chi$ condition for $A$, with Proposition \ref{prop:3855}, tell us 
that the complex $\mrm{R} \Ga_{\m}(A)$ has cofinite cohomology modules over 
$A$. By the same token, the left $\chi$ condition for $A^{\mrm{op}}$
tells us that the complex $\mrm{R} \Ga_{\m^{\mrm{op}}}(A)$ has cofinite 
cohomology modules over $A^{\mrm{op}}$. Using the isomorphism (\ref{eqn:3980}) 
we conclude that 
$P_A \in \dcat{D}^{\mrm{b}}_{(\mrm{cof}, \mrm{cof})}(A^{\mrm{en}}, \mrm{gr})$.
Now Graded Matlis Duality (Theorem \ref{thm:4045}) says that the complex $R_A$ 
satisfies 
\begin{equation} \label{eqn:4582}
R_A = (P_A)^* \in 
\dcat{D}^{\mrm{b}}_{(\mrm{f}, \mrm{f})}(A^{\mrm{en}}, \mrm{gr}) .
\end{equation}
This is condition (i) of Definition \ref{dfn:3716}. 

\medskip \noindent Step 2. 
Because the functor $\mrm{R} \Ga_{\m}$ has finite cohomological dimension, 
Van den Bergh's Local Duality (Corollary \ref{cor:3901})
says that the functor 
$\opn{RHom}_A \bigl( -, R \bigr) \lb \cong \mrm{R} \Ga_{\m}(-)^*$
also has finite cohomological dimension. This means that the complex $R_A$ has 
finite graded-injective dimension over $A$. Similarly, because 
$\mrm{R} \Ga_{\m^{\mrm{op}}}$ has finite cohomological dimension, the complex 
$R_A$ has finite graded-injective dimension over $A^{\mrm{op}}$. 
This is condition (ii) of Definition \ref{dfn:3716}.

\medskip \noindent Step 3.  
Now we shall prove that $R_A$ has the derived Morita property on the $A$ side. 
We have these isomorphisms in the category 
$\dcat{D}(A^{\mrm{en}}, \mrm{gr})$~:
\begin{equation} \label{eqn:3994}
\begin{aligned}
& \opn{RHom}^{}_A(R_A, R_A) = 
\opn{RHom}^{}_A \bigl( \mrm{R} \Ga_{\m}(A)^*, \mrm{R} \Ga_{\m}(A)^*
\bigr)  \\
& \quad \cong^{1}
\opn{RHom}^{}_{A^{\mrm{op}}} \bigl( \mrm{R} \Ga_{\m}(A), \mrm{R} 
\Ga_{\m}(A) \bigr) \\
& \quad \cong^2 
\opn{RHom}^{}_{A^{\mrm{op}}} 
\bigl( \mrm{R} \Ga_{\m^{\mrm{op}}}(A), \mrm{R} \Ga_{\m^{\mrm{op}}}(A) \bigr) 
\\ & \quad 
\cong^3 \opn{RHom}^{}_{A^{\mrm{op}}} 
\bigl( \mrm{R} \Ga_{\m^{\mrm{op}}}(A), A \bigr) 
\\ & \quad 
\cong^4 \opn{RHom}^{}_{A} 
\bigl( A^*, \mrm{R} \Ga_{\m^{\mrm{op}}}(A)^* \bigr) 
\\ & \quad 
\cong^5 \opn{RHom}^{}_{A} 
\bigl( A^*, \mrm{R} \Ga_{\m}(A)^* \bigr) 
\\ & \quad 
\cong^6 (A^{*})^* \cong^7 A \, .
\end{aligned}
\end{equation}
The isomorphisms $\cong^{1}$, $\cong^{4}$ and $\cong^{7}$ are from graded 
Matlis duality (Theorem \ref{thm:4045}); 
the isomorphisms $\cong^{2}$ and $\cong^{5}$ are by the symmetry isomorphism 
$\ep_A$; the isomorphism $\cong^{3}$ is due to Proposition \ref{prop:3990}; and 
the isomorphism $\cong^{6}$ is due to Van den Bergh's Local Duality (Corollary 
\ref{cor:3901}). 
These isomorphisms respect the derived homothety morphisms 
through $A^{\mrm{op}}$ (in lines 1, 5 and 6), the derived homothety morphisms
through $A$ (in lines 2, 3 and 4); and the canonical 
isomorphisms $A \to (A^*)^*$ and $A \to A$ in line 7. 
(Compare to the proof of Theorem \ref{thm:3681}, and the use of Lemma 
\ref{lem:4361} in it.)
The conclusion is that the derived homothety morphism
$\opn{hm}^{\mrm{R}}_{R, A^{\mrm{op}}} : A \to \opn{RHom}_A(R_A, R_A)$
is an isomorphism. 
This establishes the derived Morita property on the $A$ side for $R_A$. 

\medskip \noindent Step 4.
Replacing $A$ with $A^{\mrm{op}}$, using the symmetry isomorphism $\ep_A$, the 
same calculation as in step 3 shows that $R_A$ has the derived Morita property 
on the $A^{\mrm{op}}$ side. 
Thus condition (iii) of Definition \ref{dfn:3716} is verified. We conclude that 
$R_A$ is a graded NC dualizing complex. 

\medskip \noindent Step 5.  
In this step we prove that $R_A$ is balanced. 
By formula (\ref{eqn:4582}), the derived Morita property on the $A$ side, and 
Van den Bergh's Local Duality (Corollary \ref{cor:3901}), there are isomorphisms
\begin{equation} \label{eqn:4581}
A \underset{\cong}{\xar{\opn{hm}^{\mrm{R}}_{R_A, A^{\mrm{op}}}}} 
\opn{RHom}_{A}(R_A, R_A) \underset{\cong}{\xar{\de_{R_A}}}
\bigl( \mrm{R} \Ga_{\m}(R_A) \bigr)^* 
\end{equation}
in $\dcat{D}(A^{\mrm{en}}, \mrm{gr})$.
See the isomorphism (\ref{eqn:4580}) with $M := R_A = (P_A)^*$.
After we apply the functor $(-)^*$ to (\ref{eqn:4581}), Graded Matlis Duality 
(Theorem \ref{thm:4045}) says that the morphism 
\[ \be := 
\bigl( \de_{R_A} \circ \opn{hm}^{\mrm{R}}_{R_A, A^{\mrm{op}}} \bigr)^* : 
\Ga_{\m}(R_A) \to A^* \]
in $\dcat{D}(A^{\mrm{en}}, \mrm{gr})$ is an isomorphism. This is a 
balancing isomorphism for $R_A$, so item (B2) of Definition 
\ref{dfn:3714} is satisfied.

Similarly, the Derived Morita property on the $A^{\mrm{op}}$ side, 
Van den Bergh's Local Duality for $A^{\mrm{op}}$, and the symmetry isomorphism 
$\ep_A$ give us these isomorphisms 
\[ A \underset{\cong}{\xar{\opn{hm}^{\mrm{R}}_{R_A, A}}} 
\opn{RHom}_{A^{\mrm{op}}}(R_A, R_A) \underset{\cong}{\xar{\de_{R_A}}}
\bigl( \mrm{R} \Ga_{\m^{\mrm{op}}}(R_A) \bigr)^* \]
in $\dcat{D}(A^{\mrm{en}}, \mrm{gr})$.
By Graded Matlis Duality, this yields an isomorphism 
$\mrm{R} \Ga_{\m^{\mrm{op}}}(R_A) \lb \cong A^*$. We see that 
$R_A$ has weakly symmetric derived torsion; and hence, by Theorem 
\ref{thm:3720}, $R_A$ has symmetric derived torsion. This is condition (B1) of 
Definition \ref{dfn:3714}.
\end{proof}

\begin{cor}[Two Equivalent Properties] \label{cor:4585} 
Under Setup \tup{\ref{setup:3925}}, the two properties below are 
equivalent\tup{:}
\begin{enumerate}
\rmitem{i}  The graded ring $A$ satisfies the $\chi$ condition, and it has 
finite local cohomological dimension. 
\index{Algebraically graded ring! chi@$\chi$ condition on connected}
\index{Algebraically graded ring! connected {\indash} of finite local 
cohomological dimension}

\rmitem{ii} The graded ring $A$ has a balanced dualizing complex.
\index{Dualizing complex! balanced}
\end{enumerate}
\end{cor}

\begin{proof}
The implication (i) $\Rightarrow$ (ii) is the Van den Bergh Existence 
Theorem \ref{thm:3713}. 
The reverse implication is Corollary \ref{cor:4655}.
\end{proof}

Here are several other consequences of Theorem \ref{thm:3713}. 
Recall that a dualizing complex $R_A$ over $A$ gives rise to the 
associated duality functors 
\begin{equation} \label{eqn:5100}
D_A, D_{A^{\mrm{op}}} : \dcat{D}(A^{\mrm{en}}, \mrm{gr})^{\mrm{op}} \to
\dcat{D}(A^{\mrm{en}}, \mrm{gr}) ,
\end{equation}
namely 
$D_A = \opn{RHom}_{A}(-, R_A)$ and 
$D_{A^{\mrm{op}}} = \opn{RHom}_{A^{\mrm{op}}}(-, R_A)$.

\begin{thm} \label{thm:4055} 
Under Setup \tup{\ref{setup:3925}}, let $R_A$ be a balanced dualizing complex 
over $A$, with associated duality 
functors $D_A$ and $D_{A^{\mrm{op}}}$.
\begin{enumerate}
\item There is an isomorphism 
$D_A \cong D_{A^{\mrm{op}}}$
between the triangulated functors in \tup{(\ref{eqn:5100})}.
% $\dcat{D}^{\mrm{b}}_{(\mrm{f}, \mrm{f})}(A^{\mrm{en}}, 
% \mrm{gr})^{\mrm{op}} \to \dcat{D}(A^{\mrm{en}}, \mrm{gr})$.
 
\item If 
$M \in \dcat{D}^{\mrm{b}}_{(\mrm{f}, \mrm{f})}(A^{\mrm{en}}, \mrm{gr})$
then $D_A(M)$ belongs to  
$\dcat{D}^{\mrm{b}}_{(\mrm{f}, \mrm{f})}(A^{\mrm{en}}, \mrm{gr})$.

\item The functor 
\[ D_A : \dcat{D}^{\mrm{b}}_{(\mrm{f}, \mrm{f})}(A^{\mrm{en}}, 
\mrm{gr})^{\mrm{op}} \to
\dcat{D}^{\mrm{b}}_{(\mrm{f}, \mrm{f})}(A^{\mrm{en}}, \mrm{gr}) \]
is an equivalence of triangulated categories, and it is its own quasi-inverse. 
\end{enumerate}
\end{thm}

\begin{proof} \mbox{}

\smallskip \noindent
(1) By Theorems \ref{thm:3713} and \ref{thm:3710} there is a canonical 
isomorphism 
$R_A \cong (P_A)^*$ in $\dcat{D}(A^{\mrm{en}}, \mrm{gr})$. 

Take a complex 
$M \in \dcat{D}^{\mrm{b}}_{(\mrm{f}, \mrm{f})}(A^{\mrm{en}}, \mrm{gr})$.
We know that $A$ has the $\chi$ condition and finite local cohomological 
dimension. By Proposition \ref{prop:3712} the complex $M$ has symmetric derived 
torsion. This means that there's an isomorphism 
$\ep_M : \mrm{R} \Ga_{\m}(M) \iso \mrm{R} \Ga_{\m^{\mrm{op}}}(M)$
in $\dcat{D}(A^{\mrm{en}}, \mrm{gr})$. 
By Theorem \ref{thm:3750}  and its opposite version (namely the theorem applied 
to the ring $A^{\mrm{op}}$ instead of $A$), this translates to an isomorphism 
\begin{equation} \label{eqn:5101}
P_A \ot^{\mrm{L}}_{A} M \cong M \ot^{\mrm{L}}_{A} P_A
\end{equation}
in $\dcat{D}(A^{\mrm{en}}, \mrm{gr})$.
Applying $(-)^*$ to  (\ref{eqn:5101}) we obtain an isomorphism 
\begin{equation} \label{eqn:4080}
(P_A \ot^{\mrm{L}}_{A} M)^* \cong (M \ot^{\mrm{L}}_{A} P_A)^*
\end{equation}
in $\dcat{D}(A^{\mrm{en}}, \mrm{gr})$.
Now by Hom-tensor adjunction we get isomorphisms
\begin{equation} \label{eqn:4081}
\begin{aligned}
& (P_A \ot^{\mrm{L}}_{A} M)^* = 
\opn{Hom}_{\K}(P_A \ot^{\mrm{L}}_{A} M , \K) 
\\
& \quad 
\cong 
\opn{RHom}_{A} \bigl( M, \opn{RHom}_{\K}(P_A, \K) \bigr) 
\\
& \quad 
\cong \opn{RHom}_{A} \bigl( M, (P_A)^* \bigr) 
\cong \opn{RHom}_{A}(M, R_A) = D_A(M) .
\end{aligned} 
\end{equation}
Similarly there are isomorphisms
\begin{equation} \label{eqn:4082}
\begin{aligned}
& (M \ot^{\mrm{L}}_{A} P_A)^* = 
\opn{Hom}_{\K}(M \ot^{\mrm{L}}_{A} P_A , \K) 
\\
& \quad 
\cong 
\opn{RHom}_{A^{\mrm{op}}} \bigl( M, \opn{RHom}_{\K}(P_A, \K) \bigr) 
\\
& \quad 
\cong \opn{RHom}_{A^{\mrm{op}}} \bigl( M, (P_A)^* \bigr) 
\cong \opn{RHom}_{A^{\mrm{op}}}(M, R_A) = D_{A^{\mrm{op}}}(M) .
\end{aligned}
\end{equation}
All these isomorphisms are functorial in the argument $M$.
By combining the isomorphisms (\ref{eqn:4080}), (\ref{eqn:4081}) and 
(\ref{eqn:4082}) we get an isomorphism of functors 
$D_A \cong D_{A^{\mrm{op}}}$. 

\medskip \noindent
(2)  By Theorem \ref{thm:3820} we know that 
$D_A(M)$ belongs to 
$\dcat{D}^{\mrm{b}}_{(.., \mrm{f})}(A^{\mrm{en}}, \mrm{gr})$.
The same theorem, applied to the ring $A^{\mrm{op}}$, tells us that 
$D_{A^{\mrm{op}}}(M)$ belongs to 
$\dcat{D}^{\mrm{b}}_{(\mrm{f}, ..)}(A^{\mrm{en}}, \mrm{gr})$.
Thus 
$D_A(M) \cong D_{A^{\mrm{op}}}(M) \in 
\dcat{D}^{\mrm{b}}_{(\mrm{f}, \mrm{f})}(A^{\mrm{en}}, \mrm{gr})$.

\medskip \noindent
(3) We already know, by Theorem \ref{thm:3820}, that 
$D_{A^{\mrm{op}}} \circ D_A \cong \opn{Id}$
as triangulated functors from 
$\dcat{D}^{\mrm{b}}_{(\mrm{f}, \mrm{f})}(A^{\mrm{en}}, \mrm{gr})$
to itself. By item (2) the essential image of 
$D_A$ is 
$\dcat{D}^{\mrm{b}}_{(\mrm{f}, \mrm{f})}(A^{\mrm{en}}, \mrm{gr})$. 
And by item (1) the functors $D_A$ and $D_{A^{\mrm{op}}}$ are isomorphic 
on this triangulated category. 
\end{proof}

\begin{lem} \label{lem:4060}
Under Setup \tup{\ref{setup:3925}}, assume $A$ satisfies the $\chi$ condition 
and has finite local cohomological dimension. 
Then there are inclusions 
$\dcat{D}^{}_{\mrm{f}}(A, \mrm{gr}) \sub
\dcat{D}^{}(A, \mrm{gr})_{\mrm{com}}$
and
$\dcat{D}^{}_{\mrm{cof}}(A, \mrm{gr}) \sub
\dcat{D}^{}(A, \mrm{gr})_{\mrm{tor}}$.
\end{lem}

\begin{proof}
The inclusion 
$\dcat{D}_{\mrm{cof}}^{}(A, \mrm{gr}) \sub
\dcat{D}^{}_{\mrm{tor}}(A, \mrm{gr})$
is trivial. And Proposition \ref{prop:3965} says that 
$\dcat{D}^{}_{\mrm{tor}}(A, \mrm{gr}) = 
\dcat{D}^{}(A, \mrm{gr})_{\mrm{tor}}$.

Now to the complete complexes. 
Take a complex 
$M \in \dcat{D}^{}_{\mrm{f}}(A, \mrm{gr})$.
Then there are isomorphisms 
\begin{equation} \label{eqn:4025}
\begin{aligned}
& \opn{ADC}_{\m}(M) = \opn{RHom}_{A}(P_A, M) 
\\ & \quad
\cong^1 \opn{RHom}_{A^{\mrm{op}}} \bigl( M^*, (P_A)^* \bigr) 
= \opn{RHom}_{A^{\mrm{op}}}(M^*, R_A) 
\\ & \quad
\cong^2 \opn{RHom}_{A^{\mrm{op}}} \bigl( M^*, \mrm{R} \Ga_{\m^{\mrm{op}}}(R_A) 
\bigr) 
\\ & \quad
\cong^3 \opn{RHom}_{A^{\mrm{op}}}(M^*, A^*) 
\cong (M^*)^* \cong^4 M
\end{aligned} 
\end{equation}
in $\dcat{D}(A, \mrm{gr})$. 
The isomorphism $\cong^1$ is due to Theorem \ref{thm:3800}, that applies 
because $M, P_A \in \dcat{C}_{\mrm{dwf}}(A, \mrm{gr})$. 
The isomorphism $\cong^2$ comes from Proposition \ref{prop:3990}, and it 
applies since 
$M^* \in \dcat{D}^{}_{\mrm{cof}}(A^{\mrm{op}}, \mrm{gr}) \sub 
\dcat{D}^{}_{\mrm{tor}}(A^{\mrm{op}}, \mrm{gr})$
and 
$R_A \in \dcat{D}^{\mrm{b}}(A^{\mrm{op}}, \mrm{gr})$. 
The isomorphism $\cong^3$ is because $R_A$ is a balanced dualizing complex.  
The isomorphism $\cong^4$ is by Proposition \ref{prop:4591}, that applies 
because $M \in \dcat{C}_{\mrm{dwf}}(A, \mrm{gr})$. 
We see that $M$ is in the essential image of the functor $\opn{ADC}_{\m}$. 
According to Theorem \ref{thm:3792}(1) it follows that 
$M \in \dcat{D}^{\mrm{b}}(A, \mrm{gr})_{\mrm{com}}$. 
\end{proof}

The next theorem is a refinement of the NC graded MGM Equivalence (Theorem 
\ref{thm:3792}). Given another graded ring $B$, there are triangulated functors 
\[ \mrm{R} \Ga_{\m}, \, \opn{ADC}_{\m} : 
\dcat{D}(A \ot B^{\mrm{op}}, \mrm{gr}) \to 
\dcat{D}(A \ot B^{\mrm{op}}, \mrm{gr}) . \]
They are the derived $\m$-torsion and the abstract derived $\m$-adic 
completion, respectively. 

\begin{thm} \label{thm:4070}
Under Setup \tup{\ref{setup:3925}}, assume $A$ satisfies the $\chi$ condition 
and has finite local cohomological dimension. Let $B$ be some graded ring. 
\begin{enumerate}
\item If 
$M \in \dcat{D}^{\mrm{b}}_{(\mrm{f}, ..)}(A \ot B^{\mrm{op}}, \mrm{gr})$
then 
$\mrm{R} \Ga_{\m}(M) \in 
\dcat{D}^{\mrm{b}}_{(\mrm{cof}, ..)}(A \ot B^{\mrm{op}}, \mrm{gr})$.

\item For 
$N \in \dcat{D}^{\mrm{b}}_{(\mrm{cof}, ..)}(A \ot B^{\mrm{op}}, \mrm{gr})$
there is an isomorphism 
$\opn{ADC}_{\m}(N) \cong \lb D_{A^{\mrm{op}}}(N^*)$ 
in $\dcat{D}(A \ot B^{\mrm{op}}, \mrm{gr})$,
and it is functorial in $N$. 

\item If 
$N \in \dcat{D}^{\mrm{b}}_{(\mrm{cof}, ..)}(A \ot B^{\mrm{op}}, \mrm{gr})$
then 
$\opn{ADC}_{\m}(N) \in 
\dcat{D}^{\mrm{b}}_{(\mrm{f}, ..)}(A \ot B^{\mrm{op}}, \mrm{gr})$.

\item The functor 
\[ \mrm{R} \Ga_{\m} : 
\dcat{D}^{\mrm{b}}_{(\mrm{f}, ..)}(A \ot B^{\mrm{op}}, \mrm{gr}) \to 
\dcat{D}^{\mrm{b}}_{(\mrm{cof}, ..)}(A \ot B^{\mrm{op}}, \mrm{gr}) \]
is an equivalence of triangulated categories, with quasi-inverse 
$\opn{ADC}_{\m}$. 
\end{enumerate}
\end{thm}

\begin{proof} \mbox{}

\smallskip \noindent
(1) The ring $B$ is irrelevant, so we may omit it. 
The $\chi$ condition and finiteness of local cohomology imply that 
$\mrm{R} \Ga_{\m}(M) \in \dcat{D}^{\mrm{b}}_{\mrm{cof}}(A, \mrm{gr})$.

\medskip \noindent 
(2) Take a complex 
$N \in \dcat{D}^{\mrm{b}}_{\mrm{cof}}(A \ot B^{\mrm{op}}, \mrm{gr})$.
We have these isomorphisms 
\[ \begin{aligned}
& \opn{ADC}_{\m}(N) = \opn{RHom}_{A}(P_A, N) 
\cong^{\dag} \opn{RHom}_{A^{\mrm{op}}} \bigl( N^*, (P_A)^* \bigr) 
\\ & \quad
= \opn{RHom}_{A^{\mrm{op}}}(N^*, R_A) 
= D_{A^{\mrm{op}}}(N^*) 
\end{aligned} \]
in $\dcat{D}(A\ot B^{\mrm{op}}, \mrm{gr})$.
The isomorphism $\cong^{\dag}$ is due to Theorem \ref{thm:3800}, that applies 
because 
$N \in \dcat{C}_{\mrm{dwf}}(A\ot B^{\mrm{op}}, \mrm{gr})$ and 
$P_A \in \dcat{C}_{\mrm{dwf}}(A^{\mrm{en}}, \mrm{gr})$. 

\medskip \noindent 
(3) Again we can forget about $B$. 
Take $N \in  \dcat{D}^{\mrm{b}}_{\mrm{cof}}(A, \mrm{gr})$.
Then  
$N^* \in \dcat{D}_{\mrm{f}}(A^{\mrm{op}}, \mrm{gr})$,
so by Theorem \ref{thm:3820} we have 
$D_{A^{\mrm{op}}}(N^*) \in \dcat{D}^{\mrm{b}}_{\mrm{f}}(A, \mrm{gr})$.
But by item (2) we know that 
$\opn{ADC}_{\m}(N) \cong D_{A^{\mrm{op}}}(N^*)$. 
            
\medskip \noindent 
(4) Items (1) and (3) tell us that the functors 
$\mrm{R} \Ga_{\m}$ and $\opn{ADC}_{\m}$ send the categories 
$\dcat{D}^{\mrm{b}}_{(\mrm{f}, ..)}(A \ot B^{\mrm{op}}, \mrm{gr})$
and 
$\dcat{D}^{\mrm{b}}_{(\mrm{cof}, ..)}(A \ot B^{\mrm{op}}, \mrm{gr})$
to each other, respectively. 
Lemma \ref{lem:4060} says that the NC MGM Equivalence (Theorem \ref{thm:3792}) 
applies here, namely that there is an isomorphism of triangulated functors 
$\opn{ADC}_{\m} \circ \, \mrm{R} \Ga_{\m} \cong \opn{Id}$
from 
$\dcat{D}^{\mrm{b}}_{(\mrm{f}, ..)}(A \ot B^{\mrm{op}}, \mrm{gr})$
to itself, and an isomorphism of triangulated functors 
$\mrm{R} \Ga_{\m} \circ \opn{ADC}_{\m} \cong \opn{Id}$
from
$\dcat{D}^{\mrm{b}}_{(\mrm{cof}, ..)}(A \ot B^{\mrm{op}}, \mrm{gr})$
to itself.
\end{proof}

In Theorem \ref{thm:4051} we saw that there is an isomorphism 
$\ep : \mrm{R} \Ga_{\m} \iso \mrm{R} \Ga_{\m^{\mrm{op}}}$
of triangulated functors 
$\dcat{D}^{\mrm{b}}_{(\mrm{f}, \mrm{f})}(A^{\mrm{en}}) \to 
\dcat{D}(A^{\mrm{en}}, \mrm{gr})$.
The next corollary produces a similar isomorphism for the abstract derived 
completion functors 
\begin{equation} \label{eqn:4058}
\opn{ADC}_{\m}, \opn{ADC}_{\m^{\mrm{op}}} : \dcat{D}(A^{\mrm{en}}, \mrm{gr}) 
\to 
\dcat{D}(A^{\mrm{en}}, \mrm{gr}) , 
\end{equation}
where we define
$\opn{ADC}_{\m^{\mrm{op}}} := \opn{RHom}_{A^{\mrm{op}}}(P_A, -)$. 

\begin{cor} \label{cor:4080}   
Under Setup \tup{\ref{setup:3925}}, assume $A$ satisfies the $\chi$ condition 
and has finite local cohomological dimension. Then there is an isomorphism 
$\ga : \lb \opn{ADC}_{\m} \iso \opn{ADC}_{\m^{\mrm{op}}}$
of triangulated functors 
$\dcat{D}^{\mrm{b}}_{(\mrm{cof}, \mrm{cof})}(A^{\mrm{en}}, \mrm{gr}) \to 
\dcat{D}(A^{\mrm{en}}, \mrm{gr})$.
\end{cor}

\begin{proof}
Take a complex 
$N \in \dcat{D}^{\mrm{b}}_{(\mrm{cof}, \mrm{cof})}(A^{\mrm{en}}, \mrm{gr})$.
By Theorem \ref{thm:4070}(2), there is an isomorphism 
$\opn{ADC}_{\m}(N) \cong D_{A^{\mrm{op}}}(N^*)$
in $\dcat{D}^{\mrm{b}}(A^{\mrm{en}}, \mrm{gr})$.
Similarly (replacing $A$ with $A^{\mrm{op}}$) there is an isomorphism 
$\opn{ADC}_{\m^{\mrm{op}}}(N) \cong D_{A}(N^*)$
in $\dcat{D}^{\mrm{b}}(A^{\mrm{en}}, \mrm{gr})$.
Now 
$N^* \in \dcat{D}^{\mrm{b}}_{(\mrm{f}, \mrm{f})}(A^{\mrm{en}}, \mrm{gr})$,
so by Theorem \ref{thm:4055} we have 
$D_{A^{\mrm{op}}}(N^*) \cong D_{A}(N^*)$.
Since we relied on functorial isomorphisms, this is also a functorial 
isomorphism.   
\end{proof}

\begin{cor} \label{cor:4070}
Assume $A$ satisfies the $\chi$ condition and has finite local cohomological 
dimension. Then the functor 
\[ \mrm{R} \Ga_{\m} : 
\dcat{D}^{\mrm{b}}_{(\mrm{f}, \mrm{f})}(A^{\mrm{en}}, \mrm{gr}) \to
\dcat{D}^{\mrm{b}}_{(\mrm{cof}, \mrm{cof})}(A^{\mrm{en}}, \mrm{gr}) \]
is an equivalence of triangulated categories, with quasi-inverse
$\opn{ADC}_{\m}$. 
\end{cor}

\begin{proof}
Taking $B := A$ in Theorem \ref{thm:4070}(4) we get an equivalence 
\begin{equation} \label{eqn:4845}
\mrm{R} \Ga_{\m} : 
\dcat{D}^{\mrm{b}}_{(\mrm{f}, ..)}(A^{\mrm{en}}, \mrm{gr}) \to
\dcat{D}^{\mrm{b}}_{(\mrm{cof}, ..)}(A^{\mrm{en}}, \mrm{gr}) ,
\end{equation}
with with quasi-inverse $\opn{ADC}_{\m}$. 
From the same theorem, but switching $A$ and $A^{\mrm{op}}$, 
we get an equivalence 
\[ \mrm{R} \Ga_{\m^{\mrm{op}}} : 
\dcat{D}^{\mrm{b}}_{(.., \mrm{f})}(A^{\mrm{en}}, \mrm{gr}) \to
\dcat{D}^{\mrm{b}}_{(.., \mrm{cof})}(A^{\mrm{en}}, \mrm{gr}) , \]
with with quasi-inverse $\opn{ADC}_{\m^{\mrm{op}}}$. 
According to Theorem \ref{thm:4051} there is an isomorphism  
$\mrm{R} \Ga_{\m}(M) \cong \mrm{R} \Ga_{\m^{\mrm{op}}}(M)$
for 
$M \in \dcat{D}^{\mrm{b}}_{(\mrm{f}, \mrm{f})}(A^{\mrm{en}}, \mrm{gr})$.
Therefore, for such a complex $M$ we have 
$\mrm{R} \Ga_{\m}(M) \in 
\dcat{D}^{\mrm{b}}_{(\mrm{cof}, \mrm{cof})}(A^{\mrm{en}}, \mrm{gr})$.
Corollary \ref{cor:4080} says that 
$\opn{ADC}_{\m}(N) \cong \opn{ADC}_{\m^{\mrm{op}}}(N)$
for a complex 
$N \in \dcat{D}^{\mrm{b}}_{(\mrm{cof}, \mrm{cof})}(A^{\mrm{en}}, \mrm{gr})$; 
so for such a complex $N$ we have 
$\opn{ADC}_{\m}(N) \in 
\dcat{D}^{\mrm{b}}_{(\mrm{f}, \mrm{f})}(A^{\mrm{en}}, \mrm{gr})$.
The conclusion is that the equivalence (\ref{eqn:4845}) restricts to 
the equivalence stated in the corollary. 
\end{proof}

\begin{que} \label{que:3990}
A consequence of Theorem \ref{thm:3750} is that the complex $P_A$ has finite 
flat dimension over $A$. Is it true that the complex $P_A$ has finite 
projective dimension over $A$~? I.e., does the functor 
\[ \opn{ADC}_{\m} = \opn{RHom}_{A}(P_A, -) : \dcat{D}(A, \mrm{gr}) \to 
\dcat{D}(A, \mrm{gr}) \]
have finite cohomological dimension?

In the commutative weakly proregular case this is true -- see Examples 
\ref{exa:4595} and \ref{exa:4575}. Indeed, if $\a$ is a WPR ideal in a 
commutative ring $A$, then the cohomological dimension of the functor
$\opn{ADC}_{\a} \cong \mrm{L} \Lambda_{\a}$ 
is at most the length of a finite generating sequence of $\a$. 
\end{que}

\begin{rem} \label{rem:3857}
The proof of Theorem \ref{thm:3713} provided here is not the original proof 
from \cite{VdB}. Our proof, which relies on the NC MGM Equivalence,
seems to allow a generalization to the case of a base ring $\K$ that is not a 
field (as long as the ring $A$ is flat over $\K$). 
See \cite{VyYe} for some results in this direction. 
\end{rem}

\begin{rem} \label{rem:4082}
Q.S. Wu and J.J. Zhang \cite{WuZh}, \cite{WuZh2} studied prebalanced dualizing 
complexes over complete semilocal noncommutative rings that contain a field. 
Some of their technical results are included in Subsection 
\ref{subsec:bal-trace} below. 

An ongoing project of R. Vyas and A. Yekutieli (a continuation of \cite{VyYe}) 
aims to study balanced dualizing complexes over complete semilocal 
noncommutative rings in the arithmetic setting, namely without the presence of 
a base field. A prototypical example is the ring 
$A = \K[[G]] / I$, where $\K = \what{\Z}_p$, $G$ is a compact $p$-adic Lie 
group, $\K[[G]]$ is the noncommutative Iwasawa algebra, and 
$I \sub \K[[G]]$ is some two-sided ideal.  Because the ring $A$ could fail to 
be flat over $\K$, the balanced dualizing complex of $A$ will be an object of 
the {\em derived category of bimodules} 
$\dcat{D}(\til{A}^{\mrm{en}})$, where 
$\til{A} \to A$ is a K-flat DG ring resolution of $A$ and 
$\til{A}^{\mrm{en}} := \til{A} \ot_{\K} \til{A}^{\mrm{op}}$;
see \cite{Ye16}. 
\end{rem}

\begin{rem} \label{rem:4081}
Suppose $A$ is a commutative noetherian complete local ring, with maximal ideal 
$\m$. In the paper \cite{AlJeLi} the authors discuss two kinds of dualizing 
complexes over $A$~: the usual dualizing complexes (see Subsection 
\ref{subsec:du-cplxs}), 
that they call {\em c-dualizing complexes}, which have finite (and thus 
complete) cohomology modules; and the {\em t-dualizing complexes}, 
that have cofinite (and thus torsion) cohomology modules.
The torsion injective module $I = A^*$ (the injective hull of the residue 
field) is a t-dualizing complex in their sense. One of the results of 
\cite{AlJeLi} is that the MGM Equivalence (that they, in that early paper, 
called {\em GM Duality}) exchanges t-dualizing complexes and c-dualizing 
complexes. In particular, by applying the derived $\m$-adic completion functor 
$\mrm{L} \Lambda_{\m}$ one obtains a c-dualizing complex 
$R := \mrm{L} \Lambda_{\m}(A^*)$. 

It was observed by Vyas that the Van den Bergh Existence Theorem (Theorem 
\ref{thm:3713}) can be understood as a noncommutative variant of that very same 
result. In our noncommutative graded setting we must replace the true derived 
$\m$-adic completion functor $\mrm{L} \Lambda_{\m}$ with its abstract 
avatar $\opn{ADC}_{\m}$. The graded $A$-bimodule $A^*$ is a t-dualizing complex 
(if we adjust the definition from \cite{AlJeLi} to the NC connected graded 
setting), and indeed, as Theorem \ref{thm:3713} says, the complex
\[ R_A := \opn{ADC}_{\m}(A^*) = \opn{RHom}_{A}(P_A, A^*) = (P_A)^* \]
is a graded NC dualizing complex (as defined in Definition \ref{dfn:3716}).
\end{rem}




\mysubsection{Balanced Trace Morphisms} 
\label{subsec:bal-trace}

In this subsection we adhere to Conventions \ref{conv:3700} and \ref{conv:4560}.

\begin{dfn} \label{dfn:3765}
A graded ring homomorphism $f : A \to B$ is called {\em finite} if it makes $B$ 
into a finite graded $A$-module on both sides. 
\end{dfn}

Throughout this subsection we assume the next setup: 

\begin{setup} \label{setup:3765} 
We are given noetherian connected graded $\K$-rings $A$ and $B$, with 
augmentation ideals $\m$ and $\n$ respectively, and a finite graded
$\K$-ring homomorphism $f : A \to B$. 
We are also given a homomorphism of graded $\K$-rings  $g : C \to D$. 
\end{setup}

The rings $C$ and $D$ will play auxiliary roles, to allow us to handle 
bimodules. In practice $C$ will be 
either $A$ or $\K$, and $D$ will be either $B$ or $\K$.
The graded ring homomorphism 
$f \ot g^{\mrm{op}} : A \ot C^{\mrm{op}} \to B \ot D^{\mrm{op}}$
induces a  restriction functor
\[ \opn{Rest} = \opn{Rest}_{f \ot g^{\mrm{op}}}
 : \dcat{D}(B \ot D^{\mrm{op}}, \mrm{gr}) \to 
\dcat{D}(A \ot C^{\mrm{op}}, \mrm{gr}) . \]
This restriction functor will often remain implicit. 

\begin{thm} \label{thm:3783}
Assume Setup \tup{\ref{setup:3765}}.
\begin{enumerate}
\item There is a canonical morphism
\[ \opn{Rest} \circ \, \mrm{R} \Ga_{\n} \to 
\mrm{R} \Ga_{\m} \circ \opn{Rest} \]
of triangulated functors 
$\dcat{D}(B \ot D^{\mrm{op}}, \mrm{gr}) \to 
\dcat{D}(A \ot C^{\mrm{op}}, \mrm{gr})$.

\item Let $N \in \dcat{D}(B \ot D^{\mrm{op}}, \mrm{gr})$, and assume either 
that 
$N \in \dcat{D}^+(B \ot D^{\mrm{op}}, \mrm{gr})$ 
or that the functor $\Ga_{\m}$ has finite right cohomological dimension. 
Then the canonical morphism
$\mrm{R} \Ga_{\n}(N) \to \mrm{R} \Ga_{\m}(N)$
in $\dcat{D}(A \ot C^{\mrm{op}}, \mrm{gr})$ is an isomorphism. 
\end{enumerate}
\end{thm}

For the proof we need a few lemmas.

\begin{lem} \label{lem:3785}
For every $N \in \dcat{M}(B \ot D^{\mrm{op}}, \mrm{gr})$
there is equality 
$\Ga_{\m}(N) = \Ga_{\n}(N)$
of $(B \ot D^{\mrm{op}})$-submodules of $N$. 
\end{lem}

\begin{proof}
Take any integer $q \geq 1$,
and let 
$\m^q := \m \cdots \m$, the $q$-fold power of the ideal $\m$. 
Because $B$ is a finite $A$-module, it follows that 
$B / (\m^q \cd B) \cong (A / \m^q) \ot_A B$
is a finite module over $A / \m^q$, and hence it is a finite $\K$-module. 
This implies that the graded module $B / (\m^q \cd B)$ is concentrated in a 
finite algebraic degree interval, say $[0, i_1]$. 
So for every $i > i_1$ there is equality $(\m^q \cd B)_i = B_i$. 
Taking $q' := \opn{max}(q, i_1 + 1)$ we get 
$\n^{q'} \sub \m^q \cd B$. 

The argument above shows that for every $q \geq 1$ there exists some 
$q' \geq q$ for which $\n^{q'} \sub \m^q \cd B$. 
Trivially there is an inclusion $\m^q \cd B \sub \n^q$.
These inclusions yield 
\[ \opn{Hom}_{B} ( B / \n^{q'} , N) \sub
\opn{Hom}_{A} ( A / \m^q , N) \sub 
\opn{Hom}_{B} ( B / \n^q , N) \sub N . \]
Combined with formula (\ref{eqn:3772}) the assertion is proved. 
\end{proof}

Recall the ML condition from Definition \ref{dfn:1590}.
An inverse system $\{ N_q \}_{q \in \N}$ 
in $\dcat{M}(\K, \mrm{gr})$ is said to have the {\em trivial ML 
property}, or is called {\em prozero}, if for every $q$ there is some 
$q' \geq q$ such that the homomorphism $N_{q'} \to N_q$ is zero. See 
\cite[Definition 3.5.6]{We}. 

\begin{lem} \label{lem:3768}
Let $\{ N_q \}_{q \in \N}$ be an inverse system in 
$\dcat{M}(\K, \mrm{gr})$ that has the ML property. The following are 
equivalent\tup{:}
\begin{enumerate}
\rmitem{i} The inverse system has the trivial ML property.

\rmitem{ii} $\opn{lim}_{\lar q} N_q = 0$. 
\end{enumerate}
\end{lem}

\begin{lem} \label{lem:3820} 
Let $\{ N_q \}_{q \in \N}$ be an inverse system in 
$\dcat{M}(A, \mrm{gr})$ that has the trivial ML property, and let 
$M \in \dcat{M}(A, \mrm{gr})$. Then 
$\lim_{q \to} \opn{Hom}_{A}(N_q, M) = 0$.
\end{lem}

\begin{exer} \label{exer:3765}
Prove Lemmas \ref{lem:3768} and \ref{lem:3820}.
\end{exer}

For a finite graded $A$-module $M$, and for numbers $p, q \in \N$, we 
define \begin{equation} \label{eqn:3765}
F_{p, q}(M) := \opn{Tor}^A_p(A / \m^{q + 1}, M) = 
\opn{H}^{-p} \bigl(  (A / \m^{q + 1})  \ot^{\mrm{L}}_{A} M \bigr) 
\in \dcat{M}(A, \mrm{gr}) . 
\end{equation}
Then, fixing $p$, we define 
\begin{equation} \label{eqn:3766}
F_{p}(M) := \opn{lim}_{\lar q} F_{p, q}(M) \in \dcat{M}(A, \mrm{gr}) . 
\end{equation}
In this way we obtain functors 
\begin{equation} \label{eqn:5202}
F_{p, q}, F_{p} : \dcat{M}_{\mrm{f}}(A, \mrm{gr}) \to \dcat{M}(A, \mrm{gr}) .
\end{equation}

Lemmas \ref{lem:3765} and \ref{lem:3766} below are essentially
\cite[Lemma 6.4 and Proposition 6.5]{WuZh}. 

\begin{lem} \label{lem:3765}
Let $M$ be a finite graded $A$-module and let $p \geq 1$. Then the 
inverse system $\{ F_{p, q}(M) \}_{q \in \N}$ 
has the trivial ML property. 
\end{lem}

\begin{proof} \mbox{}

\smallskip \noindent
Step 1. Here we show that for fixed $p$, the inverse 
system $\{ F_{p, q}(M) \}_{q \in \N}$ 
satisfies the ML condition. 

Because $A$ is left noetherian and $M$ is a finite graded $A$-module, we can 
find a resolution 
$\cdots \to L^{-1} \to L^0 \to M \to 0$
in $\dcat{M}_{}(A, \mrm{gr})$,
such that each $L^{-p}$ is a finite graded-free $A$-module. 
Then 
\[ F_{p, q}(M) = 
\opn{H}^{-p} \bigl( (A / \m^{q + 1}) \ot^{\mrm{L}}_{A} M \bigr) 
\cong \opn{H}^{-p} \bigl( (A / \m^{q + 1}) \ot_{A} L \bigr) \]
is a finite $(A / \m^{q + 1})$-module. So as a graded $A$-module it is 
artinian. 
For fixed $q$, and for all $q' \geq q$, the descending chain condition on the 
submodules
$\opn{Im} \bigl( F_{p, q'}(M) \to F_{p, q}(M) \bigr) \sub F_{p, q}(M)$
tells us that this eventually stabilizes. This is the ML property. 

\medskip \noindent
Step 2. Let 
$0 \to M' \to M \to M'' \to 0$
be a short exact sequence in $\dcat{M}_{\mrm{f}}(A, \mrm{gr})$.
For each $q$ we have a long exact Tor sequence 
\begin{equation} \label{eqn:3767}
\begin{aligned}
& \cdots \to F_{1, q}(M') \to F_{1, q}(M) \to F_{1, q}(M'') 
\\
& \quad \to F_{0, q}(M') \to F_{0, q}(M) \to F_{0, q}(M'') \to 0  
\end{aligned}
\end{equation}
As $q$ varies, these long exact sequences form an inverse system. By step 1 we 
know that in every position the ML property is satisfied. Therefore by Theorem
\ref{thm:1535}, applied in each degree separately, i.e.\ to the inverse systems 
in $\dcat{M}(\K)$, the inverse limit 
\begin{equation} \label{eqn:3768}
\begin{aligned}
& \cdots \to F_{1}(M') \to F_{1}(M) \to F_{1}(M'') 
\\
& \quad \to F_{0}(M') \to F_{0}(M) \to F_{0}(M'') \to 0  
\end{aligned}
\end{equation}
is an exact sequence. 

\medskip \noindent
Step 3. For a fixed degree $k$, we have $(\m^{q + 1} \cd M)_k = 0$
when $q \gg 0$. This is because $M$ is a finite graded $A$-module, and thus it 
is bounded below in algebraic degree. It follows that the canonical homomorphism
$M_k \to (M / \m^{q + 1} \cd M)_k$ 
is bijective for $q \gg 0$. Going to the limit in $q$ we see that the canonical 
homomorphism
\[ M \to F_0(M) =  \opn{lim}_{\lar q} F_{0, q}(M) = 
\opn{lim}_{\lar q} (M / \m^{q + 1} \cd M) \]
is bijective. Therefore there is an isomorphism of functors 
$F_0 \cong \opn{Id}$. We deduce that $F_0$ is an exact functor on 
$\dcat{M}_{\mrm{f}}(A, \mrm{gr})$. 

\medskip \noindent
Step 4. We now prove, by induction on $p$, that the functor $F_p$ is zero for 
every $p \geq 1$. 

Given a module $M \in \dcat{M}_{\mrm{f}}(A, \mrm{gr})$, 
we choose a surjection $L \surj M$ from a finite graded-free module $L$. 
We get a short exact sequence
$0 \to N \to L \to M \to 0$
in $\dcat{M}_{\mrm{f}}(A, \mrm{gr})$.
Hence, by step 2, there is a long exact sequence 
\begin{equation} \label{eqn:3769}
\begin{aligned}
& \cdots \to F_{2}(N) \to F_{2}(L) \to F_{2}(M) 
\\
& \quad \xar{\de^1} F_{1}(N) \to F_{1}(L) \to F_{1}(M) 
\\
& \quad \xar{\de^0} F_{0}(N) \to F_{0}(L) \to F_{0}(M) \to 0  
\end{aligned}
\end{equation}
in $\dcat{M}(A, \mrm{gr})$;  cf.\ the exact sequence (\ref{eqn:3768}).
Now $F_{p, q}(L) = 0$ for all $q \geq 0$ and all $p > 0$, because $L$ is flat; 
so in the limit we get $F_{p}(L) = 0$ for all $p > 0$. 

Since the functor $F_0$ is exact (by step 3), we see that $\de^0 = 0$. 
Also we have $F_{1}(L) = 0$. Therefore $F_{1}(M) = 0$. 

But $M$ was an arbitrary object of 
$\dcat{M}_{\mrm{f}}(A, \mrm{gr})$, so in fact the functor $F_1$ is zero. 
This implies that in the exact sequence (\ref{eqn:3769}) 
we have $F_{1}(N) = 0$. We also have  $F_{2}(L) = 0$. 
This implies that $F_2(M) = 0$. And so on. This finishes the proof that 
$F_p = 0$ for all $p \geq 1$. 

\medskip \noindent
Step 5. Take some $M \in \dcat{M}_{\mrm{f}}(A, \mrm{gr})$.
We know that $F_p(M) = 0$ for every $p \geq 1$. By step 1 we also 
know that for every $p$ inverse system $\{ F_{p, q}(M) \}_{q \in \N}$ 
has the ML property. Lemma \ref{lem:3768} says that 
for every $p \geq 1$ the inverse system $\{ F_{p, q}(M) \}_{q \in \N}$ 
has the trivial ML property.
\end{proof}

\begin{lem} \label{lem:3766}
If $J$ is a graded-injective $B$-module, then as an $A$-module, $J$ is 
graded-$\m$-flasque. 
\end{lem}

\begin{proof}
Take some $p > 0$. We have to show that 
$\mrm{R}^p \Ga_{\m}(J) = 0$. 
Now 
\[ \mrm{R}^p \Ga_{\m}(J) 
\cong \lim_{q \to} \, \opn{Ext}^p_A \bigl( A / \m^{q + 1}, J) . \]
But 
\[ \opn{Ext}^p_A \bigl( A / \m^{q + 1}, J) = 
\opn{H}^p \bigl( \opn{RHom}_A ( A / \m^{q + 1}, J) \bigr) . \]
By derived Hom-tensor adjunction we have a canonical isomorphism 
\[ \opn{RHom}_A ( A / \m^{q + 1}, J) \cong 
\opn{RHom}_B \bigl( B \ot^{\mrm{L}}_{A} (A / \m^{q + 1}), J \bigr) . \]
in $\dcat{D}(A, \mrm{gr})$. Because $J$ is graded-injective, we have 
\[ \opn{RHom}_B \bigl( B \ot^{\mrm{L}}_{A} (A / \m^{q + 1}), J \bigr) \cong 
\opn{Hom}_B \bigl( B \ot^{\mrm{L}}_{A} (A / \m^{q + 1}), J \bigr) , \]
and also  
\[ \begin{aligned}
& \opn{H}^p \bigl( 
\opn{Hom}_B \bigl( B \ot^{\mrm{L}}_{A} (A / \m^{q + 1}), J \bigr) \cong
\opn{Hom}_B \bigl( \opn{H}^{-p} \bigl(
B \ot^{\mrm{L}}_{A} (A / \m^{q + 1}) \bigr), J \bigr) 
\\
& \quad \cong 
\opn{Hom}_B \bigl( \opn{Tor}_p^A (B, A / \m^{q + 1}), J \bigr) . 
\end{aligned} \]
Putting it all together we get 
\[ \mrm{R}^p \Ga_{\m}(J) \cong 
\lim_{q \to} \, \opn{Hom}_B \bigl( F_{p, q}(B) , J \bigr) , \]
where 
$F_{p, q}(B) := \opn{Tor}_p^A (B, A / \m^{q + 1})$, 
like in formula (\ref{eqn:3765}).
Now $B$ is a finite $A^{\mrm{op}}$-module. 
According to Lemma \ref{lem:3765}, applied to the ring $A^{\mrm{op}}$ instead 
of 
$A$, the inverse system 
$\{ F_{p, q}(M) \}_{q \in \N}$ 
has the trivial ML property. By Lemma \ref{lem:3820} we see that 
$\lim_{q \to} \opn{Hom}_B \bigl( F_{p, q}(B) , J \bigr) = 0$.
\end{proof}

\begin{proof}[Proof of Theorem \tup{\ref{thm:3783}}] \mbox{}

\smallskip \noindent
(1) For a complex 
$N \in \dcat{D}(B \ot D^{\mrm{op}}, \mrm{gr})$
we choose a K-injective resolution $\rho : N \to J$ in 
$\dcat{C}_{\mrm{str}}(B \ot D^{\mrm{op}}, \mrm{gr})$.
Then we choose a  K-injective resolution $\si : J \to I$ in 
$\dcat{C}_{\mrm{str}}(A \ot C^{\mrm{op}}, \mrm{gr})$. (We are hiding the 
restriction functor $\opn{Rest}$ here.) So 
$\si \circ \rho : N \to I$ is a K-injective resolution 
in $\dcat{C}_{\mrm{str}}(A \ot C^{\mrm{op}}, \mrm{gr})$.
We get presentations 
$\mrm{R} \Ga_{\n}(N) = \Ga_{\n}(J)$ and 
$\mrm{R} \Ga_{\m}(N) = \Ga_{\m}(I)$. 
By Lemma \ref{lem:3785} there is equality 
$\Ga_{\m}(J) = \Ga_{\n}(J)$. 
The morphism $\mrm{R} \Ga_{\n}(N) \to \mrm{R} \Ga_{\m}(N)$
in $\dcat{D}(A \ot C^{\mrm{op}}), \mrm{gr})$
we are looking for is the one represented by the homomorphism 
$\Ga_{\m}(\si) : \Ga_{\n}(J) = \Ga_{\m}(J) \to \Ga_{\m}(I)$
in $\dcat{C}_{\mrm{str}}(A \ot C^{\mrm{op}}, \mrm{gr})$. 
It does not depend on the choices made. 

\medskip \noindent
(2) Now we choose the K-injective resolution $\rho : N \to J$ in 
$\dcat{C}_{\mrm{str}}(B \ot D^{\mrm{op}}, \mrm{gr})$ more carefully: 
we take the semi-graded-cofree resolution from Theorem \ref{thm:3740}(2).
Then the complex $J$ is K-injective in $\dcat{C}(B \ot D^{\mrm{op}}, \mrm{gr})$, 
each $J^p$ is injective in $\dcat{M}(B \ot D^{\mrm{op}}, \mrm{gr})$, and 
$\opn{inf}(J) = \opn{inf}(\opn{H}(N))$. 
According to Lemma \ref{lem:3766}, $J$ is a complex of graded-$\m$-flasque 
modules. By Lemma \ref{lem:3885}, under either of the two conditions, 
the complex $J$ is graded-$\m$-flasque. Therefore the homomorphism 
$\Ga_{\m}(\si) : \Ga_{\m}(J) \to \Ga_{\m}(I)$
in $\dcat{C}_{\mrm{str}}(A \ot C^{\mrm{op}}, \mrm{gr})$ is a quasi-isomorphism. 
\end{proof}

\begin{cor} \label{cor:4030}
If the functor $\Ga_{\m}$ has finite right cohomological dimension, then 
the functor $\Ga_{\n}$ has finite right cohomological dimension.
\end{cor}

\begin{proof}
Part 2 of Theorem \ref{thm:3783} implies that the right cohomological dimension 
of $\Ga_{\n}$ is at most the right cohomological dimension of $\Ga_{\m}$.  
\end{proof}

\begin{cor} \label{cor:4031}
If the ring $A$ satisfies the left $\chi$ condition, then so does the ring $B$. 
\end{cor}

\begin{proof}
A graded $B$-module $N$ is finite iff it is finite as a graded $A$-module, and 
the same for cofiniteness. Let $N \in \dcat{M}_{\mrm{f}}(B, \mrm{gr})$. 
By Proposition \ref{prop:3855} the graded $A$-modules 
$\opn{H}^p_{\m}(N)$ are all cofinite.
But by part 2 of Theorem \ref{thm:3783} we know that 
$\opn{H}^p_{\n}(N) \cong \opn{H}^p_{\m}(N)$ as graded $A$-modules. 
So the graded $B$-modules $\opn{H}^p_{\n}(N)$ are all cofinite.
Again using Proposition \ref{prop:3855} we see that $B$ satisfies the left 
$\chi$ condition.
\end{proof}

\begin{cor} \label{cor:4032}
If the ring $A$ has a balanced dualizing complex $(R_A, \be_A)$, then the ring 
$B$ has a balanced dualizing complex $(R_B, \be_B)$.
\end{cor}

\begin{proof}
According to Corollary \ref{cor:4585}, $A$ satisfies the $\chi$ condition, and 
it has finite local cohomological dimension. Corollaries \ref{cor:4030} and 
\ref{cor:4031}, applied to the finite graded ring homomorphisms
$A \to B$ and $A^{\mrm{op}} \to B^{\mrm{op}}$, tell us that 
$B$ satisfies the $\chi$ condition, and it has finite local cohomological 
dimension. Again using Corollary \ref{cor:4585}, now in the reverse direction, 
we conclude that $B$ has a balanced dualizing complex.
\end{proof}

The graded ring homomorphism $f : A \to B$ induces, by applying the functor 
$(-)^* = \opn{Hom}_{\K}(-, \K)$, a homomorphism
$f^* : B^* \to A^*$ in $\dcat{M}(A^{\mrm{en}}, \mrm{gr})$. 

\begin{dfn} \label{dfn:4030}
Under Setup \tup{\ref{setup:3765}}, assume $A$ and $B$ have balanced dualizing 
complexes $(R_A, \be_A)$ and $(R_B, \be_B)$ respectively. A 
{\em balanced trace morphism} 
\index{Dualizing complex! balanced trace morphism of balanced {\indash}es}
\index{1-Trf@$\opn{tr}_f$}
is a morphism 
\[ \opn{tr}_f = \opn{tr}_{B / A} : R_B \to R_A \]
in $\dcat{D}(A^{\mrm{en}}, \mrm{gr})$,
such that the diagram 
\[ \UseTips \xymatrix @C=10ex @R=6ex {
\mrm{R} \Ga_{\m}(R_B)
\ar[r]^{ \mrm{R} \Ga_{\m}(\opn{tr}_f) }
\ar[d]_{\be_B}^{\cong}
&
\mrm{R} \Ga_{\m}(R_A)
\ar[d]^{\be_A}_{\cong}
\\
B^*
\ar[r]^{f^*}
&
A^*
} \]
in $\dcat{D}(A^{\mrm{en}}, \mrm{gr})$ is commutative. 
Here we use the canonical isomorphism 
$\mrm{R} \Ga_{\n}(R_B) \lb \cong \mrm{R} \Ga_{\m}(R_B)$
from Theorem \tup{\ref{thm:3783}(2)}.
\end{dfn}

In Subsection \ref{subsec:adjs-NC} we studied backward morphisms for objects in 
derived categories. This makes sense also for complexes of graded bimodules. 
Namely, given complexes
$M \in \dcat{D}(A \ot C^{\mrm{op}}, \mrm{gr})$
and
$N \in \dcat{D}(B \ot D^{\mrm{op}}, \mrm{gr})$, 
we can talk about a {\em backward morphism} 
$\th : N \to M$ in $\dcat{D}(A \ot C^{\mrm{op}}, \mrm{gr})$. 
The backward morphism $\th$ induces two morphisms by adjunction. 
First there is the morphism 
$\opn{badj}^{\mrm{R}}_{A}(\th) : N \to \opn{RHom}_{A}(B, M)$ 
in $\dcat{D}(B \ot C^{\mrm{op}}, \mrm{gr})$; this is the {\em derived backward 
adjunction morphism on the $A$ side}. The second is 
$\opn{badj}^{\mrm{R}}_{C^{\mrm{op}}}(\th) : N \to 
\opn{RHom}_{C^{\mrm{op}}}(D, M)$
in $\dcat{D}(A \ot D^{\mrm{op}}, \mrm{gr})$; this is the {\em derived backward 
adjunction morphism on the $C^{\mrm{op}}$ side}.
The constructions are just like those in Theorem \ref{thm:3192}(2), 
but now using the graded bimodule resolutions from Subsection 
\ref{gr-res-der-fun}.

\begin{dfn} \label{dfn:4031}
Consider a backward morphism $\th : N \to M$ 
in \lb $\dcat{D}(A \ot C^{\mrm{op}}, \mrm{gr})$. 
\begin{enumerate}
\item $\th$ is called a {\em nondegenerate backward morphism on the $A$ side} 
if the morphism $\opn{badj}^{\mrm{R}}_{A}(\th)$ is an isomorphism. 

\item $\th$ is called a {\em nondegenerate backward morphism on the 
$C^{\mrm{op}}$ side} if the morphism 
$\opn{badj}^{\mrm{R}}_{C^{\mrm{op}}}(\th)$ 
is an isomorphism. 

\item  $\th$ is said to be a {\em nondegenerate backward morphism on both 
sides} if it is nondegenerate both on the $A$ side and on the 
$C^{\mrm{op}}$ side. 
\end{enumerate}
\end{dfn}

\begin{exa} \label{exa:4585}
The backward morphism $f^* : B^* \to A^*$ 
in $\dcat{D}(A^{\mrm{en}}, \mrm{gr})$ is a nondegenerate backward morphism 
on both sides. 
\end{exa}

\begin{thm} \label{thm:4030}
The following hold in the situation of Definition \tup{\ref{dfn:4030}}.
\begin{enumerate}
\item There exists a unique 
balanced trace morphism 
$\opn{tr}_f : R_B \to R_A$.

\item The balanced trace morphism $\opn{tr}_f$ is a nondegenerate backward 
morphism  on both sides.
\end{enumerate}
\end{thm}

We shall need several lemmas for the proof. In analogy to Definition 
\ref{dfn:3750}, the dedualizing complex of $B$ is 
$P_B := \mrm{R} \Ga_{\n}(B) \cong \mrm{R} \Ga_{\n^{\mrm{op}}}(B)$ 
in $\dcat{D}(B^{\mrm{en}}, \mrm{gr})$.

\begin{lem} \label{lem:4035}
There are canonical isomorphisms 
$P_B \cong B \ot^{\mrm{L}}_{A} P_A$ and
$P_B \cong P_A \ot^{\mrm{L}}_{A} B$
in $\dcat{D}(B \ot A^{\mrm{op}}, \mrm{gr})$ and 
$\dcat{D}(A \ot B^{\mrm{op}}, \mrm{gr})$ respectively.
\end{lem}

\begin{proof}
According to Theorems \ref{thm:3783}(2) and \ref{thm:3750} there are 
isomorphisms
$P_B = \mrm{R} \Ga_{\n}(B) \cong \mrm{R} \Ga_{\m}(B) \cong 
P_A \ot^{\mrm{L}}_{A} B$
in $\dcat{D}(A \ot B^{\mrm{op}}, \mrm{gr})$. The same theorems, but applied to 
the opposite rings, give us isomorphisms
$P_B \cong \mrm{R} \Ga_{\n^{\mrm{op}}}(B) \cong 
\mrm{R} \Ga_{\m^{\mrm{op}}}(B) \cong B \ot^{\mrm{L}}_{A} P_A$
in $\dcat{D}(B \ot A^{\mrm{op}}, \mrm{gr})$.
\end{proof}

Like Definition \ref{dfn:3790}, the abstract derived $\n$-adic completion 
functor is 
\[ \opn{ADC}_{\n} := \opn{RHom}_{B}(P_B, -) : 
\dcat{D}(B \ot D^{\mrm{op}}, \mrm{gr}) \to 
\dcat{D}(B \ot D^{\mrm{op}}, \mrm{gr}) . \]

\begin{lem} \label{lem:4036}
For $N \in \dcat{D}(B \ot D, \mrm{gr})$
there is an isomorphism 
$\opn{ADC}_{\n}(N) \cong \opn{ADC}_{\m}(N)$
in $\dcat{D}(A \ot D^{\mrm{op}}, \mrm{gr})$.
It is functorial in $N$. 
\end{lem}

\begin{proof}
We have isomorphisms 
\[ \begin{aligned}
& \opn{ADC}_{\n}(N) = \opn{RHom}_{B}(P_B, N) \cong^1 
\opn{RHom}_{B}(B \ot^{\mrm{L}}_{A} P_A, N) 
\\
& \qquad 
\cong^2 \opn{RHom}_{A}(P_A, N) = \opn{ADC}_{\m}(N)
\end{aligned} \]
in $\dcat{D}(A \ot D^{\mrm{op}}, \mrm{gr})$. 
The isomorphism $\cong^1$ comes from Lemma \ref{lem:4035}, and the 
isomorphism $\cong^2$ is by adjunction. 
\end{proof}

In the next lemma we take $C = D := \K$. 

\begin{lem} \label{lem:4065}
Let $N \in \dcat{D}^{+}_{\mrm{f}}(B, \mrm{gr})$,
let $M \in \dcat{D}^{+}_{\mrm{f}}(A, \mrm{gr})$,
and let 
$\th : N \to M$ be a backward morphism in 
$\dcat{D}(A, \mrm{gr})$. The following conditions are 
equivalent\tup{:}
\begin{enumerate}
\rmitem{i} $\th$ is a nondegenerate backward morphism on the $A$ side.  

\rmitem{ii} The morphism 
$\mrm{R} \Ga_{\m}(\th) : \mrm{R} \Ga_{\m}(N) \to \mrm{R} \Ga_{\m}(M)$
in $\dcat{D}(A, \mrm{gr})$ is a nondegenerate backward morphism on the $A$ 
side. 
\end{enumerate}
\end{lem}

\begin{proof}
Let us choose bounded below injective resolutions 
$N \to J$ and $M \to I$ in 
$\dcat{C}_{\mrm{str}}(B, \mrm{gr})$ and 
$\dcat{C}_{\mrm{str}}(A, \mrm{gr})$ respectively. Then the backward morphism 
$\th$ is represented by a homomorphism 
$\til{\th} : J \to I$ 
in $\dcat{C}_{\mrm{str}}(A, \mrm{gr})$. 

The morphism 
$\opn{badj}^{\mrm{R}}_{A}(\th)  : N \to \opn{RHom}_{A}(B, M)$
in $\dcat{D}(B, \mrm{gr})$ is represented by the homomorphism 
$\opn{badj}_{A}(\til{\th}) : J \to \opn{Hom}_{A}(B, I)$
in $\dcat{C}_{\mrm{str}}(B, \mrm{gr})$.
The formula for it it this: 
$\opn{badj}_{A}(\til{\th})(j)(b) = \til{\th}(b \cd j)$ for $j \in J$ 
and $b \in B$. 
Now $\opn{Hom}_{A}(B, I)$ is a K-graded-injective complex over $B$, by 
adjunction. By Lemma \ref{lem:3766} the complexes $J$ and 
$\opn{Hom}_{A}(B, I)$ are graded-$\m$-flasque. So the morphism 
\[ \mrm{R} \Ga_{\m} \bigl( \opn{badj}_{A}^{\mrm{R}}(\th) \bigr) : 
\mrm{R} \Ga_{\m}(N) \to \mrm{R} \Ga_{\m} \bigl( \opn{RHom}_{A}(B, M) \bigr) \]
is represented by 
\begin{equation} \label{eqn:4850}
\Ga_{\m}(\opn{badj}_{A}(\til{\th})) : \Ga_{\m}(J) \to  
\Ga_{\m} \bigl( \opn{Hom}_{A}(B, I) \bigr) . 
\end{equation}

The morphism $\mrm{R} \Ga_{\m}(\th)$ is represented by the homomorphism
$\Ga_{\m}(\til{\th}) : \Ga_{\m}(J) \to \Ga_{\m}(I)$.
Since the complex $\Ga_{\m}(I)$ is K-graded-injective over $A$ (by 
stability, see Theorem \ref{thm:3990}), it follows that 
$\opn{badj}_A^{\mrm{R}} \bigl( \mrm{R} \Ga_{\m}(\th) \bigr)$
is represented by 
\begin{equation} \label{eqn:4851}
\opn{badj}_A(\Ga_{\m}(\til{\th})) : \Ga_{\m}(J) \to
\opn{Hom}_A(B, \Ga_{\m}(I)) . 
\end{equation}

A calculation, using the fact that $B$ is a finite $A$-module, shows that 
\begin{equation} \label{eqn:4853}
 \Ga_{\m} \bigl( \opn{Hom}_{A}(B, I) \bigr) = 
\opn{Hom}_{A} \bigl( B, \Ga_{\m}(I) \bigr) 
\end{equation}
as subcomplexes of $\opn{Hom}_{A}(B, I)$. 
Let us consider the diagram 
\begin{equation} \label{eqn:4854}
\UseTips \xymatrix @C=11ex @R=6ex {
\Ga_{\m}(J)
\ar[r]^(0.38){ \opn{badj}_A(\Ga_{\m}(\til{\th})) }
\ar[dr]_(0.45){ \Ga_{\m}(\opn{badj}_A(\til{\th})) \ \ }
&
\opn{Hom}_{A} \bigl( B, \Ga_{\m}(I) \bigr)
\ar[d]^{=}
\ar[dr]^{\sub} 
\\
&
\Ga_{\m} \bigl( \opn{Hom}_{A}(B, I) \bigr)
\ar[r]^{\sub} 
&
\opn{Hom}_{A}(B, I)
}
\end{equation}
in $\dcat{C}_{\mrm{str}}(A, \mrm{gr})$. It is commutative: the triangle on the 
right is just the equality (\ref{eqn:4853}). And an element 
$j \in \Ga_{\m}(J)$ goes by all paths to the homomorphism 
$b \mapsto \til{\th}(b \cd j)$ in $\opn{Hom}_{A}(B, I)$. 

Passing to the derived category, (\ref{eqn:4854}) gives the 
commutative diagram
\begin{equation} \label{eqn:4855}
\UseTips \xymatrix @C=14ex @R=6ex {
\mrm{R} \Ga_{\m}(N)
\ar[r]^(0.4){ \opn{badj}_A^{\mrm{R}} (\mrm{R} \Ga_{\m}(\th)) }
\ar[dr]_(0.45){ \mrm{R} \Ga_{\m} (\opn{badj}_A^{\mrm{R}}(\th)) \ \ }
&
\opn{RHom}_{A} \bigl( B, \mrm{R} \Ga_{\m}(M) \bigr)
\ar[d]^{\cong}
\\
&
\mrm{R} \Ga_{\m} \bigl( \opn{RHom}_{A}(B, M) \bigr) 
}
\end{equation}
in $\dcat{D}(A, \mrm{gr})$.

Since both $N$ and $\opn{RHom}_{A}(B, M)$ belong to 
$\dcat{D}^+_{\mrm{f}}(A, \mrm{gr})$, by Lemma \ref{lem:4060} they belong to 
$\dcat{D}^{+}(A, \mrm{gr})_{\mrm{com}}$.
By Theorem \ref{thm:3792} (the MGM Equivalence) we know that the morphism
$\opn{badj}_A^{\mrm{R}}(\th)$ is an isomorphism if and only if 
the morphism \lb $\mrm{R} \Ga_{\m}(\opn{badj}_A^{\mrm{R}}(\th))$ is an 
isomorphism. Diagram (\ref{eqn:4855}) says that this happens if and only if 
$\opn{badj}_A^{\mrm{R}} \bigl( \mrm{R} \Ga_{\m}(\th) \bigr)$
is an isomorphism. Going back to the definition, we conclude that $\th$ is 
nondegenerate if and only if $\mrm{R} \Ga_{\m}(\th)$ is nondegenerate. 
\end{proof}

\begin{proof}[Proof of Theorem \tup{\ref{thm:4030}}] \mbox{}

\smallskip \noindent
(1) According to Theorems \ref{thm:3713} and \ref{thm:3710} there is an 
isomorphism 
$R_A \cong (P_A)^* \cong \opn{ADC}_{\m}(A^*)$
in $\dcat{D}(A^{\mrm{en}}, \mrm{gr})$. 
The same theorems, together with Lemma \ref{lem:4036}
(taking $C := A$ and $D := B$), say that 
$R_B \cong (P_B)^* \cong \opn{ADC}_{\n}(B^*) \cong \opn{ADC}_{\m}(B^*)$
in $\dcat{D}(A^{\mrm{en}}, \mrm{gr})$. 

The equivalence in Corollary \ref{cor:4070} says that there are 
isomorphisms 
\begin{equation} \label{eqn:4067}
\mrm{R} \Ga_{\m}(R_B) \cong (\mrm{R} \Ga_{\m} \circ \opn{ADC}_{\m})(B^*)
\cong B^*
\end{equation}
and 
\begin{equation} \label{eqn:4068}
\mrm{R} \Ga_{\m}(R_A) \cong (\mrm{R} \Ga_{\m} \circ \opn{ADC}_{\m})(A^*)
\cong A^*
\end{equation}
in $\dcat{D}(A^{\mrm{en}}, \mrm{gr})$.
And that there is a unique morphism 
$\th : R_B \to R_A$
in $\dcat{D}(A^{\mrm{en}}, \mrm{gr})$ such that the diagram 
\begin{equation} \label{eqn:4066}
\UseTips \xymatrix @C=8ex @R=6ex {
\mrm{R} \Ga_{\m}(R_B)
\ar[r]^{\mrm{R} \Ga_{\m}(\th)}
\ar[d]_{\cong}
&
\mrm{R} \Ga_{\m}(R_A)
\ar[d]^{\cong}
\\
B^*
\ar[r]^(0.5){f^*}
&
A^*
} 
\end{equation}
in $\dcat{D}(A^{\mrm{en}}, \mrm{gr})$, in which the vertical isomorphisms are 
(\ref{eqn:4067}) and (\ref{eqn:4068}), is commutative. 

But to get a balanced trace morphism $\opn{tr}_f$ we need a commutative diagram 
like (\ref{eqn:4066}) in which the vertical isomorphisms are $\be_B$ and 
$\be_A$ respectively. Because the automorphisms of $A^*$ and $B^*$ in 
$\dcat{D}(A^{\mrm{en}}, \mrm{gr})$
are multiplication by nonzero elements of $\K$, there is a unique 
$c \in \K^{\times}$ such that 
$\opn{tr}_f := c \cd \th : R_B \to R_A$
is a balanced trace morphism.

\medskip \noindent 
(2) To see that $\opn{tr}_f$ is a nondegenerate backward morphism on the 
$A$ side we can forget the $A^{\mrm{op}}$-structure.
So we can view $\opn{tr}_f$ as a backward morphism
$\opn{tr}_f : R_B \to R_A$ in $\dcat{D}(A, \mrm{gr})$.
Now
$\mrm{R} \Ga_{\m}(\opn{tr}_f) = \be_{A}^{-1} \circ f^* \circ \be_B$.
In Example \ref{exa:4585} we saw that $f^* : B^* \to A^*$ is a nondegenerate 
backward morphism in $\dcat{D}(A, \mrm{gr})$. The morphisms $\be_A$ and $\be_B$ 
are isomorphisms. Therefore 
\[ \mrm{R} \Ga_{\m}(\opn{tr}_f) : \mrm{R} \Ga_{\m}(R_B) \to 
\mrm{R} \Ga_{\m}(R_A) \]
is a nondegenerate backward morphism in $\dcat{D}(A, \mrm{gr})$. 
According to Lemma \ref{lem:4065} the morphism $\opn{tr}_f$
is a nondegenerate backward morphism in $\dcat{D}(A, \mrm{gr})$. 

As for the $A^{\mrm{op}}$ side: the symmetry from Corollary \ref{cor:4080} 
gives isomorphisms 
$R_A \cong \opn{ADC}_{\m^{\mrm{op}}}(A^*)$
and
$R_B \cong \opn{ADC}_{\m^{\mrm{op}}}(B^*)$
in $\dcat{D}(A^{\mrm{en}}, \mrm{gr})$. 
Going over the constructions in the proof of item (1) above, we see that
there is equality 
$\mrm{R} \Ga_{\m^{\mrm{op}}}(\opn{tr}_f) = 
d \cd \be_{A}^{-1} \circ f^* \circ \be_B$
for some $d \in \K^{\times}$.
Therefore $\mrm{R} \Ga_{\m^{\mrm{op}}}(\opn{tr}_f)$ 
is a nondegenerate backward morphism 
$\dcat{D}(A^{\mrm{op}}, \mrm{gr})$.
Now Lemma \ref{lem:4065}, transcribed to the ring $A^{\mrm{op}}$, 
says that $\opn{tr}_f$ is a 
nondegenerate backward morphism in
$\dcat{D}(A^{\mrm{op}}, \mrm{gr})$.
\end{proof}

\begin{cor} \label{cor:4065}
Let $A \xar{f} B \xar{g} C$ be finite homomorphisms in 
$\catt{Rng}_{\mrm{gr}} \centover \K$, 
and assume these rings have balanced dualizing 
complexes $(R_A, \be_A)$, $(R_B, \be_B)$ and $(R_C, \be_C)$
respectively. 
Then 
$\opn{tr}_{g \circ f} = \opn{tr}_f \circ \opn{tr}_g$
as morphisms $R_C \to R_A$ in $\dcat{D}(A^{\mrm{en}}, \mrm{gr})$. 
\end{cor}

\begin{proof}
The follows from the uniqueness in Theorem \ref{thm:4030} and the fact that 
$(g \circ f)^* = f^* \circ g^*$ as homomorphisms $C^* \to A^*$ in 
$\dcat{M}(A, \mrm{gr})$. 
\end{proof}

We end this subsection with a partial converse to Corollaries 
\ref{cor:4030} and \ref{cor:4031}.
Recall that a central element $a \in A$ is called {\em regular} if it is not a 
zero divisor, i.e.\ the only element $b \in A$ such that $a \cd b = 0$ is 
$b = 0$. 

\begin{thm} \label{thm:4250}
\index{Algebraically graded ring! chi@$\chi$ condition on connected}
\index{Algebraically graded ring! connected {\indash} of finite local 
cohomological dimension}
Let $A$ be a noetherian connected graded ring, and let $a \in A$ be a regular
homogeneous central element of positive degree. Define the connected graded 
ring $B := A / (a)$. If $B$ satisfies the $\chi$ condition and it has finite 
local cohomological dimension, then $A$ also  satisfies the $\chi$ condition 
and it has finite local cohomological dimension.
\end{thm}

This is \cite[Theorem 8.8]{ArZh}. The original proof does not use 
derived categories, and is quite different from the proof below. 

\begin{proof}
By op-symmetry it suffices to prove that $A$ satisfies the left $\chi$ 
condition and that the functor $\mrm{R} \Ga_{\m}$ has finite cohomological 
dimension. 

Let $d \in \N$ be the local cohomological dimension of $B$, and let 
$i \geq 1$ be the degree of the element $a$. 
In view of Propositions \ref{prop:3855} and \ref{prop:4050}, all we need to 
prove is that the following two conditions hold for every 
$M \in \dcat{M}_{\mrm{f}}(A, \mrm{gr})$~:
\begin{itemize}
\rmitem{i} $\opn{H}^p_{\m}(M)$ is a cofinite $A$-module for every $p$.

\rmitem{ii} $\opn{H}^p_{\m}(M) = 0$ for every $p > d + 1$. 
\end{itemize}

Let us denote the augmentation ideal of $B$ by $\n$. 
The short exact sequence 
$0 \to A \xar{a \cd (-)} A(i) \to B(i) \to 0$
in $\dcat{M}(A^{\mrm{en}}, \mrm{gr})$ is viewed as a distinguished triangle
%\begin{equation} \label{4250}
$A \xar{a \cd (-)} A(i) \to B(i) \xar{ \, \triangle \, }$
%\end{equation}
in $\dcat{D}(A^{\mrm{en}}, \mrm{gr})$. 
Applying the functor $(-) \ot^{\mrm{L}}_{A} M$ to it, we obtain a 
distinguished triangle 
\begin{equation} \label{4251}
M \xar{a \cd (-)} M(i) \to N \xar{ \, \triangle \, }
\end{equation}
in $\dcat{D}(A, \mrm{gr})$, where by definition 
$N := B(i) \ot^{\mrm{L}}_{A} M \in \dcat{D}(B, \mrm{gr})$.
Note that 
$\opn{H}^{p}(N) \in \dcat{M}_{\mrm{f}}(B, \mrm{gr})$
for all $p$, and $\opn{H}^{p}(N) = 0$ unless $-1 \leq p \leq 0$. 

Next we apply the functor $\mrm{R} \Ga_{\m}$ to (\ref{4251}), and use Theorem 
\ref{thm:3783}, to obtain this distinguished triangle 
\begin{equation} \label{4252}
\mrm{R} \Ga_{\m}(M) \xar{a \cd (-)} \mrm{R} \Ga_{\m}(M(i)) \to 
\mrm{R} \Ga_{\n}(N) \xar{ \, \triangle \, }
\end{equation}
in $\dcat{D}(A, \mrm{gr})$.
Taking cohomologies in (\ref{4252}) we obtain this exact sequence 
\begin{equation} \label{4253}
\opn{H}^{p - 1}_{\n}(N) \to 
\opn{H}^{p}_{\m}(M) \xar{a \cd (-)}
\opn{H}^{p}_{\m}(M(i)) \to  
\opn{H}^{p}_{\n}(N) 
\end{equation}
in $\dcat{M}(A, \mrm{gr})$. 
Let use write 
\[ K^p := \opn{Ker} \bigl( \opn{H}^{p}_{\m}(M) \xar{a \cd (-)}
\opn{H}^{p}_{\m}(M(i)) \bigr) \sub \opn{H}^{p}_{\m}(M) . \]
Since $a \in \m$, there is an inclusion 
\begin{equation} \label{4254}
\opn{Soc} \bigl( \opn{H}^{p}_{\m}(M) \bigr) =
\opn{Hom}_{A} \bigl( \K, \opn{H}^{p}_{\m}(M) \bigr) \sub K^p . 
\end{equation}

For $p > d + 1$ we have 
$\opn{H}^{p - 1}_{\n}(N) = 0$ and 
$\opn{H}^{p}_{\n}(N) = 0$. 
The exact sequence (\ref{4253}) says that the homomorphism 
$a \cd (-) : \opn{H}^{p}_{\m}(M) \to \opn{H}^{p}_{\m}(M(i))$
is bijective. So $K^p = 0$, and by (\ref{4254}) we see that 
$\opn{Soc} \bigl( \opn{H}^{p}_{\m}(M) \bigr) = 0$. 
But $\opn{H}^{p}_{\m}(M)$ is an $\m$-torsion module, so according to 
Proposition \ref{prop:4811} we get $\opn{H}^{p}_{\m}(M) = 0$.
This establishes condition (ii). 

Finally take any $p$.
By the $\chi$ condition for $B$ and by Proposition \ref{prop:3855} we know 
that $\opn{H}^{p - 1}_{\n}(N)$ is a cofinite graded $B$-module; and thus it is 
a cofinite graded $A$-module. According to the NC Graded Matlis Duality 
(Theorem 
\ref{thm:4045}) the cofinite graded $A$-modules are the artinian objects in the 
abelian category $\dcat{M}(A, \mrm{gr})$.
From the exact sequence (\ref{4253}) we obtain a surjection 
$\opn{H}^{p - 1}_{\n}(N) \to K^p$.
Therefore $K^p$ is a cofinite graded $A$-module. 
The inclusion (\ref{4254}) says that 
$\opn{Soc} \bigl( \opn{H}^{p}_{\m}(M) \bigr)$ 
is a cofinite graded $A$-module; and thus it is a finite graded $\K$-module. 
Hence, by Corollary \ref{cor:4270}, 
$\opn{H}^{p}_{\m}(M)$ is a cofinite graded $A$-module.
This is condition (i). 
\end{proof}





\cleardoublepage
\mysection{Rigid Noncommutative Dualizing Complexes} 
\label{sec:rigid-DC-NC} 

\AYcopyright 

In this section of the book we are going to work with noncommutative (namely 
not necessarily commutative) rings. We shall often use the 
abbreviation ``NC'' for ``noncommutative''. 
The goal is to introduce {\em rigid NC dualizing complexes} in the sense of M. 
Van den Bergh \cite{VdB}, and to prove their existence and uniqueness under 
certain conditions. 


\mysubsection{Noncommutative Dualizing Complexes}
\label{subsec:NC-DC}

Here we will define NC dualizing complexes over a NC ring $A$. 
The definition is almost identical to that in the NC graded setting (see 
Subsection \ref{subsec:graded-NCDC}). We will prove 
that the NC dualizing complexes over $A$ are 
parameterized by the derived Picard group of $A$. The content of this 
subsection is adapted from the paper \cite{Ye4}.

Let $\K$ be a commutative ring. 
Recall (from Definition \ref{dfn:4125}) that a {\em central $\K$-ring} is a 
ring $A$ equipped with a ring homomorphism $\K \to \opn{Cent}(A)$, where 
$\opn{Cent}(A)$ is the center of $A$. 
In more traditional texts, a central $\K$-ring is called an {\em associative 
unital $\K$-algebra}. The category of central $\K$-rings, with $\K$-ring 
homomorphisms, is denoted by $\catt{Rng} \centover \K$. 

Here is the convention that is in force throughout this section:

\begin{conv} \label{conv:3670} 
There is a  base field $\K$. All rings are by default central 
$\K$-rings, and all homomorphisms between $\K$-rings are over $\K$; namely we 
work within the category $\catt{Rng} \centover \K$.
All bimodules are $\K$-central, and all additive functors are $\K$-linear.
We use the notation $\ot$ for $\ot_{\K}$. For a ring $A$ we write 
$A^{\mrm{en}} := A \ot A^{\mrm{op}}$, the enveloping ring of $A$. 
\end{conv}

See Remark \ref{rem:3691} regarding the assumption that the base ring 
$\K$ is a field.    

Recall that a NC ring $A$ is called noetherian if it
is both left noetherian and right noetherian.
See Remark \ref{rem:3901} regarding the failure of the noetherian property to 
be preserved under finitely generated ring extensions. 

Let $A$ be a ring. As before, $\dcat{M}(A)$ is the category of left 
$A$-modules. The category of right $A$-modules is $\dcat{M}(A^{\mrm{op}})$,
and the category of $A$-bimodules is $\dcat{M}(A^{\mrm{en}})$.
Given another ring $B$, the category of $A$-$B$-bimodules is 
$\dcat{M}(A \ot B^{\mrm{op}})$.
These are $\K$-linear abelian categories. 

The derived category of $\dcat{M}(A)$ is 
$\dcat{D}(A) := \dcat{D}(\dcat{M}(A))$; its objects are the 
complexes of (left) $A$-modules. 
The other derived categories are  
$\dcat{D}(A^{\mrm{op}}) := \dcat{D}(\dcat{M}(A^{\mrm{op}}))$,
$\dcat{D}(A^{\mrm{en}}) := \dcat{D}(\dcat{M}(A^{\mrm{en}}))$
and 
$\dcat{D}(A \ot B^{\mrm{op}}) := \dcat{D}(\dcat{M}(A \ot B^{\mrm{op}}))$.

Consider the canonical ring anti-automorphism 
$\opn{op} : B \to B^{\mrm{op}}$, that is the identity on the underlying 
$\K$-module. There is a canonical ring isomorphism 
\begin{equation} \label{eqn:4358}
A \ot B^{\mrm{op}} \iso B^{\mrm{op}} \ot A ,  \quad  
a \ot \opn{op}(b) \mapsto \opn{op}(b) \ot a .
\end{equation}
The induced isomorphism on the bimodule categories 
$\dcat{M}(A \ot B^{\mrm{op}}) \iso \lb \dcat{M}(B^{\mrm{op}} \ot A)$
is the identity on the underlying $\K$-modules. Indeed, for 
$M \in \dcat{M}(A \ot B^{\mrm{op}})$, $m \in M$ and 
$a \ot \opn{op}(b) \in A \ot B^{\mrm{op}}$, we have 
$(\opn{op}(b) \ot a) \cd m = a \cd m \cd b = (a \ot \opn{op}(b)) \cd m$.
Now take $B = A$ in (\ref{eqn:4358}). Since 
$(A^{\mrm{op}})^{\mrm{op}} = A$, there is a canonical ring isomorphism 
\begin{equation} \label{eqn:4450}
A^{\mrm{en}} = A \ot A^{\mrm{op}} \iso 
A^{\mrm{op}} \ot A = (A^{\mrm{op}})^{\mrm{en}} 
\end{equation}
and an induced isomorphism on the bimodule categories 
\begin{equation} \label{eqn:4355}
\dcat{M}(A^{\mrm{en}}) \iso \dcat{M}((A^{\mrm{op}})^{\mrm{en}}) .
\end{equation}
All these categorical relations pass to derived categories. 

The ring $A^{\mrm{en}}$ is isomorphic to its opposite ring. 
There is a ring isomorphism 
\begin{equation} \label{eqn:4875}
\opn{trop} : A^{\mrm{en}} \iso (A^{\mrm{en}})^{\mrm{op}}
\end{equation}
with this formula: for $a_1, a_2 \in A$ it sends the element 
$a_1 \ot \opn{op}(a_2) \in A^{\mrm{en}}$ to 
\begin{equation} \label{eqn:4876}
\opn{trop}(a_1 \ot \opn{op}(a_2)) := 
\opn{op}(a_2 \ot \opn{op}(a_1)) \in (A^{\mrm{en}})^{\mrm{op}} .
\end{equation}
On the underlying $\K$-module $A \ot A$ this is just the transposition 
$a_1 \ot a_2 \mapsto a_2 \ot a_1$, and hence the name ``transpose opposite'' 
and the acronym ``trop''. 

Usually an $A$-$A$-bimodule $M$ is seen as a left 
$A^{\mrm{en}}$-module, with this action of an element 
$a_1 \ot \opn{op}(a_2) \in A^{\mrm{en}}$ on an element $m \in M$~:
\begin{equation} \label{eqn:4887}
(a_1 \ot \opn{op}(a_2)) \cd m = 
a_1 \cd m \cd a_2 .
\end{equation}
But using the ring isomorphism $\opn{trop}$ from formula (\ref{eqn:4875}) we 
can also consider $M$ as a left $(A^{\mrm{en}})^{\mrm{op}}$-module, and thus as 
a right $A^{\mrm{en}}$-module. Explicitly, the right action of an element 
$a_2 \ot \opn{op}(a_1) \in A^{\mrm{en}}$ on an element $m \in M$ is:
\begin{equation} \label{eqn:4878}
\begin{aligned}
&
m \cd (a_2 \ot \opn{op}(a_1)) =^{\heartsuit}
\opn{op}(a_2 \ot \opn{op}(a_1)) \cd m
\\ & \quad 
=^{\dag} (a_1 \ot \opn{op}(a_2)) \cd m =^{\diamondsuit} a_1 \cd m \cd a_2 . 
\end{aligned}
\end{equation}
In equality $=^{\heartsuit}$ we move from a right action by $A^{\mrm{en}}$
to the left action by $(A^{\mrm{en}})^{\mrm{op}}$, and in 
equality $=^{\dag}$ we made use of formula (\ref{eqn:4876})
to move from $(A^{\mrm{en}})^{\mrm{op}}$ to $A^{\mrm{en}}$.
Finally, the equality $=^{\diamondsuit}$ is formula (\ref{eqn:4887}). 
This game of musical chairs will be played in Subsections 
\ref{subsec:Fil-exit-R-NCDC} and \ref{subsec:RNCDC-uniq}.

Let $A$ and $B$ be rings. The canonical ring homomorphisms 
\[ A \to A \ot B^{\mrm{op}} \lar B^{\mrm{op}} \]
induce restriction functors 
\begin{equation} \label{eqn:3680}
\dcat{D}(A) \xleftarrow{\opn{Rest}_A} \dcat{D}(A \ot B^{\mrm{op}}) 
\xar{\opn{Rest}_{B^{\mrm{op}}}} \dcat{D}(B^{\mrm{op}}) . 
\end{equation}
See the commutative diagram \ref{eqn:3536}. We shall usually keep 
these restriction functors implicit, and instead use terminology like in 
Definition \ref{dfn:3446}.

In Definition \ref{dfn:2135} we introduced this notation: for a ring $A$ 
we denote by $\dcat{M}_{\mrm{f}}(A)$ the full subcategory of 
$\dcat{M}(A)$ on the finite (i.e.\ finitely generated) $A$-modules.
And $\dcat{D}_{\mrm{f}}(A)$ is the full subcategory of 
$\dcat{D}(A)$ on the complexes of $A$-modules $M$ such that 
$\opn{H}^i(M) \in \dcat{M}_{\mrm{f}}(A)$ for every $i$. 

If $A$ is left noetherian, then $\dcat{M}_{\mrm{f}}(A)$ is a thick abelian 
subcategory of $\dcat{M}(A)$; and hence $\dcat{D}_{\mrm{f}}(A)$ is a full 
triangulated subcategory of $\dcat{D}(A)$. 
As usual we can combine indicators: 
$\dcat{D}_{\mrm{f}}^{\star}(A) := \dcat{D}_{\mrm{f}}(A) \cap 
\dcat{D}^{\star}(A)$,
where $\star$ is some boundedness indicator ($+$, $-$, $\mrm{b}$ or
$\bra{\mrm{empty}}$). 

Now to finiteness of $A$-$B$-bimodules. The expression  
$\dcat{M}_{(\mrm{f}, ..)}(A \ot B^{\mrm{op}})$
denotes the full subcategory of $\dcat{M}(A \ot B^{\mrm{op}})$
on the bimodules that are finite over $A$; and the expression 
$\dcat{M}_{(.., \mrm{f})}(A \ot B^{\mrm{op}})$
denotes the full subcategory on the bimodules that are finite over 
$B^{\mrm{op}}$. Of course we define 
\[ \dcat{M}_{(\mrm{f}, \mrm{f})}(A \ot B^{\mrm{op}}) :=
\dcat{M}_{(\mrm{f}, ..)}(A \ot B^{\mrm{op}}) \cap 
\dcat{M}_{(.., \mrm{f})}(A \ot B^{\mrm{op}}) . \]
The same notation shall apply to derived categories. 

To a complex of bimodules 
$M \in \dcat{D}(A \ot B^{\mrm{op}})$ 
we associated two derived homothety morphisms in Subsection 
\ref{subsec:tilting}. Let us recall them. 
The derived homothety morphism through $B^{\mrm{op}}$ is the morphism 
\begin{equation} \label{equ:4336}
\opn{hm}^{\mrm{R}}_{M, B^{\mrm{op}}} : B \to \opn{RHom}_{A}(M, M)  
\end{equation}
in $\dcat{D}(B^{\mrm{en}})$; and the derived homothety morphism through $A$ 
is the morphism 
\begin{equation} \label{equ:4337}
\opn{hm}^{\mrm{R}}_{M, A} : A \to \opn{RHom}_{B^{\mrm{op}}}(M, M) 
\end{equation}
in $\dcat{D}(A^{\mrm{en}})$. 
According to Definition \ref{3455}, the complex $M$ has the  
noncommutative derived Morita property on the $A$ side (resp.\ on the 
$B^{\mrm{op}}$ side) if the morphism 
$\opn{hm}^{\mrm{R}}_{M, B^{\mrm{op}}}$
(resp.\ $\opn{hm}^{\mrm{R}}_{M, A}$) is an isomorphism. 

\begin{dfn}[\cite{Ye4}] \label{dfn:3675} 
Let $A$ be a noetherian ring. A 
{\em noncommutative dualizing complex over $A$} 
\index{Dualizing complex! noncommutative}
is an object 
$R \in \dcat{D}^{\mrm{b}}(A^{\mrm{en}})$ satisfying the 
following three conditions:
\begin{itemize}
\rmitem{i} For every $p$ the 
$A$-bimodule $\opn{H}^p(R)$ is finite over $A$ and over $A^{\mrm{op}}$. 

\rmitem{ii} The complex $R$ has finite injective 
dimension over $A$ and over $A^{\mrm{op}}$. 

\rmitem{iii} The complex of bimodules $R$ has the noncommutative derived Morita 
property on both sides. 
\end{itemize}
\end{dfn}

Note that conditions (i) and (ii) are one-sided, namely they refer separately 
to $\opn{Rest}_{A}(R) \in \dcat{D}(A)$ and 
to  $\opn{Rest}_{A^{\mrm{op}}}(R) \in \dcat{D}(A^{\mrm{op}})$; 
whereas condition (iii) is more complicated, and we will study it further in 
Lemma \ref{lem:4361} below. Also note that condition (i) can be stated as ``$R$ 
belongs to $\dcat{D}_{(\mrm{f}, \mrm{f})}(A^{\mrm{en}})$''. 

Before going on, we need to say something about the op-symmetry of  
Definition \ref{dfn:3675}. Given a complex 
$R \in \dcat{D}(A^{\mrm{en}})$, by formula (\ref{eqn:4355}) the complex $R$ is 
also an object of $\dcat{D}((A^{\mrm{op}})^{\mrm{en}})$, and we could ask 
whether $R$ is a dualizing complex over the ring $A^{\mrm{op}}$. The answer is 
positive of course, because all three conditions are op-symmetric. We record 
this fact as the next proposition. 

\begin{prop} \label{prop:4355}
If $R$ is a NC dualizing complex over $A$, then $R$ is also a 
NC dualizing complex over $A^{\mrm{op}}$. 
\end{prop}

\begin{exa} \label{exa:3695} 
A NC ring $A$ is called {\em regular} if it has finite global cohomological 
dimension on both sides. 
This is very similar to the graded definition (see Definition \ref{dfn:3920}).
Let us recall what it means: there is some 
natural number $d$ such that $\opn{Ext}^i_A(M, N) = 0$ for all $i > d$ and 
all $M, N \in \dcat{M}(A)$; and also  
$\opn{Ext}^i_{A^{\mrm{op}}}(M', N') = 0$ for all $i > d$ and 
all $M', N' \in \dcat{M}(A^{\mrm{op}})$. 
The smallest such number $d$ is the global cohomological dimension of $A$. 

Such rings are easy to find; for instance, if $\K$ is a field, and $\g$ is 
finite Lie algebra over $\K$, then the {\em universal enveloping ring}
$A := \opn{U}(\g)$ is regular; the global cohomological dimension of $A$ is 
$d := \opn{rank}_{\K}(\g)$.  

A weaker condition is this: $A$ is called {\em Gorenstein} if $A$ has finite 
injective dimension both as a left module and as a right module over itself. 
Namely there is some natural number $d$ such that 
$\opn{Ext}^i_A(M, A) = 0$ for all $i > d$ and all $M \in \dcat{M}(A)$; and 
also $\opn{Ext}^i_{A^{\mrm{op}}}(M', A) = 0$ for all $i > d$ and 
all $M' \in \dcat{M}(A^{\mrm{op}})$. 

Assume $A$ is Gorenstein. Then the complex of bimodules 
$R := A$ satisfies condition (ii) of Definition \ref{dfn:3675}. 
Conditions (i) and (iii) of this definition are automatically true. We conclude 
that $R := A$ is a NC dualizing complex over itself. 
\end{exa}

\begin{dfn} \label{dfn:3680}
Let $A$ be a noetherian ring, let $R$ be a NC dualizing complex over $A$, and 
let $B$ be a second ring. 
The {\em duality functors} 
\index{Duality functor}
associated to $R$ are the triangulated functors 
\[ D_A : \dcat{D}(A \ot B^{\mrm{op}})^{\mrm{op}} \to 
\dcat{D}(B \ot A^{\mrm{op}}), \quad 
D_A := \opn{RHom}_A(-, R)   \]
and 
\[ D_{A^{\mrm{op}}} : \dcat{D}(B \ot A^{\mrm{op}})^{\mrm{op}} 
\to \dcat{D}(A \ot B^{\mrm{op}}), \quad 
D_{A^{\mrm{op}}} := \opn{RHom}_{A^{\mrm{op}}}(-, R) .  \]
\end{dfn}

Notice that the expressions $D_A$ and $D_{A^{\mrm{op}}}$ leave $R$ and $B$ 
implicit; this is because the dualizing complex $R$ is taken to be fixed, and 
the second ring $B$ is less important. 

The NC derived Hom-evaluation morphisms were already discussed in Subsection 
\ref{subsec:graded-NCDC}, in the algebraically graded context. The nongraded 
definition is very similar. Given a complex 
$M \in \dcat{D}(A \ot B^{\mrm{op}})$, there is a morphism 
\begin{equation} \label{equ:4335}
\opn{ev}^{\mrm{R, R}}_{M} : M \to D_{A^{\mrm{op}}}(D_{A}(M)) = 
\opn{RHom}_{A^{\mrm{op}}} \bigl( \opn{RHom}_{A}(M, R), R \bigr)
\end{equation}
in $\dcat{D}(A \ot B^{\mrm{op}})$, which is functorial in $M$. If we choose a 
K-injective resolution $R \to I$ in $\dcat{D}(A^{\mrm{en}})$, then we have an 
isomorphism 
\[ D_{A^{\mrm{op}}}(D_{A}(M)) \cong 
\opn{Hom}_{A^\mrm{op}} \bigl( \opn{Hom}_{A}(M, I), I \bigr) \]
in $\dcat{D}(A \ot B^{\mrm{op}})$, and the morphism 
$\opn{ev}^{\mrm{R, R}}_{M}$ is represented by the Hom-evaluation homomorphism
\[ \opn{ev}_{M, I} : M \to 
\opn{Hom}_{A^\mrm{op}} \bigl( \opn{Hom}_{A}(M, I), I \bigr)  \]
in $\dcat{C}_{\mrm{str}}(A \ot B^{\mrm{op}})$. 
For this we use the results of Subsection \ref{subsec:flat-dg-rng}, that apply 
because $A$ and $B$ are flat over $\K$. 

In the special case when $B = A$ and $M = A$, the morphism 
$\opn{ev}^{\mrm{R, R}}_{A}$ recovers the derived 
homothety morphism through $A$, see (\ref{equ:4337}); i.e.\ 
\begin{equation} \label{eqn:4356}
\opn{hm}^{\mrm{R}}_{R, A} = \opn{ev}^{\mrm{R, R}}_{A} : 
A \to \opn{RHom}_{A^{\mrm{op}}}(R, R)
\end{equation}
in $\dcat{D}(A^{\mrm{en}})$. 
This formula makes implicit use of the left co-unitor isomorphism
$\opn{lcu} : \opn{RHom}_{A}(A, R) \iso R$ 
in $\dcat{D}(A^{\mrm{en}})$. 
Likewise for the homothety morphism $\opn{hm}^{\mrm{R}}_{R, A^{\mrm{op}}}$.

\begin{thm} \label{thm:3680}
Let $A$ be a noetherian ring, let $R$ be a NC dualizing complex over $A$, let 
$B$ be a second ring, and let $D_A$ and $D_{A^{\mrm{op}}}$ be the  
duality functors associated to $R$. Let $\star$ be a boundedness indicator, and 
let $M \in \dcat{D}^{\star}_{(\mrm{f}, ..)}(A \ot B^{\mrm{op}})$.
Then the following hold\tup{:}
\begin{enumerate}
\item The complex $D_A(M)$ belongs to 
$\dcat{D}^{-\star}_{(.., \mrm{f})}(B \ot A^{\mrm{op}})$, where $-\star$ 
is the reversed boundedness indicator. 

\item The derived Hom-evaluation morphism 
$\opn{ev}^{\mrm{R, R}}_{M} : M \to D_{A^{\mrm{op}}}(D_{A}(M))$
in $\dcat{D}(A \ot B^{\mrm{op}})$ is an isomorphism. 
\end{enumerate}
\end{thm}

\begin{proof}
The proof is the same as that of Theorem \ref{thm:3820}, with the obvious 
modifications (i.e.\ neglecting the algebraic grading). 
It all boils down to Theorems \ref{thm:2160} and \ref{thm:2135}.
\end{proof}

\begin{cor} \label{cor:3680}
Under the assumptions of Theorem \tup{\ref{thm:3680}}, 
the functor
\[ D_{A} : \dcat{D}^{\star}_{(\mrm{f}, ..)}(A \ot B^{\mrm{op}})^{\mrm{op}} \to 
\dcat{D}^{-\star}_{(.., \mrm{f})}(B \ot A^{\mrm{op}}) \]
is an equivalence of triangulated categories, with quasi-inverse
$D_{A^{\mrm{op}}}$. 
\end{cor}

\begin{proof}
First we note that the op-symmetry (Proposition \ref{prop:4355}) implies that  
Theorem \ref{thm:3680} is true after exchanging $A$ with $A^{\mrm{op}}$. 
Thus 
$\opn{ev}^{\mrm{R, R}}_{N} : N \to D_{A}(D_{A^{\mrm{op}}}(N))$
is an isomorphism for every 
$N \in \dcat{D}^{-\star}_{(.., \mrm{f})}(B \ot A^{\mrm{op}})$. 
Now the assertion is clear. 
\end{proof}

We need a better understanding of the derived homothety morphisms. 
Suppose $M \in \dcat{C}(A \ot B^{\mrm{op}})$. The right action of $B$ on $M$ is 
the homothety ring homomorphism through $B^{\mrm{op}}$,
$\opn{hm}_{M, B^{\mrm{op}}} : B^{\mrm{op}} \to
\opn{End}_{\dcat{C}_{\mrm{str}}(A)}(M)$.
Namely for elements $b \in B$ and $m \in M^i$ we have 
$\opn{hm}_{M, B^{\mrm{op}}}(\opn{op}(b))(m) = m \cd b \in M^i$.
By composing $\opn{hm}_{M, B^{\mrm{op}}}$ with the localization functor 
$\opn{Q} : \dcat{C}_{\mrm{str}}(A) \to \dcat{D}(A)$
we obtain the ring homomorphism 
\begin{equation} \label{eqn:4357}
\opn{hm}_{M, B^{\mrm{op}}}^{\dcat{D}} := 
\opn{Q} \circ \opn{hm}_{M, B^{\mrm{op}}} : 
B^{\mrm{op}} \to \opn{End}_{\dcat{D}(A)}(M) .
\end{equation}

\begin{lem} \label{lem:4361}
Let $M \in \dcat{D}(A \ot B^{\mrm{op}})$.
The three conditions below are equivalent. 
\begin{itemize}
\rmitem{i} $M$ has the derived Morita property on the $A$ side, i.e.\ 
\tup{(\ref{equ:4336})} is an isomorphism.

\rmitem{ii} For every $p \neq 0$ the module 
$\opn{Hom}_{\dcat{D}(A)}(M, M[p])$
is zero, and the ring homomorphism 
$\opn{hm}_{M, B^{\mrm{op}}}^{\dcat{D}} : B^{\mrm{op}} \to
\opn{End}_{\dcat{D}(A)}(M)$
is bijective. 

\rmitem{iii} For every $p \neq 0$ the module
$\opn{H}^p \bigl( \opn{RHom}_{A}(M, M) \bigr)$
is zero, and the $B^{\mrm{op}}$-module 
$\opn{H}^0 \bigl( \opn{RHom}_{A}(M, M) \bigr)$
is free, with basis the element 
$\opn{id}_M$. 
\end{itemize}
\end{lem}

\begin{proof}
This is clear from the canonical isomorphisms 
\[ \opn{Hom}_{\dcat{D}(A)} \bigl( M, M[p] \bigr) \cong 
\opn{H}^p \bigl( \opn{RHom}_{A}(M, M) \bigr) \]
of Corollary \ref{cor:2120}. 
\end{proof}

\begin{lem} \label{lem:4362}
Let $T \in \dcat{D}(A^{\mrm{en}})$, and consider the triangulated functor 
\[ G := T \ot^{\mrm{L}}_{A} (-) : 
\dcat{D}(A \ot B^{\mrm{op}}) \to \dcat{D}(A \ot B^{\mrm{op}}) . \]
For every $M \in \dcat{D}(A \ot B^{\mrm{op}})$  
the diagram 
\[ \UseTips \xymatrix @C=12ex @R=6ex {
B^{\mrm{op}}
\ar[r]^(0.4){ \opn{hm}_{M, B^{\mrm{op}}}^{\dcat{D}} }
\ar[dr]_(0.45){ \opn{hm}_{G(M), B^{\mrm{op}}}^{\dcat{D}} }
&
\opn{End}_{\dcat{D}(A)}(M)
\ar[d]^{G}
\\
&
\opn{End}_{\dcat{D}(A)}(G(M))
} \]
of ring homomorphisms is commutative. 
\end{lem}

\begin{proof}
This is because the action of $B^{\mrm{op}}$ on $M$ is categorical, i.e.\ by
endomorphisms in $\dcat{D}(A)$. 
\end{proof}

\begin{lem} \label{lem:4363}
Let $R \in \dcat{D}(A^{\mrm{en}})$, and consider the triangulated functor 
\[ D := \opn{RHom}_{A}(-, R) : 
\dcat{D}(A \ot B^{\mrm{op}})^{\mrm{op}} \to \dcat{D}(B \ot A^{\mrm{op}}) . \]
For every $M \in \dcat{D}(A \ot B^{\mrm{op}})$  
the diagram 
\[ \UseTips \xymatrix @C=12ex @R=6ex {
B^{\mrm{op}}
\ar[r]^(0.4){ \opn{hm}_{M, B^{\mrm{op}}}^{\dcat{D}} }
\ar[dr]_(0.45){ \opn{Op}(\opn{hm}_{D(M), B}^{\dcat{D}}) \quad }
&
\opn{End}_{\dcat{D}(A)}(M)
\ar[d]^{D}
\\
&
\opn{End}_{\dcat{D}(A^{\mrm{op}})}(D(M))^{\mrm{op}}
} \]
of ring homomorphisms is commutative. 
\end{lem}

\begin{proof}
Same as the previous lemma, only here $D$ is contravariant. 
\end{proof}

Tilting DG bimodules were defined in Definition \ref{dfn:3455}. Here we call 
them {\em tilting complexes}. 
Given a tilting complex $T \in \dcat{D}(A^{\mrm{en}})$, we know that its 
quasi-inverse $T^{\vee}$ satisfies
$T^{\vee} \cong \opn{RHom}_{A}(T, A) \cong  \opn{RHom}_{A^{\mrm{op}}}(T, A)$ 
in $\dcat{D}(A^{\mrm{en}})$. See Corollary \ref{cor:3455} with $B = A$,
and use the ring isomorphism (\ref{eqn:4450}). 

\begin{lem} \label{lem:4350}
Let $T$ be a tilting complex complex over $A$, with 
quasi-inverse $T^{\vee}$. Then there is an isomorphism 
$T \ot^{\mrm{L}}_{A} (-) \cong \opn{RHom}_{A}(T^{\vee}, -)$
of triangulated functors from 
$\dcat{D}(A \ot B^{\mrm{op}})$ to itself. 
\end{lem}

\begin{proof}
Let's write $S := T^{\vee}$. Then 
$T \cong S^{\vee} \cong \opn{RHom}_{A}(S, A)$. 
We know that $S$ is perfect on the $A$ side (see Corollary \ref{cor:3505}),
so by Theorem \ref{thm:3400} we have 
\begin{align*}
&
T \ot^{\mrm{L}}_{A} M \cong \opn{RHom}_{A}(S, A) \ot^{\mrm{L}}_{A} M 
\\ & \quad 
\cong \opn{RHom}_{A}(S, A \ot^{\mrm{L}}_{A} M) \cong 
\opn{RHom}_{A}(T^{\vee}, M)  . \qedhere
\end{align*} 
\end{proof}

\begin{thm}[\cite{Ye4}] \label{thm:3681}
Let $A$ be a noetherian ring, and let $R$ be a NC dualizing complex over $A$. 
\index{Dualizing complex! noncommutative}
\index{Tilting complex!}
\begin{enumerate}
\item Given a tilting complex $T$ over $A$, the complex
$R' := T \ot^{\mrm{L}}_{A} R$ 
is also a NC dualizing complex over $A$. 

\item Given a second NC dualizing complex $R'$ over $A$, the complex 
$T := \lb \opn{RHom}_A(R, R')$ 
is a tilting complex over $A$, and 
$R' \cong T \ot^{\mrm{L}}_{A} R$
in $\dcat{D}(A^{\mrm{en}})$. 

\item If $T$ is a tilting complex over $A$, and if
$T \ot^{\mrm{L}}_{A} R \cong R$ 
in $\dcat{D}(A^{\mrm{en}})$,  
then $T \cong A$ in $\dcat{D}(A^{\mrm{en}})$.
\end{enumerate}
\end{thm}

\begin{proof} \mbox{}

\smallskip \noindent
(1) We need to verify that $R'$ satisfies conditions (i)-(iii) of Definition 
\ref{dfn:3675}. 

Consider the equivalence of triangulated categories 
$G := T \ot^{\mrm{L}}_{A} (-) : \dcat{D}(A) \to \dcat{D}(A)$.
The functor $G$ has finite cohomological dimension (see Theorem \ref{thm:3400}, 
Corollary \ref{cor:3505} and Proposition \ref{prop:3415}). Now 
$G(A) = T$, and  
$T \in \dcat{D}^{\mrm{b}}_{\mrm{f}}(A)$
by Theorem \ref{thm:3415}, so Theorem \ref{thm:4600}
tells us that $G(M) \in \dcat{D}^{\mrm{b}}_{\mrm{f}}(A)$
for every $M \in \dcat{D}^{\mrm{b}}_{\mrm{f}}(A)$. 
Taking $M := R$ we see that 
$R' = G(R) \in \dcat{D}^{\mrm{b}}_{\mrm{f}}(A)$. 

Let $T^{\vee}$ be the quasi-inverse of $T$. By Lemma \ref{lem:4350}, with 
$B = A$, we see that 
\begin{equation} \label{eqn:4365}
R' = T \ot^{\mrm{L}}_{A} R \cong \opn{RHom}_{A}(T^{\vee}, R) = 
D_A(T^{\vee})
\end{equation}
in $\dcat{D}(A^{\mrm{en}})$. The tilting complex $T^{\vee}$ belongs to 
$\dcat{D}^{\mrm{b}}_{\mrm{f}}(A)$, so according to Theorem \ref{thm:3680}(1) 
with $B = \K$, we see that 
$D_A(T^{\vee}) \in \dcat{D}^{\mrm{b}}_{\mrm{f}}(A^{\mrm{op}})$. 
Thus $R' \in \dcat{D}^{\mrm{b}}_{\mrm{f}}(A^{\mrm{op}})$.
We have now verified condition (i) of Definition \ref{dfn:3675} for $R'$. 

To prove that $R'$ has finite injective dimension over $A$, we need to find a 
uniform bound on the cohomology of $\opn{RHom}_{A}(M, R')$, for all 
$M \in \dcat{M}(A)$. This is the same as finding a natural number $d$ such 
that $\opn{Hom}_{\dcat{D}(A)}(M, R'[i]) = 0$ for all 
$M \in \dcat{M}(A)$ and $\abs{i} > d$. 
We shall use the equivalence of triangulated categories 
\[ G^{\vee} := T^{\vee} \ot^{\mrm{L}}_{A} (-) \cong 
\opn{RHom}_{A}(T, -) : \dcat{D}(A) \to \dcat{D}(A) . \]
Since $R \cong G^{\vee}(R')$, we obtain a $\K$-module isomorphism
\[ G^{\vee} : \opn{Hom}_{\dcat{D}(A)} \bigl( M, R'[i] \bigr) \iso
\opn{Hom}_{\dcat{D}(A)} \bigl( G^{\vee}(M), R[i] \bigr)  \]
for every $i$. The functor $G^{\vee}$ has finite cohomological dimension, and 
$R$ has finite injective dimension; their sum $d$ serves as the required bound. 

Now let us prove that $R'$ has finite injective dimension over $A^{\mrm{op}}$.
Take a module $M \in \dcat{M}(A^{\mrm{op}})$. 
We have these isomorphisms for every integer $i$~:
\[ \begin{aligned}
& 
\opn{H}^i \bigl( \opn{RHom}_{A^{\mrm{op}}}(M, R') \bigr) \cong^{1}
\opn{Hom}_{\dcat{D}(A^{\mrm{op}})} \bigl( M, R'[i] \bigr)
\\ & \qquad 
\cong^{2} 
\opn{Hom}_{\dcat{D}(A^{\mrm{op}})} \bigl( M, D_A(T^{\vee})[i] \bigr)
\\ & \qquad 
\cong^{3} 
\opn{Hom}_{\dcat{D}(A)} \bigl( T^{\vee}, D_{A^{\mrm{op}}}(M)[i] \bigr)
\\ & \qquad 
\cong^{4} \opn{H}^i  \bigl(
\opn{RHom}_{A} \bigl( T^{\vee}, \opn{RHom}_{A^{\mrm{op}}}(M , R) \bigr)
\bigr)
\\ & \qquad
\cong^{5} \opn{H}^i  \bigl( 
T \ot^{\mrm{L}}_{A} \opn{RHom}_{A^{\mrm{op}}}(M, R) \bigr) 
\\ & \qquad
= \opn{H}^i  \bigl( (G \circ D_{A^{\mrm{op}}})(M) \bigr) . 
\end{aligned} \]
The justifications are as follows: 
\begin{itemize}
\item[$\cong^{1}$~:] This is by Corollary \ref{cor:2120}.

\item[$\cong^{2}$~:] This is by formula (\ref{eqn:4365}).

\item[$\cong^{3}$~:] This relies on Corollary \ref{cor:3680}.

\item[$\cong^{4}$~:] Again we use Corollary \ref{cor:2120}.

\item[$\cong^{5}$~:] It is due to Lemma \ref{lem:4350}.
\end{itemize}
The functors $G$ and $D_{A^{\mrm{op}}}$ have finite cohomological dimensions, 
so we have a uniform bound on the cohomology of 
$(G \circ D_{A^{\mrm{op}}})(M)$.
This completes the proof that $R'$ satisfies condition (ii) of Definition 
\ref{dfn:3675}. 

To prove that $R'$ has the derived Morita property on the $A$ side we 
use the fact that $R$ has the derived Morita property on the $A$ side, and we
invoke Lemma \ref{lem:4361} with $B := A$ and $M := R$. Since $R' = G(R)$, for 
every $i \in \Z$ we have an isomorphism of $\K$-modules 
\[ G : \opn{Hom}_{\dcat{D}(A)} \bigl( R, R[i] \bigr) \iso
\opn{Hom}_{\dcat{D}(A)} \bigl( R', R'[i] \bigr) , \]
and this is zero when $i \neq 0$. For $i = 0$ we use Lemma \ref{lem:4362}:
there is a commutative diagram of rings
\[ \UseTips \xymatrix @C=12ex @R=6ex {
A^{\mrm{op}}
\ar[r]^(0.4){\opn{hm}_{A^{\mrm{op}}, R}^{\dcat{D}}}_{\cong}
\ar[dr]_(0.45){\opn{hm}_{A^{\mrm{op}}, R'}^{\dcat{D}}}
&
\opn{End}_{\dcat{D}(A)}(R)
\ar[d]^{G}_{\cong}
\\
&
\opn{End}_{\dcat{D}(A)}(R')
} \]
and the vertical arrow $G$ is an isomorphism. 
Therefore $\opn{hm}_{A^{\mrm{op}}, R'}^{\dcat{D}}$ is a ring isomorphism. 

Finally, to prove that $R'$ has the derived Morita property on the 
$A^{\mrm{op}}$ side we use the fact that the tilting complex 
$T^{\vee}$ has the derived Morita property on the $A$ side 
(by Corollary \ref{cor:3505}), together with Lemma \ref{lem:4361}, applied to 
the complexes 
$T^{\vee} \in \dcat{M}(A \ot A^{\mrm{op}})$
and $R' \in \dcat{M}(A^{\mrm{op}} \ot A)$.
Recall from (\ref{eqn:4365}) that 
$R' \cong D_A(T^{\vee})$ in $\dcat{D}(A^{\mrm{en}})$.
By Corollary \ref{cor:3680}, for every $i \in \Z$ we 
have an isomorphism of $\K$-modules 
\[ D_A : \opn{Hom}_{\dcat{D}(A)} \bigl( T^{\vee}, T^{\vee}[i] \bigr) \iso
\opn{Hom}_{\dcat{D}(A^{\mrm{op}})} \bigl( R', R'[i] \bigr) , \]
and this is zero when $i \neq 0$. For $i = 0$ we use Lemma \ref{lem:4363}:
there is a commutative diagram of rings
\[ \UseTips \xymatrix @C=12ex @R=6ex {
A^{\mrm{op}}
\ar[r]^(0.4){\opn{hm}_{A^{\mrm{op}}, T^{\vee}}^{\dcat{D}}}_{\cong}
\ar[dr]_(0.45){ \opn{Op}(\opn{hm}_{A, R'}^{\dcat{D}}) \ }
&
\opn{End}_{\dcat{D}(A)}(T^{\vee})
\ar[d]^{D_A}_{\cong}
\\
&
\opn{End}_{\dcat{D}(A^{\mrm{op}})}(R')^{\mrm{op}}
} \]
It follows that 
$\opn{hm}_{A, R'}^{\dcat{D}} : A \to \opn{End}_{\dcat{D}(A^{\mrm{op}})}(R')$
is a ring isomorphism. 

\medskip \noindent 
(2) The proof of this item is a minor modification of the proof of Theorem 
\ref{thm:3780}. Indeed, an ungraded version of Lemma \ref{lem:3917} holds here, 
and we use Theorem \ref{thm:3680} instead of Theorem \ref{thm:3820}. 

\medskip \noindent 
(3) Assume that $T \ot^{\mrm{L}}_{A} R \cong R$ in $\dcat{D}(A^{\mrm{en}})$.
Of course $R = D_A(A)$. Formula (\ref{eqn:4365}) says that 
$D_A(A) \cong D_A(T^{\vee})$ in $\dcat{D}(A^{\mrm{en}})$.
By Corollary \ref{cor:3680} we conclude that 
$A \cong T^{\vee}$ in $\dcat{D}(A^{\mrm{en}})$.
But then $A \cong T$  in $\dcat{D}(A^{\mrm{en}})$.
\end{proof}

The op-symmetry that gave rise to Proposition \ref{prop:4355} implies that 
several variations of Theorem \ref{thm:3681} are true. We wish to write down 
only one of them, that will be needed later. 

\begin{cor} \label{cor:4365}
Let $A$ be a noetherian ring, and let $R$ and $R'$ be NC dualizing complexes 
over $A$. Then there is a tilting complex $T'$ such that 
$R' \cong R \ot^{\mrm{L}}_{A} T'$
in $\dcat{D}(A^{\mrm{en}})$. 
\end{cor}

\begin{proof}
This is a trivial (yet confusing) consequence of item (2) of  
Theorem \ref{thm:3681}, when viewing $R$ and $R'$ as NC dualizing complexes 
over the ring $A^{\mrm{op}}$. 
\end{proof}

Recall that a left action of a group $G$ on a nonempty set $X$ is called {\em 
simply transitive} if there is one $G$-orbit in $X$, and the stabilizer of 
every point $x \in X$ is trivial. I.e.\ for every (or equivalently, for some) 
point $x \in X$ the action function 
$G \times \{ x \} \to X$, $(g, x) \mapsto g(x)$, is bijective. 

\begin{cor} \label{cor:3690}
Let $A$ be a ring, and assume it has at least one NC dualizing complex. 
\index{1-DPic(A)@$\opn{DPic}_{\K}(A)$}
\index{Picard group! noncommutative derived}
Then the left action of the group $\opn{DPic}_{\K}(A)$ on the set of 
isomorphism classes of NC dualizing complexes, induced by the bifunctor
$(T, R) \mapsto T \ot^{\mrm{L}}_{A} R$,
is simply transitive.  
\end{cor}

\begin{proof}
By Theorem \ref{thm:3681}(1) this is a well-defined action:
$T \ot^{\mrm{L}}_{A} R$ is a NC dualizing complex. 
The action is transitive by Theorem \ref{thm:3681}(2), and it has trivial 
stabilizers by Theorem \ref{thm:3681}(3).
\end{proof}

\begin{rem} \label{rem:3696}
Suppose $A$ is a noetherian commutative $\K$-ring.
A noncommutative dualizing complex over $A$ is not the same as a commutative
dualizing complex over $A$, in the sense of Definition \ref{dfn:2155}. Indeed, 
there is a $\K$-ring 
homomorphism 
$A^{\mrm{en}} \to A$, and this induces a restriction functor 
$\dcat{D}(A) \to \dcat{D}(A^{\mrm{en}})$. 
If $R \in \dcat{D}(A)$ is a commutative dualizing complex over $A$, then
its image in $\dcat{D}(A^{\mrm{en}})$ is a noncommutative dualizing complex 
over $A$. However, according to Theorem \ref{thm:3681}(1), 
given any tilting complex $T$ over $A$, the complex 
$R' := T \ot^{\mrm{L}}_{A} R$
is also a NC dualizing complex over $A$. And $R'$ can very easily fail to be in 
the essential image of $\dcat{D}(A)$; e.g.\ by taking 
$T := A(\psi)$, the twist by an automorphism $\psi$ of the trivial bimodule 
$A$. 
\end{rem}

\begin{rem} \label{rem:3691}
In case the base ring $\K$ is not a field, but $A$ is a {\em flat noetherian 
$\K$-ring}, then the definitions and results in this subsection are 
still valid. 

In case the ring $A$ is noetherian, but is not flat over the base ring $\K$, we 
can still define NC dualizing complexes over it. For that we choose a 
K-flat DG ring resolution $\til{A} \to A$ over $\K$; see 
Theorem \ref{thm:4310} and Proposition \ref{prop:4305}. 
A NC dualizing complex over $A$ is then a complex
$R \in \dcat{D}(\til{A}^{\mrm{en}})$ 
that satisfies the three conditions of Definition \ref{dfn:3675}.
As explained in Remark \ref{rem:1420}, the triangulated category 
$\dcat{D}(\til{A}^{\mrm{en}})$ is independent of the resolution 
$\til{A}$, up to a canonical equivalence. This means that the dualizing complex 
$R$ is also independent of the resolution.

It is expected that Theorems \ref{thm:3680} and \ref{thm:3681} will still hold 
in this general situation; but we did not verify this assertion. 

However, flatness alone is not sufficient for many of the results 
in the following subsections (mainly in Subsection 
\ref{subsec:Fil-exit-R-NCDC}, in which we rely on results from Section 
\ref{sec:BDC}, that assumed Convention \ref{conv:3700}.) Therefore we decided 
to assume $\K$ is a field in the current subsection as well.
\end{rem}




\mysubsection{Rigid NC DC: Definition and Uniqueness}
\label{subsec:RNCDC-uniq}

In this subsection we introduce rigid NC dualizing complexes, and prove their 
uniqueness. The material is mostly from Van den Bergh's paper \cite{VdB}, with 
some improvements coming from \cite{Ye4}. 
We continue with Convention \ref{conv:3670}; in particular, $\K$ is a base 
field, and all rings are central over $\K$.

Let $A$ be a ring, and let 
$M, N \in \dcat{M}(A^{\mrm{en}})$.
The $\K$-module $M \ot N$ has two $A$-module structures and two 
$A^{\mrm{op}}$-module structures, all commuting with each other. 
We organize them into an {\em outside action} and an {\em inside 
action} of the ring $A^{\mrm{en}}$. Here are the formulas: given elements 
$m \in M$, $n \in N$ and $a_1, a_2 \in A$, we let
\begin{equation} \label{eqn:4215}
(a_1 \ot \opn{op}(a_2)) \, \cd^{\mrm{out}} \, (m \ot n) := 
(a_1 \cd m) \ot (n \cd a_2) 
\end{equation}
and
\begin{equation} \label{eqn:4216}
(a_1 \ot \opn{op}(a_2)) \, \cd^{\mrm{in}} \, (m \ot n) := 
(m \cd a_2) \ot (a_1 \cd n) . 
\end{equation}
In these formulas we are using the canonical anti-automorphism
$\opn{op} : A \to A^{\mrm{op}}$. 
We shall denote the two copies 
of $A^{\mrm{en}}$ that act on $M \ot N$ by
\begin{equation} \label{eqn:4460}
A^{\mrm{en, out}} = A^{\mrm{out}} \ot A^{\mrm{op, out}}
\end{equation}
and
\begin{equation} \label{eqn:4461}
A^{\mrm{en, in}} = A^{\mrm{in}} \ot A^{\mrm{op, in}} . 
\end{equation}
Namely the ring $A^{\mrm{en, out}}$ acts on $M \ot N$ by the outside action
(\ref{eqn:4215}), 
and the ring $A^{\mrm{en, in}}$ acts by the inside action (\ref{eqn:4216}). 

Let us define the ring 
\begin{equation} \label{eqn:4411}
A^{\mrm{four}} := A^{\mrm{en, out}} \ot A^{\mrm{en, in}}  .
\end{equation}
There are ring homomorphisms 
\begin{equation} \label{eqn:4500}
A^{\mrm{en, out}} \to A^{\mrm{four}} \lar \, A^{\mrm{en, in}}  .
\end{equation}
With this notation we have 
$M \ot N \in \dcat{M}(A^{\mrm{four}})$.
The action of the ring $A^{\mrm{four}}$ on the module $M \ot N$ is this, 
explicitly:  an element 
\[ \bigl( a_1 \ot \opn{op}(a_2) \bigr) \ot 
\bigl( a_3 \ot \opn{op}(a_4) \bigr) \in A^{\mrm{four}}\]
acts on an element 
$m \ot n \in M \ot N$ by
\[ \bigl( (a_1 \ot \opn{op}(a_2)) \ot (a_3 \ot \opn{op}(a_4)) \bigr) 
\cd (m \ot n) = 
(a_1 \cd m \cd a_4) \ot (a_3 \cd n \cd a_2) . \]

In some situations we prefer to view the inside action of $A^{\mrm{en}}$ on
$M \ot N$ as a right action. This is possible by the isomorphism 
$\opn{trop}$ from formula (\ref{eqn:4875}). 
To be explicit, for elements $m \ot n \in M \ot N$ and 
$a_3 \ot \opn{op}(a_4) \in A^{\mrm{en, in}}$
the action is as follows, using formula (\ref{eqn:4878}):
\begin{equation} \label{eqn:4877}
\begin{aligned}
&
(m \ot n) \, \cd^{\mrm{in}} \, (a_3 \ot \opn{op}(a_4)) = 
(a_4 \ot \opn{op}(a_3))\, \cd^{\mrm{in}} \, (m \ot n)
\\ & \quad 
= (m \cd a_3) \ot (a_4 \cd n) . 
\end{aligned}
\end{equation}

\begin{exa} \label{exa:5125}
Suppose we take $M = N := A \in \dcat{M}(A^{\mrm{en}})$.
Then $P := M \ot N = A \ot A$ can also be viewed as 
$P \cong A^{\mrm{en}}$, and then in it an 
$A^{\mrm{en}}$-bimodule. The left $A^{\mrm{en}}$-module structure on 
$P \cong A^{\mrm{en}}$ is this: for a module element $b_1 \ot b_2 \in P$ and a 
ring element $a_1 \ot \opn{op}(a_2) \in A^{\mrm{en}}$ we have 
\[ \bigl( a_1 \ot \opn{op}(a_2) \bigr) \cd (b_1 \ot b_2) = 
(a_1 \cd b_1) \ot (b_2 \cd a_2) ; \]
so this is the outside action on $P = A \ot A$. The right $A^{\mrm{en}}$-module 
structure on $P \cong A^{\mrm{en}}$ is this: for a ring element 
$a_3 \ot \opn{op}(a_4) \in A^{\mrm{en}}$ we have 
\[ (b_1 \ot b_2) \cd \bigl( a_3 \ot \opn{op}(a_4) \bigr) = 
(b_1 \cd a_3) \ot (a_4 \cd b_2) ; \]
so this is the inside action on $P = A \ot A$, cf.\ (\ref{eqn:4877}). 
\end{exa}

Given bimodules $L, M, N \in \dcat{M}(A^{\mrm{en}})$,
we denote by
$\opn{Hom}_{A^{\mrm{en, out}}}(L, M \ot N)$
the $\K$-module of homomorphisms $\phi : L \to M \ot N$ 
that are $A^{\mrm{en}}$-linear for the outside action of $A^{\mrm{en}}$ on 
$M \ot N$. This is an 
$A^{\mrm{en}}$-module by the inside action on $M \ot N$. 
Namely 
\[ \opn{Hom}_{A^{\mrm{en, out}}}(L, M \ot N)
\in \dcat{M}(A^{\mrm{en, in}}) = \dcat{M}(A^{\mrm{en}}) . \]

All these four-module operations extend in an obvious way to complexes, and 
can be derived. 

\begin{dfn} \label{dfn:4215}
Let $A$ be a ring and $M \in \dcat{D}(A^{\mrm{en}})$. 
The {\em NC square of $M$} 
\index{Squaring operation! noncommutative}
\index{1-SqA@$\opn{Sq}_{A}(M)$}
is the complex
\[ \opn{Sq}_{A}(M) := 
\opn{RHom}_{A^{\mrm{en, out}}}(A, M \ot M) \in 
\dcat{D}(A^{\mrm{en, in}}) = \dcat{D}(A^{\mrm{en}}) . \]
\end{dfn}

Of course the square is relative to the base field $\K$, which is implicit in 
the formulas. 

The square of the complex $M$ is calculated as follows. First note that $M$ is 
K-flat over $\K$, since $\K$ is a field; so 
$M \ot M = M \ot^{\mrm{L}}_{\K} M \in \dcat{D}(A^{\mrm{four}})$.
We choose a K-injective resolution 
$M \ot M \to I$ in $\dcat{C}_{\mrm{str}}(A^{\mrm{four}})$. 
The complex $I$ is unique up to a homotopy equivalence, that itself is unique 
up to homotopy (in other words, $I$ is unique up to a unique isomorphism in 
$\dcat{K}(A^{\mrm{four}})$). Because of the flatness of the ring homomorphisms 
(\ref{eqn:4500}), the complex 
$I$ is K-injective over $A^{\mrm{en, out}}$, and 
\begin{equation} \label{eqn:4370}
\opn{Sq}_{A}(M) = \opn{Hom}_{A^{\mrm{en, out}}}(A, I) \in 
\dcat{D}(A^{\mrm{en, in}}) = \dcat{D}(A^{\mrm{en}}) .
\end{equation}

\begin{dfn} \label{dfn:4370}
Let $\phi : M \to N$ be a morphism in $\dcat{D}(A^{\mrm{en}})$. 
The {\em NC square of $\phi$} 
\index{1-SqA@$\opn{Sq}_{A}(\phi)$}
is the morphism
\[ \opn{Sq}_{A}(\phi) := 
\opn{RHom}_{A^{\mrm{en, out}}}(\opn{id}_A, \phi \ot \phi) :
\opn{Sq}_{A}(M) \to \opn{Sq}_{A}(N) \]
in $\dcat{D}(A^{\mrm{en}})$.
\end{dfn}

The square of the morphism $\phi$ is calculated as follows. Choose a 
K-projective resolution $P \to M$ in $\dcat{C}_{\mrm{str}}(A^{\mrm{en}})$. 
Then $\phi$ is represented by a homomorphism 
$\til{\phi} : P \to N$ in 
$\dcat{C}_{\mrm{str}}(A^{\mrm{en}})$. The morphism 
$\phi \ot \phi : M \ot M \to N \ot N$ 
in $\dcat{D}(A^{\mrm{four}})$ is represented by the homomorphism 
$\til{\phi}  \ot \til{\phi} : P \ot P \to N \ot N$ 
in $\dcat{C}_{\mrm{str}}(A^{\mrm{four}})$.
Next we choose K-injective resolutions 
$P \ot P \to I$ and $N \ot N \to J$ in 
$\dcat{C}_{\mrm{str}}(A^{\mrm{four}})$. 
The homomorphism $\til{\phi}  \ot \til{\phi}$ induces a homomorphism 
$\til{\psi} : I \to J$ 
in $\dcat{C}_{\mrm{str}}(A^{\mrm{four}})$. Note that $\til{\psi}$ is unique up 
to homotopy. Finally
\begin{align*}
& \opn{Sq}_{A}(\phi) = 
\opn{Q} \bigl( \opn{Hom}_{A^{\mrm{en, out}}}(\opn{id}_A, \til{\psi}) \bigr) :
\\ & \quad 
\opn{Sq}_{A}(M) = \opn{Hom}_{A^{\mrm{en, out}}}(A, I)
\to \opn{Sq}_{A}(N) = \opn{Hom}_{A^{\mrm{en, out}}}(A, J) 
\end{align*}
% \[ \opn{Sq}_{A}(\phi) = 
% \opn{Q} \bigl( \opn{Hom}_{A^{\mrm{en, out}}}(\opn{id}_A, \til{\psi}) \bigr) :
% \opn{Sq}_{A}(M) \to \opn{Sq}_{A}(N)  \]
in $\dcat{D}(A^{\mrm{en, in}}) = \dcat{D}(A^{\mrm{en}}) $.

\begin{dfn} \label{dfn:4216}
Let $A$ be a ring. A {\em NC rigid complex over $A$} 
\index{Rigid complex! noncommutative}
is a pair $(M, \rho)$, where $M \in \dcat{D}(A^{\mrm{en}})$, and 
$\rho : M \iso \opn{Sq}_{A}(M)$
is an isomorphism in $\dcat{D}(A^{\mrm{en}})$, called a {\em NC rigidifying 
isomorphism}. 
\end{dfn}

\begin{dfn} \label{dfn:4372} 
Let $(M, \rho)$ and $(N, \si)$ be NC rigid complexes over $A$. A
{\em NC rigid morphism over $A$} 
\[ \phi : (M, \rho) \to (N, \si) \]
is a morphism $\phi : M \to N$ 
in $\dcat{D}(A^{\mrm{en}})$, such that the diagram 
\[ \UseTips \xymatrix @C=8ex @R=6ex {
M
\ar[r]^(0.35){\rho}
\ar[d]_{\phi}
&
\opn{Sq}_{A}(M)
\ar[d]^{\opn{Sq}_{A}(\phi)}
\\
N
\ar[r]^(0.35){\si}
&
\opn{Sq}_{A}(N)
} \]
in $\dcat{D}(A^{\mrm{en}})$ is commutative. 
\end{dfn}

\begin{dfn} \label{dfn:4380}
Let $A$ be a ring. The category of {\em NC rigid complexes over $A$} 
is the category whose objects are the NC rigid complexes from Definition 
\ref{dfn:4216}, and whose morphisms are the NC rigid morphisms  Definition 
\ref{dfn:4372}.
\end{dfn}

\begin{dfn}[Van den Bergh \cite{VdB}] \label{dfn:4375}
Let $A$ be a noetherian ring. A {\em rigid NC dualizing complex over $A$} 
\index{Rigid complex! noncommutative {\indash} dualizing {\indash}}
\index{Dualizing complex! noncommutative rigid}
is a NC rigid complex $(R, \rho)$, such that $R$ is a NC dualizing complex over 
$A$, in the sense of Definition \ref{dfn:3675}.
\end{dfn}

Again, we remind that the notion of rigidity is relative to the base field 
$\K$. 

Recall that the center of the ring $A$ is denoted by $\opn{Cent}(A)$. 
Of course $\opn{Cent}(A) = \opn{Cent}(A^{\mrm{op}})$. 
Given a bimodule $M \in \dcat{M}(A^{\mrm{en}})$, there are two ring 
homomorphisms
\[ \opn{chm}_{M, A}, \, \opn{chm}_{M, A^{\mrm{op}}} : 
\opn{Cent}(A) \to \opn{End}_{\dcat{M}(A^{\mrm{en}})}(M) , \]
that we call the 
{\em central homotheties}%
\index{Homothety morphism! central}
through $A$ and through 
$A^{\mrm{op}}$ respectively.  In terms of elements they are 
$\opn{chm}_{M, A}(a)(m) := a \cd m$ and
$\opn{chm}_{M, A^{\mrm{op}}}(a)(m) := m \cd a$ 
for $a \in \opn{Cent}(A)$ and $m \in M$. 

\begin{exa} \label{exa:4375}
For the bimodule $M = A$, both ring homomorphisms
\[ \opn{chm}_{A, A}, \, \opn{chm}_{A, A^{\mrm{op}}} : 
\opn{Cent}(A) \to \opn{End}_{\dcat{M}(A^{\mrm{en}})}(A) \]
are isomorphisms, and they are equal. 
\end{exa}

The ring homomorphisms $\opn{chm}_{M, A}$ and $\opn{chm}_{M, A^{\mrm{op}}}$ 
extend to complexes, namely for  $M \in \dcat{C}(A^{\mrm{en}})$ 
there are ring homomorphisms 
\[ \opn{chm}_{M, A}, \, \opn{chm}_{M, A^{\mrm{op}}} : 
\opn{Cent}(A) \to \opn{End}_{\dcat{C}_{\mrm{str}}(A^{\mrm{en}})}(M) . \]
By postcomposing them with the categorical localization functor $\opn{Q}$ we 
obtain ring homomorphisms 
\[ \opn{chm}^{\dcat{D}}_{M, A}, \, \opn{chm}^{\dcat{D}}_{M,  A^{\mrm{op}}} : 
\opn{Cent}(A) \to \opn{End}_{\dcat{D}(A^{\mrm{en}})}(M) . \]
These central homotheties were already used in Subsection 
\ref{subsec:tilting-rings}.

\begin{exa} \label{exa:4376}
For the complex of bimodules $M = A$, both ring homomorphisms
\[ \opn{chm}^{\dcat{D}}_{M, A}, \, \opn{chm}^{\dcat{D}}_{M,  A^{\mrm{op}}} : 
\opn{Cent}(A) \to \opn{End}_{\dcat{D}(A^{\mrm{en}})}(A) \]
are isomorphisms, and they are equal. This is because of the fully faithful 
embedding 
$\dcat{M}(A^{\mrm{en}}) \to \dcat{D}(A^{\mrm{en}})$ and the previous example. 
\end{exa}

\begin{lem} \label{lem:4380}
Given $R \in \dcat{D}(A^{\mrm{en}})$, consider the functor 
\[ D := \opn{Hom}_A(-, R) : \dcat{D}(A^{\mrm{en}})^{\mrm{op}} \to
\dcat{D}(A^{\mrm{en}}) . \]
Then for every $M \in \dcat{D}(A^{\mrm{en}})$ the diagram of rings 
\[ \UseTips \xymatrix @C=12ex @R=6ex {
\opn{Cent}(A)
\ar[r]^(0.4){ \opn{chm}^{\dcat{D}}_{M,  A^{\mrm{op}}}}_{}
\ar[dr]_(0.4){\opn{chm}^{\dcat{D}}_{D(M), A}}
&
\opn{End}_{\dcat{D}(A^{\mrm{en}})}(M)
\ar[d]^{D}_{}
\\
&
\opn{End}_{\dcat{D}(A^{\mrm{en}})} \bigl( D(M) \bigr)^{\mrm{op}}
} \]
is commutative. 
\end{lem}

\begin{proof}
Choose a K-injective resolution $R \to I$ in 
$\dcat{C}_{\mrm{str}}(A^{\mrm{en}})$, 
and consider the functor 
$\til{D} := \opn{Hom}_A(-, I)$. 
So $D \cong \til{D}$ as functors 
$\dcat{D}(A^{\mrm{en}})^{\mrm{op}} \to \dcat{D}(A^{\mrm{en}})$. 
Because of the contravariance of $\til{D}$, the diagram of rings 
\[ \UseTips \xymatrix @C=12ex @R=6ex {
\opn{Cent}(A)
\ar[r]^(0.4){ \opn{chm}_{M,  A^{\mrm{op}}}}_{}
\ar[dr]_(0.4){\opn{chm}_{\til{D}(M), A}}
&
\opn{End}_{\dcat{C}_{\mrm{str}}(A^{\mrm{en}})}(M)
\ar[d]^{\til{D}}_{}
\\
&
\opn{End}_{\dcat{C}_{\mrm{str}}(A^{\mrm{en}})} 
\bigl( \til{D}(M) \bigr)^{\mrm{op}}
} \]
is commutative. After applying the localization functor $\opn{Q}$ to the 
vertical arrow $\til{D}$ in the diagram above, we obtain the commutative 
diagram 
on the derived level. 
\end{proof}

\begin{lem} \label{lem:4381}
Given $M, N \in \dcat{D}(A^{\mrm{en}})$, consider the complex  
\[ L := \opn{RHom}_{A^{\mrm{en, out}}}(A, M \ot N) \in 
\dcat{D}(A^{\mrm{en, in}}) = \dcat{D}(A^{\mrm{en}}) . \]
Let $a, a', a'' \in \opn{Cent}(A)$ be elements such that 
$\opn{chm}^{\dcat{D}}_{N, A}(a) = \opn{chm}^{\dcat{D}}_{N, A^{\mrm{op}}}(a')$ 
in \lb $\opn{End}_{\dcat{D}(A^{\mrm{en}})}(N)$ 
and 
$\opn{chm}^{\dcat{D}}_{M, A}(a') = 
\opn{chm}^{\dcat{D}}_{M, A^{\mrm{op}}}(a'')$
in $\opn{End}_{\dcat{D}(A^{\mrm{en}})}(M)$.
Then \lb 
$\opn{chm}_{L, A}^{\dcat{D}}(a) = \opn{chm}^{\dcat{D}}_{L, A^{\mrm{op}}}(a'')$ 
in $\opn{End}_{\dcat{D}(A^{\mrm{en}})}(L)$.
\end{lem}

\begin{proof}
The proof resembles that of Lemma \ref{lem:3581}.
Consider the object 
$M \ot N\in \dcat{D}(A^{\mrm{four}})$. 
There are four ring homomorphisms 
\[ \opn{chm}^{\dcat{D}}_{M \ot N, B} : \opn{Cent}(A) \to 
\opn{End}_{\dcat{D}(A^{\mrm{four}})}(M \ot N) , \]
corresponding to these four options for the ring $B$,
in terms of formulas (\ref{eqn:4460}) and (\ref{eqn:4461}):
$B := A^{\mrm{out}}, A^{\mrm{op, out}}, A^{\mrm{in}} , A^{\mrm{op, in}}$. 
We know that 
\begin{equation} \label{eqn:4455}
\begin{aligned}
& \opn{chm}^{\dcat{D}}_{M \ot N, A^{\mrm{in}}}(a) = 
\opn{chm}^{\dcat{D}}_{M \ot N, A^{\mrm{op, out}}}(a') 
\\ & 
\opn{chm}^{\dcat{D}}_{M \ot N, A^{\mrm{out}}}(a') = 
\opn{chm}^{\dcat{D}}_{M \ot N, A^{\mrm{op, in}}}(a'') . 
\end{aligned}
\end{equation}

Let $M \ot N \to J$ be a K-injective resolution in 
$\dcat{C}_{\mrm{str}}(A^{\mrm{four}})$, 
so 
\[ \opn{End}_{\dcat{D}(A^{\mrm{four}})}(M \ot N) \cong 
\opn{H}^0 \bigl( \opn{End}_{\dcat{C}_{\mrm{str}}(A^{\mrm{four}})}(J) \bigr) \]
as rings. Formulas (\ref{eqn:4455}) imply that 
\begin{equation} \label{eqn:4456}
\begin{aligned}
& \opn{chm}_{J, A^{\mrm{in}}}(a) \sim  
\opn{chm}_{J, A^{\mrm{op, out}}}(a') , 
\\ & 
\opn{chm}_{J, A^{\mrm{out}}}(a') \sim
\opn{chm}_{J, A^{\mrm{op, in}}}(a'') , 
\end{aligned}
\end{equation}
where ``$\sim$'' means the homotopy relation on degree $0$ cocycles in the 
DG ring $\opn{End}_{A^{\mrm{four}}}(J)$. 

Now let us write 
$\til{L} = \opn{Hom}_{A^{\mrm{en, out}}}(A, J) \in 
\dcat{C}(A^{\mrm{en, in}})$.
The complex $\til{L}$ is K-injective over $A^{\mrm{en, in}}$, 
again by the flatness of the ring homomorphism (\ref{eqn:4500}), 
and $\til{L} \cong L$ in $\dcat{D}(A^{\mrm{en, in}})$. 
Hence there is a ring isomorphism 
\begin{equation} \label{eqn:4459}
\opn{H}^0 \bigl( \opn{End}_{A^{\mrm{en, in}}}(\til{L}) \bigr) \cong
\opn{End}_{\dcat{D}(A^{\mrm{en, in}})}(L) . 
\end{equation}
Now $A$ is a central bimodule over its center.
This implies that the two actions of $\opn{Cent}(A)$ on $\til{L}$ that 
come from $\opn{chm}_{J, A^{\mrm{out}}}$ and 
$\opn{chm}_{J, A^{\mrm{op, out}}}$ are equal. 
The two actions of $\opn{Cent}(A)$ on $\til{L}$ that 
come from $\opn{chm}_{J, A^{\mrm{in}}}$ and 
$\opn{chm}_{J, A^{\mrm{op, in}}}$
are the actions  $\opn{chm}_{\til{L}, A^{\mrm{in}}}$ and 
$\opn{chm}_{\til{L}, A^{\mrm{op, in}}}$
respectively. From formulas (\ref{eqn:4456}) we conclude that 
$\opn{chm}_{\til{L}, A^{\mrm{in}}}(a) \sim
\opn{chm}_{\til{L}, A^{\mrm{op, in}}}(a'')$. 
By formula (\ref{eqn:4459}) we see that 
$\opn{chm}_{L, A^{\mrm{in}}}^{\dcat{D}}(a) = 
\opn{chm}^{\dcat{D}}_{L, A^{\mrm{op, in}}}(a'')$
as claimed. 
\end{proof}

\begin{thm}[\cite{Ye4}] \label{thm:4375}
Let $A$ be a noetherian ring, and let $R$ be a NC dualizing complex over $A$.
\begin{enumerate}
\item The two ring homomorphisms
\[ \opn{chm}^{\dcat{D}}_{R, A}, \, \opn{chm}^{\dcat{D}}_{R, A^{\mrm{op}}} : 
\opn{Cent}(A) \to \opn{End}_{\dcat{D}(A^{\mrm{en}})}(R) \]
are bijective.

\item If $R$ is a rigid NC dualizing complex, then 
$\opn{chm}^{\dcat{D}}_{R, A} = \opn{chm}^{\dcat{D}}_{R, A^{\mrm{op}}}$.
\end{enumerate}
\end{thm}

\begin{proof} \mbox{}

\smallskip \noindent 
(1) Recall the duality functor $D_A$ from Definition \ref{dfn:3680}, with 
$B = A$. By Corollary \ref{cor:3680} it gives rise to an equivalence 
$D_{A} : \dcat{D}^{\mrm{b}}_{(\mrm{f}, ..)}(A^{\mrm{en}})^{\mrm{op}} 
\to \dcat{D}^{\mrm{b}}_{(.., \mrm{f})}(A^{\mrm{en}})$.
By Lemma \ref{lem:4380}, for every 
$M \in \dcat{D}^{\mrm{b}}_{(\mrm{f}, ..)}(A^{\mrm{en}})$
we get a commutative diagram of rings 
\[ \UseTips \xymatrix @C=12ex @R=6ex {
\opn{Cent}(A)
\ar[r]^(0.4){\opn{chm}^{\dcat{D}}_{M, A^{\mrm{op}}}}_{}
\ar[dr]_(0.4){\opn{chm}^{\dcat{D}}_{D_A(M), A}}
&
\opn{End}_{\dcat{D}(A^{\mrm{en}})}(M)
\ar[d]^{D_A}_{\cong}
\\
&
\opn{End}_{\dcat{D}(A^{\mrm{en}})} \bigl( D_A(M) \bigr)^{\mrm{op}}
} \]
Note that the vertical arrow $D_A$ is an isomorphism. 
Taking $M = R$ we have $D_A(R) \cong A$ in $\dcat{D}(A^{\mrm{en}})$. As we saw 
in Example \ref{exa:4376}, the homomorphism 
$\opn{chm}^{\dcat{D}}_{A, A}$ 
is bijective. This proves that 
$\opn{chm}^{\dcat{D}}_{R, A^{\mrm{op}}}$ is bijective. 

Likewise for $\opn{chm}^{\dcat{D}}_{R, A}$, but this time we use the 
duality functor $D_{A^{\mrm{op}}}$.

\medskip \noindent
(2) Now there is a rigidifying isomorphism 
$\rho : R \iso \opn{Sq}_{A}(R)$ in $\dcat{D}(A^{\mrm{en}})$. 
Let $f : \opn{Cent}(A) \to \opn{Cent}(A)$
be the ring automorphism 
$f := (\opn{chm}^{\dcat{D}}_{R, A^{\mrm{op}}})^{-1} \circ \, 
\opn{chm}^{\dcat{D}}_{R, A}$.
We need to prove that $f = \opn{id}$.

For this we shall use Lemma \ref{lem:4381}, with 
$M = N := R$, so $L \cong \opn{Sq}_A(R)$ in $\dcat{D}(A^{\mrm{en}})$. 
Take an element $a \in \opn{Cent}(A)$, and let 
$a' := f(a)$ and $a'' := f(a')$ in $\opn{Cent}(A)$. 
The lemma says that 
$\opn{chm}_{L, A}^{\dcat{D}}(a) = \opn{chm}^{\dcat{D}}_{L, A^{\mrm{op}}}(a'')$. 
But by the definition of $f$ and $a'$ we have  
$\opn{chm}_{R, A}^{\dcat{D}}(a) = \opn{chm}^{\dcat{D}}_{R, A^{\mrm{op}}}(a')$. 
Since $R \cong L$ in $\dcat{D}(A^{\mrm{en}})$, 
and since $\opn{chm}^{\dcat{D}}_{R, A^{\mrm{op}}}$
is bijective, we conclude that $a' = a''$. This means that $f(a) = f(f(a))$, 
so $a = f(a)$. Finally, the element $a$ was arbitrary, and hence
$f = \opn{id}$. 
\end{proof}


\begin{lem} \label{lem:4384}
Let $(R, \rho)$ and $(R', \rho')$ be rigid dualizing complexes, 
let $\phi : R \to R'$ be a morphism in $\dcat{D}(A^{\mrm{en}})$, and let
$a \in \opn{Cent}(A)$. Then 
\[ \opn{Sq}_{A} \bigl( \phi \circ \opn{chm}^{\dcat{D}}_{R, A}(a) \bigr) = 
\opn{Sq}_{A}(\phi) \circ \opn{chm}^{\dcat{D}}_{\opn{Sq}_{A}(R), A}(a^2) \]
as morphisms $\opn{Sq}_{A}(R) \to \opn{Sq}_{A}(R')$ in 
$\dcat{D}(A^{\mrm{en}})$.
\end{lem}

\begin{proof}
Because $\opn{Sq}_{A}$ is a functor, there is equality 
\[ \opn{Sq}_{A} \bigl( \phi \circ \opn{chm}^{\dcat{D}}_{R, A}(a) \bigr) =
\opn{Sq}_{A}(\phi) \circ 
\opn{Sq}_{A} \bigl( \opn{chm}^{\dcat{D}}_{R, A}(a) \bigr) .  \]
So we can assume that $R' = R$ and 
$\phi = \opn{id}_R$. We need to prove that the equality 
\begin{equation} \label{eqn:4501}
\opn{Sq}_{A} \bigl( \opn{chm}^{\dcat{D}}_{R, A}(a) \bigr) = 
\opn{chm}^{\dcat{D}}_{\opn{Sq}_{A}(R), A}(a^2) 
\end{equation}
holds for these endomorphisms of $\opn{Sq}_{A}(R)$ in $\dcat{D}(A^{\mrm{en}})$.

From Definition \ref{dfn:4370}, applied to the morphism 
$\opn{chm}^{\dcat{D}}_{R, A}(a) : R \to R$ 
in $\dcat{D}(A^{\mrm{en}})$, we see that 
\begin{equation} \label{eqn:4462}
\opn{Sq}_{A} \bigl( \opn{chm}^{\dcat{D}}_{R, A}(a) \bigr) = 
\opn{chm}^{\dcat{D}}_{\opn{Sq}_{A}(R), A^{\mrm{out}}}(a) \circ 
\opn{chm}^{\dcat{D}}_{\opn{Sq}_{A}(R), A^{\mrm{in}}}(a) 
\end{equation}
as endomorphisms of $\opn{Sq}_{A}(R)$ in 
$\dcat{D}(A^{\mrm{en}}) = \dcat{D}(A^{\mrm{en, in}})$.
But because $R$ is a rigid dualizing complex,  Theorem \ref{dfn:4370}(2) 
says that 
$\opn{chm}^{\dcat{D}}_{R, A}(a) = \opn{chm}^{\dcat{D}}_{R, A^{\mrm{op}}}(a)$.
Therefore, in formula (\ref{eqn:4462}) we can replace the action of $a$ 
through $A^{\mrm{out}}$ with its action through $A^{\mrm{op, in}}$. This gives 
\begin{equation} \label{eqn:4464}
\opn{Sq}_{A} \bigl( \opn{chm}^{\dcat{D}}_{R, A}(a) \bigr) = 
\opn{chm}^{\dcat{D}}_{\opn{Sq}_{A}(R), A^{\mrm{op, in}}}(a) \circ 
\opn{chm}^{\dcat{D}}_{\opn{Sq}_{A}(R), A^{\mrm{in}}}(a) . 
\end{equation}
Next, using Lemma \ref{lem:4381}, with $M = N := R$ and $a' = a'' := a$, 
we obtain 
\[ \opn{chm}^{\dcat{D}}_{\opn{Sq}_{A}(R), A^{\mrm{op, in}}}(a) =
\opn{chm}^{\dcat{D}}_{\opn{Sq}_{A}(R), A^{\mrm{in}}}(a) . \]
Plugging this into (\ref{eqn:4464}) we get 
\[ \opn{Sq}_{A} \bigl( \opn{chm}^{\dcat{D}}_{R, A}(a) \bigr) = 
\opn{chm}^{\dcat{D}}_{\opn{Sq}_{A}(R), A^{\mrm{in}}}(a)^2 . \]
Finally, since 
$\opn{chm}^{\dcat{D}}_{\opn{Sq}_{A}(R), A^{\mrm{in}}}$ 
is a ring homomorphism, we end up with (\ref{eqn:4501}). 
\end{proof}

In Definition \ref{dfn:4470} we introduced derived pseudo-finite DG modules
over a DG ring $A$. Recall that a DG module $M \in \dcat{D}(A)$ is called 
derived pseudo-finite if it belongs to the saturated full triangulated 
subcategory of $\dcat{D}(A)$ generated by the pseudo-finite semi-free DG 
modules. When $A$ is a ring, the pseudo-finite semi-free DG $A$-modules are 
precisely the bounded above complexes of finite free $A$-modules. 

\begin{thm}[Uniqueness, \cite{VdB}, \cite{Ye4}] \label{thm:4380} 
Let $A$ be a noetherian ring, and assume $A$ is a derived pseudo-finite complex 
over $A^{\mrm{en}}$. Suppose $(R, \rho)$ is a NC rigid dualizing complex over 
$A$. Then $(R, \rho)$ is unique, up to a unique NC rigid isomorphism.
\end{thm}

\begin{proof}
Suppose $(R', \rho')$ is another NC rigid dualizing complex over $A$.
According to Theorem \ref{thm:3681} and Corollary \ref{cor:4365}, there are 
tilting complexes $T$ and $T'$ over $A$ such that 
\begin{equation} \label{eqn:4380}
R' \cong T \ot^{\mrm{L}}_{A} R \cong R \ot^{\mrm{L}}_{A} T'
\end{equation}
in $\dcat{D}(A^{\mrm{en}})$. 
We have the following sequence of isomorphisms in $\dcat{D}(A^{\mrm{en}})$~: 
\begin{equation} \label{eqn:4381}
\begin{aligned}
&
T \ot^{\mrm{L}}_{A} R \cong^{1} R' 
\cong^2 \opn{RHom}_{A^{\mrm{en, out}}}(A, R' \ot R')
\\ & \qquad 
\cong^{1} \opn{RHom}_{A^{\mrm{en, out}}} \bigl( A, (R \ot^{\mrm{L}}_{A} T') \ot 
(T \ot^{\mrm{L}}_{A} R) \bigr)
\\ & \qquad 
\cong^3 \opn{RHom}_{A^{\mrm{en, out}}} \bigl( A, (R \ot R) 
\ot^{\mrm{L}}_{A^{\mrm{en}}} (T' \ot T) \bigr) 
\\ & \qquad 
\cong^4 \opn{RHom}_{A^{\mrm{en, out}}} ( A, R \ot R) 
\ot^{\mrm{L}}_{A^{\mrm{en}}} (T' \ot T)  
\\ & \qquad 
\cong^5 R \ot^{\mrm{L}}_{A^{\mrm{en}}} (T' \ot T) 
\cong^6 T \ot^{\mrm{L}}_{A} R \ot^{\mrm{L}}_{A} T' .
\end{aligned}
\end{equation}
Here are the explanations for the various isomorphisms:
\begin{itemize}
\item[$\cong^{1}$~:] These are due to (\ref{eqn:4380}). 

\item[$\cong^{2}$~:] This is the rigidifying isomorphism $\rho'$.

\item[$\cong^{3}$~:] This isomorphism is by rearranging the tensor factors, and 
in it the inside action of $A^{\mrm{en}}$ on $R \ot R$ is viewed as a right 
action, and this matches the left outside action of $A^{\mrm{en}}$ on 
$T' \ot T$. Cf.\ formulas (\ref{eqn:4877}) and (\ref{eqn:4215}).

\item[$\cong^{4}$~:] It is an application of the tensor-evaluation 
isomorphism (see Theorem \ref{thm:4320}). This is where we need $A$ to be a 
derived pseudo-finite complex over $A^{\mrm{en}}$ -- so that 
condition (i) of Theorem \ref{thm:4320} will hold. 
Condition (ii) that theorem holds since the complex $R \ot R$ 
has bounded cohomology. As for condition (iii) of the theorem: 
the complex $T' \ot T$ is a tilting complex over the ring $A^{\mrm{en}}$,
so it has finite flat dimension, and thus it certainly has derived bounded 
below tensor displacement. 

\item[$\cong^{5}$~:] This uses the rigidifying isomorphism $\rho$.

\item[$\cong^{6}$~:] It is a rearranging the tensor factors. Recall 
that the left action of $A^{\mrm{en}}$ on $T' \ot T$ is the outside action, 
cf.\ formula(\ref{eqn:4215}).
\end{itemize}

Let $T^{\vee}$ be the quasi-inverse of $T$. 
After applying $T^{\vee} \ot^{\mrm{L}}_{A} (-)$
to the isomorphisms (\ref{eqn:4381}) we get 
$R \cong R \ot^{\mrm{L}}_{A} T'$
in $\dcat{D}(A^{\mrm{en}})$. 
Equation (\ref{eqn:4380}) tells us that there is an isomorphism 
$\phi^{\dag} : R \iso R'$ in $\dcat{D}(A^{\mrm{en}})$. 

The isomorphism $\phi^{\dag}$ above need not be rigid. 
What we do know is that both morphisms 
$\rho' \circ \phi^{\dag}, \, \opn{Sq}_{A}(\phi^{\dag}) \circ \rho : 
R \to \opn{Sq}_{A}(R')$ 
are isomorphisms in $\dcat{D}(A^{\mrm{en}})$.
By Theorem \ref{thm:4375} the automorphisms of $R$ in $\dcat{D}(A^{\mrm{en}})$ 
are all of the form $\opn{chm}^{\dcat{D}}_{R, A}(a)$, 
for elements $a \in \opn{Cent}(A)^{\times}$.
Thus there is a unique invertible central element $a \in A$ such that  
\begin{equation} \label{eqn:4469}
\opn{Sq}_{A}(\phi^{\dag}) \circ \rho = 
\rho' \circ \phi^{\dag} \circ \opn{chm}^{\dcat{D}}_{R, A}(a) . 
\end{equation}
Define the isomorphism 
\[ \phi := \phi^{\dag} \circ \opn{chm}^{\dcat{D}}_{R, A}(a^{-1})
: R \iso R' \]
in $\dcat{D}(A^{\mrm{en}})$. 
Then we have the following equalities:
\[ \begin{aligned}
& 
\opn{Sq}_{A}(\phi) \circ \rho = 
\opn{Sq}_{A} \bigl( \phi^{\dag} \circ \opn{chm}^{\dcat{D}}_{R, A}(a^{-1}) 
\bigr) \circ \rho
\\ & \quad 
=^{\mrm{(i)}} \opn{Sq}_{A}(\phi^{\dag}) \circ 
\opn{chm}^{\dcat{D}}_{\opn{Sq}_{A}(R), A}(a^{-2}) \circ \rho
\\ & \quad 
=^{\mrm{(ii)}} \opn{Sq}_{A}(\phi^{\dag}) \circ \rho \circ 
\opn{chm}^{\dcat{D}}_{R, A}(a^{-2})
\\ & \quad 
=^{\mrm{(iii)}} \rho' \circ \phi^{\dag} \circ \opn{chm}^{\dcat{D}}_{R, A}(a)
\circ \opn{chm}^{\dcat{D}}_{R, A}(a^{-2})
\\ & \quad 
=^{\mrm{(iv)}} \rho' \circ \phi^{\dag} \circ \opn{chm}^{\dcat{D}}_{R, A}(a^{-1})
= \rho' \circ \phi . 
\end{aligned} \]
The equality $=^{\mrm{(i)}}$ is due to Lemma \ref{lem:4384}. In 
the equality $=^{\mrm{(ii)}}$ we have used the fact that $\rho$
is a $\opn{Cent}(A)$-linear morphism for the action through $A$, so $\rho$ 
commutes with $\opn{chm}^{\dcat{D}}_{R, A}(a^{-2})$.
Equality $=^{\mrm{(iii)}}$ is from (\ref{eqn:4469}). 
And equality $=^{\mrm{(iv)}}$ is because $\opn{chm}^{\dcat{D}}_{R, A}$ is a 
ring homomorphism. 
We see that $\phi : R \iso R'$ is a rigid isomorphism. The uniqueness of the 
element $a$ implies the uniqueness of the rigid isomorphism $\phi$ (by the 
same sort of calculation). 
\end{proof}

\begin{rem} \label{rem:4455}
As mentioned in Remark \ref{rem:3691}, for the purposes of Subsection 
\ref{subsec:NC-DC}, the only requirement on the base ring $\K$ is that the 
central $\K$-ring $A$ is flat over it. 

In this subsection that is not enough. Suppose that $A$ is flat over $\K$.
In order to prove the uniqueness of a 
rigid dualizing complex $R \in \dcat{D}^{\mrm{b}}(A^{\mrm{en}})$,
utilizing Theorem \ref{thm:4320}, we need that the complex 
$R \ot^{\mrm{L}}_{\K} R \in \dcat{D}(A^{\mrm{en}})$ 
will have bounded below cohomology. The only way we know to guarantee this is 
if the base ring $\K$ is {\em regular}. Notice that the regularity requirement 
is also needed in the commutative theory (see Setup \ref{set:3210} and Remark 
\ref{rem:4192}). 

Presumably the regularity of the base ring $\K$ is the only obstruction.
The lack of flatness of $A$ can be handled using a suitable DG ring resolution 
$\til{A} \to A$, as explained in Remark \ref{rem:3691}.
Thus we predict that the definitions and results of this subsection will remain 
valid if $\K$ is a regular nonzero commutative base ring, and $A$ is a 
noetherian central $\K$-ring. 
\end{rem}




\mysubsection{Interlude: Graded Rings of Laurent Type} 
\label{subsec:str-gr-rings}

Algebraically graded rings were introduced in Section 
\ref{sec:alg-gra-rings}, 
and we now return to them, and refer to them simply as graded rings. 
In this section we specialize to a particular kind of graded ring, that is 
needed for the next two subsections. 
The two main results here are Theorems \ref{thm:4510} and \ref{thm:4511}. 

We continue with Convention \ref{conv:3670}; in particular, $\K$ is a base 
field, and all rings are central over $\K$.

\begin{dfn} \label{dfn:4510}
A graded ring $\til{B}$ is called a {\em graded ring of Laurent type} 
\index{Algebraically graded ring! of Laurent type}
if there is an invertible central element $\til{c}$ in $\til{B}$ of degree $1$. 
Such an element $\til{c}$ is called a {\em uniformizer} of $\til{B}$. 
\end{dfn}

Of course if $\til{B}$ is a graded ring of Laurent type, then 
$\til{B}_i \cd \til{B}_j = \til{B}_{i + j}$
for all $i, j \in \Z$. A graded ring with this property is called a 
{\em strongly graded ring}. 

The next lemma describes the structure of graded rings of Laurent type. 
Let $\K[t, t^{-1}]$ be the ring of Laurent polynomials in the degree $1$ 
variable $t$. Given a ring $B$ (not graded), we define the graded ring of 
Laurent type 
$B[t, t^{-1}] := B \ot_{\K} \K[t, t^{-1}]$.
The uniformizer is $t$, and the degree $0$ component is $B$.  

\begin{lem} \label{lem:4510}
Let $\til{B}$ be a graded ring of Laurent type, with uniformizer $\til{c}$, and 
define $B := \til{B}_0$, the degree $0$ subring.
Then there is a unique isomorphism of graded rings 
$B[t, t^{-1}] \iso \til{B}$ 
that is the identity on $B$, and  sends $t \mapsto \til{c}$.
\end{lem}

\begin{exer} \label{exer:4511}
Prove Lemma \ref{lem:4510}.
\end{exer}

Note that for $\til{B}$ as in the lemma, there is also a surjective ring 
homomorphism $\til{B} \to B$ that sends $\til{c} \mapsto 1$. This is not a 
graded ring homomorphism. 

Recall that given a graded ring $\til{B}$, the category of graded 
$\til{B}$-modules is $\dcat{M}(\til{B}, \mrm{gr})$.
For a graded module $\til{M} \in \dcat{M}(\til{B}, \mrm{gr})$,
let us write 
$\opn{Deg}_0(\til{M}) := \til{M}_0$,
the homogeneous component of degree $0$ of $\til{M}$. 

\begin{lem} \label{lem:4390}
Assume $\til{B}$ is a graded ring of Laurent type, with $B := \til{B}_0$. 
\begin{enumerate}
\item The functor 
$\opn{Deg}_0 : \dcat{M}(\til{B}, \mrm{gr}) \to \dcat{M}(B)$
is an equivalence of abelian categories, with quasi-inverse 
$N \mapsto \til{B} \ot_{B} N \cong \K[t, t^{-1}] \ot N$.

\item There is a bifunctorial isomorphism 
\[ \opn{Deg}_0 \bigl( \opn{Hom}_{\til{B}}(\til{M} , \til{N}) \bigr) \iso 
\opn{Hom}_{B} \bigl( \opn{Deg}_0(\til{M}), \opn{Deg}_0(\til{N}) \bigr) \]
for $\til{M} , \til{N} \in \dcat{M}(\til{B}, \mrm{gr})$. It sends 
$\til{\phi} \mapsto \til{\phi}|_{\opn{Deg}_0(\til{M})}$.

\item If the graded ring $\til{B}$ is left noetherian, then the ring 
$B$ is left noetherian, and the functor 
$\opn{Deg}_0 : \dcat{M}_{\mrm{f}}(\til{B}, \mrm{gr}) \to 
\dcat{M}_{\mrm{f}}(B)$
is an equivalence  of abelian categories. 
\end{enumerate}
\end{lem}

Since $\opn{Deg}_0$ is exact, it extends to derived categories. 

\begin{lem} \label{lem:4400}
Assume $\til{B}$ is a graded ring of Laurent type, with $B := \til{B}_0$. 
\begin{enumerate}
\item For every boundedness condition $\star$, the functor 
$\opn{Deg}_0 : \dcat{D}^{\star}(\til{B}, \mrm{gr}) \to \dcat{D}^{\star}(B)$
is an equivalence of triangulated categories. 

\item If the graded ring $\til{B}$ is left noetherian, then 
$\opn{Deg}_0 : \dcat{D}^{\star}_{\mrm{f}}(\til{B}, \mrm{gr}) \to 
\dcat{D}^{\star}_{\mrm{f}}(B)$
is an equivalence. 
\end{enumerate}
\end{lem}

\begin{exer} \label{exer:4445}
Prove Lemmas \ref{lem:4390} and \ref{lem:4400}.
\end{exer}

\begin{lem} \label{lem:4401}
Let $\til{B}$ and $\til{C}$ be graded rings, with 
$B := \til{B}_0$ and $C := \til{C}_0$. Assume $\til{B}$ is 
a graded ring of Laurent type.
For $\til{M}, \til{N} \in \dcat{D}(\til{B} \ot \til{C}^{\mrm{op}}, \mrm{gr})$ 
there is an isomorphism 
\[ \opn{Deg}_0 \bigl( \opn{RHom}_{\til{B}}(\til{M}, \til{N}) \bigr) \iso  
\opn{RHom}_{B} \bigl( \opn{Deg}_0(\til{M}), 
\opn{Deg}_0(\til{N}) \bigr) \]
in $\dcat{D}(C^{\mrm{en}})$, which is functorial in these complexes.
\end{lem}

\begin{proof}
Choose a semi-graded-free resolution $\til{P} \to \til{M}$ over 
$\til{B} \ot \til{C}^{\mrm{op}}$. 
Because $\til{C}$ is a graded-free $\K$-module, it follows that the complex 
$\til{P}$ is semi-graded-free over $\til{B}$. 
Since $B = \opn{Deg}_0(\til{B})$, the complex $\opn{Deg}_0(\til{P})$ is 
semi-free over the ring $B$; so 
$\opn{Deg}_0(\til{P}) \to \opn{Deg}_0(\til{M})$
is a semi-free resolution over $B$. 
We get these isomorphisms: 
\[ \begin{aligned}
&
\opn{Deg}_0 \bigl( \opn{RHom}_{\til{B}}(\til{M}, \til{N}) \bigr) 
\cong 
\opn{Deg}_0 \bigl( \opn{Hom}_{\til{B}}(\til{P}, \til{N}) \bigr) 
\\ & \quad 
\cong^{\dag} \opn{Hom}_{B} \bigl( \opn{Deg}_0(\til{P}), \opn{Deg}_0(\til{N}) 
\bigr)
%\\ \quad &
\cong \opn{RHom}_{B} \bigl( \opn{Deg}_0(\til{M}), \opn{Deg}_0(\til{N}) \bigr) 
\end{aligned} \]
in $\dcat{D}(C^{\mrm{en}})$. 
The isomorphism $\cong^{\dag}$ comes from Lemma \ref{lem:4390}(2). 
\end{proof}

\begin{lem} \label{lem:4505}
Let $\til{B}$ be a graded ring, and let $\til{c}$ be a central homogeneous 
element of $\til{B}$ of positive degree. 
Then the central element $\til{c} - 1$ is regular in $\til{B}$.
\end{lem}

\begin{proof}
Take a nonzero element $\til{b} \in \til{B}$,
with homogeneous component decomposition 
$\til{b} = \til{b}_{1} + \cdots + \til{b}_{r}$, 
see (\ref{eqn:4260}). So either $r = 1$ and $\til{b} = \til{b}_{1}$
is homogeneous, or $r \geq 2$, and then
$\opn{deg}(\til{b}_{1}) < \opn{deg}(\til{b}_{i})$ for all 
$2 \leq i \leq r$. Then the homogeneous component decomposition
of $\til{b} \cd (\til{c} - 1)$ is 
\[ \til{b} \cd (\til{c} - 1) = - \til{b}_{1} + 
(\tup{higher degree terms}) . \]
The element $-\til{b}_{1}$ is nonzero, and therefore 
$\til{b} \cd (\til{c} - 1) \neq 0$. 
\end{proof}

Derived pseudo-finite complexes were introduced in Definition 
\ref{dfn:4470}.

\begin{lem} \label{lem:4406}
Let $A, B, C$ be rings; let 
$M_1, M_2 \in \dcat{D}(A)$; let $N_1 \in \dcat{D}(B)$; and let 
$N_2 \in \dcat{D}(B \ot C)$.
There is a morphism 
\[ \th : \opn{RHom}_{A}(M_1, M_2) \ot \opn{RHom}_{B}(N_1, N_2) \to 
\opn{RHom}_{A \ot B}(M_1 \ot N_1, M_2 \ot N_2) \]
in $\dcat{D}(C)$, which is functorial in these complexes. 

If $M_1$ is derived pseudo-finite over $A$, $N_1$ is derived pseudo-finite 
over $B$, and the complexes $M_2$ and $N_2$ have bounded below cohomologies, 
then $\th$ is an isomorphism. 
\end{lem}

\begin{proof}
Choose semi-free resolutions $P_1 \to M_1$ and $Q_1 \to N_1$, over $A$ and $B$ 
respectively. Then 
$P_1 \ot Q_1 \to M_1 \ot N_1$
is a semi-free resolution over $A\ot B$; see Proposition \ref{prop:4636}.
The morphism $\th$ is represented by the obvious homomorphism 
\[ \til{\th} : \opn{Hom}_{A}(P_1, M_2) \ot \opn{Hom}_{B}(Q_1, N_2) \to 
\opn{Hom}_{A \ot B}(P_1 \ot Q_1, M_2 \ot N_2) \]
in $\dcat{C}_{\mrm{str}}(C)$. 

Now assume that $\opn{H}(M_2)$ and $\opn{H}(N_2)$ are bounded below. By 
replacing  $M_2$ and $N_2$ with suitable smart truncations, we can assume these
are bounded below complexes. We fix them.

If $P_1$ and $Q_1$ are pseudo-finite semi-free, then $\til{\th}$ is an 
isomorphism in $\dcat{C}_{\mrm{str}}(C)$, by the usual calculation; 
see Lemma \ref{lem:4735}. 

If we fix $N_1$, then the complexes $M_1$ for which $\th$ is an isomorphism 
form a saturated full triangulated subcategory of $\dcat{D}(A)$. 
Likewise, if we fix $M_1$, then the complexes $N_1$ for which $\th$ is an 
isomorphism form a saturated full triangulated subcategory of 
$\dcat{D}(B)$. See Proposition \ref{prop:4603}. 
Therefore, as in the proof of Theorem \ref{thm:4320},
$\th$ is an isomorphism whenever $M_1$ is derived pseudo-finite over $A$ and 
$N_1$ is derived pseudo-finite over $B$.
\end{proof}

The inside and outside actions from Subsection \ref{subsec:RNCDC-uniq} make 
sense also in the graded setting. 
Here is the first main result of this subsection. 

\begin{thm} \label{thm:4510} 
Let $\til{B}$ be a graded ring of Laurent type, with $B := \til{B}_0$.
Assume that $B$ is a derived pseudo-finite complex over $B^{\mrm{en}}$. 
Given complexes 
$\til{M}, \til{N} \in \dcat{D}^{+}(\til{B}^{\mrm{en}},  \mrm{gr})$,
there is an isomorphism 
\[ \begin{aligned}
&
\opn{Deg}_0 \bigl( \opn{RHom}_{\til{B}^{\mrm{en, out}}}
(\til{B}, \til{M} \ot \til{N}) \bigr) 
\\ & \quad 
\iso \opn{RHom}_{B^{\mrm{en, out}}} \bigl( B, 
\opn{Deg}_0(\til{M}) \ot \opn{Deg}_0(\til{N}) \bigr)[-1]
\end{aligned} \]
in $\dcat{D}(B^{\mrm{en, in}}) = \dcat{D}(B^{\mrm{en}})$, 
which is functorial in these complexes.
\end{thm}

\begin{proof}
Say $\til{c}$ is a uniformizer of $\til{B}$. 
The ring $\til{B}^{\mrm{en}}  = \til{B} \ot \til{B}^{\mrm{op}}$ is actually
$\Z^2$-graded. Taking the degree $0$ component is for the total degree,
so the ring 
$\opn{Deg}_0(\til{B}^{\mrm{en}})$
retains a $\Z$-grading, that we shall call the {\em hidden grading}. For the 
hidden grading there is a graded ring isomorphism 
\begin{equation} \label{eqn:4440}
B^{\mrm{en}}[s, s^{-1}] = B^{\mrm{en}} \ot \K[s, s^{-1}] \iso 
\opn{Deg}_0(\til{B}^{\mrm{en}}) , 
\end{equation}
where $s$ is a variable of hidden degree $1$, and it goes to the element 
$\til{c} \ot \til{c}^{-1} \in \opn{Deg}_0(\til{B}^{\mrm{en}})$.

Similarly for complexes: given 
$\til{M}, \til{N} \in \dcat{D}(\til{B}^{\mrm{en}},  \mrm{gr})$,
there is a canonical isomorphism 
\begin{equation} \label{eqn:4441}
\opn{Deg}_0(\til{M}) \ot \opn{Deg}_0(\til{N}) \ot \K[s, s^{-1}] \iso 
\opn{Deg}_0(\til{M} \ot \til{N})
\end{equation}
in $\dcat{D} \bigl( B^{\mrm{four}}[s, s^{-1}] \bigr)$. 
We are neglecting the hidden grading from this stage onward in the proof. 

According to Lemma \ref{lem:4401} -- that applies because 
$\til{B}^{\mrm{en}}$ is a graded ring of Laurent type -- 
there is an isomorphism 
\begin{equation} \label{eqn:4442}
\opn{Deg}_0 \bigl( \opn{RHom}_{\til{B}^{\mrm{en, out}}}
(\til{B}, \til{M} \ot \til{N}) \bigr) \cong 
\opn{RHom}_{\opn{Deg}_0(\til{B}^{\mrm{en, out}})}
\bigl( B, \opn{Deg}_0(\til{M} \ot \til{N}) \bigr) 
\end{equation}
in $\dcat{D} \bigl( \opn{Deg}_0(\til{B}^{\mrm{en, in}}) \bigr)$.
After applying the restriction functor 
$\dcat{D} \bigl( \opn{Deg}_0(\til{B}^{\mrm{en, in}}) \bigr) \to 
\dcat{D}(B^{\mrm{en, in}}) =  \dcat{D}(B^{\mrm{en}})$,
(\ref{eqn:4442}) becomes an isomorphism in $\dcat{D}(B^{\mrm{en}})$. 
We then have the following isomorphisms in $\dcat{D}(B^{\mrm{en}})$~: 
\begin{equation} \label{eqn:4443}
\begin{aligned}
& 
\opn{RHom}_{\opn{Deg}_0(\til{B}^{\mrm{en, out}})}
\bigl( B, \opn{Deg}_0(\til{M} \ot \til{N}) \bigr)
\\ & \quad 
\cong^{\dag} \opn{RHom}_{B^{\mrm{en, out}} \ot \, \K[s, s^{-1}]}
\Bigl( B \ot \K, 
\bigl( \opn{Deg}_0(\til{M}) \ot \opn{Deg}_0(\til{N}) \bigr) \ot \K[s, s^{-1}] 
\Bigr)
\\ & \quad 
\cong^{\ddag} \opn{RHom}_{B^{\mrm{en, out}}} \bigl( B, 
\opn{Deg}_0(\til{M}) \ot \opn{Deg}_0(\til{N}) \bigr) \ot 
\opn{RHom}_{\K[s, s^{-1}]} \bigl( \K, \K[s, s^{-1}] \bigr) . 
\end{aligned}
\end{equation}
The isomorphism $\cong^{\dag}$ comes from formulas 
(\ref{eqn:4440}) and (\ref{eqn:4441}).
The isomorphism $\cong^{\ddag}$ is by Lemma \ref{lem:4406}; it 
justified because $B$ is a derived pseudo-finite complex over $B^{\mrm{en}}$, 
and of course $\K$ is a derived pseudo-finite complex over the noetherian ring 
$\K[s, s^{-1}]$.

The element $s - 1 \in \K[s, s^{-1}]$ is regular, so we 
have the Koszul resolution 
\[ 0 \to \K[s, s^{-1}] \xar{(s - 1) \cd (-)} \K[s, s^{-1}] \to \K \to 0 . \]
We use it to calculate
\[ \opn{RHom}_{\K[s, s^{-1}]} \bigl( \K, \K[s, s^{-1}] \bigr) \cong \K[-1] \]
in $\dcat{D}(\K)$. 
Plugging this into (\ref{eqn:4443}) finishes the proof. 
\end{proof}

From here until the end of this subsection we
consider a connected graded ring $\til{A}$, with a degree $1$
central element $\til{c} \in \til{A}$. Define the ring
$A := \til{A} / (\til{c} - 1)$,
and the canonical ring surjection $f : \til{A} \to A$.
Of course the ring $A$ is not graded. 

Let $\til{A}_{\til{c}}$ be the localization  of $\til{A}$ with respect to 
$\til{c}$, i.e.\ inverting the homogeneous multiplicatively closed set 
$\{ \til{c}^{\, i} \}_{i \in \N}$. 
Because the image of $\til{c}$ in $\til{A}_{\til{c}}$ is an invertible central 
element of degree $1$, the graded ring $\til{A}_{\til{c}}$ is a
graded ring of Laurent type, with uniformizer $\til{c}$. We let 
$\opn{inc} : \opn{Deg}_0(\til{A}_{\til{c}}) = (\til{A}_{\til{c}})_0 \to 
\til{A}_{\til{c}}$ 
be the inclusion of the degree $0$ component of $\til{A}_{\til{c}}$.
There is also the ring homomorphism 
$f_{\til{c}} : \til{A}_{\til{c}} \to A$, 
$f_{\til{c}}(\til{c}) = 1$. 

\begin{lem} \label{lem:4385}
\mbox{}
\begin{enumerate}
\item The ring homomorphism 
\[ g := f_{\til{c}} \circ \opn{inc} : 
\opn{Deg}_0(\til{A}_{\til{c}}) = (\til{A}_{\til{c}})_0 \to A \]
is an isomorphism. 

\item The ring isomorphism $g$ extends to a graded ring isomorphism 
$\til{g} : \til{A}_{\til{c}} \iso A[t, t^{-1}]$
that sends $\til{c} \mapsto t$. 
\end{enumerate}
\end{lem}

The situation is shown in this commutative diagram of rings: 
\[ \UseTips \xymatrix @C=8ex @R=6ex {
\til{A}
\ar@{->>}@(ur,ul)[rr]^{f}
\ar[r]
&
\til{A}_{\til{c}}
\ar@{->>}[r]^{f_{\til{c}}}
&
*+{A}
\\
&
*++{(\til{A}_{\til{c}})_0}
\ar@{>->}[u]^{\mrm{inc}}
\ar@{>->>}[ur]_{g}
} \]

\begin{proof}
In this proof we work with the graded ring of Laurent type 
$\til{B} := \til{A}_{\til{c}}$ and its uniformizer $\til{c}$.

\medskip \noindent 
(1) The homomorphism $f : \til{A} \to A$ is surjective, so every nonzero 
$a \in A$ can be written as 
$a = f(\til{b})$ for some nonzero $\til{b} \in \til{A}$. 
Consider the homogeneous component decomposition 
\begin{equation}  \label{eqn:4867}
\til{b} = \til{b}_{1} + \cdots + \til{b}_{r}
\end{equation}
of $\til{b}$, as in (\ref{eqn:4260}), with $r \geq 1$. Define 
$a' := \sum_{i = 1}^{r} \, \til{b}_i \cdot 
\til{c}^{\, - \opn{deg}(\til{b}_i)} \in (\til{A}_{\til{c}})_0$.
Then $a = g(a')$. We see that $g$ is surjective.

Next we look at 
\begin{equation}  \label{eqn:4868}
\opn{Ker}(g) = \opn{Ker}(f_{\til{c}}) \cap (\til{A}_{\til{c}})_0 
\sub (\til{A}_{\til{c}})_0 . 
\end{equation}
The ideal $\opn{Ker}(f_{\til{c}})$ is generated by the central element 
$\til{c} - 1$.
Consider a nonzero element $\til{b} \cd (\til{c} - 1) \in \til{A}_{\til{c}}$.
Let (\ref{eqn:4867}) be the homogeneous component decomposition of the nonzero 
element $\til{b}$. Then the homogeneous component decomposition
of $\til{b} \cd (\til{c} - 1)$ is 
$\til{b} \cd (\til{c} - 1) = -\til{b}_{1} + \cdots + \til{b}_{r} \cd \til{c}$.
This is not a homogeneous element, and hence it cannot lie inside 
$(\til{A}_{\til{c}})_0$. From formula (\ref{eqn:4868}) we conclude that 
$\opn{Ker}(g) = 0$. 
Thus $g$ is a ring isomorphism. 

\medskip \noindent 
(2) Let 
\begin{equation} \label{eqn:4476}
h : A \iso (\til{A}_{\til{c}})_0
\end{equation}
be the inverse of $g$. It extends to a graded ring homomorphism
$\til{h} : A[t, t^{-1}] \to  \til{A}_{\til{c}}$,
$t \mapsto \til{c}$. 
This is an isomorphism, because $\til{A}_{\til{c}}$ is a graded ring of Laurent 
type. The isomorphism $\til{g}$ is the inverse of $\til{h}$.
\end{proof}

\begin{exer} \label{exer:4475}
Show that the functor
$\opn{Ind}_f : \dcat{M}(\til{A}, \mrm{gr}) \to \dcat{M}(A)$,
$\opn{Ind}_f = A \ot_{\til{A}} (-)$, 
is exact.
\end{exer}

Here is the second main result of the subsection. 
Connected graded $\K$-rings were defined in Definition \ref{dfn:3703}. 

\begin{thm} \label{thm:4511}   
Let $\til{A}$ be a noetherian connected graded ring, let
$\til{c} \in \til{A}$ be a homogeneous central element of degree $1$, and let 
$A := \til{A} / (\til{c} - 1)$. Then\tup{:}
\begin{enumerate}
\item The ring $A$ is noetherian.

\item There is an isomorphism 
$A^{\mrm{en}} \ot^{\mrm{L}}_{\til{A}^{\mrm{en}}} \til{A} \cong A \oplus A[1]$
in $\dcat{D}(A^{\mrm{en}})$. 

\item The bimodule $A$ is a derived pseudo-finite complex over the ring 
$A^{\mrm{en}}$.
\end{enumerate}
\end{thm}

\begin{proof} \mbox{}

\smallskip \noindent
(1) By Theorem \ref{thm:4245} the ring $\opn{Ungr}(\til{A})$ is noetherian. 
Since there is a surjection of rings $\opn{Ungr}(\til{A}) \to A$, it follows 
that $A$ is noetherian. 

\medskip \noindent 
(2) We have these isomorphisms 
\[ \begin{aligned}
&
A^{\mrm{en}} \ot^{\mrm{L}}_{\til{A}^{\mrm{en}}} \til{A}
\cong^1 A \ot^{\mrm{L}}_{\til{A}} \til{A} \ot^{\mrm{L}}_{\til{A}} A
\\[0.3em] & \quad 
\cong^2 A \ot^{\mrm{L}}_{\til{A}} A
\cong^3 A \ot^{\mrm{L}}_{\til{A}} \til{A}_{\til{c}} 
\ot^{\mrm{L}}_{\til{A}_{\til{c}}} \! A 
\cong^4 A \ot^{\mrm{L}}_{\til{A}_{\til{c}}} \! A 
\end{aligned} \]
in $\dcat{D}(A^{\mrm{en}})$. 
The isomorphism $\cong^1$ is a rearrangement of the derived tensor factors. 
The isomorphism $\cong^2$ is from the left unitor isomorphism
$\til{A} \ot^{\mrm{L}}_{\til{A}} A \cong A$ in 
$\dcat{D}(\til{A} \ot A^{\mrm{op}})$. The isomorphism $\cong^3$
is because the ring homomorphism $\til{A} \to A$ factors through  
$\til{A}_{\til{c}}$, and 
$A \cong \til{A}_{\til{c}} \ot^{\mrm{L}}_{\til{A}_{\til{c}}} \! A$
in $\dcat{D}(\til{A} \ot A^{\mrm{op}})$. And, finally, the  
isomorphism $\cong^4$ is because 
$A \ot^{\mrm{L}}_{\til{A}} \til{A}_{\til{c}} \cong A$ in 
$\dcat{D}(\til{A} \ot \til{A}_{\til{c}}^{\mrm{op}})$

Now by Lemma \ref{lem:4505} 
we know that ${\til{c}} - 1$ is a regular central element in 
$\til{A}_{\til{c}}$. So there is a short exact sequence 
\begin{equation} \label{eqn:4510}
0 \to \til{A}_{\til{c}} \xar{\, (\til{c} - 1) \cd (-) \, } 
\til{A}_{\til{c}} \xar{ \, f_{\til{c}} \, } A \to  0
\end{equation}
in $\dcat{M}(\til{A}_{\til{c}} \ot A^{\mrm{op}})$. 
Here $\til{A}_{\til{c}}$ is an $A^{\mrm{op}}$-module via the ring homomorphism 
$h : A \to \til{A}_{\til{c}}$ from formula (\ref{eqn:4476}). We view the short 
exact sequence (\ref{eqn:4510}) as a quasi-isomorphism $\til{P} \to A$ in 
$\dcat{C}_{\mrm{str}}(\til{A}_{\til{c}} \ot A^{\mrm{op}})$,
where 
\[ \til{P} := \bigl( \cdots \to 0 \to 
\til{A}_{\til{c}} \xar{\, (\til{c} - 1) \cd (-) \, } 
\til{A}_{\til{c}} \to  0 \to \cdots \bigr) \]
is a complex concentrated in cohomological degrees $-1, 0$. 
Because $\til{P}$ is semi-free over $\til{A}_{\til{c}}$, this allows us to 
calculate: 
\[ A \ot^{\mrm{L}}_{\til{A}_{\til{c}}} \! A \cong A 
\ot_{\til{A}_{\til{c}}} \! \til{P}  \cong 
\bigl( \cdots \to 0 \to A \xar{ \,0 \, } A \to 0 \to \cdots \bigr) \cong A[1] 
\oplus A \]
in $\dcat{D}(A^{\mrm{en}})$.

\medskip \noindent 
(3) By Proposition \ref{prop:4266} there is a quasi-isomorphism 
$\til{Q} \to \til{A}$ in 
$\dcat{C}_{\mrm{str}}(\til{A}^{\mrm{en}}, \mrm{gr})$
from a nonpositive complex $\til{Q}$ of finite graded-free 
$\til{A}^{\mrm{en}}$-modules. We can forget the grading now, and just view 
$\til{Q} \to \til{A}$ as a free resolution of the module $\til{A}$ over the 
ring $\til{A}^{\mrm{en}}$. We get 
$P := A^{\mrm{en}} \ot_{\til{A}^{\mrm{en}}} \til{Q}  \cong 
A^{\mrm{en}} \ot^{\mrm{L}}_{\til{A}^{\mrm{en}}} \til{A}$ 
in $\dcat{D}(A^{\mrm{en}})$. The complex $P$ is pseudo-finite semi-free over 
$A^{\mrm{en}}$.
By item (2) the complex $A$ is a direct summand of $P$
in $\dcat{D}(A^{\mrm{en}})$, and hence it is a derived 
pseudo-finite complex over $A^{\mrm{en}}$.
\end{proof}




\mysubsection{Graded Rigid NC DC}
\label{subsec:GrR-NC-DC}

Here we make a bridge between {\em balanced dualizing complexes} and 
{\em rigid dualizing complexes}. This is done by introducing an intermediate 
kind of object: the {\em graded rigid dualizing complex}. All is in the 
noncommutative setting of course. 

The results of this subsection were originally in \cite{VdB}, with a few 
improvements later in \cite{YeZh1}. Here we give a much more detailed 
discussion, and some corrections. 

We continue with Convention \ref{conv:3670}. 
In particular $\K$ is a base field, and all rings are by default central over 
$\K$. Throughout this subsection we assume the following setup: 

\begin{setup} \label{set:4510}
We are given a noetherian connected graded central $\K$-ring $\til{A}$ and
a central homogeneous element $t \in \til{A}$ of degree $1$. 
We define the central $\K$-ring  
$A := \til{A} / (t - 1)$.
\end{setup}

Note that the ring $A$ is not graded. Also the rings
$\til{A}^{\mrm{en}}$ and $A^{\mrm{en}}$ need not be 
noetherian. However: 

\begin{prop} \label{prop:4515}
Under Setup \tup{\ref{set:4510}}, the ring $A$ is noetherian, and $A$ is a 
derived pseudo-finite complex over $A^{\mrm{en}}$.
\end{prop}

\begin{proof} 
Use Theorem \ref{thm:4511}. 
\end{proof}

The concepts of outside and inside actions from the beginning of Subsection 
\ref{subsec:RNCDC-uniq} make sense in the graded setting. Thus we have a 
connected graded ring 
$\til{A}^{\mrm{four}} = \til{A}^{\mrm{en, out}} \ot \til{A}^{\mrm{en, in}}$. 
For a complex
$\til{M} \in \dcat{D}(\til{A}^{\mrm{en}}, \mrm{gr})$,
its tensor product with itself $\til{M} \ot \til{M}$
is an object of 
$\dcat{D}(\til{A}^{\mrm{four}}, \mrm{gr})$,
and 
\begin{equation} \label{eqn:4860}
\opn{RHom}_{\til{A}^{\mrm{en, out}}} ( \til{A}, \til{M} \ot \til{M} )
\in \dcat{D}(\til{A}^{\mrm{en, in}}, \mrm{gr}) = 
\dcat{D}(\til{A}^{\mrm{en}}, \mrm{gr}) . 
\end{equation}

Recall that 
$\dcat{D}_{(\mrm{f},..)}(\til{A}^{\mrm{en}}, \mrm{gr})$
(resp.\ $\dcat{D}_{(.., \mrm{f})}(\til{A}^{\mrm{en}}, \mrm{gr})$)
is the full subcategory of \lb 
$\dcat{D}(\til{A}^{\mrm{en}}, \mrm{gr})$
on the complexes whose cohomology bimodules are finite over $\til{A}$
(resp.\  $\til{A}^{\mrm{op}}$), and 
\[ \dcat{D}_{(\mrm{f}, \mrm{f})}(\til{A}^{\mrm{en}}, \mrm{gr}) = 
\dcat{D}_{(\mrm{f},..)}(\til{A}^{\mrm{en}}, \mrm{gr}) \cap 
\dcat{D}_{(.., \mrm{f})}(\til{A}^{\mrm{en}}, \mrm{gr}) . \]

\begin{dfn} \label{dfn:4391} 
Under Setup \tup{\ref{set:4510}}, a 
{\em graded rigid NC dualizing complex}%
\index{Dualizing complex! graded rigid NC} 
over $\til{A}$ is a pair $(\til{R}, \til{\rho})$, where 
$\til{R} \in \dcat{D}^{\mrm{b}}_{(\mrm{f}, \mrm{f})}
(\til{A}^{\mrm{en}}, \mrm{gr})$
is a graded NC dualizing complex over $\til{A}$, in the sense of Definition 
\ref{dfn:3716}; and 
\[ \til{\rho} : \til{R} \iso 
\opn{RHom}_{\til{A}^{\mrm{en, out}}} ( \til{A}, \til{R} \ot \til{R} ) \]
is an isomorphism in $\dcat{D}(\til{A}^{\mrm{en}}, \mrm{gr})$,
see (\ref{eqn:4860}) with $\til{M} := \til{R}$.  
\end{dfn}

Balanced dualizing complexes were introduced in Definition \ref{dfn:3714}.

\begin{thm} \label{thm:4400}
\index{Dualizing complex! graded rigid NC}
\index{Dualizing complex! balanced}
Under Setup \tup{\ref{set:4510}}, if $\til{R}$ is a balanced NC dualizing 
complex over $\til{A}$, then $\til{R}$ is a graded rigid NC dualizing complex 
over $\til{A}$.
\end{thm}

We are a bit sloppy in stating the theorem; the proper way to state it is this:
if $(\til{R}, \til{\be})$ is a balanced dualizing complex over $\til{A}$, then 
there exists an isomorphism $\til{\rho}$ such that $(\til{R}, \til{\rho})$ is a 
graded rigid NC dualizing complex over $\til{A}$.

We need a lemma first. 
Recall (from Definition \ref{dfn:3723}) that a graded $\til{A}$-module 
$\til{M} = \bigoplus_i \til{M}_i$ is called degreewise finite over $\K$ if each 
homogeneous component $\til{M}_i$ is a finite $\K$-module. We denoted by 
$\dcat{M}_{\mrm{dwf}}(\til{A}, \mrm{gr})$
the full subcategory of $\dcat{M}(\til{A}, \mrm{gr})$ on the degreewise finite
graded modules. This is a thick abelian subcategory of 
$\dcat{M}(\til{A}, \mrm{gr})$, closed under subobjects and quotients.
Then we denoted by 
$\dcat{D}_{\mrm{dwf}}(\til{A}, \mrm{gr})$
the full subcategory of $\dcat{D}(\til{A}, \mrm{gr})$
on the complexes $\til{M}$ such that 
$\opn{H}^q(\til{M}) \in \dcat{M}_{\mrm{dwf}}(\til{A}, \mrm{gr})$
for all $q$. This is a full triangulated subcategory.

\begin{lem} \label{lem:4402}
Let 
$\til{M}, \til{N} \in 
\dcat{D}^{\mrm{b}}_{(\mrm{f}, ..)}(\til{A}^{\mrm{en}}, \mrm{gr})$.
Then\tup{:}
\begin{enumerate}
\item The complexes $\til{M}$ and $\til{M} \ot \til{N}$ satisfy 
$\til{M} \in \dcat{D}^{\mrm{b}}_{\mrm{dwf}}(\til{A}^{\mrm{en}}, \mrm{gr})$
and 
$\til{M} \ot \til{N} \in 
\dcat{D}^{\mrm{b}}_{\mrm{dwf}}(\til{A}^{\mrm{four}}, \mrm{gr})$. 

\item The canonical morphism 
$\th : \til{M}^* \ot \til{N}^* \to (\til{M} \ot \til{N})^*$ 
in $\dcat{D}^{\mrm{b}}_{\mrm{dwf}}(\til{A}^{\mrm{four}}, \mrm{gr})$
is an isomorphism. 
\end{enumerate}
\end{lem}

\begin{proof} \mbox{}

\smallskip \noindent
(1) For every $q$ the $\til{A}$-module $\opn{H}^q(\til{M})$ is finite, so
it is the image of a finite direct sum of algebraic degree shifts of $\til{A}$. 
Since $\til{A} \in \dcat{M}_{\mrm{dwf}}(\K, \mrm{gr})$,
and since the subcategory
$\dcat{M}_{\mrm{dwf}}(\K, \mrm{gr})$
is closed under quotients inside $\dcat{M}(\K, \mrm{gr})$, we see that 
$\opn{H}^q(\til{M}) \in \dcat{M}_{\mrm{dwf}}(\K, \mrm{gr})$. 
We conclude that 
\[ \opn{H}^q(\til{M}) \in \dcat{M}_{\mrm{dwf}}(\K, \mrm{gr}) \cap 
\dcat{M}(\til{A}^{\mrm{en}}, \mrm{gr}) 
= \dcat{M}_{\mrm{dwf}}(\til{A}^{\mrm{en}}, \mrm{gr}) . \]

Let $[q_0, q_1]$ be a finite integer interval that contains 
$\opn{con}(\opn{H}(\til{M}))$. 
For every $p \in \Z$ there is an isomorphism 
\[ \opn{H}^p(\til{M} \ot \til{N}) \cong 
\bigoplus_{q \in [q_0, q_1]} \, 
\bigl( \opn{H}^q(\til{M}) \ot \opn{H}^{p - q}(\til{N}) \bigr) \]
in $\dcat{M}(\K, \mrm{gr})$. 
Each $\opn{H}^q(\til{M})$ and $\opn{H}^{p - q}(\til{N})$
is a finite graded $\til{A}$-module, so it is degreewise finite and bounded 
below (in algebraic degree). Therefore 
\[ \opn{H}^q(\til{M}) \ot \opn{H}^{p - q}(\til{N}) \in 
\dcat{M}_{\mrm{dwf}}(\K, \mrm{gr}) , \]
and hence so is $\opn{H}^p(\til{M} \ot \til{N})$. 
This means that 
$\til{M} \ot \til{N} \in 
\dcat{D}^{\mrm{b}}_{\mrm{dwf}}(\til{A}^{\mrm{four}}, \mrm{gr})$. 

\medskip \noindent 
(2) We start by giving an explicit formula for the canonical morphism $\th$. 
There is a canonical homomorphism 
\begin{equation} \label{eqn:4485}
\til{\th} : \til{M}^* \ot \til{N}^* \to (\til{M} \ot \til{N})^*
\end{equation}
in $\dcat{C}_{\mrm{str}}(A^{\mrm{four}}, \mrm{gr})$; 
it is 
\[ \til{\th}(\phi \ot \psi)(m \ot n) := 
(-1)^{j \cd k} \cd \phi(m) \ot \psi(n) \in \K \]
for $\phi \in (\til{M}^*)^i$, $\psi \in (\til{N}^*)^j$, 
$m \in \til{M}^k$ and $n \in \til{N}^l$.  
Then $\th = \opn{Q}(\til{\th})$ in 
$\dcat{D}(A^{\mrm{four}}, \mrm{gr})$.
We need to show that under the given finiteness and boundedness conditions, 
$\th$ is an isomorphism. For that we can forget the $A^{\mrm{four}}$-module 
structures, and view $\th$ as a morphism in 
$\dcat{D}(\K, \mrm{gr})$.

But in $\dcat{D}(\K, \mrm{gr})$ there are canonical isomorphisms 
$\til{M} \cong \opn{H}(\til{M})$ and 
$\til{N} \cong \opn{H}(\til{N})$.
This means that we can assume that the differentials of the complexes $\til{M}$ 
and $\til{N}$ are zero. We know that for every
$\sbmat{p \\[0.1em] i} \in \sbmat{\Z \\[0.1em] \Z} = \Z^2$,
in the notation from Subsection \ref{subsec:alg-gr-mods},
the $\K$-modules $\til{M}^p_i$ and  $\til{N}^p_i$ are finite. 
Also there are uniform bounds on vanishing: there are integers 
$p_0, p_1, i_0$ such that 
$\til{M}^p_i = \til{N}^p_i = 0$
unless $p_0 \leq p \leq p_1$ and $i_0 \leq i$. 
An easy calculation shows that for these objects $\til{M}$ and $\til{N}$, the 
homomorphism $\til{\th}$ of (\ref{eqn:4485}) is bijective. 
\end{proof}

\begin{proof}[Proof of Theorem \tup{\ref{thm:4400}}]
We need to produce a graded rigidifying isomorphism $\til{\rho}$ for the graded 
NC dualizing complex $\til{R}$. 

Consider the complex 
$\til{R} \in \dcat{D}^{\mrm{b}}_{(\mrm{f}, \mrm{f})}(\til{A}^{\mrm{en}}, 
\mrm{gr})$.
According to Lemma \ref{lem:4402}(1) we know that 
\begin{equation} \label{eqn:4410}
\til{R} \ot \til{R} \in 
\dcat{D}^{\mrm{b}}_{\mrm{dwf}}(\til{A}^{\mrm{four}}, \mrm{gr})  
\end{equation}
and 
\begin{equation} \label{eqn:4412}
\til{A} \in \dcat{D}^{\mrm{b}}_{\mrm{dwf}}(\til{A}^{\mrm{en}}, \mrm{gr}) . 
\end{equation}
There is a sequence of isomorphisms in 
$\dcat{D}(\til{A}^{\mrm{en}}, \mrm{gr})$~:
\begin{equation} \label{eqn:4413}
\begin{aligned}
& \opn{RHom}_{\til{A}^{\mrm{en, out}}}(\til{A}, \til{R} \ot \til{R})
\cong^{1}
\opn{RHom}_{\til{A}^{\mrm{en, out}}} \bigl( (\til{R} \ot \til{R})^*, \til{A}^* 
\bigr)
\\
& \qquad \cong^{2}
\opn{RHom}_{\til{A}^{\mrm{en, in}}}(\til{R}^* \ot \til{R}^*, \til{A}^*)
\\
& \qquad \cong^{3}
\opn{RHom}_{\til{A}^{\mrm{en, in}}}(\til{P} \ot \til{P}, \til{A}^*) 
\\
& \qquad \cong^{4}
\opn{RHom}_{\til{A}} \bigl( \til{P} , 
\opn{RHom}_{\til{A}^{\mrm{op}}} (\til{P}, \til{A}^*) \bigr) 
\\
& \qquad \cong^{5}
\opn{RHom}_{\til{A}} (\til{P}, \til{R} ) \cong^{6} \til{R} . 
\end{aligned}
\end{equation}
The explanations for the isomorphisms are:
\begin{itemize}
\item[$\cong^{1}$~:] By Theorem \ref{thm:3800}, that applies by formulas 
(\ref{eqn:4410}) and (\ref{eqn:4412}).  

\item[$\cong^{2}$~:] This is by Lemma \ref{lem:4402}(2). Note that the outside 
action changes to an inside action. 

\item[$\cong^{3}$~:] Here 
$\til{P} := \mrm{R} \Ga_{\til{\m}}(\til{A})$, where $\til{\m}$ is the 
augmentation ideal of $\til{A}$. According to Theorems \ref{thm:3713} and 
\ref{thm:3710} there is an isomorphism 
$\til{R} \cong \til{P}^*$
in $\dcat{D}(\til{A}^{\mrm{en}}, \mrm{gr})$; dualizing we get 
an isomorphism $\til{R}^* \cong \til{P}$ in 
$\dcat{D}(\til{A}^{\mrm{en}}, \mrm{gr})$. 

\item[$\cong^{4}$~:] This is one of the standard Hom-tensor identities. 

\item[$\cong^{5}$~:] This is another instance of the isomorphism 
$\til{P}^* \cong \til{R}$. 

\item[$\cong^{6}$~:] By Lemma \ref{lem:4060}, the restriction of 
$\til{R}$ to $\dcat{D}(\til{A}, \mrm{gr})$ satisfies  
$\til{R} \in \dcat{D}_{\mrm{f}}(\til{A}, \mrm{gr}) \lb 
\sub \dcat{D}(\til{A}, \mrm{gr})_{\mrm{com}}$,
so $\til{R}$ is isomorphic, in $\dcat{D}(\til{A}^{\mrm{en}}, \mrm{gr})$,
to its abstract derived completion 
$\opn{ADC}_{\til{\m}}(\til{R}) = \opn{RHom}_{\til{A}} (\til{P}, \til{R})$.
See Subsection \ref{subsec:NC-MGM}. 
\end{itemize}

The composition of the isomorphisms in (\ref{eqn:4413}) is the graded 
rigidifying isomorphism $\til{\rho}$. 
\end{proof}

The localization homomorphism $\til{A} \to \til{A}_t$
induces a flat homomorphism of graded rings  
$\til{A}^{\mrm{en}} \to (\til{A}_t)^{\mrm{en}} =
\til{A}_t \ot \til{A}_t^{\mrm{op}}$.
This homomorphism can also be viewed as a localization of 
$\til{A}^{\mrm{en}}$ with respect to the degree $1$ central elements 
$t \ot 1$ and $1 \ot t$. 
There is a corresponding induction functor 
\begin{equation} \label{eqn:4420}
\opn{Ind}_{(\til{A}_t)^{\mrm{en}}}
: \dcat{D}(\til{A}^{\mrm{en}}, \mrm{gr}) \to 
\dcat{D} \bigl( (\til{A}_t)^{\mrm{en}}, \mrm{gr} \bigr) .
\end{equation}
Here is what it does on objects: given a complex 
$\til{M} \in \dcat{D}(\til{A}^{\mrm{en}}, \mrm{gr})$,
we have 
\[ \opn{Ind}_{\til{A}^{\mrm{en}}}(\til{M}) = 
(\til{A}_t)^{\mrm{en}} \ot_{\til{A}^{\mrm{en}}} \til{M} \cong 
\til{A}_t \ot_{\til{A}} \til{M} \ot_{\til{A}} \til{A}_t . \]
There is a ring homomorphism  
\[ \opn{inc} : (\til{A}_t)^{\mrm{en}}_0 = 
\opn{Deg}_0 \bigl( (\til{A}_t)^{\mrm{en}} \bigr)
\to (\til{A}_t)^{\mrm{en}} , \]
the inclusion of the degree $0$ component. And there is a functor 
\begin{equation} \label{eqn:4421}
\opn{Deg}_0 : \dcat{D} \bigl( (\til{A}_t)^{\mrm{en}}, \mrm{gr} \bigr) \to 
\dcat{D} \bigl( (\til{A}_t)^{\mrm{en}}_0 \bigr) . 
\end{equation}

We have a ring isomorphism 
$h : A \iso (\til{A}_t)_0$; see formula (\ref{eqn:4476}). It gives rise to 
ring homomorphisms 
\[ A^{\mrm{en}} \xar[\cong]{h^{\mrm{en}}}
(\til{A}_t)_0 \ot (\til{A}_t^{\mrm{op}})_0 
\xar[\sub]{\opn{inc}} 
(\til{A}_t \ot \til{A}_t^{\mrm{op}})_0 = (\til{A}_t)^{\mrm{en}}_0 . \]
Relative to this composed ring homomorphism we have the restriction functor 
\begin{equation} \label{eqn:4416}
\opn{Rest}_{A^{\mrm{en}}} : 
\dcat{D} \bigl( (\til{A}_t)^{\mrm{en}}_0 \bigr) \to 
\dcat{D}(A^{\mrm{en}}) . 
\end{equation}
Composing the functors from (\ref{eqn:4420}), (\ref{eqn:4421}) and 
(\ref{eqn:4416}) we get a triangulated functor 
\begin{equation} \label{eqn:4422}
\opn{Rest}_{A^{\mrm{en}}} \circ \opn{Deg}_0 \circ 
\opn{Ind}_{(\til{A}_t)^{\mrm{en}}} :
\dcat{D}(\til{A}^{\mrm{en}}, \mrm{gr}) \to \dcat{D}(A^{\mrm{en}}) . 
\end{equation}

\begin{thm} \label{thm:4395}
Under Setup \tup{\ref{set:4510}}, let 
$\til{R} \in \dcat{D}(\til{A}^{\mrm{en}}, \mrm{gr})$ be a 
graded rigid NC dualizing complex. Define the complex 
\[ R := (\opn{Rest}_{A^{\mrm{en}}} \circ \opn{Deg}_0 \circ 
\opn{Ind}_{(\til{A}_t)^{\mrm{en}}}) (\til{R})[-1] \in \dcat{D}(A^{\mrm{en}}) . 
\]
Then $R$ is a rigid NC dualizing complex over $A$.
\end{thm}

Once more we were a bit sloppy. The proper statement is this: 
we are given a graded rigid NC dualizing complex $(\til{R}, \til{\rho})$; 
then the complex $R$ defined above admits a rigidifying isomorphism $\rho$, 
such that $(R, \rho)$ is a rigid NC dualizing complex over $A$. The proof will 
come after two lemmas.  

\begin{lem} \label{lem:4393}
Assume that $\til{R}$ is a graded rigid dualizing complex over $\til{A}$.
Then for every $p$ the graded bimodule $\opn{H}^p(\til{R})$ is 
central over $\opn{Cent}(\til{A})$. 
\end{lem}

\begin{proof}
We know that the central homotheties 
\[ \opn{chm}^{}_{\til{A}, \til{A}} ,  \, 
\opn{chm}^{}_{\til{A}, \til{A}^{\mrm{op}}} :
\opn{Cent}(\til{A}) \to 
\opn{End}_{\dcat{M}(\til{A}^{\mrm{en}}, \mrm{gr})}(\til{A}) \]
are equal graded ring isomorphism.
Therefore the graded version of Theorem \ref{thm:4375} -- that is proved the 
same way -- tells us that the ring homomorphisms
\[ \opn{chm}^{\dcat{D}}_{R, \til{A}} , \,
\opn{chm}^{\dcat{D}}_{R, \til{A}^{\mrm{op}}} :
\opn{Cent}(\til{A}) 
\to \opn{End}_{\dcat{D}(\til{A}^{\mrm{en}}, \mrm{gr})}(\til{R})  \]
are equal isomorphisms.
  
The left $\opn{Cent}(\til{A})$-module structure on $\opn{H}^p(\til{R})$
coincides with the categorical action $\opn{chm}^{\dcat{D}}_{R, \til{A}}$; 
and likewise from the right. We see that the left and right 
$\opn{Cent}(\til{A})$-module structure on $\opn{H}^p(\til{R})$
are the same, so it is a central bimodule. 
\end{proof}

\begin{lem} \label{lem:4395}
Assume that $\til{R}$ is a graded rigid dualizing complex over $\til{A}$.
Then the homomorphisms 
\[ \tag{l} \til{A}_t \ot_{\til{A}} \til{R} \to 
\til{A}_t \ot_{\til{A}} \til{R} \ot_{\til{A}} \til{A}_t \]
and
\[ \tag{r} \til{R} \ot_{\til{A}} \til{A}_t \to 
\til{A}_t \ot_{\til{A}} \til{R} \ot_{\til{A}} \til{A}_t \]
in $\dcat{C}_{\mrm{str}}(\til{A}^{\mrm{en}}, \mrm{gr})$
are quasi-isomorphisms. 
\end{lem}

\begin{proof}
We shall only treat (l); the homomorphism (r) is treated similarly. 
Let $\til{C} := \opn{Cent}(\til{A})$. Then 
$\til{A}_t = \til{A} \ot_{\til{C}} \til{C}_t = 
\til{C}_t \ot_{\til{C}} \til{A}$.
Hence (l) is true if (and only if) the homomorphism
$\til{C}_t \ot_{\til{C}} \til{R} \to 
\til{C}_t \ot_{\til{C}} \til{R} \ot_{\til{C}} \til{C}_t$
is a quasi-isomorphism. 

Because of flatness of localization, for every $p$ there is a commutative 
diagram 
\begin{equation} \label{eqn:5127}
\UseTips \xymatrix @C=6ex @R=6ex {
\opn{H}^p \bigl( \til{C}_t \ot_{\til{C}} \til{R} \bigr)
\ar[r]
\ar[d]_{\cong}
&
\opn{H}^p \bigl( \til{C}_t \ot_{\til{C}} \til{R} \ot_{\til{C}} \til{C}_t \bigr)
\ar[d]_{\cong}
\\
\til{C}_t \ot_{\til{C}} \opn{H}^p (\til{R})
\ar[r]
&
\til{C}_t \ot_{\til{C}} \opn{H}^p (\til{R}) \ot_{\til{C}} \til{C}_t 
} 
\end{equation}
in $\dcat{M}(\til{A}^{\mrm{en}}, \mrm{gr})$, with vertical isomorphisms.
By Lemma \ref{lem:4393}, $\opn{H}^p (\til{R})$ is a central $\til{C}$ 
bimodule. And localization satisfies 
$\til{C}_t = \til{C}_t \ot_{\til{C}} \til{C}_t$.
This implies that the bottom arrow in diagram (\ref{eqn:5127}) is bijective. 
Hence the top arrow is also bijective. 
\end{proof}

\begin{proof}[Proof of Theorem \tup{\ref{thm:4395}}]
The proof is divided into several steps. 

\smallskip \noindent
Step 1. This is the easiest step: we prove that for every $q$ the bimodule 
$\opn{H}^q(R)$ is a finite module over $A$ and over $A^{\mrm{op}}$. 
In fact, we are going to treat the left module structure of $\opn{H}^q(R)$ 
only; the right module structure is treated the same way, by replacing $A$ 
with $A^{\mrm{op}}$.

Let us write
\[ \til{R}_t := \opn{Ind}_{(\til{A}_t)^{\mrm{en}}}(\til{R}) \cong 
\til{A}_t \ot_{\til{A}} \til{R} \ot_{\til{A}} \til{A}_t \in 
\dcat{D} \bigl( (\til{A}_t)^{\mrm{en}}, \mrm{gr} \bigr) . \]
According to Lemma \ref{lem:4395} there is an isomorphism 
$\til{R}_t \cong \til{A}_t \ot_{\til{A}} \til{R}$ 
in $\dcat{D}(\til{A}_t, \mrm{gr})$. 
Since $\til{R} \in \dcat{D}_{\mrm{f}}(\til{A}, \mrm{gr})$, 
it follows that 
$\til{R}_t \in \dcat{D}_{\mrm{f}}(\til{A}_t, \mrm{gr})$. 
We identify $(\til{A}_t)_0$ with $A$ using the canonical isomorphism $h$
from (\ref{eqn:4476}). 
By Lemma \ref{lem:4400}(2) we get 
$\opn{Deg}_0(\til{R}_t) \in \dcat{D}_{\mrm{f}} \bigl( (\til{A}_t)_0 \bigr) 
= \dcat{D}_{\mrm{f}}(A)$.
Hence 
\begin{equation} \label{eqn:4426}
R = \opn{Deg}_0(\til{R}_t)[-1] \in \dcat{D}_{\mrm{f}}(A) . 
\end{equation}
We see that $\opn{H}^q(R)$ is a finite $A$-module for all $q$. 

\medskip \noindent
Step 2. Here we prove that $R$ has finite injective dimension over $A$ and over 
$A^{\mrm{op}}$. Again, we only examine the left side; the right side is done 
similarly, just replacing $A$ with $A^{\mrm{op}}$.

Because $A$ is noetherian, it suffices to find a uniform bound for the 
cohomology of $\opn{RHom}_{A}(M, R)$, for all $M \in \dcat{M}_{\mrm{f}}(A)$. 
Take such an $A$-module $M$. 
By Lemma \ref{lem:4390}(3) there exists a graded module  
$\til{M} \in \dcat{M}_{\mrm{f}}(\til{A}, \mrm{gr})$
such that, letting
$\til{M}_t := \til{A}_t \ot_{\til{A}} M$, we have 
\begin{equation} \label{eqn:4425}
M \cong \opn{Deg}_0(\til{M}_t) 
\end{equation} 
in $\dcat{M}_{\mrm{f}}(A)$. There are isomorphisms  
\begin{equation} \label{eqn:4427}
\begin{aligned}
&
\opn{RHom}_{\til{A}}(\til{M}, \til{R}) \ot_{\til{A}} \til{A}_t
\cong^1 \opn{RHom}_{\til{A}}(\til{M}, \til{R} \ot_{\til{A}} \til{A}_t)
\\
& \quad 
\cong^2 \opn{RHom}_{\til{A}}(\til{M}, \til{R}_t)
\cong^3 \opn{RHom}_{\til{A}_t}(\til{M}_t, \til{R}_t)
\end{aligned}
\end{equation}
in $\dcat{D}(\til{A}_t^{\mrm{op}}, \mrm{gr})$.
The isomorphism $\cong^1$ is an instance of the graded tensor-evaluation 
isomorphism (Theorem \ref{thm:4535}), that is valid because $\til{A}$ is 
noetherian, $\til{M} \in \lb \dcat{M}_{\mrm{f}}(\til{A}, \mrm{gr})$, 
$\til{R} \in \dcat{D}^{\mrm{b}}(\til{A}^{\mrm{en}}, \mrm{gr})$, 
and $\til{A}_t$ is flat over $\til{A}$. 
The isomorphism $\cong^2$ is by Lemma \ref{lem:4395}. 
And the isomorphism $\cong^3$ is adjunction for the ring homomorphism 
$\til{A} \to \til{A}_t$.
 
Next, according to Lemma \ref{lem:4401}, and using the isomorphisms 
(\ref{eqn:4425}) and (\ref{eqn:4426}), there is an isomorphism 
\[ \begin{aligned}
&
\opn{Deg}_0 \bigl( \opn{RHom}_{\til{A}_t}(\til{M}_t, \til{R}_t) \bigr) \cong
\opn{RHom}_{\opn{Deg}_0(\til{A}_t)} 
\bigl( \opn{Deg}_0(\til{M}_t), \opn{Deg}_0(\til{R}_t) \bigr)
\\ & \quad 
\cong \opn{RHom}_{A}(M, R)[1] 
\end{aligned} \]
in $\dcat{D}(A^{\mrm{op}})$. 
Combining this with formula (\ref{eqn:4427}), 
we see that the injective dimension of $R$ over $A$ is at most one more than 
the graded injective dimension of $\til{R}$ over $\til{A}$; and this is known 
to be finite.

\medskip \noindent
Step 3. In this step we prove that $R$ has the derived Morita property on both 
sides. As before, we only prove this property on the $A$ side; the 
$A^{\mrm{op}}$ side is done similarly. 

The isomorphisms (\ref{eqn:4427}), but for the complex of bimodules $\til{R}$ 
instead of the for the module $\til{M}$, become isomorphisms 
\begin{equation} \label{eqn:4428}
\begin{aligned}
&
\opn{RHom}_{\til{A}}(\til{R}, \til{R}) \ot_{\til{A}} \til{A}_t
\cong^1 \opn{RHom}_{\til{A}}(\til{R}, \til{R} \ot_{\til{A}} \til{A}_t)
\\
& \quad 
\cong^2 \opn{RHom}_{\til{A}}(\til{R}, \til{R}_t)
\cong^3 \opn{RHom}_{\til{A}_t}(\til{A}_t \ot_{\til{A}} \til{R}, \til{R}_t)
\\
& \quad 
\cong^4 \opn{RHom}_{\til{A}_t}(\til{R}_t, \til{R}_t)
\end{aligned}
\end{equation}
in $\dcat{D}(\til{A}_t^{\mrm{op}}, \mrm{gr})$.
The last isomorphism $\cong^4$ is another case of Lemma \ref{lem:4395}.
Define 
$R' := \opn{Deg}_0(\til{R}_t) \in \dcat{D}_{\mrm{f}}(A)$;
so $R = R'[-1]$. 
We now apply $\opn{Deg}_0$ to the last object in (\ref{eqn:4428}), obtaining the
isomorphisms 
\begin{equation} \label{eqn:4431}
\begin{aligned}
&
\opn{Deg}_0 \bigl( \opn{RHom}_{\til{A}_t}(\til{R}_t, \til{R}_t) \bigr)
\\
& \quad 
\cong^{\tup{(i)}} \opn{RHom}_{\opn{Deg}_0(\til{A}_t)} 
\bigl( \opn{Deg}_0(\til{R}_t), \opn{Deg}_0(\til{R}_t) \bigr)
\\
& \quad 
\cong^{\tup{(ii)}} \opn{RHom}_{A}(R', R')
\cong^{\tup{(iii)}} \opn{RHom}_{A}(R, R)
\end{aligned}
\end{equation}
in $\dcat{D}(A^{\mrm{op}})$.
The isomorphism $\cong^{\tup{(i)}}$ is due to Lemma \ref{lem:4401}, with 
$\til{B} = \til{C}:= \til{A}_t$. 
The isomorphism $\cong^{\tup{(ii)}}$ is simply the definition of $R'$, and 
the isomorphism $\cong^{\tup{(iii)}}$ is because 
$R = R'[-1]$. 

We know that $\til{R}$ has the graded derived NC Morita property on the 
$\til{A}$ side. Using Lemma \ref{lem:4361} (that is true also in the graded 
situation) we see that \lb 
$\opn{H}^q \bigl( \opn{RHom}_{\til{A}}(\til{R}, \til{R}) \bigr) = 0$ 
for all $q \neq 0$. Because the functors $\opn{H}^q$ and $\opn{Deg}_0$ commute 
with each other, the isomorphisms (\ref{eqn:4428}) and (\ref{eqn:4431}) tell us 
that 
\begin{equation} \label{eqn:4433}
\opn{H}^q \bigl( \opn{RHom}_{A}(R, R) \bigr) = 0
\end{equation}
for all $q \neq 0$.

For $q = 0$ we know that 
$\opn{H}^0 \bigl( \opn{RHom}_{\til{A}}(\til{R}, \til{R}) \bigr)$
is a graded-free $\til{A}^{\mrm{op}}$-module with basis $\opn{id}_{\til{R}}$. 
This is by Lemma \ref{lem:4361}, applied in the graded situation. 
The isomorphisms (\ref{eqn:4428}) send the element
$\opn{id}_{\til{R}} \ot \, 1$ to $\opn{id}_{\til{R}_t}$; and therefore \lb 
$\opn{H}^0 \bigl( \opn{RHom}_{\til{A}_t}(\til{R}_t, \til{R}_t) \bigr)$
is a graded-free $\til{A}_t^{\mrm{op}}$-module with basis 
$\opn{id}_{\til{R}_t}$.
The element $\opn{id}_{\til{R}_t}$ has degree $0$, and hence  
\[ \opn{H}^0 \bigl(  \opn{Deg}_0 \bigl( \opn{RHom}_{\til{A}_t}(\til{R}_t, 
\til{R}_t) \bigr) \bigr) \cong 
\opn{Deg}_0 \bigl( \opn{H}^0 \bigl(   \opn{RHom}_{\til{A}_t}(\til{R}_t, 
\til{R}_t) \bigr) \bigr) \]
is a free $A^{\mrm{op}}$-module with basis $\opn{id}_{\til{R}_t}$.
The isomorphisms (\ref{eqn:4431}) send the element 
$\opn{id}_{\til{R}_t}$ to $\opn{id}_{R}$; so 
$\opn{H}^0 \bigl( \opn{RHom}_{A}(R, R) \bigr)$ 
is a free $A^{\mrm{op}}$-module with basis $\opn{id}_{R}$.
Again invoking Lemma \ref{lem:4361}, this fact, with formula (\ref{eqn:4433}), 
say that $R$ has the derived NC Morita property on the $A$ side.  

\medskip \noindent
Step 4. Here we produce a rigidifying isomorphism $\rho$ for $R$, starting with 
a given graded rigidifying isomorphism 
\[ \til{\rho} : \til{R} \iso 
\opn{RHom}_{\til{A}^{\mrm{en, out}}} ( \til{A}, \til{R} \ot \til{R} ) \]
in $\dcat{D}(\til{A}^{\mrm{en}}, \mrm{gr})$.  
Here is a list of isomorphisms in 
$\dcat{D} \bigl( (\til{A}_t)^{\mrm{en}}, \mrm{gr} \bigr)$.
\begin{equation} \label{eqn:4435}
\begin{aligned}
& \til{R}_t \cong^1 
\opn{RHom}_{\til{A}^{\mrm{en, out}}} ( \til{A}, \til{R} \ot \til{R})
\ot_{\til{A}^{\mrm{en}}} (\til{A}_t)^{\mrm{en}}
\\
& \quad \cong^{2} 
\opn{RHom}_{\til{A}^{\mrm{en, out}}} \bigl( \til{A}, 
(\til{R} \ot \til{R}) \ot_{\til{A}^{\mrm{en}}} (\til{A}_t)^{\mrm{en}} \bigr)
\\
& \quad \cong^{3} 
\opn{RHom}_{\til{A}^{\mrm{en, out}}} 
\bigl( \til{A}, (\til{R} \ot_{\til{A}} \til{A}_t) \ot
(\til{A}_t \ot_{\til{A}} \til{R}) \bigr) 
\\
& \quad \cong^{4} 
\opn{RHom}_{\til{A}^{\mrm{en, out}}} (\til{A}, \til{R}_t \ot \til{R}_t) 
\\
& \quad \cong^{5} 
\opn{RHom}_{(\til{A}_t)^{\mrm{en, out}}} (\til{A}_t, \til{R}_t \ot \til{R}_t) .
\end{aligned}
\end{equation}
The justifications are:
\begin{itemize}
\item[$\cong^{1}$~:] Here we apply the functor 
$\opn{Ind}_{(\til{A}_t)^{\mrm{en}}}$ to the isomorphism $\til{\rho}$. The 
inside action of $\til{A}^{\mrm{en}}$ on $\til{R} \ot \til{R}$ is treated as 
a right action. 

\item[$\cong^{2}$~:] This is by the graded tensor-evaluation 
isomorphism (Theorem \ref{thm:4535}). It holds because $\til{A}$ is
a derived graded-pseudo-finite complex over $\til{A}^{\mrm{en}}$,
see Corollary \ref{cor:4531}; 
$\til{R} \ot \til{R} \in \dcat{D}^{\mrm{b}}(\til{A}^{\mrm{four}}, \mrm{gr})$;
and $(\til{A}_t)^{\mrm{en}}$ is flat over $\til{A}^{\mrm{en}}$.

\item[$\cong^{3}$~:] This is a rearrangement of derived tensor factors. 

\item[$\cong^{4}$~:] By Lemma \ref{lem:4395}. 

\item[$\cong^{5}$~:] This is adjunction for the graded ring homomorphism 
$\til{A}^{\mrm{en}} \to (\til{A}_t)^{\mrm{en}}$, with the fact that 
$\til{A}_t = \til{A}_t \ot_{\til{A}} \til{A} \ot_{\til{A}} \til{A}_t$ as graded 
rings.
\end{itemize}

Now we apply $\opn{Deg}_0$ to (\ref{eqn:4435}), 
and we obtain the first isomorphism below in $\dcat{D}(A^{\mrm{en}})$.
\begin{equation} \label{eqn:4436}
\begin{aligned}
& 
R' = \opn{Deg}_0(\til{R}_t) \cong^{\tup{(i)}} 
\opn{Deg}_0 \bigl( 
\opn{RHom}_{(\til{A}_t)^{\mrm{en, out}}} (\til{A}_t, 
\til{R}_t \ot \til{R}_t) \bigr) 
\\ & \quad 
\cong^{\tup{(ii)}} \opn{RHom}_{A^{\mrm{en, out}}} (A, R' \ot R')[-1]
\end{aligned}
\end{equation}
The isomorphism $\cong^{\tup{(ii)}}$ is by Theorem \ref{thm:4510}, with 
$\til{B} := \til{A}_t$, that we are allowed to use because $A$ is a derived 
pseudo-finite complex over $A^{\mrm{en}}$; see Theorem \ref{thm:4511}. 
Plugging in $R = R'[-1]$ we get 
$R \cong \opn{RHom}_{A^{\mrm{en, out}}}(A, R \ot R)$
in $\dcat{D}(A^{\mrm{en}})$. This is the rigidifying isomorphism $\rho$.
\end{proof}




\mysubsection{Filtered Rings and Existence of Rigid NC DC}
\label{subsec:Fil-exit-R-NCDC}

In this subsection we give M. Van den Bergh's proof of the existence of NC 
rigid dualizing complexes (Theorem \ref{thm:4396}). This proof appeared in his 
seminal paper \cite{VdB}, and then there were some improvements in 
\cite{YeZh1}. The proof here is much more detailed. 

We continue with Convention \ref{conv:3670}.
Specifically, $\K$ is a base field, and all rings are central over $\K$. 

Earlier in the book (in Section \ref{sec:exist-resol}) we encountered 
filtrations of DG 
modules. Here we are interested in filtrations of rings. By a {\em filtration 
of a ring $A$} we mean a collection 
$\bigl\{ F_j(A) \bigr\}_{j \geq -1}$ of $\K$-submodules 
\[ F_{-1}(A) \sub F_{0}(A) \sub F_{1}(A) \sub \cdots \sub A \]
such that 
$F_{-1}(A) = 0$, $\bigcup_j F_j(A) = A$, 
$1_A \in F_{0}(A)$ and 
$F_{j}(A) \cd F_{k}(A) \sub F_{j + k}(A)$ 
for all $j, k$. 

Given a filtration $F = \bigl\{ F_j(A) \bigr\}_{j \geq -1}$ of the ring $A$, 
we write
$\opn{Gr}^F_j(A) := \lb F_{j}(A) / F_{j - 1}(A)$ 
for $j \geq 0$. 
The {\em associated graded ring} 
$\opn{Gr}^F(A) := \bigoplus_{j \geq 0} \opn{Gr}^F_j(A)$ 
is a nonnegative algebraically graded ring. 

A filtration $F$ on the ring $A$ gives rise to another graded ring: it is the  
{\em Rees ring} 
$\opn{Rees}^F(A) := \bigoplus_{j \geq 0} F_j(A) \cd t^j \sub A[t]$. 
Here $t$ is a central variable of degree $1$, that we call the Rees parameter. 
There is a ring isomorphism 
$A \cong \opn{Rees}^F(A) / (t - 1)$
and a graded ring isomorphism
$\opn{Gr}^F(A) \cong \opn{Rees}^F(A) / (t)$.
We often use the abbreviations 
$\bar{A} := \opn{Gr}^F(A)$ and 
$\til{A} := \opn{Rees}^F(A)$. 
This is consistent with the notation used earlier in the section: 
$\til{A}$ is a graded ring, $t$ is a degree $1$ central regular element in 
it, and $A = \til{A} / (t - 1)$.

For an element $c \in \K$ we write 
$\opn{pr}_{c} : \til{A} \to \til{A} / (t - c)$ 
for the canonical surjective ring homomorphism. 
The homomorphisms $\opn{pr}_{0}$ and $\opn{pr}_{1}$ are 
displayed in this diagram of rings:
\[ \UseTips \xymatrix @C=2ex @R=5ex {
&
\til{A} = \opn{Rees}^F(A)
\ar[dl]_(0.6){\opn{pr}_1}
\ar[dr]^(0.6){\opn{pr}_0}
\\
*+{A}
&
&
*+{\bar{A} = \opn{Gr}^F(A)}
} \] 
The filtration $F$ can be recovered from the Rees ring as follows: 
\begin{equation} \label{eqn:4506}
F_j(A) = \opn{pr}_{1} (\til{A}_j) 
= \opn{pr}_{1} \Bigl( \bigoplus\nolimits_{i = 0}^j \, \til{A}_i \Bigr) \sub A . 
\end{equation}

\begin{dfn}[\cite{YeZh1}] \label{dfn:4445}
A filtration $F$ of the ring $A$ is called a 
{\em noetherian connected filtration} 
\index{Filtration on a ring! noetherian connected}
if the graded ring $\opn{Gr}^F(A)$ is a noetherian connected graded 
ring (Definitions \ref{dfn:3703} and \ref{dfn:3706}). 
\end{dfn}

Let $\til{B}$ be a connected graded ring. Recall from Definition \ref{dfn:4532} 
that a complex $M \in \dcat{D}(\til{B}, \mrm{gr})$ is said to be 
derived graded-pseudo-finite if it belongs to the saturated full 
triangulated subcategory of $\dcat{D}(\til{B}, \mrm{gr})$
generated by the bounded above complexes $P$ such that each $P^i$ is a 
finite graded-free $\til{B}$-module. 

\begin{prop} \label{prop:4395}
Assume the ring $A$ admits a noetherian connected filtration $F$. Then\tup{:}
\begin{enumerate}
\item The Rees ring $\til{A} := \opn{Rees}^F(A)$ is a noetherian connected 
graded ring, and $\til{A}$ is a derived graded-pseudo-finite complex  
over $\til{A}^{\mrm{en}}$.

\item The ring $A$ is noetherian, and $A$ is a derived pseudo-finite complex 
over $A^{\mrm{en}}$.
\end{enumerate}
\end{prop}

\begin{proof} \mbox{}

\smallskip \noindent 
(1) Let $\bar{A} := \opn{Gr}^F(A)$, which is a noetherian connected 
graded ring. Since 
$\til{A}_0 = \opn{Rees}^F_0(A) = F_0(A) = \opn{Gr}^F_0(A) = \bar{A}_0 = \K$
and 
$\til{A}_i = \opn{Rees}^F_i(A) \cong F_i(A) \cong 
\bigoplus_{0 \leq j \leq i} \bar{A}_j$ 
as ungraded $\K$-modules, we see that $\til{A}$ is connected graded. 
Next, because $\bar{A} = \til{A} / (t)$, by Theorem \ref{thm:4240} we deduce 
that $\til{A}$ is noetherian. 
By Proposition \ref{prop:4266}, $\til{A}$ is a derived graded-pseudo-finite
complex over $\til{A}^{\mrm{en}}$.

\medskip \noindent 
(2) By item (1)  the graded ring $\til{A}$ is connected and noetherian. 
Since $A = \til{A} / (t - 1)$, it follows that $A$ is noetherian. And according 
to Theorem \ref{thm:4511}(3), $A$ is a derived pseudo-finite complex over 
$A^{\mrm{en}}$.  
\end{proof}

\begin{thm}[Van den Bergh Existence, \cite{VdB}] \label{thm:4396}
\index{Dualizing complex! balanced}
\index{Dualizing complex! noncommutative rigid}
\index{Filtration on a ring! noetherian connected}
Let $A$ be a ring. Assume $A$ admits a noetherian connected filtration 
$F$, such that the connected graded ring 
$\opn{Gr}^F(A)$ has a balanced NC dualizing complex.
Then the ring $A$ has a rigid NC dualizing complex $(R, \rho)$. Moreover, 
$(R, \rho)$ is unique up to a unique rigid isomorphism.
\end{thm}

\begin{proof}
As before we write $\til{A} := \opn{Rees}^F(A)$ and $\bar{A} := \opn{Gr}^F(A)$.
These are noetherian connected graded rings, and 
$\bar{A} = \til{A} / (t)$. Since $\bar{A}$ admits a balanced NC dualizing 
complex, Corollary \ref{cor:4585} says that $\bar{A}$ satisfies the $\chi$ 
condition and it has finite local cohomological dimension. 
By Theorem \ref{thm:4250} the graded ring $\til{A}$ also satisfies the $\chi$ 
condition and it has finite local cohomological dimension.
Using Corollary \ref{cor:4585} for $\til{A}$, we conclude that it has a 
balanced NC dualizing complex $\til{R}$. 

According to Theorem \ref{thm:4400}, $\til{R}$ is a graded rigid NC dualizing 
complex over $\til{A}$. Theorem \ref{thm:4395} then says that 
\[ R := (\opn{Rest}_{A^{\mrm{en}}} \circ \opn{Deg}_0 \circ 
\opn{Ind}_{(\til{A}_t)^{\mrm{en}}}) (\til{R})[-1] \in \dcat{D}(A^{\mrm{en}}) \]
is a rigid NC dualizing complex over $A$; i.e.\ there exists a rigidifying 
isomorphism $\rho$ for $R$. 

By Proposition \ref{prop:4395}(2), $A$ is a derived pseudo-finite complex over
$A^{\mrm{en}}$. Then, according to Theorem \ref{thm:4380}, the rigid dualizing 
complex $(R, \rho)$ is unique up to a unique rigid isomorphism.
\end{proof}

Here is a result we won't prove here. It is needed for the example that comes 
after it. 

\begin{prop} [{\cite[Proposition 6.13]{YeZh1}}]  \label{prop:4505}
Let $A$ be a commutative ring, and let $f : A \to B$ be a finite central ring 
homomorphism. If $A$ has a noetherian connected filtration $F$, then 
there exists a noetherian connected filtration $G$ of $B$, such that 
$f(F_j(A)) \sub G_j(B)$ for all $j$, and the graded ring homomorphism 
$\opn{Gr}(f) : \opn{Gr}^F(A) \to \opn{Gr}^G(B)$ 
is finite. 
\end{prop}

\begin{exa} \label{exa:4506}
Suppose the ring $B$ is finite over its center $\opn{Cent}(B)$, and 
$\opn{Cent}(B)$ is a finitely generated $\K$-ring. 
Then we can find a finite central homomorphism $f : A \to B$, where
$A =  \K[t_1, .., t_n]$ is a commutative polynomial ring. 
The grading on $A$ with $\opn{deg}(t_i) = 1$ gives rise to a noetherian 
connected filtration $F$, by the formula 
$F_j(A) := \bigoplus_{k \leq j} A_k$. Of course 
$\opn{Gr}^F(A) \cong A$. 

Now we can use Proposition \ref{prop:4505}
to conclude that $B$ has a noetherian connected filtration $G$ such that 
$\opn{Gr}^F(A) \to \opn{Gr}^G(B)$ is finite. According to Example \ref{exa:3990}
and Corollary \ref{cor:4032}, the 
graded ring $\opn{Gr}^G(B)$ has a balanced dualizing complex. 
Therefore, by Theorem \ref{thm:4396}, $B$ has a rigid NC dualizing complex. 
\end{exa}

Sometimes rings come with filtrations of the following type. 

\begin{dfn} \label{dfn:4515}
A {\em central filtration of finite type} of a ring $A$ is a filtration 
$G = \bigl\{ G_j(A) \bigr\}_{j \geq -1}$
of $A$, such that the graded ring 
$\opn{Gr}^G(A) = \bigoplus_{j \geq 0} \opn{Gr}^G_j(A)$
is finite over its center 
$C = \bigoplus_{j \geq 0} C_j  := \opn{Cent} \bigl( \opn{Gr}^G(A) \bigr)$, 
and the commutative graded ring $C$ is finitely generated over $\K$.
\end{dfn}

In \cite{YeZh2.5} such a filtration was called a ``differential 
filtration of finite type'', but this name sounds too confusing when we also 
talk about DG rings. The reason for the name ``differential'' will be apparent 
in Example \ref{exa:4515} below. 

It is not hard to see, using Theorem \ref{thm:4240}, 
that if $A$ admits a central filtration of finite type, then 
it is noetherian. However, much more is true, as the next theorem shows.  
This is a result from \cite{YeZh2.5} that we present without a 
proof. A slightly weaker version of this theorem was in the paper
\cite{McSt} of J.C. McConnell and J.T. Stafford, which 
the authors attributed to J. Bernstein. 

\begin{thm}[Two Filtrations, {\cite[Theorem 3.1]{YeZh2.5},
\cite{McSt}}] \label{thm:4505}
Let $A$ be a ring that admits a central filtration of finite type $G$.
Then $A$ admits a noetherian connected filtration $F$, such that the 
noetherian connected graded ring $\opn{Gr}^F(A)$ 
is a commutative finitely generated $\K$-ring.
\end{thm}

As explained in Example \ref{exa:3991}, the 
ring $\opn{Gr}^F(A)$ has a balanced dualizing complex. Therefore, using 
Theorems \ref{thm:4505} and \ref{thm:4396}, we conclude that:

\begin{cor} \label{cor:4515}
If $A$ admits a central filtration of finite type, then $A$ has a 
rigid NC dualizing complex $(R, \rho)$, and it is unique up to 
a unique rigid isomorphism.
\end{cor}

\begin{exa} \label{exa:4515}
Assume $\K$ has characteristic $0$. Let $C$ be a smooth commutative 
$\K$-ring of pure dimension $n$, and let $A := \mcal{D}(C)$, the ring of 
differential operators. 
Consider the filtration
$G = \bigl\{ G_j(A) \bigr\}_{j \geq -1}$
of $A$ by order of operators. For this filtration we have 
$G_0(A) = C$ and 
$G_1(A) = C \oplus \mcal{T}(C)$,
where $\mcal{T}(C)$ is the module of derivations.
Then $\opn{Gr}^G(A)$ is a commutative ring, and there is a canonical 
graded ring isomorphism 
$\opn{Gr}^G(A) \cong \opn{Sym}_{C}(\mcal{T}(C))$,
the symmetric tensor ring on the finitely generated projective $C$-module 
$\mcal{T}(C)$. This is a commutative finite type $\K$-ring, so $G$ is a central 
filtration of finite type.
Corollary \ref{cor:4515} says that $A$ has a rigid NC dualizing complex 
$(R, \rho)$. In \cite{Ye4.5} it was shown that $R \cong A[2 n]$.
\end{exa}

We end the subsection with three remarks. 

\begin{rem} \label{rem:4870}
Here is a brief history of the early days NC dualizing complexes (during the 
1990's). The first paper on this topic was \cite{Ye1} by A. Yekutieli from 
1992, in which dualizing complexes over NC connected graded rings (always over a 
base field $\K$) were introduced (see Definition \ref{dfn:3716}).
This material was part of the Ph.D. thesis of Yekutieli, and was inspired by 
ideas of his advisor M. Artin on a noncommutative analogue of Serre Duality. 
The paper \cite{Ye1} also contained the definition of balanced dualizing 
complexes (Definition \ref{dfn:3714}), a proof of NC Local Duality (Theorem 
\ref{thm:4616}), a proof of the existence of balanced dualizing complexes 
(Theorem \ref{thm:3781}), and some particular calculations (Example 
\ref{exa:3780}). 

The paper \cite{YeZh0} by Yekutieli and J.J. Zhang from 1997 used the results 
of \cite{Ye1} to establish NC Serre Duality on noncommutative projective 
schemes, in the sense of Artin and Zhang \cite{ArZh}. In the same paper we 
also find Corollary \ref{cor:4655}, saying that the existence of a balanced 
dualizing complex over the noetherian connected ring $A$ implies that $A$ 
satisfies the $\chi$ condition and has finite local cohomological dimension. 

The paper \cite{VdB} by M. Van den Bergh from 1997 begins by 
establishing the important converse to Corollary \ref{cor:4655}, namely that 
a noetherian connected ring $A$, which satisfies the $\chi$ condition and has 
finite local cohomological dimension, admits a balanced dualizing complex 
(Theorem \ref{thm:3713}). Thus we have Corollary \ref{cor:4585} on the 
``two equivalent properties''. Van den Bergh went on to define rigid dualizing 
complexes, to prove their uniqueness (Theorem \ref{thm:4380}), and to 
prove existence of rigid dualizing complexes, in the presence of suitable 
filtrations and balanced dualizing complexes (Theorem \ref{thm:4396}).
The subsequent paper \cite{VdB2} by Van den Bergh utilized the rigidity 
formalism to deduce what is now called Van den Bergh Duality 
(Theorem \ref{thm:4860} below). 

At about the same time that Van den Bergh wrote \cite{VdB}, Yekutieli worked on 
the paper \cite{Ye4}, in which the derived Picard group $\opn{DPic}(A)$ of a NC 
ring $A$ was introduced. The fact that any two NC dualizing complexes $R, R'$ 
over $A$ are related by $R' \cong T \ot^{\mrm{L}}_{A} R$ for some tilting 
complex $T$ (Theorem \ref{thm:3681}, proved in \cite{Ye4}) 
was used by Van den Bergh in his proof of the uniqueness of rigid dualizing 
complexes in \cite{VdB}.

The paper \cite{YeZh1} by Yekutieli and Zhang from 1999 built on Van den 
Bergh's ideas. The uniqueness results of Van den Bergh were sharpened, and were 
upgraded to functoriality (see Remark \ref{rem:4380}). The Auslander condition 
on dualizing complexes (see Remark \ref{rem:3821}) was also introduced and 
studied in \cite{YeZh1}.

It is hard to track the progress of the theory of NC dualizing complexes beyond 
the year 2000. This theory has merged with, or has been absorbed 
into, the more active research on Calabi-Yau rings and categories
(see next subsection), with its ties 
to birational algebraic geometry and mathematical physics. 
\end{rem}

\begin{rem} \label{rem:4380}
This remark is about the functorial properties of NC rigid dualizing 
complexes. In Remark \ref{rem:4191} we talked about the rigid trace 
homomorphism and the rigid localization homomorphism, between rigid residue 
complexes over commutative rings. These are actual homomorphisms of complexes. 
For NC rings there are usually no residue complexes (because of the pathological 
nature of injective modules in the NC setting). 
Still we can consider rigid trace morphisms and rigid localization morphisms 
between NC rigid dualizing complexes. These are morphisms in 
derived categories. Below is an outline. 

Suppose $f : A \to B$ is a finite ring homomorphism, and the rigid dualizing 
complexes $(R_A, \rho_A)$ and $(R_B, \rho_B)$ exist. A {\em rigid trace 
morphism} 
$\opn{tr}_f : (R_B, \rho_B) \to (R_A, \rho_A)$ 
is a backward morphism $\opn{tr}_f : R_B \to R_A$ in $\dcat{D}(A^{\mrm{en}})$, 
which is rigid (the obvious 
variant of Definition \ref{dfn:4372}), and is nondegenerate on both sides 
(see Definition \ref{dfn:4031}). In \cite{YeZh1} the existence and uniqueness 
of the rigid trace morphism exists was proved, when $f : A \to 
B$ is finite centralizing, and $A$ admits a noetherian connected filtration $F$ 
such that $\opn{Gr}^F(A)$ has a balanced dualizing complex. The proof is along 
the same lines as the proof of Theorem \ref{thm:4396} above, combined 
with the existence of the balanced trace morphism, established in Theorem 
\ref{thm:4030}. Among the applications of the rigid trace: (1) It was used 
in \cite{Ye4.5} to calculate the rigid dualizing complex of 
$A = \mrm{U}(\g)$, see Example \ref{exa:4505} below. (2) In 
\cite{YeZh1.5} it was used to calculate multiplicities of indecomposable 
injectives in minimal injective resolutions. 

Now consider a localization ring homomorphism $g : A \to A'$. Assume the rigid 
dualizing complexes $(R_A, \rho_A)$ and $(R_{A'}, \rho_{A'})$ exist. A {\em 
rigid localization morphism} 
$\opn{q}_g : (R_A, \rho_A) \to (R_{A'}, \rho_{A'})$ 
is a forward morphism $\opn{q}_g : R_A \to R_{A'}$ in 
$\dcat{D}(A^{\mrm{en}})$, which is rigid (it respects 
the two rigidifying isomorphisms), and is nondegenerate on both sides. 
Existence and uniqueness of the rigid localization morphism was proved in 
\cite[Theorem 6.2]{YeZh2.5}; it is related to Ore conditions. One application 
of the rigid localization morphism is in the concept of {\em homological 
transcendence degree} \cite{YeZh1.7} -- the ring $A$ is an 
integral domain which has a rigid dualizing complex (say because of the 
existence of a good filtration on it), and $D = A'$ is the division ring of 
fractions of $A$. The homological transcendence degree of $D$ is the unique 
number $n$ such that the rigid dualizing complex $R_D$ satisfies 
$\opn{H}^{-n}(R_D) \neq 0$. 
\end{rem}

\begin{rem} \label{rem:3821}
In this remark we discuss the {\em Auslander property} of NC dualizing 
complexes. Around 1990, J.E. Bj\"ork \cite{Bjo} and T. Levasseur \cite{Lev}  
had studied NC regular noetherian rings with the Auslander property. This 
property was later generalized in the papers \cite{Ye2.5}
and \cite{YeZh1} by A. Yekutieli and J.J. Zhang. Below is a quick explanation 
of the meaning and applications of this concept. 

Consider a NC noetherian ring $A$ with a NC dualizing complex $R$. 
Given a module $M \in \dcat{M}_{\mrm{f}}(A)$, its grade with respect to $R$ is 
the number 
\[ \opn{grade}_R(M) := 
\opn{\inf} \bigl( \{ q \mid \opn{Ext}^q_A(M, R) \neq 0 \} \bigr) \in 
\Z \cup  \{ \infty \} . \]
Note that $\opn{Ext}^q_A(M, R) \in \dcat{M}_{\mrm{f}}(A^{\mrm{op}})$.
The number $\opn{grade}_R(N)$ for $N \in \dcat{M}_{\mrm{f}}(A^{\mrm{op}})$
is defined similarly. The complex $R$ is called {\em Auslander} if for every 
$M \in \dcat{M}_{\mrm{f}}(A)$, every $q$, and 
every $A^{\mrm{op}}$-submodule 
$N \sub \opn{Ext}^q_A(M, R)$, the inequality $\opn{grade}_R(N) \geq q$
holds; and the same inequality holds with $A$ and $A^{\mrm{op}}$ 
exchanged. The original definition used by Bj\"ork and Levasseur, for a regular 
ring $A$, was for the particular dualizing complex $R := A$. 

In \cite{YeZh1} it was proved that in many situations the rigid dualizing 
complex $R$ from Theorem \ref{thm:4396} has the Auslander property. For 
instance, \cite[Theorem 0.7]{YeZh1} implies that if $A$ has a noetherian 
connected filtration $F$ such that $\opn{Gr}^F(A)$ is commutative, then the 
rigid dualizing complex $R$ of $A$ is Auslander. 

Having an Auslander rigid dualizing complex gives rise to the {\em canonical 
dimension} of $A$-modules. For $M \in \dcat{M}_{\mrm{f}}(A)$ its canonical 
dimension is $\opn{Cdim}(M) := - \opn{grade}_R(M)$.
This can be considered as a NC version of the rigid dimension from 
Section \ref{sec:dual-cplx-comm-rng}.
Indeed, if $A$ is a finite type commutative $\K$-ring, with commutative rigid 
dualizing complex $R$, then for a prime ideal $\p \sub A$ and the $A$-module 
$M := A / \p$ we have 
$\opn{Cdim}(M) = \opn{rig{.}dim}_{\K}(\p) = \opn{dim}(A / \p)$,
where the second number is from Definition \ref{dfn:2537}, and the 
third number is the Krull dimension of the ring $A / \p$. 

As shown in \cite{YeZh1}, several structural results for NC rings 
can be deduced from the presence of an Auslander rigid dualizing complex; for 
instance, the catenarity of $\opn{Spec}(A)$, see \cite[Theorem 0.1]{YeZh1}.
Moreover, in \cite[Section 7]{YeZh2.5} a link was made 
between Auslander rigid dualizing 
complexes and perverse t-structures: the canonical dimension of modules defines 
a perverse t-structure on $\dcat{D}^{\mrm{b}}_{\mrm{f}}(A)$, and under the 
duality functor $D := \opn{RHom}_A(-, R)$ this t-structure goes to the standard 
t-structure on $\dcat{D}^{\mrm{b}}_{\mrm{f}}(A^{\mrm{op}})$. 
\end{rem}




\mysubsection{Twisted Calabi-Yau Rings}
\label{subsec:tw-CY}

In this final subsection we study a class of regular rings that is of great 
importance in recent ring theory. 
We continue with Convention \ref{conv:3670}.
Specifically, $\K$ is a base field, and all rings are central over $\K$. 
For a ring $A$, its enveloping ring is 
$A^{\mrm{en}} = A \ot A^{\mrm{op}}$. 

\begin{dfn} \label{dfn:4860}
A ring $A$ is called {\em homologically smooth over $\K$} if $A$ is a perfect 
complex over the ring $A^{\mrm{en}}$.
\end{dfn}

This definition was proposed by M. Kontsevich, around the year 2000, cf.\ 
\cite[Erratum]{VdB2}.
Clearly if $A$ is homologically smooth over $\K$, then $A$ is a derived 
pseudo-finite complex over $A^{\mrm{en}}$. 
Recall (from Example \ref{exa:3695}) that the ring $A$ is called regular if it 
has finite global cohomological dimension. 

\begin{prop} \label{prop:4860}
If $A$ is homologically smooth over $\K$, then $A$ is a regular ring. 
\end{prop}

\begin{proof}
By Theorem \ref{thm:3415} there is a quasi-isomorphism 
$\rho : P \to A$ in $\dcat{C}_{\mrm{str}}(A^{\mrm{en}})$, where $P$ 
is a bounded complex of finite projective $A^{\mrm{en}}$-modules. 
Let $d$ be the amplitude of $P$. 
Given an $A$-module $M$, define $P_M := P \ot_A M$, and let 
$\rho_M : P_M \to M$ be the homomorphism in
$\dcat{C}_{\mrm{str}}(A)$ induced by $\rho$. Since $\rho$ is a homotopy 
equivalence in $\dcat{C}_{\mrm{str}}(A^{\mrm{op}})$, it follows that $\rho_M$ is 
a quasi-isomorphism. Since $P^i$ is a projective $A^{\mrm{en}}$-module, 
and $M$ is a free $\K$-module, we see that $P_M^i = P^i \ot_A M$ is a 
projective $A$-module (often infinitely generated; cf.\ Example \ref{exa:5125}). 
The conclusion is that the 
projective dimension of $M$ is at most $d$. 

A similar consideration shows that every $M' \in \dcat{M}(A^{\mrm{op}})$
has projective dimension at most $d$. So the global homological 
dimension of $A$ is at most $d$. 
\end{proof}

\begin{rem} \label{rem:4860}
Suppose $A$ is a finite type commutative $\K$-ring. It can be shown that $A$ is 
homologically smooth over $\K$ if and only if it is smooth over $\K$, in the 
sense of \cite{EGA-IV}. Cf.\ \cite[Proposition 3.24]{Mil}.
\end{rem}

Given a $\K$-ring automorphism $\mu$ of $A$, we define the
{\em $\mu$-twisted bimodule} $A(\mu)$ as follows: as a left $A$-module,
$A(\mu)$ is the free $A$-module $A$; and the right $A$-module action is 
twisted by $\mu$. To be explicit, let us denote by $a \cd b$ the 
multiplication of the ring $A$. Then the left action of a ring element 
$a \in A$ on a bimodule element $b \in A(\mu)$ is 
$a \, \cd^{\mu} \, b := a \cd b \in A(\mu)$; and the right action is 
$b \, \cd^{\mu} \, a := b \cd \mu(a) \in A(\mu)$.
In fancier notation,  a ring element 
$a_1 \ot \opn{op}(a_2) \in A^{\mrm{en}}$
acts on a bimodule element $b \in A(\mu)$ from the left like this:
$(a_1 \ot \opn{op}(a_2)) \, \cd^{\mu} \, b = a_1 \cd b \cd \mu(a_2) 
\in A(\mu)$.
The graded version of this twisting already appeared in Definition 
\ref{dfn:3780}. 
If $\nu$ is another automorphism of $A$, then 
$A(\mu) \ot_A A(\nu) \cong A(\mu \circ \nu)$
in $\dcat{M}(A^{\mrm{en}})$. 

It is possible to twist other bimodules: 
given $M \in \dcat{M}(A^{\mrm{en}})$ we define 
\begin{equation} \label{eqn:4896}
M(\mu) := M \ot_A A(\mu) \in \dcat{M}(A^{\mrm{en}}) . 
\end{equation}
Let us make explicit the left and right actions of the 
ring $A^{\mrm{en}}$ on the module $M(\mu)$, using formula (\ref{eqn:4878}). 
The original left and right actions of an element $a \in A$ on an element
$m \in M$ are denoted by $a \cd m$ and $m \cd a$. Then an element 
$(a_1 \ot \opn{op}(a_2)) \in A^{\mrm{en}}$ acts on an element 
$m \in M(\mu)$ from the left and the right  by 
\begin{equation} \label{eqn:4904}
\begin{aligned}
& 
(a_1 \ot \opn{op}(a_2)) \, \cd^{\mu} \, m = 
a_1 \cd m \cd \mu(a_2) ,
\\
&
m \, \cd^{\mu} \, (a_1 \ot \opn{op}(a_2)) = 
a_2 \cd m \cd \mu(a_1) . 
\end{aligned}
\end{equation}

\begin{dfn} \label{dfn:4490}
\index{Ring! Calabi-Yau}
\index{Ring! twisted Calabi-Yau}
\index{Ring! Nakayama automorphism}
A ring $A$ is called a {\em twisted CY ring} if it has these 
two properties:
\begin{itemize}
\rmitem{i} The ring $A$ is homologically smooth over $\K$. 
 
\rmitem{ii} There is a natural number $n$, called the {\em dimension of $A$}, 
and a ring automorphism $\nu$ of $A$, called the {\em Nakayama automorphism}, 
such that 
\[ \opn{Ext}^i_{A^{\mrm{en}}}(A, A^{\mrm{en}}) \cong
\begin{cases}
A(\nu) & \tup{if} \ \ i = n 
\\
0 & \tup{if} \ \ i \neq n 
\end{cases} 
\]
as $A$-bimodules. 
\end{itemize}
The ring $A$ is called a {\em CY ring} if the third property also holds:
\begin{itemize}
\rmitem{iii} The automorphism $\nu$ is inner, so that $A(\nu) \cong A$ as 
$A$-bimodules.
\end{itemize}
\end{dfn}

In property (ii) the $A^{\mrm{en}}$-module 
$\opn{Ext}^i_{A^{\mrm{en}}}(A, A^{\mrm{en}})$
should be viewed as a special case of 
$\opn{Ext}^i_{B}(M, N) \in \dcat{M}(C^{\mrm{op}})$ 
for $M \in \dcat{M}(B)$ and $N \in \dcat{M}(B \ot C^{\mrm{op}})$,
where we take $B, C := A^{\mrm{en}}$, $M := A$ and 
$N := A^{\mrm{en}}$, and we use the isomorphism (\ref{eqn:4875}) to switch
from a right $A^{\mrm{en}}$-module to a left $A^{\mrm{en}}$-module. 

The letters ``CY'' are an abbreviation for ``Calabi-Yau''. 
Most texts use the term ``CY algebra'' of course. See Remark 
\ref{rem:4505} for a discussion of the background of this definition. 
Clearly this is a relative notion: it pertains to $A$ as a central $\K$-ring. 
It is easy to see that the Nakayama automorphism 
$\nu$ is unique up to composition with an inner automorphism. 

Here is another definition, taken from \cite{Ye4.5}. 

\begin{dfn} \label{dfn:4875}
Let $A$ be a noetherian ring with rigid dualizing complex $R$. If 
$R \cong A(\mu)[n]$ in $\dcat{D}(A^{\mrm{en}})$ for some 
automorphism $\mu$ and some integer $n$, then $\mu$ is 
called the {\em dualizing automorphism} of $A$. 
\end{dfn}

Like the Nakayama automorphism, the dualizing automorphism is unique up to 
composition with an inner automorphism.  

\begin{prop} \label{prop:4900}  
Let $A$ be a noetherian ring which is homologically smooth over $\K$, let $n$ 
be a natural number, and let $\nu$ be an automorphism of $A$. 
The following two conditions are equivalent. 
\begin{enumerate}
\rmitem{i} $A$ is a twisted CY ring of dimension $n$ with Nakayama automorphism 
$\nu$. 

\rmitem{ii} The complex $R := A(\nu^{-1})[n]$ is a rigid dualizing complex over 
$A$. 
\end{enumerate}
\end{prop}

Thus the dualizing automorphism is the inverse of the Nakayama automorphism.

\begin{proof}
The key formula is this: for $M, N \in \dcat{D}(A^{\mrm{en}})$ there are 
isomorphisms
\begin{equation} \label{eqn:4890}
\begin{aligned}
& 
\opn{RHom}_{A^{\mrm{en}}}(A, M \ot N) \cong^{1}
\opn{RHom}_{A^{\mrm{en}}} \bigl( A, A^{\mrm{en}} 
\ot^{\mrm{L}}_{A^{\mrm{en}}} (M \ot N) \bigr) 
\\ & \quad
\cong^{2} \opn{RHom}_{A^{\mrm{en}}}(A, A^{\mrm{en}}) 
\ot^{\mrm{L}}_{A^{\mrm{en}}} (M \ot N)
\\ & \quad 
\cong^3 N \ot^{\mrm{L}}_{A} \opn{RHom}_{A^{\mrm{en}}}(A, A^{\mrm{en}}) 
\ot^{\mrm{L}}_{A} M 
\end{aligned}
\end{equation}
in $\dcat{D}(A^{\mrm{en}})$.
The isomorphism $\cong^1$ is the left unitor. 
The isomorphism $\cong^2$ is by derived tensor-evaluation, 
and it exists because 
$A$ is a perfect complex over $A^{\mrm{en}}$; see Theorem \ref{thm:3400}.
The isomorphism $\cong^3$ is because the left $A^{\mrm{en}}$ action on 
$M \ot N$ is the outside action, cf.\ formula (\ref{eqn:4215}). 

\medskip \noindent 
(i) $\Rightarrow$ (ii): By Example \ref{exa:3695} the complex $R$ is dualizing. 
What is needed is a rigidifying isomorphism for it. We are given that 
$\opn{RHom}_{A^{\mrm{en}}}(A, A^{\mrm{en}}) \cong A(\nu)[-n]$ 
in $\dcat{D}(A^{\mrm{en}})$. 
Plugging this and $M = N := R$ in (\ref{eqn:4890}) gives 
\[ \begin{aligned}
&
\opn{RHom}_{A^{\mrm{en}}}(A, R \ot R) \cong
A(\nu^{-1})[n] \ot^{\mrm{L}}_{A} A(\nu)[-n] \ot^{\mrm{L}}_{A} A(\nu^{-1})[n]
\\ & \quad 
\cong A(\nu^{-1} \circ \nu \circ \nu^{-1})[n - n + n] = 
A(\nu^{-1})[n] = R , 
\end{aligned} \]
as required. 

\medskip \noindent 
(ii) $\Rightarrow$ (i):
Now we are given that 
$\opn{RHom}_{A^{\mrm{en}}}(A, R \ot R) \cong R$ 
in $\dcat{D}(A^{\mrm{en}})$. Substituting 
$M = N := R$ in (\ref{eqn:4890}) gives 
\[ A(\nu^{-1})[n] \cong A(\nu^{-1})[n] \ot^{\mrm{L}}_{A}
\opn{RHom}_{A^{\mrm{en}}}(A, A^{\mrm{en}}) \ot^{\mrm{L}}_{A} A(\nu^{-1})[n] \]
in $\dcat{D}(A^{\mrm{en}})$. Therefore 
$A(\nu^{})[-n] \cong \opn{RHom}_{A^{\mrm{en}}}(A, A^{\mrm{en}})$ 
in $\dcat{D}(A^{\mrm{en}})$.
This is property (iii) of Definition \ref{dfn:4490}.
\end{proof}

\begin{exa} \label{exa:4895}
This is a continuation of Example \ref{exa:4515}: $\K$ has characteristic $0$, 
$C$ is a smooth commutative $\K$-ring of pure dimension $n$, and 
$A = \mcal{D}(C)$, the ring of differential operators.
It is known that $A$ is homologically smooth -- in fact, in this case 
$A^{\mrm{en}}$ is a regular noetherian ring, of dimension $4 n$. 
According to \cite[Theorem 2.6]{Ye4.5} the rigid dualizing complex of $A$ is
$R = A[2 n]$. Hence $A$ is a CY ring of dimension $2 n$.
\end{exa}

AS regular graded rings were introduced in Definition \ref{dfn:3782}. 

\begin{thm} \label{thm:4490}
\index{Ring! twisted Calabi-Yau}
\index{Algebraically graded ring! AS regular}
Let $A$ be a ring that admits some noetherian connected filtration $F$, such 
that the graded ring $\opn{Gr}^F(A)$ is AS regular of dimension $n$. 
Then $A$ is a noetherian twisted CY ring of dimension $n$, with a Nakayama 
automorphism $\nu$ that respects the filtration $F$. 
\end{thm}

\begin{proof}
Let $\til{A} := \opn{Rees}^F(A)$ and $\bar{A} := \opn{Gr}^F(A)$.
We are given that $\bar{A}$ is AS regular of dimension $n$. According to 
Theorem \ref{thm:4546}(2), $\til{A}$ is AS regular of dimension 
$n + 1$. 

Consider the minimal graded-free resolution 
$\til{Q} \to \til{A}$ of $\til{A}$ over $\til{A}^{\mrm{en}}$.
We know, from Theorem \ref{thm:4548}  and Proposition \ref{prop:4266}, that 
each $\til{Q}^i$ is a 
finite graded-free 
$\til{A}^{\mrm{en}}$-module, and $\til{Q}^i = 0$ for $i < -n -1$. 
Therefore the complex 
\[ P := A^{\mrm{en}} \ot_{\til{A}^{\mrm{en}}} \til{Q} \cong 
A^{\mrm{en}} \ot^{\mrm{L}}_{\til{A}^{\mrm{en}}} \til{A} \in 
\dcat{D}(A^{\mrm{en}}) \]
is perfect. But according to Theorem \ref{thm:4511}(2), $A$ is a direct summand 
of $P$ in $\dcat{D}(A^{\mrm{en}})$. Thus $A$ is a perfect complex over 
$A^{\mrm{en}}$. This is property (ii) of Definition \ref{dfn:4490}.

By Corollary \ref{cor:3780} the graded ring $\til{A}$
has a balanced dualizing complex
$\til{R} := \til{A}(\til{\mu}, l)[n + 1]$,
where $\til{\mu}$ is a graded ring automorphism of $\til{A}$, and $l$ is an 
integer. Theorem \ref{thm:4400} says that $\til{R}$ is a graded rigid dualizing 
complex over $\til{A}$. 

Lemma \ref{lem:4393} says that for every $p$ the $\til{A}$-bimodule 
$\opn{H}^{p}(\til{R})$ is a central bimodule over $\opn{Cent}(\til{A})$.
But for $p = -n - 1$ we have 
$\opn{H}^{-n - 1}(\til{R}) \cong \til{A}(\til{\mu}, l)$,
and this implies that $\til{\mu}$ acts trivially on the element 
$t \in \opn{Cent}(\til{A})$. Hence $\til{\mu}$ extends to an 
automorphism of $\til{A}_t$, and there is an isomorphism
\begin{equation} \label{eqn:4507}
\opn{Ind}_{(\til{A}_t)^{\mrm{en}}} (\til{R}) \cong 
\til{A}_t(\til{\mu}, l)[n + 1]
\end{equation}
in $\dcat{D} \bigl( (\til{A}_t)^{\mrm{en}}, \mrm{gr} \bigr)$. 
Also there is an induced automorphism $\mu$ of the ring 
$A$, and formula (\ref{eqn:4506}) shows that $\mu$ respects the filtration $F$. 

According to Theorem \ref{thm:4395} the noetherian ring $A$ 
has a rigid NC dualizing complex 
\[ R := (\opn{Rest}_{A^{\mrm{en}}} \circ \opn{Deg}_0 \circ 
\opn{Ind}_{(\til{A}_t)^{\mrm{en}}}) (\til{R})[-1] \in \dcat{D}(A^{\mrm{en}}) . 
\]
From (\ref{eqn:4507}) we see that $R \cong A(\mu)[n]$ 
in $\dcat{D}(A^{\mrm{en}})$. Now use Proposition \ref{prop:4900} with 
$\nu := \mu^{-1}$. 
\end{proof}

\begin{exa} \label{exa:4505}
Let $\g$ be a finite Lie algebra over $\K$, with 
$\opn{rank}_{\K}(\g) = n$, and let 
$A := \opn{U}(\g)$ be its universal enveloping algebra. 
The ring $A$ is noetherian twisted CY of dimension $n$. This was worked out in 
the paper \cite{Ye4.5}, and is outlined below. 

The ring $A$ has a standard noetherian connected filtration $F$, such that 
$\bar{A} = \opn{Gr}^F(A) \cong \K[t_1, \ldots, t_n]$,
the commutative polynomial ring in $n$ variables, all of degree $1$. 
The ring $\bar{A}$ is AS regular of dimension $n$, as explained in
Example \ref{exa:3990}.
The Rees ring here is the homogeneous  universal enveloping ring from Example 
\ref{exa:4075}.

Theorem \ref{thm:4490} says that $A$ is twisted CY of dimension $n$, with a 
Nakayama automorphism $\nu$ that respects the filtration $F$. 
In case $\g$ is an abelian Lie algebra (i.e.\ the Lie bracket is zero), then  
$\nu$ is the identity automorphism. When $\g$ is semi-simple, 
Van den Bergh \cite{VdB2} proved that $\nu$ is trivial too. So in these extreme 
cases the ring $A = \opn{U}(\g)$ is CY (untwisted). 

The general case was done in \cite{Ye4.5}. Consider the character (i.e.\ 
rank $1$ representation) $\bwedge^n(\g)$ of $\g$, the $n$-th exterior power of 
the adjoint representation. This can be seen as a Lie algebra 
homomorphism $\ep : \g \to \K$. Then the dualizing automorphism $\mu$ is the 
unique ring automorphism of $A = \opn{U}(\g)$ such that 
$\mu(v) = v - \ep(v) \cd 1_A$
for $v \in \g \sub F_1(A) \sub A$.
The Nakayama automorphism is $\nu = \mu^{-1}$. 

The comultiplication of $A$ makes the tensor product of $A$-bimodules into a 
bimodule. This allows us to express $A(\mu)$ like this:
$A(\mu) \cong A \ot \bwedge^n(\g)$,
where the left action of $A$ on $\bwedge^n(\g)$ is trivial, and the right 
action is the coadjoint action.
\end{exa}

Given an $A$-bimodule $M$, its {\em $i$-th Hochschild homology} is the 
$\K$-module \lb 
$\opn{HH}_i(A, M) := \opn{Tor}_i^{A^{\mrm{en}}}(A, M)$,  
and its {\em $i$-th Hochschild cohomology} is 
$\opn{HH}^i(A, M) := \lb \opn{Ext}^i_{A^{\mrm{en}}}(A, M)$. 

\begin{thm}[Van den Bergh Duality, \cite{VdB2}] \label{thm:4860}
Let $A$ be an $n$-dimensional twisted CY ring, with Nakayama automorphism 
$\nu$. For every $M \in \dcat{M}(A^{\mrm{en}})$ and 
$0 \leq i \leq n$ there is an isomorphism 
\[ \opn{HH}^i(A, M) \cong \opn{HH}_{n - i} \bigl( A, M(\nu) \bigr) \] 
in $\dcat{M}(\K)$, and it is functorial in $M$. 
\end{thm}

The theorem does not require the rings $A$ nor $A^{\mrm{en}}$ to be noetherian.
We need a lemma first. 

\begin{lem} \label{lem:4895}
For a complex $M \in \dcat{D}(A^{\mrm{en}})$, and an automorphism $\mu$ of 
$A$, there is an isomorphism 
$A \ot^{\mrm{L}}_{A^{\mrm{en}}} M(\mu) \cong 
A(\mu) \ot^{\mrm{L}}_{A^{\mrm{en}}} M$
in $\dcat{D}(\K)$. This isomorphism is functorial in $M$. 
\end{lem}

\begin{proof}
It suffices to produce a functorial isomorphism 
$\psi : A \ot_{A^{\mrm{en}}} M(\mu) \iso A(\mu) \ot_{A^{\mrm{en}}} M$
for $M \in \dcat{M}(A^{\mrm{en}})$. 
We define $\psi$ on elements $b \in A$ and $m \in M$ by the formula
$\psi(b \ot m) := \mu(b) \ot m$.
It must be verified that the relations coming from 
the action of the ring $A^{\mrm{en}}$ are respected by $\psi$. 
Take a ring element $a_1 \ot \opn{op}(a_2) \in A^{\mrm{en}}$. Then in the 
module $A \ot_{A^{\mrm{en}}} M(\mu)$ we have equalities
\begin{equation} \label{eqn:4901}
\begin{aligned}
& 
(a_2 \cd b \cd a_1) \ot m = (b \cd (a_1 \ot \opn{op}(a_2))) \ot m = 
\\ & \quad
= b \ot ((a_1 \ot \opn{op}(a_2)) \, \cd^{\mu} \, m) = 
b \ot (a_1 \cd m \cd \mu(a_2)) . 
\end{aligned}
\end{equation}
See formula (\ref{eqn:4904}). In the module $A(\mu) \ot_{A^{\mrm{en}}} M$ 
we have equalities
\begin{equation} \label{eqn:4900}
\begin{aligned} 
& 
\mu(a_2 \cd b \cd a_1) \ot m = 
(\mu(a_2) \cd \mu(b) \cd \mu(a_1)) \ot m = 
\\ & \quad
= (\mu(b) \, \cd^{\mu} \, (a_1 \ot \opn{op}(\mu(a_2)))) \ot m 
% ? (b \cd (a_1 \ot \opn{op}(a_2))) \ot m = 
\\ & \quad
= \mu(b) \ot ((a_1 \ot \opn{op}(\mu(a_2))) \cd m) 
\\ & \quad
= \mu(b) \ot (a_1 \cd m \cd \mu(a_2))  .
\end{aligned}
\end{equation}
The isomorphism $\psi$ sends the first (resp.\ last) term in (\ref{eqn:4901}) 
to the first  (resp.\ last) term in (\ref{eqn:4900}). 
\end{proof}

\begin{proof}[Proof of Theorem \tup{\ref{thm:4860}}]
We are given an isomorphism 
\begin{equation} \label{eqn:4862}
A(\nu)[-n] \cong 
\opn{RHom}_{A^{\mrm{en}}} (A, A^{\mrm{en}})
\end{equation}
in $\dcat{D}(A^{\mrm{en}})$.
Consider these isomorphisms
\begin{equation} \label{eqn:4863}
\begin{aligned}
& A \ot^{\mrm{L}}_{A^{\mrm{en}}} M(\nu)[-n] \cong^{1} 
A(\nu)[-n] \ot^{\mrm{L}}_{A^{\mrm{en}}} M 
\\ & \quad 
\cong^{2} \opn{RHom}_{A^{\mrm{en}}} (A, A^{\mrm{en}}) 
\ot^{\mrm{L}}_{A^{\mrm{en}}} M
\\ & \quad 
\cong^{3} \opn{RHom}_{A^{\mrm{en}}} (A, A^{\mrm{en}}
\ot^{\mrm{L}}_{A^{\mrm{en}}} M)
\cong^4 \opn{RHom}_{A^{\mrm{en}}} (A, M)
\end{aligned}
\end{equation}
in $\dcat{D}(\K)$.
The first isomorphism is by Lemma \ref{lem:4895}.
The isomorphism $\cong^{2}$ comes from (\ref{eqn:4862}). 
The isomorphism $\cong^{3}$ is by derived tensor-evaluation, and it holds 
because $A$ is perfect over $A^{\mrm{en}}$, see Theorem 
\ref{thm:3400}. The isomorphism $\cong^4$ is clear. 

Taking $\opn{H}^i(-)$ in (\ref{eqn:4863}) we obtain an isomorphism 
\[ \opn{Tor}_{n - i}^{A^{\mrm{en}}} \bigl( A, M(\nu) \bigr) \cong 
\opn{Ext}^i_{A^{\mrm{en}}} (A, M) \]
in $\dcat{M}(\K)$. Functoriality is clear. 
\end{proof}

\begin{rem} \label{rem:4505}
Here is a quick discussion of the evolution of the {\em Calabi-Yau} concept. 
In algebraic geometry, an $n$-dimensional smooth proper scheme $X$ over a field 
$\K$ has the {\em canonical sheaf} $\bsym{\om}_X := \Om^n_{X / \K}$; this is a 
rank $1$ locally free $\OO_X$-module. 
In the paper \cite{BoKa0} from 1990, A.I. Bondal and M.M. Kapranov
introduced the {\em Serre functor}
\begin{equation} \label{eqn:4866}
\opn{S} := \bsym{\om}_X[n] \ot_{\OO_X} (-) . 
\end{equation}
This is an auto-equivalence of the derived category 
$\dcat{D}^{\mrm{b}}_{\mrm{c}}(X) = 
\dcat{D}^{\mrm{b}}_{\mrm{c}}(\cat{Mod} \OO_X)$
which admits a bifunctorial isomorphism 
\begin{equation} \label{eqn:4865}
\opn{Hom}_{\dcat{D}^{\mrm{b}}_{\mrm{c}}(X)}( \MM, \NN) \cong 
\opn{Hom}_{\dcat{D}^{\mrm{b}}_{\mrm{c}}(X)}( \NN, \opn{S}(\MM) )^* .
\end{equation}
Here $(-)^* := \opn{Hom}_{\K}(-, \K)$. 
This is an upgrade of the familiar Serre Duality for $X$, which is the 
special case where $\MM = \OO_X$ and $\NN$ is the translation of a locally free 
$\OO_X$-module. 

The Serre functor makes sense on any $\K$-linear triangulated category 
$\cat{D}$ that has suitable finiteness properties (``smooth'' and ``proper''). 
It is a triangulated auto-equivalence $\opn{S}$ of $\cat{D}$ admitting a 
bifunctorial isomorphism
\begin{equation} \label{eqn:4872}
\opn{Hom}_{\cat{D}}(M, N) \cong \opn{Hom}_{\cat{D}}(N, \opn{S}(M))^* 
\end{equation}
in $\dcat{M}(\K)$ for $M, N \in \cat{D}$. In \cite{BoKa0} it was proved that if 
a Serre functor exists, then it is unique. In many situations (where this makes 
sense) the Serre functor is realized by tensoring with the rigid dualizing 
complex. This is true in the geometric context above, in which the rigid 
dualizing complex of the scheme $X$ is $\bsym{\om}_X[n]$; see formula 
(\ref{eqn:4866}). It is also true for a regular finite $\K$-ring $A$, as 
explained in Example \ref{exa:4870} below. 

Back in algebraic geometry, a smooth proper scheme $X$ is called Calabi-Yau 
if the canonical sheaf $\bsym{\om}_X$ is trivial. (Examples are abelian 
varieties and K3 surfaces.) The CY condition implies that the Serre functor 
$\opn{S}$ is isomorphic to $\opn{T}^n$, the $n$-th power of the 
translation functor of $\cat{D}  = \dcat{D}^{\mrm{b}}_{\mrm{c}}(X)$. 
Thus the dimension of $X$ is precisely the ratio between the 
auto-equivalences $\opn{S}$ and $\opn{T}$. 

Around the year 1994, M. Konstevich began discussing {\em Homological Mirror 
Symmetry}, which is a vast and deep program (still not fully understood) 
relating the algebraic geometry of a CY scheme $X$ (as above) with the {\em 
symplectic geometry} of its mirror partner manifold $Y$.  See \cite{HMS}
for more information on this program. 

Among the ideas arising in the discussion 
by Kontsevich was that of a {\em Calabi-Yau category}. This is a smooth 
proper $\K$-linear triangulated category $\cat{D}$, whose translation functor 
$\opn{T}$ and Serre functor $\opn{S}$ satisfy the numerical relation 
$\opn{S} \cong \opn{T}^n$ for some integer $n$, called the {\em CY dimension} 
of $A$. More generally, if the relation $\opn{S}^m \cong \opn{T}^n$ holds for 
some $m, n \in \Z$, then $\cat{D}$ is called  {\em fractionally CY}
of dimension $(n, m)$. By abuse of notation, such a ring $A$ is sometimes 
called fractionally CY of dimension $n / m$. 

The concept of CY triangulated categories later showed up in various 
mathematical areas, including:
\begin{itemize}
\item Algebraic geometry (see \cite{Kuz}). 

\item Conformal field theories (see \cite{Cost}).

\item Representation theory of algebras and cluster categories
(see  \cite{CCS}, \cite{BMRRT}, \cite{Kel4}, \cite{Kel2}, \cite{Lad} and 
Example \ref{exa:4870} below).
\end{itemize}

In the preprint \cite{Ginz} from 2006, V. Ginzburg defined {\em Calabi-Yau 
algebras}. At about the same time K. Brown and J.J. Zhang \cite{BrZh} 
introduced {\em rigid Gorenstein algebras}. These are both precisely 
the CY rings from Definition \ref{dfn:4490} above.
Eventually the CY terminology was adopted by ring theorists. 

In case $A$ is a noetherian connected graded ring, then it is twisted CY in the 
graded sense if and only if it is AS regular; see \cite{ReRoZh}.
More results on graded CY rings can be found in \cite{ReRo} and its references.
\end{rem}

We finish the subsection (and the book) with an example of a fractionally CY 
category. 

\begin{exa} \label{exa:4870}
Suppose $A$ is a finite $\K$-ring. (Traditional texts would say that 
$A$ is a finite dimensional $\K$-algebra.)  
We claim that the rigid dualizing complex of $A$ is $R := A^{*}$, the 
$\K$-linear dual of $A$. That $R \in \dcat{D}(A^{\mrm{en}})$ is a NC dualizing 
complex over $A$ is very easy to see (cf.\ Definition \ref{dfn:3675}). 
As for rigidity: in this case it is just linear algebra. Indeed, 
$R \ot R = A^* \ot A^* \cong (A^{\mrm{en}})^*$ 
in $\dcat{M}(A^{\mrm{four}})$, and hence there are canonical isomorphisms
\[ \opn{Sq}_A(R) = \opn{RHom}_{A^{\mrm{en}}}(A, R \ot R) \cong 
\opn{Hom}_{A^{\mrm{en}}}(A, (A^{\mrm{en}})^*) \cong A^* = R \]
in $\dcat{D}(A^{\mrm{en}})$.

Now let's assume that $A$ is regular, i.e.\ it has finite global 
cohomological dimension. (If the base field $\K$ is perfect, one can show 
that $A$ is homologically smooth over it; see \cite{Hap2} and \cite{BGMS}.)
Then, according to Theorems \ref{thm:3416} and \ref{thm:3400}, 
every object of $\dcat{D}^{\mrm{b}}_{\mrm{f}}(A)$ is a compact object of 
$\dcat{D}(A)$. This includes $R$. So $R$ satisfies condition (iii) of 
Theorem \ref{thm:3455}, and therefore it is a tilting complex over $A$. 
We see that the triangulated functor 
\[ \opn{S} := R \ot^{\mrm{L}}_{A} (-) : \dcat{D}^{\mrm{b}}_{\mrm{f}}(A)
\to \dcat{D}^{\mrm{b}}_{\mrm{f}}(A) \]
is an auto-equivalence. 

It turns out that $\opn{S}$ is a Serre functor of 
$\dcat{D}^{\mrm{b}}_{\mrm{f}}(A)$. Let us prove this. 
For $M, N \in \dcat{D}^{\mrm{b}}_{\mrm{f}}(A)$
we have canonical isomorphisms
\[ \begin{aligned}
& \opn{RHom}_{A}(N, \opn{S}(M)) = \opn{RHom}_{A}(N, A^* \ot^{\mrm{L}}_{A} M)
\\ & \quad 
\cong^{\dag}  \opn{RHom}_{A}(N, A^*) \ot^{\mrm{L}}_{A} M \cong 
N^* \ot^{\mrm{L}}_{A} M 
\end{aligned} \]
in $\dcat{D}(\K)$. 
The isomorphism $\cong^{\dag}$ is by derived tensor-evaluation, and it holds 
because $N$ is a compact object of $\dcat{D}(A)$. 
Therefore we get functorial isomorphisms 
\begin{equation} \label{eqn:4874}
\begin{aligned}
& \opn{RHom}_{A}(N, \opn{S}(M))^* \cong (N^* \ot^{\mrm{L}}_{A} M)^* 
\cong \opn{RHom}_{\K}( N^* \ot^{\mrm{L}}_{A} M, \K) 
\\ & \quad 
\cong^{\heartsuit} \opn{RHom}_{A}(M, \opn{RHom}_{\K}(N^*, \K)) 
\cong \opn{RHom}_{A}(M, N) 
\end{aligned}
\end{equation}
in $\dcat{D}(\K)$. 
The isomorphism $\cong^{\heartsuit}$ is by adjunction for the ring homomorphism 
$\K \to A$. Taking $\opn{H}^0$ in (\ref{eqn:4874}) we obtain a functorial 
isomorphism 
\[ \opn{Hom}_{\dcat{D}^{\mrm{b}}_{\mrm{f}}(A)}(N, \opn{S}(M))^*
 \cong \opn{Hom}_{\dcat{D}^{\mrm{b}}_{\mrm{f}}(A)}(M, N)  \]
in $\dcat{M}(\K)$. 

Finally we specialize to the setting of Example \ref{exa:4180}.
So $A$ is the ring of upper triangular $n \times n$ matrices over $\K$ for some 
$n \geq 2$. As calculated by A. Yekutieli and E. Kreines in \cite{Ye4}, there 
is an isomorphism 
\[ \underbrace{R \ot^{\mrm{L}}_{A} \cdots \ot^{\mrm{L}}_{A} R}_{n + 1} 
\cong A[n - 1] \]
in $\dcat{D}(A^{\mrm{en}})$. Since the translation functor of
$\dcat{D}^{\mrm{b}}_{\mrm{f}}(A)$ is 
$\opn{T} \cong A[1] \ot^{\mrm{L}}_{A} (-)$, we obtain an isomorphism of 
functors $\opn{S}^{n + 1} \cong \opn{T}^{n - 1}$. 
The conclusion is that the category 
$\dcat{D}^{\mrm{b}}_{\mrm{f}}(A)$ is fractionally CY of dimension 
$(n - 1, n + 1)$.
\end{exa}






